\documentclass[11pt]{article}
\usepackage[T1]{fontenc}
\usepackage{lmodern,microtype,amsmath,amssymb,amsthm,mathtools}
\usepackage[a4paper,margin=29mm]{geometry}
\usepackage[round,authoryear]{natbib}
\usepackage{booktabs,array,longtable}
\usepackage{xurl}
\usepackage[hidelinks]{hyperref}
\makeatletter
\renewcommand{\l@subsection}{\@dottedtocline{2}{1.5em}{3.2em}}
\makeatother
\newcommand{\workbenchurl}{https://github.com/KokunoYumeto/schemes-and-moduli-workbench}
\newcommand{\recordurl}{https://github.com/KokunoYumeto/mathematics-workbenches}
\newcommand{\F}{\mathbf F}
\newcommand{\Q}{\mathbf Q}
\newcommand{\Z}{\mathbf Z}
\newcommand{\Ocal}{\mathcal O}
\newcommand{\D}{\mathbb D}
\newcommand{\Gal}{\operatorname{Gal}}
\newcommand{\Tr}{\operatorname{Tr}}
\newcommand{\Nm}{\operatorname{N}}
\newcommand{\GL}{\operatorname{GL}}
\newcommand{\ord}{\operatorname{ord}}
\newcommand{\nuMin}{\nu_{\min}}
\newtheorem{theorem}{Theorem}[section]
\newtheorem{proposition}[theorem]{Proposition}
\newtheorem{lemma}[theorem]{Lemma}
\newtheorem{corollary}[theorem]{Corollary}
\theoremstyle{definition}
\newtheorem{definition}[theorem]{Definition}
\newtheorem{question}[theorem]{Question}
\theoremstyle{remark}
\newtheorem{remark}[theorem]{Remark}
\newtheorem{example}[theorem]{Example}
\numberwithin{equation}{section}
\title{Ramification thresholds for Grothendieck's inertia-character question,\protect\\ after Bondarko and Drinfeld}
\author{}
\date{23 November 2026}
\begin{document}
\maketitle
\begin{abstract}
Grothendieck asked whether finite torsion determines the special inertia character of a potentially good group. For the full height-two Barsotti--Tate filtration class over complete mixed-characteristic fields containing $\mu_p$, we determine the character index at every absolute ramification index $e$ prime to $p$: for $p\ge5$, it is $\lceil(1+\log_p e)/2\rceil$. Thus index $r\ge2$ occurs exactly when $p^{2r-3}<e<p^{2r-1}$, with $(p-1)\mid e$; the minimal cyclotomic index has value one. For $n\ge2$, actual $p^n$-torsion recovers the character when $e\le p^{n-1}$. For two one-slope models, isomorphic $p^n$-torsion determines the special character when $n\ge v_p(e)+2$; a comparison of different Newton polygons forces $e>p^{2n-2-v_p(e)}$. The mixed comparison threshold is exactly $e>p^{2n-3-v_p(e)}$. For $r\ge0$ and $2\le k\le(p^2-1)/2$, the index at $e=kp^r(p-1)$ is $r+2$ when $p\nmid k$ or $k=p$, and $r+3$ for the other multiples of $p$. At the next coefficient and $r\ge1$, any unequal-character comparison must have equal Hasse parameters, opposite tangent scalars and exchanged canonical lines in split residual representations. At that boundary, the finite image after the common tame good extension cannot lie in a principal nonsplit Cartan, and a tame quadratic self-twist with irreducible intertwiner reduction is excluded. At arbitrary indices, the two primitive Frobenius chains control the uniform bound. The prime-to-$p$ mechanism is central $p^2$ rigidity, proved by computing and transporting an exact lower break between actual point fields. A pointwise graph valuation bound allows successive canonical quotients while retaining two torsion levels, and excludes every residual branch at the first wild index. An equivariant construction with prescribed degree-$p$ stages supplies actual lower examples at every permitted mixed index. For elliptic curves, a separate theorem gives the full rational $p^n$-torsion criterion for every tame defect and permitted index, including its dependence on $v_p(e/(p-1))$ and the least level forcing good reduction. Equivariant deformation and transverse specialization produce actual realizations; wild characteristic-three defects are treated separately. In every dimension, a potentially ordinary abelian variety with unramified $p$-torsion has good reduction for odd $p$, without a roots-of-unity assumption. The height-one character index is zero for every prime. Explicit Hesse specializations give full rational three-torsion and exact variable defects, including twelve; a residue-field obstruction excludes additive potentially good full $p$-torsion over $\Q_p(\zeta_p)$ for $p\ge5$. The classical logarithmic cyclotomic extension gives an actual filtered height-two group that is not potentially Hodge--Tate, clarifying the scope of the filtration condition. The mathematics and text were produced and checked by GPT-6 Astra (ChatGPT Pro mode); no human has checked this paper.
\end{abstract}
\tableofcontents
\section{Finite torsion, central rigidity and ramification}
This section states the complete prime-to-$p$ character index, the mixed threshold at every index, and infinite exact wild bands in the full height-two filtration class. The prime-to-$p$ comparison uses a degree-$p$ extension between actual point fields above a canonical quotient. Central $p^2$ rigidity and repeated quotients give its sharp upper range. Primitive coordinates give uniform recovery without a restriction on $v_p(e)$, while the pointwise graph bound controls canonical iteration. Integral constructions with prescribed degree-$p$ stages attain the mixed thresholds. A comparison of stable Frobenius spans treats all one-slope cases in the remaining proposed wild range. At the first coefficient beyond the exact bands, an actual comparison cannot have its finite image contained in a principal nonsplit Cartan after the common tame good extension. A separate problem asks which indices allow an additive elliptic curve to have full rational torsion.

Throughout, $K$ is a complete discretely valued field of characteristic zero with perfect residue field of characteristic $p>0$. We normalize $v_K(K^\times)=\Z$ and write $e_K=v_K(p)$. A statement of existence at absolute ramification index $e$ allows the field $K$ to vary with $e$. The constructions use residue field $\overline{\F}_p$; they make no assertion for every field of that index.

For an elliptic curve, \emph{full rational $p^n$-torsion} means that all of $E[p^n](\overline K)$ lies in $E(K)$. If $E$ is potentially good, its \emph{semistability defect} is the degree of its minimal good-reduction extension after passage to the completed maximal unramified base. The character-recovery problem concerns a Barsotti--Tate group $M/K$ admitting a finite filtration whose graded pieces acquire good reduction. Over a common finite Galois extension where those pieces have good models, inertia acts on their rational covariant Dieudonn\'e modules. The sum of these finite characters, evaluated on elements of order prime to $p$ in the finite inertia quotient, will be denoted $\theta_M$. This is the character defined in Grothendieck's \emph{Cristaux} manuscripts, cote~9, sheet~11.

On sheet~12 Grothendieck asks, for elliptic curves with isomorphic $p$-torsion,
\begin{quote}
``est-il vrai que les fonctions traces $\phi_{E_1}$ et $\phi_{E_2}$ sur le groupe d'inertie [\ldots] sont \'egales ?''
\end{quote}
The manuscript restricts the comparison to elements of order prime to $p$ and explicitly assumes that $K$ contains the $p$th roots of unity. The quotation is from archive sheet~12, which is PDF page~13; the character definition occupies archive sheet~11, PDF page~12 \citep[pp.~11--13]{Grothendieck1969}. We study the refined question of the least higher truncation sufficient for that character, within the stated filtration class.

We use actual finite torsion throughout: an isomorphism $M[p^n]\simeq M'[p^n]$ is an isomorphism of finite group schemes over $K$, equivalently of their finite Galois modules. It is not an isomorphism of rational Tate modules. In every character-index assertion we restrict the base class by $\mu_p\subset K$. Let $\nuMin(p,e,h)$ be the least nonnegative integer $\nu$ for which this character is determined uniformly, among allowed height-$h$ groups over bases of index $e$, by $M[p^{\nu+1}]$. Thus the index $\nu=2$ refers to $p^3$-torsion.

The allowed filtration condition is the monodromy condition defined on archive sheet~10, PDF page~11: a finite composition series of Barsotti--Tate groups whose graded pieces acquire good reduction after finite separable extension. Sheet~11, PDF page~12, defines its sum character \citep[sheets~10--11]{Grothendieck1969}. Proposition~\ref{prop:filtered-non-ht} shows that this condition can hold when the whole Tate representation is not potentially Hodge--Tate. The results below comparing an elliptic curve to its full rational torsion do not by themselves determine $\nuMin$.

\subsection{The complete prime-to-$p$ index}
The first theorem determines the character index at every absolute ramification index prime to $p$, in the full height-two filtration class defined above.

\begin{theorem}[The complete prime-to-$p$ character index]\label{thm:complete-prime-index}
Let $p\ge5$ and let $e$ be a positive multiple of $p-1$ with $p\nmid e$. Consider complete discretely valued characteristic-zero fields with perfect residue field of characteristic $p$, absolute index $e$, and $\mu_p\subset K$. In the height-two Barsotti--Tate filtration class,
\begin{equation}\label{eq:complete-prime-index}
 \nuMin(p,e,2)=\left\lceil\frac{1+\log_p e}{2}\right\rceil.
\end{equation}
Equivalently, $\nuMin(p,p-1,2)=1$, and for every integer $r\ge2$,
\[
 \nuMin(p,e,2)=r\qquad\text{if}\qquad p^{2r-3}<e<p^{2r-1}.
\]
Put $r=\nuMin(p,e,2)$. The upper statement is uniform over every field in the specified base class: $p^{r+1}$-torsion determines the special prime-to-$p$ inertia character. At every such $e$ there are actual allowed groups over some field with residue $\overline{\F}_p$, isomorphic $p^r$-torsion and different special characters. These lower examples are Barsotti--Tate groups; no assertion that every pair comes from elliptic curves is made.
\end{theorem}

The strict endpoints cause no ambiguity because $p\nmid e$. In particular, the complete intervals for indices two, three and four are
\[
 p<e<p^3,\qquad p^3<e<p^5,\qquad p^5<e<p^7,
\]
respectively, with the same divisibility and base-field restrictions. The proof combines an upper comparison theorem with actual lower constructions at every index.

\begin{theorem}[Iterated supersingular recovery]\label{thm:ss-iterated-range}
Let $p\ge5$, $n\ge3$, and let $K$ be a complete discretely valued field of characteristic zero with perfect residue field of characteristic $p$, containing $\mu_p$. Put $e=e_K$ and assume $p\nmid e$. If
\begin{equation}\label{eq:ss-iterated-range}
 e<\frac{(p-1)p^{3n-3}}{p^n-1},
\end{equation}
then any isomorphism $G[p^n]\simeq H[p^n]$ between height-two Barsotti--Tate groups with potentially supersingular reduction implies $\theta_G=\theta_H$. Since $(p-1)\mid e$, the numerical condition is equivalently
\[
 e\le(p-1)\bigl(p^{2n-3}+p^{n-3}\bigr).
\]
\end{theorem}

\begin{corollary}[The full filtration range]\label{cor:complete-prime-recovery}
Under the same base hypotheses, for every $n\ge3$ one has
\[
 e<p^{2n-3}\quad\Longrightarrow\quad\nuMin(p,e,2)\le n-1.
\]
\end{corollary}

Section~\ref{sec:canonical-iteration} proves these statements by repeatedly taking a canonical quotient of one model and the quotient by its transported kernel in the other. Each step loses exactly one torsion level. The graph valuation bound ensures that the first model enters the central range while at least $p^2$-torsion remains, so Theorem~\ref{thm:central-p2-rigidity} applies. The transported kernel need not be canonical in the partner. The sharp mixed obstruction then supplies both the full-class upper threshold and the matching lower examples.

\subsection{Uniform recovery and exact wild bands}
Primitive coordinates give a uniform comparison without a restriction on the $p$-part of the absolute ramification index; a compatible cyclic-point construction gives exact mixed thresholds and infinitely many sharp full-class values.

\begin{theorem}[A uniform range at every index]\label{thm:all-index-upper}
Let $p\ge5$ and $n\ge2$. In the standing height-two filtration class, if
\[
                           e_K\le p^{n-1},
\]
then isomorphic actual $p^n$-torsion forces equality of the special inertia characters. This holds over every standing complete field of that index, with no restriction on $v_p(e_K)$. Equivalently,
\[
 \nuMin(p,e,2)\le\left\lceil\log_p e\right\rceil
\]
for every positive multiple $e$ of $p-1$.
\end{theorem}

\begin{theorem}[The mixed threshold at every index]\label{thm:mixed-all-indices}
Let $p\ge5$, $n\ge2$ and let $e$ be a positive multiple of $p-1$. Put $v=v_p(e)$. Over some standing field of exact index $e$, a whole potentially supersingular height-two group and a group filtered by two potentially good height-one groups of total dimension one have isomorphic actual $p^n$-torsion and different special characters if and only if
\[
                              e>p^{2n-3-v}.
\]
The existence field may always have residue field $\overline{\F}_p$. A negative exponent is interpreted in $\Q_{>0}$.
\end{theorem}

\begin{theorem}[Infinite exact wild bands]\label{thm:wild-bands-exact}
For every $p\ge5$, integer $r\ge0$ and integer $2\le k\le(p^2-1)/2$,
\[
 \nuMin\bigl(p,kp^r(p-1),2\bigr)=
 \begin{cases}
 r+2,&p\nmid k\ \text{or }k=p,\\
 r+3,&p\mid k,\ k\ge2p.
 \end{cases}
\]
In the first case, $p^{r+3}$-torsion recovers the special character over every standing field of that exact index, while some such field supports two actual permitted groups with isomorphic $p^{r+2}$-torsion and different special characters. In the second case, these two levels are $p^{r+4}$ and $p^{r+3}$, respectively. In particular,
\[
 \begin{gathered}
 \nuMin\bigl(p,kp(p-1),2\bigr)=3\quad(2\le k\le2p-1),\\
 \nuMin\bigl(p,p^s(p-1),2\bigr)=s+1\quad(s\ge1).
 \end{gathered}
\]
At $k=2p$, this gives
\[
 \nuMin\bigl(p,2p^{r+1}(p-1),2\bigr)=r+3.
\]
\end{theorem}

Section~\ref{sec:one-slope-all-indices} proves the uniform theorem by comparing two Frobenius chains and the larger exact band by matching residual characters and transporting a canonical quotient. Section~\ref{sec:all-wild-indices} gives a separate graph proof of the sufficient range $e\le(p-1)p^{n-2}$, proves the exact mixed criterion and constructs the lower examples for the bands. The mixed criterion determines exactly when an actual group of one reduction type can have the same truncation as a group of the other type. Comparisons between two supersingular groups outside the proved ranges remain a separate question, so the infinite exact bands do not constitute a formula for every later wild index.


\subsection{Primitive comparisons at arbitrary indices}
The next results distinguish Newton polygons without assuming that the absolute index is prime to $p$. Let $G,H$ be whole potentially supersingular height-two groups over a standing base $K$, with $p\ge5$, $\mu_p\subset K$ and $v=v_p(e_K)$. After completed unramified base change, use the common tame good field $L/K$ of degree $p^2-1$ and normalize $v_L(L^\times)=\Z$. A one-slope model has all nonzero $p$-point values equal to $e_K$. A two-slope model has an integer parameter $z$ with
\[
 0<z<pe_K,\qquad
 (p+1)e_K-z\quad(p-1\text{ points}),\qquad
 z/p\quad(p^2-p\text{ points}).
\]

\begin{theorem}[One-slope recovery]\label{thm:one-slope-all-indices}
If both good models have one first-level Newton slope and
\[
 n\ge v_p(e_K)+2,\qquad G[p^n]\simeq H[p^n],
\]
then $\theta_G=\theta_H$. There is no upper bound on $e_K$.
\end{theorem}

\begin{theorem}[An obstruction to different polygons]\label{thm:mixed-polygon-obstruction}
Let $n\ge1$. If $G[p^n]\simeq H[p^n]$, one model has one slope and the other has two slopes with parameter $z$, then
\begin{equation}\label{eq:mixed-polygon-obstruction}
 v_p(z)\ge\max\{1,2n-1-v_p(e_K)\},\qquad
 e_K>p^{\,2n-2-v_p(e_K)}.
\end{equation}
This assertion does not require unequal special characters.
\end{theorem}

\begin{corollary}[The remaining two-slope comparison]\label{cor:all-index-two-slope-reduction}
Let $p\ge5$, $n\ge2$, and let $M,M'$ belong to the full standing height-two filtration class over $K$. Suppose
\[
 e_K\le p^{\,2n-3-v_p(e_K)},\qquad M[p^n]\simeq M'[p^n].
\]
Their total dimensions and finite determinant factors agree. If their special characters differ, both groups are whole potentially supersingular groups and both good models have two first-level Newton slopes. In every other structural or polygon case the characters agree.
\end{corollary}

Define $\nu_{\mathrm{SS}}(p,e)$ to be the least nonnegative integer $r$ such that actual $p^{r+1}$-torsion recovers the special character uniformly among pairs of whole potentially supersingular height-two groups over the standing fields of index $e$.

\begin{theorem}[The full index and the supersingular index]\label{thm:full-supersingular-maximum}
For every $p\ge5$ and positive multiple $e$ of $p-1$, put $v=v_p(e)$. Then
\[
 \nuMin(p,e,2)=\max\left\{\nu_{\mathrm{SS}}(p,e),
       \left\lceil\frac{1+\log_p e+v}{2}\right\rceil\right\}.
\]
The index on the left concerns the full filtration class. The identity does not evaluate $\nu_{\mathrm{SS}}$ at every wild index.
\end{theorem}

\begin{theorem}[Reduction to a central comparison]\label{thm:all-index-central-reduction}
Under the hypotheses of Corollary~\ref{cor:all-index-two-slope-reduction}, suppose the special characters differ. For the two two-slope models one has
\[
 \max\{v_p(z_1),v_p(z_2)\}\ge v+2,\qquad v=v_p(e_K).
\]
Choose a model $G$ attaining this bound. There are $c$ successive quotients of $G$ by its canonical $p$-lines, paired with quotients of the other model by their transported lines, with
\[
 0\le c\le n-v-2,\qquad r=n-c\ge v+2,
\]
which give an unequal-character actual $p^r$-torsion comparison. The chosen terminal model has parameter $z_c=p^c z_G$ satisfying
\[
 e_K<z_c<pe_K,\qquad
 v_p\bigl((p+1)e_K-z_c\bigr)=v_p(e_K).
\]
The partner is whole potentially supersingular and may have either Newton polygon. The original numerical bound at level $n$ is not asserted at the smaller level $r$.
\end{theorem}

Section~\ref{sec:one-slope-all-indices} proves these assertions. The scaled torsion graph has $2n$ primitive Frobenius dimensions. A one-slope coordinate has a two-dimensional stable Frobenius span, while a two-slope coordinate has a one-dimensional span; this distinction controls overlaps of their nonzero weights. The primitive correspondence is classical \citep[\S\S2.1--2.2, pp.~3--4]{Koch2015}. The two-chain divisibility argument also proves Theorem~\ref{thm:all-index-upper}. At a failed canonical-line matching, Lemma~\ref{os-lem:canonical-matching} supplies an additional truncated valuation-sum inequality. For $v=0$, central $p^2$ rigidity excludes the terminal comparison; for $v>0$, the corresponding assertion with a $p$-divisible high numerator remains open.

\subsection{The first wild character index}
The graph valuation bound also gives an exact full-class answer at the first absolute ramification index divisible by $p$ under the cyclotomic hypothesis.

\begin{theorem}[The general first wild index]\label{thm:first-wild-general}
Let $p\ge5$. For the height-two Barsotti--Tate filtration class over complete discretely valued characteristic-zero fields with perfect residue field, $\mu_p\subset K$ and $e_K=p(p-1)$, one has
\[
                         \nuMin(p,p(p-1),2)=2.
\]
Thus $p^3$-torsion determines the special prime-to-$p$ inertia character over every such field. There are actual unequal-character comparisons at level $p^2$ over a field of that exact index with residue field $\overline{\F}_p$. The lower examples consist of a whole potentially supersingular group and a group filtered by two potentially good height-one groups.
\end{theorem}

Section~\ref{sec:first-wild-general} gives the full proof. For two whole supersingular groups the graph estimate would require
\[
p(p-1)(p^3-1)\ge(p-1)p^4,
\]
whose left side is smaller by $p(p-1)$. The mixed obstruction and determinant comparison settle the other filtration types. An invariant Eisenstein specialization of the cyclic-point construction attains the lower level.


\subsection{The smallest character index}
The next theorem determines the least torsion level at the smallest possible absolute index under the cyclotomic hypothesis.

\begin{theorem}[Minimal cyclotomic ramification]\label{thm:minimal-index}
Let $p\ge5$. For the height-two filtration class over complete discretely valued characteristic-zero fields with perfect residue field, $\mu_p\subset K$ and $e_K=p-1$,
\[
                         \nuMin(p,p-1,2)=1.
\]
Thus $p^2$-torsion determines the special prime-to-$p$ inertia character over every such field. Over $K=\operatorname{Frac}W(\overline{\F}_p)(\zeta_p)$ there are two whole potentially supersingular groups with isomorphic $p$-torsion and different special characters.
\end{theorem}

\begin{corollary}[The exact index-two range]\label{cor:exact-index-two}
Let $p\ge5$ and let $e$ be a positive multiple of $p-1$ with $p\nmid e$. In the same height-two filtration class, with $\mu_p$ in every base field,
\[
              \nuMin(p,e,2)=2\quad\Longleftrightarrow\quad p<e<p^3.
\]
Equivalently $e=t(p-1)$ with $2\le t\le p^2+p+1$ and $p\nmid t$. In particular, at $p=5$ the index is one at $e=4$, and two at every multiple of four from $8$ through $124$ that is prime to five.
\end{corollary}

Section~\ref{sec:minimal-index} proves the theorem by combining the graph valuation bound with an explicit Lubin--Tate twist. The twist changes the characteristic-zero descent trace while Frobenius on $\F_{p^2}$ identifies its first-level representation. The sharp mixed examples then identify both ends of the index-two range. These statements give an exact answer in a Barsotti--Tate class; they do not assert that every lower example is the torsion of an elliptic curve.

\subsection{Central rigidity and prime-to-$p$ recovery}
For the next two theorems, $G$ and $H$ are height-two Barsotti--Tate groups with potentially supersingular reduction, and $p\ge5$. Pass to the completed maximal unramified base, without changing $e=e_K$, and take the common tame good-reduction extension $L/K$ of degree $m=p^2-1$. Normalize $v_L(L^\times)=\Z$. For a two-slope good model, its normalized Hasse value $z$ is the positive integer specified by the nonzero $p$-torsion valuations
\[
 A=(p+1)e-z\quad(p-1\text{ points}),\qquad
 B=z/p\quad(p^2-p\text{ points}),\qquad 0<z<pe.
\]
Equivalently, its deformation Hasse valuation is $h=(p-1)z$. The points of larger valuation form the canonical subgroup $C$. We call the range $e<z<pe$ \emph{central}; it is exactly the range $A<p^2B$ in which every $p$-division point of a nonzero point of $C$ has valuation $A/p^2$. These conventions and their independence of integral parameter changes are justified in Lemma~\ref{rc-lem:newton}.

\begin{theorem}[Central $p^2$ rigidity]\label{thm:central-p2-rigidity}
Let $K$ be a complete discretely valued field of characteristic zero with perfect residue field of characteristic $p\ge5$, and suppose $\mu_p\subset K$. Let $G,H/K$ be height-two Barsotti--Tate groups with potentially supersingular reduction. Suppose the good model of $G$ over the common tame field $L$ has normalized Hasse value $z_G$ satisfying
\[
 e_K<z_G<pe_K,\qquad p\nmid A_G=(p+1)e_K-z_G.
\]
Then every isomorphism $G[p^2]\simeq H[p^2]$ over $K$ implies
$\theta_G=\theta_H$. There is no restriction on $v_p(e_K)$, and no Newton-slope or Hasse-divisibility assumption on $H$.
\end{theorem}

The proof retains two fields attached to a point $Q$ with $pQ\in C\setminus\{0\}$:
\[
 F=L(Q),\qquad F_0=L(q(Q)),\qquad q:G\longrightarrow G/C.
\]
Their degrees over $L$ are $p^2$ and $p$, respectively, and $F/F_0$ is cyclic with exact lower break $(p^2-1)A_G$. A finite torsion isomorphism identifies both fields through their point stabilizers. The lower-break inequality for the transported coordinate then opposes the strict residual-break inequality of a quotient with a $p$-divisible Hasse value. Section~\ref{sec:central-ramification} proves every field-degree, normality and transport assertion used here.

\begin{theorem}[Recovery with one canonical quotient]\label{thm:ss-prime-range}
Let $p\ge5$ and $n\ge3$. Let $K$ be a complete discretely valued field of characteristic zero with perfect residue field of characteristic $p$, containing $\mu_p$, and put $e=e_K$. Suppose $p\nmid e$ and
\begin{equation}\label{eq:ss-prime-range}
 e<\frac{(p-1)p^{2n}}{p^n-1}.
\end{equation}
For any two height-two Barsotti--Tate groups $G,H/K$ with potentially supersingular reduction, an isomorphism $G[p^n]\simeq H[p^n]$ implies $\theta_G=\theta_H$.
\end{theorem}

An invariant augmentation-ideal estimate gives, in a hypothetical unequal-character comparison, the pointwise lower bound $(p-1)/(p^n-1)$ on every nonzero torsion coordinate. For a two-slope model this forces $z/p^{2n-1}\ge(p-1)/(p^n-1)$. Under \eqref{eq:ss-prime-range}, a $p$-divisible Hasse value is either already central or becomes central after one canonical quotient. Theorem~\ref{thm:central-p2-rigidity} then applies. The one-slope case is treated without an upper bound on $e$. This proves Theorem~\ref{thm:ss-prime-range} in Section~\ref{sec:central-ramification}.

\begin{corollary}[Consequences of one canonical quotient]\label{cor:prime-recovery-intro}
For $p\ge5$, a positive multiple $e$ of $p-1$ prime to $p$, and the base class $\mu_p\subset K$, one has
\[
 \nuMin(p,e,2)\le n-1
 \quad\begin{cases}
 n=3,&e<p^3,\\
 n\ge4,&e\le(p-1)(p^n+1).
 \end{cases}
\]
In particular, the following exact values hold:
\begin{align*}
 \nuMin(p,e,2)&=2&&\text{if }p<e<p^3,\\
 \nuMin(p,e,2)&=3&&\text{if }p^3<e\le(p-1)(p^4+1),\\
 \nuMin(p,e,2)&=4&&\text{if }p^5<e\le(p-1)(p^5+1).
\end{align*}
The filtration class and index convention are those fixed above.
\end{corollary}

Theorem~\ref{thm:prime-index-upper} and Corollary~\ref{cor:prime-index-intervals} prove the full-class range $e<p^n$ by recovery of canonical lines. Section~\ref{sec:central-ramification} gives the displayed one-quotient bounds, while Section~\ref{sec:canonical-iteration} proves the complete ranges of Theorem~\ref{thm:complete-prime-index} by iteration. The two comparisons use different ranges of canonical quotient geometry. At $n=3$, the exact mixed obstruction requires the full-class restriction $e<p^3$, even though the supersingular-only range is larger. Section~\ref{sec:primitive-grading} gives a further primitive-function reduction and excludes all three of its numerical configurations by Theorem~\ref{thm:central-p2-rigidity}.

\subsection{Full rational torsion with tame elliptic descent}
The next three results determine exactly which absolute ramification indices admit an additive potentially good elliptic curve with full rational torsion and a prescribed tame defect. The field is allowed to vary; the existence assertions hold with algebraically closed residue field.

\begin{theorem}[Tame full torsion at level one]\label{thm:tame-level-one}
Let $p\ge5$. There exist a complete discretely valued characteristic-zero field $K$ with perfect residue field of characteristic $p$ and $e_K=e$, and an elliptic curve $E/K$ with full rational $p$-torsion, additive potentially good reduction and semistability defect $m$, if and only if
\begin{equation}\label{eq:tame-first-criterion}
 m\in\{2,3,4,6\},\qquad m\mid p+1,\qquad p-1\mid e,
 \qquad e\ge\frac{p^2-1}{m}.
\end{equation}
The residue field may be $\overline{\F}_p$. Consequently the set of possible indices is
\[
 \left\{e\in(p-1)\Z_{>0}:e\ge\frac{p^2-1}{m_{\max}(p)}\right\},
 \qquad
 m_{\max}(p)=
 \begin{cases}
 2,&p\equiv1\pmod{12},\\
 6,&p\equiv5\pmod{12},\\
 4,&p\equiv7\pmod{12},\\
 6,&p\equiv11\pmod{12}.
 \end{cases}
\]
For $p=3$ and a specified tame defect, $m=2$ occurs exactly at even indices $e\ge4$, and $m=4$ occurs at every positive even index. This last assertion concerns tame inertia only.
\end{theorem}

\begin{theorem}[Every tame defect and every level]\label{thm:tame-full-level}
Fix $p\ge5$, $n\ge1$, and $m\in\{2,3,4,6\}$ with $m\mid p+1$. Let $e$ be a positive multiple of $p-1$, and put $s=v_p(e/(p-1))$. There exist a complete discretely valued field $K$ with $e_K=e$ and perfect residue field of characteristic $p$, and an elliptic curve over $K$ with full rational $p^n$-torsion and exact semistability defect $m$, if and only if
\begin{equation}\label{eq:tame-full-criterion}
 s\ge n-1,\qquad
 e\ge\frac{p^2-1}{m}\,
       p^{\max\{2n-2,\,4n-4-s\}}.
\end{equation}
Every permitted index is attained with residue field $\overline{\F}_p$.
The same assertion holds for $p=3$ and specified tame defect $m=2$ or $4$. The wild characteristic-three defects are treated separately in Theorem~\ref{thm:three-range}.
\end{theorem}

\begin{corollary}[The least level forcing good reduction]\label{cor:tame-forcing-level}
Let $p\ge5$ and $e\in(p-1)\Z_{>0}$. If $m_{\max}(p)e<p^2-1$, put $N=0$. Otherwise set
\[
 s=v_p(e/(p-1)),\qquad
 L=\left\lfloor\log_p\frac{m_{\max}(p)e}{p^2-1}\right\rfloor,
 \qquad
 N=1+\min\left\{s,\left\lfloor\frac L2\right\rfloor,
                         \left\lfloor\frac{L+s}{4}\right\rfloor\right\}.
\]
Among potentially good elliptic curves over all complete fields of absolute index $e$ with perfect residue field of characteristic $p$, the least positive level $n$ at which full rational $p^n$-torsion forces good reduction is $N+1$. For $p=3$, the same formula with $m_{\max}=4$ gives the least level within the class with tame inertia image.
\end{corollary}

Section~\ref{sec:tame-full-torsion} proves necessity from the smallest and largest exact-order formal torsion valuations. For sufficiency it uses the regular full-level deformation surface: a tame action scales both basis-point parameters, a divisorial valuation prescribes their extrema, and an invariant transverse section gives an actual elliptic curve over a complete field of exactly the required index. This last step restores an algebraically closed residue field after an intermediate valuation with residue field $\overline{\F}_p(z)$.

The basic full-division-field bound
\[
 D=(p^2-1)p^{2n-2}
\]
is already the lower bound in \citet[Theorem~1.1, p.~2 of the author version]{Smith2023}, in the setting of a good supersingular elliptic curve at a prime of a number field. Our local proof includes arbitrary perfect residue fields. The assertions above further specify the exact tame defect, the dependence on $s=v_p(e/(p-1))$, and realization at every permissible index; no priority claim is made for the basic bound or the torsion-valuation classification.


\subsection{Wild full torsion in residue characteristic three}
The following theorem treats the wild defects at $p=3$. Together with Theorem~\ref{thm:tame-full-level}, it gives the complete existential answer for additive elliptic curves in this characteristic.

\begin{theorem}\label{thm:three-range}
Let $e=2u$ be a positive even integer, put $s=v_3(u)$, and let $n\ge1$. For each $m\in\{3,6,12\}$, an elliptic curve with additive reduction, exact semistability defect $m$, and full rational $3^n$-torsion exists over some complete discretely valued field of absolute ramification $e$ and perfect residue field if and only if
\begin{equation}\label{eq:wild-criterion}
 s\ge n-1,\qquad
 me\ge 8\,3^{\max\{2n-2,\,4n-4-s\}}.
\end{equation}
The field can be chosen with residue field $\overline{\F}_3$.

Allowing every nontrivial defect, the largest possible level is
\begin{equation}\label{eq:maximum-three}
 N(e)=1+\min\left\{s,\left\lfloor\frac{L+1}{2}\right\rfloor,
                         \left\lfloor\frac{L+s+1}{4}\right\rfloor\right\},
 \qquad L=\lfloor\log_3u\rfloor.
\end{equation}
Among potentially good elliptic curves, full rational $3^{N(e)+1}$-torsion forces good reduction, and that level is least.
\end{theorem}

The theorem describes how much ramification is needed to hide additive reduction from full rational torsion. The smallest absolute index admitting full rational $27$-torsion with additive reduction is $162$. The index $216$ also admits such a curve, of defect twelve; its significance is that the tame construction does not attain that level there. We give the parameters of both examples in Section~\ref{sec:three}.

The proof combines two ingredients. First, rational full level structure gives an exact sum of the valuations of primitive torsion coordinates. A wild automorphism forces the largest such valuation to be divisible by three. Second, one must realize the numerical possibility by a deformation carrying the complete descent action. Regularity of a level ring alone does not provide this realization. We construct an invariant equation on a weighted chart and prove that it has a unique simple branch; that branch is the required discrete valuation ring.

\subsection{Ordinary reduction and boundary examples}
The following results distinguish unramified torsion of an ordinary abelian variety from the broader filtration condition used in the character problem.

\begin{theorem}[Ordinary reduction from unramified $p$-torsion]\label{thm:ordinary-unramified-p}
Let $K$ be complete discretely valued of characteristic zero with perfect residue field of odd characteristic $p$. Let $A/K$ be an abelian variety of any dimension with potentially good ordinary reduction. If $A[p]$ is unramified, then $A$ has good reduction. Neither $\mu_p\subset K$ nor a condition on the degree of a polarization is required.
\end{theorem}

\begin{corollary}[Defect two with full rational $p$-torsion]\label{cor:defect-two-full-p}
For every odd prime $p$, there are a complete field $K$ with residue field $\overline{\F}_p$, absolute ramification index $(p^2-1)/2$ and $\mu_p\subset K$, and an elliptic curve $E/K$ with full rational $p$-torsion and additive potentially good supersingular reduction of exact semistability defect two. Its special inertia involution has trace $-2$.
\end{corollary}

\begin{proposition}[An explicit wild cubic]\label{prop:explicit-cubic-defect-six}
Over $K=\Q_3(\zeta_3)$, the elliptic curve
\[
                         E:\ y^2+y=x^3
\]
has full rational $3$-torsion, additive potentially supersingular reduction and exact semistability defect six, also after unramified base change. Its finite special inertia image is cyclic of order six, and its involution has trace $-2$.
\end{proposition}

\begin{proposition}[Good graded pieces without potential Hodge--Tateness]\label{prop:filtered-non-ht}
For every odd $p$ and $K=\Q_p(\zeta_p)$, there is a height-two Barsotti--Tate group $G/K$ in an exact sequence
\[
0\longrightarrow(\Q_p/\Z_p)_K\longrightarrow G
 \longrightarrow(\Q_p/\Z_p)_K\longrightarrow0
\]
whose rational Tate representation is not potentially Hodge--Tate and hence not potentially semistable. The special character defined by its two good graded pieces is the constant character $2$.
\end{proposition}

Section~\ref{sec:early-reduction} proves these assertions. Corollary~\ref{cor:defect-two-full-p} is the defect-two specialization of Theorem~\ref{thm:tame-level-one}; Proposition~\ref{prop:explicit-cubic-defect-six} is proved by an explicit good model and its faithful descent action. The representation in Proposition~\ref{prop:filtered-non-ht} is the classical logarithmic cyclotomic example displayed by \citet[Section~II.1.2, p.~9 of the 2002 author version]{Berger2004}. The proposition constructs its actual Barsotti--Tate group and identifies what it says about the filtration class; it makes no new claim in Sen theory.

\subsection{Height one and full first-level elliptic torsion}
The first layer already recovers the height-one character, while elliptic full torsion has additional constraints that depend on the particular base field.

\begin{proposition}[The height-one index]\label{prop:height-one-index}
Let $p$ be any prime and let $K$ satisfy the standing complete-field hypotheses, with $\mu_p\subset K$. For every height-one group $M$ in the allowed Barsotti--Tate filtration class, its special character on prime-to-$p$ inertia is determined by $M[p]$. Thus
\[
                         \nuMin(p,e_K,1)=0.
\]
This includes $p=2$.
\end{proposition}

\begin{theorem}[A full-three-torsion curve of defect twelve]\label{thm:explicit-hesse-twelve}
Put
\[
 K=\Q_3(i,\zeta),\qquad i^2=-1,\quad \zeta^2+\zeta+1=0,
 \qquad \pi=\zeta-1,\qquad t=1+i\pi.
\]
Then $e_K=2$ and the residue field is $\F_9$. The elliptic curve
\[
 E:\ y^2=x^3-27t(t^3+8)x+54(t^6-20t^3-8)
\]
has full rational $3$-torsion and additive potentially supersingular reduction of exact semistability defect twelve. The same exact defect holds after unramified base change.
\end{theorem}

\begin{theorem}[A variable Hesse defect]\label{thm:hesse-variable-defect}
Let $K$ be a complete discretely valued field of characteristic zero with algebraically closed residue field of characteristic three, containing $\mu_3$ and of absolute ramification index $e$. Choose a uniformizer $\varpi$ and an integer $r$ such that
\[
                         0<2r<e,\qquad 3\nmid e+r.
\]
Put $t=1+\varpi^r$. The elliptic curve
\[
 E_t:\quad y^2=x^3-27t(t^3+8)x+54(t^6-20t^3-8)
\]
has full rational $3$-torsion and potentially supersingular reduction of exact semistability defect
\[
                             m=\frac{12}{\gcd(4,e+r)}.
\]
There is an explicit good-reduction extension of this degree, and the ramification degree of every good-reduction extension is divisible by $m$.
\end{theorem}

\begin{corollary}[Every first-level defect in characteristic three]\label{cor:three-first-level-table}
Allow the complete field to vary, with arbitrary perfect residue field of characteristic three. An additive potentially good elliptic curve with full rational $3$-torsion and specified defect $m$ exists at exactly the indices in the following table:
\[
\begin{array}{c|c}
 m&\text{absolute ramification indices}\\\hline
 2&e\in2\Z,\ e\ge4\\
 3&e\in2\Z,\ e\ge4\\
 4&e\in2\Z,\ e\ge2\\
 6&e\in2\Z,\ e\ge2\\
 12&e\in2\Z,\ e\ge2.
\end{array}
\]
There are no other nontrivial defects. In particular every positive even index is attained. Over the particular field $\Q_3(\zeta_3)$ there is an example of defect six.
\end{corollary}

\begin{corollary}[The cyclotomic base field]\label{cor:cyclotomic-full-p}
For $p\ge5$, no elliptic curve defined over $\Q_p(\zeta_p)$ has both full rational $p$-torsion and additive potentially good reduction. For $p=3$ such a curve exists.
\end{corollary}

Section~\ref{sec:first-level-refinements} proves these statements. The Hesse family and its preservation of symplectic $3$-torsion structures are classical; see \citet[\S1, Theorem~1.1, p.~2]{AnemaTopTuijp2018}, over perfect fields of characteristic different from three. Here the generic field has characteristic zero. The proofs supply the points and the exact local good-reduction fields for the displayed curves. The first-level table is a direct specialization of the complete tame and wild existence criteria. The fixed cyclotomic field requires an additional residue-field argument at $p=5$: an existential example at the same absolute index over a different field would not settle that corollary. The classifications of \citet{OzekiYoshida2026} concern curves defined over $\Q_p$; they are neighboring context and are not used as a proof of this assertion for curves defined over $\Q_p(\zeta_p)$.


\subsection{A cyclic point and a cyclotomic obstruction}
The next results concern inputs to the comparison problem. They are stated here in their independent field-theoretic and formal-group forms.

\begin{theorem}\label{thm:cyclic-bound}
Let $p\ge5$, let $R$ be a complete discrete valuation ring with algebraically closed residue field of characteristic $p$, and put $a=v_R(p)$ and $v=v_p(a)$. If a one-dimensional formal group of special height two has an $R$-rational point of exact order $p^n$, with $n\ge2$, then
\begin{equation}\label{eq:cyclic-bound}
 a>(p^2-1)p^{2n-3-v}.
\end{equation}
A negative exponent on the right is interpreted in $\Q_{>0}$.
\end{theorem}

This inequality records the number of times a rational inverse chain can follow the larger Newton slope. Each such step consumes one factor of $p$ in $a$. The valuations themselves are the canonical-subgroup valuations described by \citet[Theorem~4.4 and Corollaries~4.5--4.6, pp.~10--11 of the author version]{Smith2023}. Section~\ref{sec:cyclic} gives the argument for arbitrary height-two formal groups and derives the strict inequality, including its equality boundary.

\begin{theorem}\label{thm:cyclotomic-exact}
Let $p\ge3$, and suppose that $K$ has algebraically closed residue field, $e_K=p(p-1)$, $\mu_p\subset K$, and $\mu_{p^2}\not\subset K$. Choose a Cohen coefficient field $K_0\subset K$, put $K_1=K_0(\mu_p)$, and write $d=v_K(\mathfrak D_{K/K_1})$. The last upper ramification break of $K(\mu_{p^3})/K$ is exactly
\begin{equation}\label{eq:exact-second-break}
 \delta_2=(p-1)(2p+1)-d\ge p(p-1).
\end{equation}
\end{theorem}

This identity follows by taking the normal closure of the degree-$p$ extension $K/K_1$ and applying Herbrand's theorem. It supplies the determinant obstruction in the formal comparison proofs. A separate argument with differents gives the following bound, without computing that normal closure.

\begin{theorem}\label{thm:cyclotomic-bound}
Let $p\ge3$. Suppose that $K$ has algebraically closed residue field, $e_K=p(p-1)$, $\mu_p\subset K$, and $\mu_{p^2}\not\subset K$. Then $C=K(\mu_{p^3})$ is cyclic of degree $p^2$ over $K$. If $\delta_1<\delta_2$ are its upper ramification breaks, in the classical convention without a conductor shift, then
\begin{equation}\label{eq:cyclotomic-bound}
 \delta_1+p\delta_2\ge p(p^2-p-1).
\end{equation}
Consequently $\delta_1\le p-1$ implies
\begin{equation}\label{eq:second-break}
 \delta_2\ge p^2-p-2+\frac1p.
\end{equation}
\end{theorem}

The inequality follows by comparing differents in two towers over a Cohen field. Both arguments apply to infinite perfect residue fields. A proposed common torsion field has to contain the cyclotomic extension through its determinant; an upper bound on the torsion field's last break can therefore contradict~\eqref{eq:exact-second-break}.

\subsection{Sharp comparison thresholds}
The following special cases illustrate the character thresholds at small ramification indices. Throughout the following statement, the class defining $\nuMin$ includes the potentially good graded filtrations of Section~1.

\begin{theorem}\label{thm:thresholds-intro}
Let $p\ge5$ and assume $\mu_p\subset K$.
\begin{enumerate}
\item For $n\ge2$ and a positive multiple $e$ of $p-1$ prime to $p$, a potentially supersingular height-two group and a group filtered by two potentially good height-one groups of total dimension one can have isomorphic $p^n$-torsion and different special characters over some field of index $e$ if and only if $e>p^{2n-3}$.
\item Among indices $e$ divisible by $p(p-1)$ with $v_p(e)=1$, such a comparison at level $p^3$ exists if and only if $e>p^2$.
\item Over all complete fields with perfect residue field of absolute index $2p(p-1)$, the least uniform character index is
\[
                         \nuMin(p,2p(p-1),2)=3.
\]
\item At the first wild index, $\nuMin(p,p(p-1),2)=2$ for every $p\ge5$.
\item If $p\nmid e$ and $p^3<e<p^4$, then $\nuMin(p,e,2)=3$.
\end{enumerate}
\end{theorem}

Part~(1) is Corollary~\ref{cor:sharp-mixed-prime}, part~(2) is Corollary~\ref{cor:sharp-mixed-wild}, part~(3) is Theorem~\ref{thm:second-wild-exact}, part~(4) is Theorem~\ref{thm:first-wild-general}, and part~(5) is Corollary~\ref{cor:prime-index-intervals}. The existence assertions allow the base to vary; part~(3) is a uniform recovery assertion. Their proofs explain why a cyclic point can conceal a special character: quotienting by the point produces a finite extension of rank-one characters, and Kummer theory lifts exactly that extension to a filtered Barsotti--Tate group. For the upper bound, the common determinant forces the cotangent weights of the graph quotient to agree modulo $p-1$. This strengthens the coefficient estimate in Bondarko's integral kernel argument.

\subsection{The first unresolved wild coefficient}
The next theorem narrows the first coefficient beyond the exact wild bands to one necessary configuration of actual torsion representations.

\begin{theorem}[The first coefficient beyond the exact bands]
\label{thm:first-boundary-normal-form}
Let $p\ge5$ and $r\ge1$, and put
\[
 k_0=\frac{p^2+1}{2},\qquad b_0=\frac{p-1}{2},\qquad
 e=k_0(p-1)p^r,\qquad n=r+3.
\]
Suppose that two groups in the standing height-two filtration class
over a field of absolute ramification index $e$ have isomorphic
actual $p^n$-torsion and different special inertia characters.
Then both are whole potentially supersingular groups with two-slope
good models. In the common tame normalization their Hasse parameters
and special tangent scalars satisfy
\[
 z_1=z_2=b_0p^n,\qquad \frac{\alpha_2}{\alpha_1}=-1.
\]
Both residual $\mathbf F_p$-representations over $K$ are split.
Their canonical and complementary lines have distinct inertia
characters, and the given isomorphism sends each canonical line
to the complementary line of the other representation.
\end{theorem}

This is a necessary form for an assumed unequal-character comparison.
It does not assert that this form occurs or determine the character
index at $e=k_0(p-1)p^r$.

Section~\ref{sec:first-boundary-normal-form} proves the result using the general two-chain congruence and the actual quotient markings.

\begin{corollary}[Point fields at the first boundary]
\label{cor:boundary-point-field-degrees}
Under the assumptions of
Theorem~\ref{thm:first-boundary-normal-form}, pass to the completed
maximal unramified base and the common tame good field $L$.
For every exact $p^j$-point $Q$ of either group whose $p^{j-1}$-multiple is a
nonzero canonical point, write $[L(Q):L]=p^d$. Then
\[
 d\ge\max\{0,2j-1-r\}\quad(2\le j<n),\qquad
 d\ge r+5\quad(j=n).
\]
Each relative step at levels $r+1,r+2,r+3$ has degree $p^2$.
In particular every canonical $p^{n-1}$-point field has degree
at least $p^{r+3}$. If $r=1$ and $n=4$, the degrees at levels
$p^2,p^3,p^4$ are respectively
\[
 p^2,\qquad p^4,\qquad p^6.
\]
\end{corollary}

Proposition~\ref{prop:proper-canonical-point-fields} proves a
more general proper-level bound for an exchanged central
comparison. The point-field restrictions above do not settle
whether the boundary comparison occurs.

\begin{corollary}[Nonsplit Cartan exclusion at the first boundary]
\label{cor:boundary-nonsplit-cartan-image}
Let $p\ge5$, $r\ge1$, $n=r+3$, and
\[
                    e=\frac{p^2+1}{2}(p-1)p^r.
\]
Suppose that two groups in the standing height-two filtration class
over a field of absolute index $e$ have isomorphic actual
$p^n$-torsion and different special inertia characters.
Pass to the completed maximal unramified base and its common
tame good extension $L$ of degree $p^2-1$.
For each group, the finite image of
$\operatorname{Gal}(\overline L/L)$ on its $p^n$-torsion is not
contained in any conjugate of
\[
                    (1+p\mathcal O)/(1+p^n\mathcal O),
\]
where $\mathcal O$ is the unramified quadratic
$\mathbf Z_p$-algebra and $\mathcal O/p^n\mathcal O$ acts on
$(\mathbf Z/p^n\mathbf Z)^2$ through its regular representation.
\end{corollary}

This is a necessary restriction on an assumed actual comparison.
It does not exclude the remaining non-Cartan images or determine
the boundary character index. Section~\ref{sec:boundary-nonsplit-image-proof}
proves the restriction by constructing an actual invariant-ring map
under the additional image-containment assumption.

\begin{corollary}[The irreducible quadratic self-twist case]
\label{cor:boundary-nonsplit-self-twist}
Under the assumptions of
Corollary~\ref{cor:boundary-nonsplit-cartan-image}, the finite
representation of either group admits no tame quadratic self-twist
with an invertible intertwiner whose reduction has irreducible
quadratic characteristic polynomial over $\mathbf F_p$.
\end{corollary}


The proof in Section~\ref{sec:boundary-nonsplit-image-proof} covers
arbitrary perfect residue fields. The split-intertwiner case is not
part of this corollary.


\subsection{Relation to earlier work and scope}
Drinfeld's regularity and level-parameter theorems are the foundation of the full-torsion constructions. For the tame criterion we use the original \citet[Propositions~4.2--4.3, pp.~570--573]{Drinfeld1974}; the primitive factorization used in the wild characteristic-three proof is stated in \citet[Theorem~2.1.2, p.~437; Proposition~2.2.3, pp.~439--441]{CarterStrauch2021}. The equivalence between elliptic and Barsotti--Tate deformations is the Serre--Tate theorem, in the form of \citet[Theorem~1.2.1]{Katz1981}. The divisibility of a discriminant for full rational $3$-torsion follows from the Hesse-pencil theorem of \citet[Theorem~1.1, p.~2; Proposition~5.1, pp.~4--5]{AnemaTopTuijp2018}.

Proposition~\ref{prop:divisorial-full-cartan} computes the full principal nonsplit Cartan quotient of the universal height-two deformation at every level $p^n$, for odd $p$ and $n\ge2$, with monomial weights $(p^2+p+1,p^2)$. Its ramification index and purely inseparable residue degree are both $p^{n-1}$, and its graded fixed field is $k(X^{p^{n-1}},Y^{p^{n-1}})$. Proposition~\ref{prop:divisorial-last-cartan} gives the separate last-layer invariants, while Proposition~\ref{prop:divisorial-scalar-level-four} computes the distinct full scalar subgroup at level four. At $p=5,n=4$, a further tame quotient of ramification index $12$ with unchanged residue field has absolute index exactly $6500$ (Corollary~\ref{cor:divisorial-full-cartan-6500}). This statement concerns the specified divisorial valuation with imperfect residue field; it does not identify its curve specializations with arbitrary boundary comparisons.

The local rings of this Cartan construction have a stronger property. Put $\pi=x^{1-p}y^p$, $t=x^{p^2}/y^{p^2+p+1}$ and $B_c=(R_n[\pi,t,t^{-1}])_{(\pi,t-[c])}$ for $c\in k^*$. Proposition~\ref{prop:cartan-torus-regularity} proves that $B_c^C$ and its reduced vertical divisor are regular, with divisor ring $k[t^{p^{n-1}},t^{-p^{n-1}}]_{(t^{p^{n-1}}-c^{p^{n-1}})}$. Its local equation $\rho$ satisfies $p=u\rho^{(p^4-1)p^{n-1}}$ for a unit $u$. Thus every local mixed-characteristic specialization of these charts has absolute index divisible by $(p^4-1)p^{n-1}$. At $p=5,n=4$, this is $78000$; an additional field extension of ramification index at most $12$ gives base index at least $6500$. The proof uses the principal-augmentation criterion of \citet[Theorem~2 and Proposition~5, pp.~64--69]{KiralyLutkebohmert2013}. Corollary~\ref{cor:boundary-nonsplit-cartan-image} constructs the invariant-ring map under an actual nonsplit-Cartan image hypothesis and contradicts this divisibility for every first-boundary index. Lemma~\ref{lem:scalar-compatible-frame} supplies the coefficient-compatible universal frame over arbitrary algebraically closed residue extensions; its proof derives the required scalar extension degree by degree from \citet[Propositions~4.1--4.3, pp.~570--573]{Drinfeld1974}. If the torus residue is transcendental over $\overline{\mathbf F}_p$, the map from a closed chart need not be local, but the exact unit identity still gives the contradiction. The image hypothesis is not forced for every surviving comparison.

Proposition~\ref{prop:cartan-orbit-data} shows that the maximal primitive-point orbit sizes do not bound the full image by a Cartan: determinant-constrained congruence groups have the same orbits and are noncommutative. Lemma~\ref{lem:quadratic-self-twist-cartan} gives the stronger Cartan containment supplied by an actual quadratic self-twist, with a distinct split case. Corollary~\ref{cor:boundary-nonsplit-self-twist} excludes its irreducible-intertwiner case for an actual first-boundary comparison.

The local divisor class is an additional constraint on specialization. Lemma~\ref{lem:cartier-index-contact} shows that a single vertical multiplicity $N$ and Cartier index $i$ force $N/\gcd(N,e_K)\mid i$ along a local mixed-characteristic curve. Example~\ref{ex:noncartier-quotient-specialization} applies the known cyclic quotient of \citet[Example~5.10]{LiedtkeMatsumoto2018} to show why normality and generic \'etaleness alone do not preserve a divisorial index bound. Proposition~\ref{prop:boundary-specialization-contact} separately forces contact divisible by $5$ for a specified upstairs chart carrying an actual boundary comparison. Example~\ref{ex:fierce-transverse-specialization} shows how a downstairs transverse curve may acquire such contact.

For finite flat group schemes, \citet[Theorems~A--B, p.~661]{Bondarko2006} proves an almost-fullness theorem for the generic-fibre functor and finite torsion criteria for good reduction. These are uniform extension results; they do not on their own give an exact least character index in height two. \citet[Theorem~1.1, p.~1 of the preprint]{Hattori2008} determines tame characters of ramification graded pieces of a finite flat group scheme by its Abbes--Saito filtration. That theorem concerns generic finite representations and their ramification; the character here is the finite descent character of good special graded pieces. Recovering the latter from a given truncation requires an additional argument.

The full-division-field lower bound and higher-torsion valuation classification of \citet[Theorem~1.1, p.~2; Theorem~4.4 and Corollaries~4.5--4.6, pp.~10--11 of the author version]{Smith2023}, and the reduction restrictions for a rational point of order $p^n$ in \citet[Theorems~1.1--1.3, pp.~1--2]{Melistas2022}, address closely related local questions. The numerical lower bound $(p^2-1)p^{2n-2}$ is therefore credited to Smith. The mixed criterion in Theorem~\ref{thm:mixed-all-indices} additionally constructs actual groups of two different reduction types at every admissible absolute index, using a full integral descent action. The exact fixed-defect, every-index assertions~\eqref{eq:tame-full-criterion} and~\eqref{eq:wild-criterion} additionally require an equivariant elliptic realization and the refined valuation of the base index. More recently, \citet{OzekiYoshida2026} determine first-level point fields and ramification breaks for elliptic curves over $\Q_p$; its Section~3.4 also treats bounds for higher torsion. Those results concern fixed elliptic curves over $p$-adic fields, whereas our existence theorem varies the base and permits an algebraically closed residue field.

At the level of entire Galois images, \citet[Theorem~1.3, p.~2]{BisattEtAl2026} classify the $p$-adic images of elliptic curves over $\Q_p$, for $p>7$, whose mod-$p$ image lies in a nonsplit Cartan normalizer; in the higher semistability-defect case their conclusion has an explicit additional bound relating $p$ to the valuation of the $j$-invariant. Their fixed base field and residual-image hypothesis differ from the torsion-character recovery problem over the ramified bases considered here.

The Newton inputs to Theorem~\ref{thm:central-p2-rigidity} agree, for elliptic formal groups, with the canonical-subgroup valuations of \citet[Theorem~4.4 and Corollaries~4.5--4.6]{Smith2023}; Lemma~\ref{rc-lem:newton} derives the required formulas for height-two formal models. The elementary lower-ramification estimate in Lemma~\ref{cr-lem:valuation} belongs to the classical framework of \citet[Chapter~IV, \S1]{Serre1979}. The comparison here retains the cyclic extension $L(Q)/L(q(Q))$, its exact break, and its identification under a specified $p^2$-torsion isomorphism. In contrast, the first-level classifications of \citet[\S\S2.1, 2.3]{OzekiYoshida2026} concern point fields of elliptic curves over $\Q_p$, while their higher-level Proposition~3.25 bounds the last upper break of the full division field. Those stated results do not provide the transported relative-field comparison proved here. Hattori's theorem concerns the generic ramification-graded characters described above. These are comparisons of the precise statements inspected in a bounded primary-literature search, not a claim that the relative-field method or this criterion has no earlier occurrence.

All conversions between upper and lower numbering, and all different calculations, use the classical conventions of \citet[Chapters~III--IV]{Serre1979}. The full-class prime-to-$p$ table is complete. Theorems~\ref{thm:mixed-all-indices} and~\ref{thm:wild-bands-exact} give the complete mixed criterion and infinite exact wild bands. The final section identifies the remaining full-class questions outside those bands and the distinct supersingular-only and fixed-field problems. No claim of priority is made.

Recent work also studies ramification through integral lattices and actual point fields. \citet[Theorem~11.13, p.~43]{DembeleEtAl2026} bound the residual point-field breaks of nonordinary rank-two crystalline representations over ramified quadratic $2$-adic fields. The global theorem of \citet[Theorem~1.1, p.~1]{OpenAI2026} concerns modularity of odd regular two-dimensional $2$-adic representations. Their base fields, torsion levels, and conclusions differ from the odd-prime character thresholds and exact elliptic full-torsion criteria considered here. We use neither result as a proof input.

For the supersingular comparison we also prove the following elementary obstruction in Proposition~\ref{fc-thm:valuation}: if two height-two potentially supersingular groups over \(K\) have isomorphic \(p^n\)-torsion and unequal special inertia traces, then
\[
e_K(p^n-1)\geq(p-1)p^{2n-2}.
\]
No condition on \(v_p(e_K)\) is needed for this inequality. At \(n=3,e_K=p(p-1)\), its left side falls short of its right side by \(p(p-1)\), proving the supersingular upper step of Theorem~\ref{thm:first-wild-general} for every \(p\ge5\).

\section{The complete tame elliptic criterion}\label{sec:tame-full-torsion}
This section proves Theorems~\ref{thm:tame-level-one}--\ref{thm:tame-full-level} and Corollary~\ref{cor:tame-forcing-level}, including actual elliptic realizations at every allowed absolute ramification index.

\subsection{Inertia and the good formal group}
We first identify the good special fibre and the restriction imposed by a tame tangent character.

\begin{lemma}\label{tf-lem:inertia}
Let $p$ be odd, and let $E/K$ be potentially good with nontrivial tame semistability defect $m$ and full rational $p$-torsion. Its good special fibre is supersingular. If $p\ge5$, every nontrivial potentially good defect is tame and belongs to $\{2,3,4,6\}$; for $p=3$ a nontrivial tame defect belongs to $\{2,4\}$.

After completed maximal unramified base change, let $L/K$ be the minimal good extension. There is an integral formal parameter $t$ and a residue class $r\in(\Z/m\Z)^\times$ such that every nonzero rational descended $p$-power torsion point satisfies
\begin{equation}\label{tf-eq:congruence}
 v_L(t(P))\equiv r\pmod m.
\end{equation}
\end{lemma}

\begin{proof}
Completed maximal unramified base change preserves the absolute index and the inertia image. We may therefore assume the residue field $k$ is algebraically closed. We use the good-reduction criterion of \citet[Theorem~1 and Corollary~4, pp.~493--494]{SerreTate1968}. The finite inertia image acts faithfully on the good special elliptic curve $E_0$: specialization identifies prime-to-$p$ Tate modules, and an automorphism acting trivially on every $\ell$-power torsion point is the identity. Indeed a nonzero elliptic endomorphism is finite and cannot kill an infinite Tate module. The good-reduction criterion identifies the kernel with the good-reduction inertia kernel.

For $p\ge5$, a short Weierstrass equation $y^2=x^3+Ax+B$ identifies every origin-preserving automorphism with
$(x,y)\mapsto(u^2x,u^3y)$, where $u^4A=A$ and $u^6B=B$. Its group is cyclic of order $2$, $4$ or $6$, and all subgroups are tame. In characteristic three, complete the square and write
$y^2=x^3+a_2x^2+a_4x+a_6$. An automorphism has
$x\mapsto u^2x+r$, $y\mapsto u^3y$. If $a_2\ne0$, comparison of the quadratic coefficients gives $u^2=1$. If $a_2=0$, nonsingularity gives $a_4\ne0$, and comparison of linear coefficients gives $u^4=1$. A finite automorphism of order prime to $p$ with tangent multiplier one is the identity: its first nonidentity term $cT^j$ would become $hcT^j$ on taking an $h$th iterate, which cannot vanish when $p\nmid h$. Thus a tame subgroup injects into the tangent multipliers. This proves the asserted lists and the faithfulness of the tame tangent character.

Suppose $E_0$ were ordinary. The finite flat good-model group scheme $E[p]$ shows that reduction of geometric generic $p$-torsion surjects onto $E_0[p](k)$: each special point lies in the closure of generic points. All generic points are rational over $L$, so the descent automorphism acts trivially on these special points. Its action on the height-one \emph{\'etale} $p$-divisible group is a finite-order element of $\Z_p^*$ trivial modulo $p$. The group $1+p\Z_p$ is torsion-free for odd $p$, so this action is trivial. The endomorphism $\gamma-1$ kills the entire \emph{\'etale} Tate module and hence is zero. This contradicts nontrivial inertia. Therefore $E_0$ is supersingular and all geometric $p$-power torsion lies in the good formal group.

The tame extension $L/K$ is cyclic of degree $m$. Choose a uniformizer $\pi$ and generator $\sigma$ with $\sigma\pi=\xi\pi$, where $\xi\in K$ is a primitive $m$th root. This follows from the tame Kummer description over a strictly henselian DVR; units have $m$th roots by Hensel's lemma. Write the faithful tangent multiplier as $\lambda=\xi^r$, with $\gcd(r,m)=1$, using the action convention on functions. For any formal parameter $t_0$, the average
\[
 t=\frac1m\sum_{i=0}^{m-1}\lambda^{-i}\sigma^i(t_0)
\]
is again a parameter and satisfies $\sigma(t)=\lambda t$. Its special linear coefficient is a unit. Evaluation on a descended torsion section gives
$\sigma(t(P))=\lambda t(P)$. If $q=v_L(t(P))$, the residue of $t(P)/\pi^q$ is fixed by inertia, so $\xi^q=\lambda$ in $k$. This gives~\eqref{tf-eq:congruence}.
\end{proof}

\subsection{The smallest and largest division valuations}
The next calculation retains both extrema of exact-order torsion and includes every equality case.

\begin{lemma}\label{tf-lem:extrema}
Let $F$ be a one-dimensional formal group of special height two over a complete mixed-characteristic DVR, with full rational $p^n$-torsion. The residue field may be imperfect. Normalize the valuation to $\Z$, put $E=v(p)$, and let $a\ge b>0$ be the largest and smallest nonzero $p$-torsion valuations. Then
\begin{equation}\label{tf-eq:mass}
 M:=\frac{E}{p-1}=a+pb.
\end{equation}
If $b_n,A_n$ are the minimum and maximum exact $p^n$-torsion valuations, then
\begin{equation}\label{tf-eq:extrema}
 b_n=\frac b{p^{2n-2}},\qquad
 A_n=\frac{M-p^{c+1}b}{p^{2n-2-c}},\qquad
 c=\min\left\{n-1,\left\lfloor\log_p\frac{M}{(p+1)b}\right\rfloor\right\}.
\end{equation}
\end{lemma}

\begin{proof}
The positive-valuation filtration on $F[p]$ is a filtration of a two-dimensional $\F_p$-space. Indeed $F(X,Y)=X+Y$ plus terms divisible by $XY$, so addition has valuation at least the minimum, with equality at unequal valuations; multiplication by a $p$-adic unit preserves valuation. A strictly nested filtration has at most one nonzero proper subspace. Consequently, if $a>b$, exactly $p-1$ nonzero points have value $a$ and $p^2-p$ have value $b$. If $a=b$, all $p^2-1$ points have that value.

Weierstrass preparation of $[p](T)$ has distinguished degree $p^2$. Comparing its derivative at zero with the product of its nonzero roots gives~\eqref{tf-eq:mass}. The Newton polygon has vertices
\[
 (1,E),\qquad(p,H),\qquad(p^2,0),\qquad H=p(p-1)b,
\]
with the middle point collinear in the one-slope case. This follows from the root valuations and multiplicities and requires no special coordinate normal form.

Let $P\ne0$ be a torsion point of valuation $q\le a$. The lower polygon for $[p](T)-t(P)$ governing positive root valuations is the hull of
\[
 (0,q),\qquad(p,H),\qquad(p^2,0).
\]
The point $(1,E)$ lies above this hull because $E>q$; all old coefficient points lie above the old polygon, which lies above the proposed new hull. Coefficients beyond degree $p^2$ have nonnegative valuations. Thus the division values are
\begin{equation}\label{tf-eq:division}
\begin{cases}
 q/p^2\quad\text{for all }p^2\text{ preimages},&q\le p^2b,\\
 (q-H)/p\quad\text{for }p\text{ preimages, and }b
 \text{ for }p^2-p\text{ preimages},&q>p^2b.
\end{cases}
\end{equation}
At equality both descriptions coincide. This formula is used only for nonzero torsion with $q\le a$. Multiplication by $p$ strictly increases a positive nonzero formal valuation, so every higher-order torsion valuation is smaller than $a$ and the required range is preserved.

The minimum branch starting at $b$ divides by $p^2$ at every step, giving the first formula in~\eqref{tf-eq:extrema}. The largest division value is the continuous increasing function
\[
 f(q)=\begin{cases}q/p^2,&q\le p^2b,\\
 (q-p(p-1)b)/p,&q\ge p^2b.
 \end{cases}
\]
Hence $A_n=f^{n-1}(a)$. On the upper branch, after $j-1$ divisions the value is $M/p^{j-1}-pb$. The integer $c$ in~\eqref{tf-eq:extrema} counts the upper steps until the value enters $[b,p^2b]$, or until the chain ends. The remaining steps divide by $p^2$, and give the second displayed formula. If $M/((p+1)b)$ is an exact power of $p$, the two adjacent branch choices agree at their common endpoint. If $a=b$, then $c=0$ and both extrema are $b/p^{2n-2}$.
\end{proof}

\begin{proposition}\label{tf-prop:necessity}
Every curve in Theorem~\ref{thm:tame-full-level} satisfies~\eqref{eq:tame-full-criterion}. Full rational $p$-torsion with nontrivial tame defect also forces $m\mid p+1$ and $p-1\mid e$.
\end{proposition}

\begin{proof}
The Weil pairing gives $\mu_{p^n}\subset K$, and
\[
 v_K(\zeta_{p^n}-1)=\frac{e}{p^{n-1}(p-1)}\in\Z.
\]
In particular $p-1\mid e$. On the minimal good field, $E=me$ and
$M=me/(p-1)$. Lemma~\ref{tf-lem:inertia} gives $a\equiv b\equiv r\pmod m$ with $r$ invertible. Since $m\mid M=a+pb$, we obtain $(p+1)r\equiv0\pmod m$, hence $m\mid p+1$.

Put $s=v_p(e/(p-1))=v_p(M)$. Integrality of $b_n$ in Lemma~\ref{tf-lem:extrema} gives
\[
 p^{2n-2}\mid b,\qquad b\ge p^{2n-2}.
\]
Integrality of $A_n$ gives $p^{2n-2-c}\mid M-p^{c+1}b$. The second term is divisible by $p^{2n+c-1}$, so
\[
 s\ge2n-2-c\ge n-1.
\]
The definition of $c$, including its capped case, gives $M\ge(p+1)bp^c$. Therefore
\[
 M\ge(p+1)p^{2n-2+\max(0,2n-2-s)}
       =(p+1)p^{\max(2n-2,4n-4-s)}.
\]
Multiplication by $(p-1)/m$ proves~\eqref{eq:tame-full-criterion}.
\end{proof}

\subsection{Actual elliptic deformations with an exact tame action}
We next supply the elliptic family and its full-level surface before prescribing a valuation on that surface.

\begin{lemma}\label{tf-lem:surface}
Let $k=\overline{\F}_p$, and let $m$ be one of the allowed tame defects. There is a complete regular local $W(k)$-algebra $R_n$ of dimension two, flat over $W(k)$, carrying an actual good elliptic scheme with a full Drinfeld $p^n$-level structure. It has an integral formal parameter $t$ such that the coordinates $x=t(P)$ and $y=t(Q)$ of a labelled basis form a regular system of parameters. A faithful semilinear $C_m$-action fixes the level labels and $W(k)$ and satisfies
\begin{equation}\label{tf-eq:scalar-action}
 \gamma x=\zeta x,\qquad\gamma y=\zeta y,
\end{equation}
where $\zeta\in W(k)$ has exact order $m$. Its special elliptic action has exact order $m$.
\end{lemma}

\begin{proof}
First choose a supersingular elliptic $E_0/k$ with the required automorphism. For $m=3$ or $6$, the congruence $m\mid p+1$ gives $p\equiv2\pmod3$, and we take $y^2=x^3+1$. For $m=4$, take $y^2=x^3-x$, with $p\equiv3\pmod4$, including $p=3$. The usual scaling automorphisms have the required faithful tangent multipliers. These curves are supersingular: the coefficient of $x^{p-1}$ in the $(p-1)/2$ power of the cubic vanishes, respectively because $3\nmid p-1$ and because $(p-1)/2$ is odd.

For $m=2$, use $[-1]$ on any supersingular elliptic curve. Such a curve exists for every odd $p$ by an elementary Legendre calculation. Put $h=(p-1)/2$. The coefficient just used for $x(x-1)(x-\lambda)$ is $(-1)^h$ times
\[
 H(\lambda)=\sum_{i=0}^h\binom hi^2\lambda^i.
\]
This has positive degree, $H(0)=1$, and
$H(1)=\binom{p-1}{h}\ne0$ in $k$. It has a root distinct from $0,1$, hence a smooth Legendre curve with zero Hasse coefficient. To recall why the coefficient detects supersingularity, write
$y^{-1}=y^{-p}f(x)^{(p-1)/2}$ and apply Cartier to $dx/y$. The degree $3(p-1)/2<2p-1$ shows that only the coefficient of $x^{p-1}$ contributes. Cartier, with its Frobenius twist, is the differential of Verschiebung. Its vanishing makes the degree-$p$ Verschiebung inseparable. Since multiplication by $p$ is Verschiebung composed with the purely inseparable degree-$p$ Frobenius, the $p$-torsion is then connected of rank $p^2$. This is the height-two case.

Let $F_0$ be the formal group of $E_0$. The original deformation theorem of \citet[Propositions~4.2--4.3, pp.~570--573]{Drinfeld1974} gives
$R_0=W(k)[[u]]$ and a finite flat full-level ring $R_n$ which is regular of dimension two, with basis-point coordinates as parameters. The hypotheses allow the integers of any nonarchimedean local field, in particular $\Z_p$ for every odd $p$ \citep[\S1C, pp.~564--565]{Drinfeld1974}. In characteristic zero the level is full: the defining divisor identity gives all $p$-torsion roots, and successive $p$-division gives $p^{2n}$ distinct roots, since multiplication by $p$ is then \emph{\'etale}.

We explain the actual elliptic realization over the complete ring. On each Artinian quotient of $R_0$, the kernels of the universal $[p^j]$ are finite locally free of rank $p^{2j}$ by Weierstrass preparation and form a $p$-divisible group. Apply the Serre--Tate equivalence to this group and $E_0$ \citep[Theorem~1.2.1, p.~143, proof pp.~144--146]{Katz1981}. It gives compatible elliptic deformations. These algebraize directly: the section modules of $\Ocal(j[0])$, $j\ge1$, have their genus-one ranks and compatible reduction maps, since their first cohomology vanishes. Lift bases of the pole-order filtrations to choose compatible Weierstrass coordinates and coefficients. Their inverse limits define a Weierstrass equation over $R_0$ with unit discriminant. The degree-six pole-order relation uses only finitely many coefficients, so this elliptic scheme induces all the Artinian deformations. Its $p$-divisible group is the prescribed one by the compatible finite-level identifications. Pull it back to $R_n$.

The special automorphism $\gamma_0$ acts on the framed deformation functor by changing the residual framing, leaving the level labels fixed. Representability supplies actions on $R_0$ and $R_n$ fixing $W(k)$. Full faithfulness in Serre--Tate supplies the compatible elliptic isomorphisms and their cyclic cocycle identities. These extend to the algebraized Weierstrass model: the finitely many coefficients in the usual Weierstrass coordinate transformation have compatible inverse limits. Thus the action is an actual semilinear action, not only an action on formal isomorphism classes.

Average any parameter $t_0$ by
\[
 t=\frac1m\sum_{i=0}^{m-1}\zeta^{-i}\gamma^i(t_0),
\]
where $\zeta$ is the lift of the special tangent multiplier. Its special derivative is a unit, and $\gamma t=\zeta t$. Changing an integral parameter by a unit linear term preserves the regular system of basis-point parameters. As the labels are fixed, this gives~\eqref{tf-eq:scalar-action}. The special action remains $\gamma_0$ of exact order $m$.
\end{proof}

\subsection{A divisorial valuation and an invariant transverse section}
The next lemma turns the regular surface into a complete field with perfect residue field while retaining its exact index and torsion marking.

\begin{lemma}\label{tf-lem:realization}
Fix $n\ge1$ and an allowed tame defect $m$. For every integer $A\ge1$ with $A\equiv1\pmod m$, there is an elliptic curve with full rational $p^n$-torsion and exact defect $m$ over a complete field with residue $\overline{\F}_p$ and absolute index $E/m$, where
\begin{equation}\label{tf-eq:index-A}
 \frac{E}{p-1}=p^{2n-2-c}(A+p^{2c+1}).
\end{equation}
Here $0\le c<n-1$ when $p^{2c}\le A\le p^{2c+2}$, and $c=n-1$ when $A\ge p^{2n-2}$. Adjacent expressions agree at their common endpoints.
\end{lemma}

\begin{proof}
Use the surface of Lemma~\ref{tf-lem:surface}. Blow up its closed point and then the origins of the successive charts with coordinates $y,x/y^j$, for $1\le j<A$. The final exceptional divisor defines the valuation
\[
 v(y)=1,\qquad v(x)=A,
 \qquad k(v)=k(z),\quad z=x/y^A.
\]
All these charts are regular. For example the first chart is defined in the regular three-dimensional ring $R_n[t]$ by $x-yt$; at a closed point of the exceptional line this relation has a nonzero linear term, leaving regular parameters $y,t-t_0$. The same argument applies successively. The centres are invariant, and on the final chart
\[
 \gamma y=\zeta y,\qquad\gamma z=\zeta^{1-A}z=z.
\]

The minimum and maximum primitive $p^n$-torsion valuations are exactly $1,A$. Here is the cancellation check. A primitive vector $(a,b)$ has coordinate
\[
 f\equiv\bar a x+\bar b y\pmod{(x,y)^2},
 \qquad(\bar a,\bar b)\ne(0,0).
\]
If $A>1$ and $\bar b\ne0$, its value is one. If $A=1$, its leading coefficient is $\bar a z+\bar b\ne0$ in $k(z)$, again giving value one. If $A>1$ and $\bar b=0$, then $\bar a\ne0$. The level-one product factorization of $[p](T)/T$ implies $p\in(x)$: one nonzero $p$-torsion factor is $[p^{n-1}](x)$, itself divisible by $x$. Hence the complete regular quotient $R_n/(x)$ is the characteristic-$p$ DVR $k[[y]]$. Lift the image of $f$ there using Teichm\"uller coefficients to write
\[
 f=ux+g(y),\qquad \bar u=\bar a\ne0,
 \qquad \ord_y g\ge2
\]
with $g=0$ allowed. If $\ord_y g\ne A$, the smaller valuation determines $v(f)\le A$. At equality the leading coefficient is $\bar a z+c\ne0$ in $k(z)$, so $v(f)=A$. Every primitive coordinate has value at least one, and $Q,P$ attain $1,A$. This proves both extrema.

Complete the local DVR at the generic point of the exceptional divisor. It has mixed characteristic because $R_n$ is a domain flat over $W(k)$. The good special formal group remains height two over $k(z)$, and all full-level points remain rational. Lemma~\ref{tf-lem:extrema} applies, despite this intermediate imperfect residue field. Its first equation gives $b=p^{2n-2}$. Its second equation is a continuous strictly increasing piecewise affine function of $M=E/(p-1)$: on
\[
 (p+1)bp^c\le M\le(p+1)bp^{c+1}
\]
it takes the values $p^{2c}\le A\le p^{2c+2}$, with the final interval unbounded. Inverting it gives exactly~\eqref{tf-eq:index-A}. This determines the actual $v(p)$, rather than merely assigning a proposed weight to $p$.

We now specialize transversely. Near the generic point of the exceptional divisor, $p=y^E u$ with $u$ a regular unit after restriction to a nonempty open. Intersect its finitely many translates to make the open invariant. Choose a closed point $z=z_0$ there; only finitely many points of the affine exceptional line are excluded, and $k$ is infinite. Complete the regular local ring at that point and impose
\[
 z-[z_0]=0,
\]
where $[z_0]\in W(k)$ is the Teichm\"uller lift. This is an invariant regular parameter transverse to $y=0$. The quotient $B$ is a complete DVR with uniformizer $y$, residue $k$, and $p=y^E$ times a unit. In particular it has characteristic zero and the same absolute index $E$. The action descends and remains faithful on $y$.

Put $L=\operatorname{Frac}(B)$ and $K=L^{C_m}$. The invariant value group is exactly $m\Z$: the leading residue of an invariant element of valuation $q$ transforms by $\zeta^q$, forcing $m\mid q$, while $y^m$ is invariant. Thus $[L:K]=m$, the fixed ring is a complete DVR of residue $k$, and $e_K=E/m$.

The actual elliptic family pulls back to $B$, and its semilinear generic action descends by finite Galois descent to an elliptic curve over $K$. The labelled full-level sections descend individually. Their generic level remains full because $p\ne0$. The special action has faithful order $m$, so the action on every prime-to-$p$ Tate module has exact order $m$. The good-reduction criterion identifies the exact defect with $m$. Its reduction is additive: it is potentially good and not good, whereas multiplicative reduction would give negative valuation of $j$.
\end{proof}

\subsection{Solving the integer condition and the least-level problem}
We finish by choosing the monomial valuation for each allowed index and then maximizing the torsion level.

\begin{proof}[Proof of Theorem~\ref{thm:tame-full-level}]
Necessity is Proposition~\ref{tf-prop:necessity}. Suppose~\eqref{eq:tame-full-criterion} holds, and put
\[
 M=\frac{me}{p-1},\qquad b=p^{2n-2},\qquad
 c=\min\left\{n-1,\left\lfloor\log_p\frac{M}{(p+1)b}\right\rfloor\right\}.
\]
The bound gives $M\ge(p+1)bp^{\max(0,2n-2-s)}$. Because $s\ge n-1$, the last exponent is at most $n-1$, so
$c\ge\max(0,2n-2-s)$. Hence
\begin{equation}\label{tf-eq:choose-A}
 A=\frac{M}{p^{2n-2-c}}-p^{2c+1}
\end{equation}
is an integer. If $c<n-1$, its defining interval gives
$p^{2c}\le A<p^{2c+2}$; if $c=n-1$, it gives $A\ge p^{2n-2}$. In every case $A\ge1$. Since $m\mid M$, $p\nmid m$ and $p\equiv-1\pmod m$,
\[
 A\equiv-p^{2c+1}\equiv1\pmod m.
\]
Apply Lemma~\ref{tf-lem:realization}. Its good field has absolute index
\[
 E=(p-1)p^{2n-2-c}(A+p^{2c+1})=(p-1)M=me.
\]
Its fixed field therefore has exactly the prescribed index $e$, residue $\overline{\F}_p$, full rational $p^n$-torsion and exact defect $m$. This proves sufficiency for every permitted index.
\end{proof}

\begin{proof}[Proof of Theorem~\ref{thm:tame-level-one}]
At $n=1$, condition~\eqref{eq:tame-full-criterion} reduces to $e\ge(p^2-1)/m$. Lemma~\ref{tf-lem:inertia} and Proposition~\ref{tf-prop:necessity} supply the other necessary conditions. The construction proves their sufficiency. Among $2,3,4,6$, the largest divisor of $p+1$ is respectively $2,6,4,6$ for $p\equiv1,5,7,11\pmod{12}$. This gives the union over all defects. At $p=3$, the two bounds are $e\ge4$ for $m=2$ and $e\ge2$ for $m=4$, with even $e$ in both cases.
\end{proof}

\begin{proof}[Proof of Corollary~\ref{cor:tame-forcing-level}]
Write $m_*=m_{\max}(p)$. If $m_*e<p^2-1$, no nontrivial defect admits full rational $p$-torsion, so the least forcing level is one. Otherwise let $s,L$ be as in the statement. A bad potentially good curve with full rational $p^n$-torsion exists if and only if the criterion holds for $m_*$: all smaller allowable defects give stronger lower bounds, while $m_*$ itself is realized. With $r=n-1$, the exact inequalities are
\[
 r\le s,\qquad 2r\le L,\qquad4r-s\le L.
\]
The largest bad level is therefore
\[
 N=1+\min\left\{s,\left\lfloor L/2\right\rfloor,
                       \left\lfloor(L+s)/4\right\rfloor\right\}.
\]
An actual bad example exists at $N$, and its lower-level torsion is also rational. No bad example exists at $N+1$. For $p\ge5$ all potentially good nontrivial defects were covered by Lemma~\ref{tf-lem:inertia}; for $p=3$ the assertion is restricted to tame inertia and uses $m_*=4$.
\end{proof}

\section{The wild characteristic-three range}\label{sec:three}
This section proves Theorem~\ref{thm:three-range} for the wild defects and combines it with the tame criterion to obtain the largest possible characteristic-three level. Its valuation obstruction and wild invariant branch are proved in full below.

\subsection{Primitive torsion and its valuation sum}
We first isolate the valuation identity that applies before any descent condition is imposed.

\begin{lemma}\label{lem:valuation-sum}
Let $R$ be a complete discrete valuation ring with algebraically closed residue field of characteristic three. Suppose that a one-dimensional formal group of height two has full rational $3^n$-torsion over $\operatorname{Frac}(R)$. Write $d=v_R(3)$, and let $A$ and $B$ be the largest and smallest valuations of primitive $3^n$-torsion coordinates. Then $A\ge B\ge1$ and
\begin{equation}\label{eq:primitive-sum}
 d=2\,3^{n-1}\left(A+3^nB+
       \sum_{j=1}^{n-1}2\,3^{n-j-1}\min(A,3^{2j}B)\right).
\end{equation}
If $0\le r\le n-1$ and $3^{2r}B\le A\le3^{2r+2}B$, with no upper condition when $r=n-1$, then
\begin{equation}\label{eq:regional-sum}
 d=2\,3^{2n-r-2}(A+3^{2r+1}B).
\end{equation}
\end{lemma}

\begin{proof}
Choose primitive points $P,Q$ attaining $A,B$. When $A>B$, they form a basis over $\Z/3^n$: otherwise their reductions modulo three are dependent, so $Q=aP+3R_0$ for a scalar $a$ and a torsion point $R_0$. If $a$ is a unit, both summands have valuation greater than $B$; if it is not a unit, the same holds because multiplication by three strictly increases a positive formal valuation. Either alternative contradicts the choice of $Q$. When $A=B$, any basis has the required valuations.

We use Drinfeld's full-level ring with this marking. Its coordinate factorization, in \citet[Proposition~2.2.3(vii), pp.~439--441]{CarterStrauch2021}, has the following consequence:
\begin{equation}\label{eq:low-multiples}
 v_R(t(3^jQ))=3^{2j}B\qquad(0\le j<n).
\end{equation}
Indeed, after one level transition the coordinate of a primitive point is a unit times the product of its three primitive line lifts, each to the third power. Iterating $j$ times gives $3^j$ factors, each to exponent $3^j$. For the point $3^jQ$, these factors are represented by $aP+Q$ with $a\equiv0\pmod{3^{n-j}}$. Each has valuation $B$, since the contribution from $aP$ has strictly larger valuation. Specializing the universal identity proves~\eqref{eq:low-multiples} in $R$.

Primitive vectors modulo unit scalars form $\mathbf P^1(\Z/3^n)$. There are $3^n$ representatives $(a,1)$, all of valuation $B$, and the representative $(1,0)$ has valuation $A$. For each $1\le j<n$, there are $2\,3^{n-j-1}$ representatives $(1,b)$ with $v_3(b)=j$. Their valuations are $\min(A,3^{2j}B)$. If the two arguments of this minimum coincide, cancellation cannot raise the valuation: the resulting point is still primitive, whereas $A$ was maximal. The same observation handles $(a,1)$ when $A=B$.

Every line has $2\,3^{n-1}$ primitive generators, all with the same valuation. The power series $[3^n](T)/[3^{n-1}](T)$ has constant term three and its roots are precisely the primitive points. Weierstrass preparation therefore identifies $d$ with the sum of their valuations. This gives~\eqref{eq:primitive-sum}. Splitting its sum at $r$ and summing the two geometric series gives~\eqref{eq:regional-sum}. The formulas agree at a boundary $A=3^{2r+2}B$.
\end{proof}

\subsection{What a wild descent action forces}
We next explain why the largest primitive valuation has an extra factor of three when the descent group has wild order three.

\begin{lemma}\label{lem:additive-potential}
An elliptic curve over $K$ with full rational $3$-torsion and additive reduction is potentially good. After a finite good-reduction extension its special fibre is supersingular.
\end{lemma}

\begin{proof}
The Weil pairing gives $\mu_3\subset K$. If $v_K(j(E))<0$, Tate uniformization expresses $E$ as a quadratic twist of a split Tate curve. Its semisimplified three-torsion representation is
$\overline\epsilon\,\overline\chi_3\oplus\overline\epsilon$, where $\epsilon$ is the twisting character. Both $E[3]$ and $\overline\chi_3$ are trivial. Reduction modulo three is injective on $\{1,-1\}$, so $\epsilon=1$. The curve would have split multiplicative reduction, contrary to the hypothesis. Hence $j(E)$ is integral and $E$ is potentially good.

Over an algebraically closed residue field of characteristic three, the automorphism group of an ordinary elliptic curve is $\{1,-1\}$. Its nontrivial element acts nontrivially on the \emph{\'etale} three-torsion. A faithful nontrivial inertia action on an ordinary special fibre would therefore contradict full rational three-torsion. The special fibre is supersingular. We use here the good-reduction and specialization statements of \citet[Theorems~1--2 and Corollary~2 to Theorem~2, pp.~493, 496--497]{SerreTate1968}.
\end{proof}

\begin{lemma}\label{lem:wild-divisibility}
Suppose $E[3^n]$ is rational over $K$ and $E$ has semistability defect $m\in\{3,6,12\}$. Let $L/K$ be the minimal good-reduction extension after completed maximal unramified base change. In the good formal group over $\Ocal_L$, the largest primitive valuation $A$ is divisible by three.
\end{lemma}

\begin{proof}
First, the discriminant of every Weierstrass equation for $E$ is a cube in $K^\times$. To see this, use $E_*:y^2+y=x^3$, whose discriminant is $(-3)^3$ and whose three-torsion is rational over $K\supset\mu_3$. Since both three-torsion modules are constant symplectic planes, there is a symplectic isomorphism $E_*[3]\simeq E[3]$. The Hesse-pencil theorem puts $E$ in the pencil of $E_*$; the discriminant formula for that pencil is $\Delta(E_*)(\det A_t)^3$. These are exactly \citet[Theorem~1.1 and Proposition~5.1]{AnemaTopTuijp2018}. Changing a Weierstrass equation multiplies its discriminant by a twelfth power. Thus $\delta=v_K(\Delta)$ is divisible by three.

Let $\omega_K$ be a rational invariant differential and $\omega_L=u\omega_K$ an integral generator on the good model. The transformation formula for the discriminant gives
\[
 v_L(u)=\frac{m\delta}{12}.
\]
This is an integer divisible by three. Let $\sigma$ generate the wild subgroup and let $F=L^{\langle\sigma\rangle}$. Since the value group of $F^\times$ inside that of $L$ is $3\Z$, choose $b_0\in F^\times$ with $v_L(b_0)=v_L(u)$. Replacing $\omega_L$ by $(b_0/u)\omega_L$ gives a $\sigma$-invariant integral generator. Choose a formal coordinate with this differential at the origin, and make it odd by averaging under $[-1]$; two is a unit. The descent series has linear coefficient exactly one and only odd powers.

The supersingular curve is $y^2=x^3-x$. Its order-three automorphisms translate $x$ by a nonzero element of $\F_3$. Their first nonidentity term in an odd local coordinate has degree three. Hence the descent series is
\[
 t\longmapsto t+c t^3+O(t^5),\qquad c\in\Ocal_L^\times.
\]
For a rational nonzero torsion point with parameter $z$, this implies
\begin{equation}\label{eq:wild-point-displacement}
 v_L(\sigma z-z)=3v_L(z).
\end{equation}
The inverse descent convention changes $c$ by a unit and gives the same equality.

Let $b=v_L(\sigma\pi-\pi)-1$ be the positive lower break. Expanding $z$ in $\pi$ with coefficients in the fixed Cohen ring gives
\begin{equation}\label{eq:uniformizer-displacement}
 v_L(\sigma z-z)\ge v_L(z)+b,
\end{equation}
with equality if $3\nmid v_L(z)$. The lowest term gives equality in that case, and every later term changes at a strictly larger valuation. Applying the inequality to a point of minimum valuation $B$ gives $b\le2B$. If $A>B$ and $3\nmid A$, the equality at $A$ would give $b=2A>2B$, a contradiction.

It remains to consider $A=B$ and $3\nmid A$. Now $b=2A$. Write every primitive coordinate as $z=a_z\pi^A+O(\pi^{A+1})$. Comparing the leading terms of~\eqref{eq:wild-point-displacement} gives
\[
 \bar c\,a_z^3=d_0a_z,\qquad d_0\ne0,
\]
where $d_0=A\,\overline{(\sigma\pi-\pi)/\pi^{2A+1}}$ is independent of $z$. The roots of this additive polynomial form an $\F_3$-space of dimension at most one. On the other hand, for a primitive basis $P,Q$, the leading coefficient of $aP+bQ$ defines an injective $\F_3$-linear map from $\F_3^2$: a nonzero combination is primitive, so its leading coefficient cannot vanish without contradicting maximality of $A$. This contradiction proves $3\mid A$.
\end{proof}

\begin{proposition}\label{prop:necessary-wild}
Every curve in Theorem~\ref{thm:three-range} satisfies~\eqref{eq:wild-criterion}. The same numerical inequality holds for a tame defect $m=2$ or $4$.
\end{proposition}

\begin{proof}
Full rational $3^n$-torsion contains the cyclotomic extension, so $2\,3^{n-1}\mid e$ and $s\ge n-1$. The good model has $d=me$. Choose a region $r$ in Lemma~\ref{lem:valuation-sum}. For a wild defect, $A+3^{2r+1}B$ is divisible by three and $v_3(m)=1$. Equation~\eqref{eq:regional-sum} therefore implies
\[
 s\ge2n-r-2,\qquad r\ge\max(0,2n-2-s).
\]
For a tame defect, $v_3(m)=0$ and the integrality of the parenthesis gives exactly the same inequality. Since $A\ge3^{2r}B$ and $B\ge1$,
\[
 me=2\,3^{2n-r-2}(A+3^{2r+1}B)
       \ge8\,3^{2n+r-2}.
\]
Taking the smallest permitted $r$ proves the asserted condition.
\end{proof}

\subsection{A wild invariant curve in the level space}
We now construct an action-preserving one-dimensional specialization. The proof gives the complete local equation of the branch, rather than assuming that an invariant arc exists in a regular ring.

\begin{lemma}\label{lem:wild-arc}
Let $n\ge1$ and let $A\ge1$ be divisible by three. Set
\begin{equation}\label{eq:arc-degree}
 d=2\,3^{n-1}\left(A+3^n+
            \sum_{j=1}^{n-1}2\,3^{n-j-1}\min(A,3^{2j})\right).
\end{equation}
There is a complete discrete valuation ring $R$ with residue field $\overline{\F}_3$, $v_R(3)=d$, and a good supersingular elliptic curve over $R$ with full rational $3^n$-torsion and primitive valuations having extrema $A,1$.

The ring and marked curve admit a faithful semilinear $C_3$-action preserving every marked point. They admit such an action of $C_6$ if $A$ is odd, and of $C_3\rtimes C_4$ if $A\equiv1\pmod4$. The order-four generator acts by inversion on $C_3$. All actions are trivial on the residue field and induce the faithful indicated automorphisms of $y^2=x^3-x$.
\end{lemma}

\begin{proof}
Put $k=\overline{\F}_3$ and $W=W(k)$. A framed deformation of the supersingular curve has the form
\begin{equation}\label{eq:elliptic-family}
 \mathcal E_u:\quad Y^2=X^3+uX^2-X,\qquad u\in\mathfrak m.
\end{equation}
Indeed, completing the square is possible because two is invertible. The simple root lifting $X=0$ is unique; translate it to zero. A fourth root congruent to one then makes the coefficient of $X$ equal to $-1$. These operations are unique for a coordinate change reducing to the identity, leaving $u$ as the deformation parameter. The discriminant $16(u^2+4)$ is a unit. The Serre--Tate theorem identifies this deformation with the deformation of its height-two formal group.

Use the odd parameter $T=X^{-1/2}$ with sign chosen so that
\[
 Y=-T^{-3}(1+uT^2-T^4)^{1/2},\qquad
 \omega=\frac{dT}{(1+uT^2-T^4)^{1/2}}.
\]
Its logarithm begins $T-uT^3/6$, and comparison of $\log([3]T)$ with $3\log T$ gives
\begin{equation}\label{eq:cubic-coefficient}
 [3]_{T^3}=4u.
\end{equation}
This is an identity of integral coefficients, although it can be calculated after inverting three. On the special fibre, $[3](T)=-T^9$. One verification uses the three-division polynomial $3X^4+4uX^3-6X^2-1$: its tripling formula reduces to $X\mapsto X^9$, $Y\mapsto-Y^9$. For the latter sign, writing $f=X^3-X$, the characteristic-three identity
\[
 (f^2-1)(X-X^9)-f+f^5=0
\]
checks the addition formula directly.

Let $R_n$ be the universal full Drinfeld-level deformation ring. It is regular of dimension two, and the coordinates $x,y$ of a universal basis of $3^n$-torsion are regular parameters \citep[Theorem~2.1.2(ii), p.~437; Proposition~2.2.3(i), p.~439]{CarterStrauch2021}. The automorphisms of the special fibre act on the framing and hence on $R_n$, its elliptic curve, and its marking. Let $\sigma$ lift $X\mapsto X+1$, and let $\tau$ lift the automorphism for which
\[
 \tau u=-u,\qquad \tau T=iT,\qquad i^2=-1.
\]
These have orders three and four, with $\tau\sigma\tau^{-1}=\sigma^{-1}$. Their isomorphisms preserve the universal level structure. Modulo $(3,u)$ the series for $\sigma^j$, $j=0,1,2$, is
\[
 h_j(T)=\frac{T}{(1+jT^2)^{1/2}}.
\]
Because the origin is preserved, the difference of the two series lies in $(3,u)T\,W[[u,T]]$.

Set $M=3^{2n-2}$ and $H=6M$. The level-factorization theorem puts every three-torsion coordinate in $\mathfrak m_{R_n}^{M}$. In a Weierstrass factorization of $[3](T)$, its coefficient of $T^3$ contains at least six of the eight nonzero three-torsion coordinates in every term. Equation~\eqref{eq:cubic-coefficient} gives
\begin{equation}\label{eq:deformation-orders}
 u\in\mathfrak m_{R_n}^{H},\qquad
 3\in\mathfrak m_{R_n}^{8M}.
\end{equation}
In particular $H\ge6$. For $j<n$, the initial term of $[3^j]y$ is $(-1)^j y^{3^{2j}}$: the errors from $(3,u)$ have order at least $H$, whereas $3^{2j}\le M$.

Give $x,y,3$ weights $A,1,d$. The primitive coordinates have initial forms consisting of a nonzero scalar multiple of $Y$, or the lower-weight term of $X+c_jY^{3^{2j}}$, or this binomial when the weights coincide. These forms follow either from the level factorization used in Lemma~\ref{lem:valuation-sum}, or from the preceding exact special multiplication series and formal addition. At a weight equality the nonzero constants are in $\F_3$. Choose $c=i\in k$; then evaluation at $(X,Y)=(c,1)$ kills none of the initial factors because $i\ne\pm1$.

A Cohen presentation can be chosen as
\begin{equation}\label{eq:cohen-relation}
 R_n\simeq W[[X_0,Y_0]]/(3-P(X_0,Y_0)),
\end{equation}
where $P$ is a unit times the product of all primitive coordinates. The sum of their weights is $d$ by~\eqref{eq:arc-degree}. Its initial form has nonzero value $v_0$ at $(c,1)$. Each individual factor has weight less than $d$. Therefore a term in a lift of that factor which is divisible by three has larger weight than the factor, and cannot affect the weight-$d$ initial product. The coefficient of the linear term three in the relation is a unit; the relation generates the Cohen kernel.

Introduce
\[
 x=\pi^A Z,\qquad y=\pi,\qquad3=\pi^d V,
\]
and complete at $Z=c$, $V=v_0$, $\pi=0$. The chart has relations
\begin{equation}\label{eq:chart-relations}
 3=\pi^d V,\qquad V=P(\pi^A Z,\pi)/\pi^d.
\end{equation}
The quotient on the right is an integral convergent series. To verify this also for coefficients divisible by three, expand coefficients in Teichm\"uller digits and replace each $3^b$ by $\pi^{db}V^b$. All resulting exponents of $\pi$ are nonnegative. Its reduction is a polynomial in $Z$, independent of $V$, with value $v_0$ at $c$. Thus the second relation has derivative one in $V$ and eliminates $V$. The chart is a regular complete local ring of dimension two, with parameters $\pi,Z-c$. On it, $u\in\pi^H$ and $3\in\pi^d$, where $d\ge8M>H$.

The wild action extends to this chart before imposing any additional equation. The ratios $\sigma x/x$ and $\sigma y/y$ are units and are one modulo the maximal ideal. Hence
\[
 \sigma Z=(\sigma x/x)(\sigma y/y)^{-A}Z,\qquad
 \sigma V=(\sigma y/y)^{-d}V
\]
preserve its centre and its defining relations.

In the fraction ring of the chart, set $z_x=x^{-2}$, $z_y=y^{-2}$ and define the exact invariants
\begin{align}
 D_y&=\Tr_{\langle\sigma\rangle}(z_y^2),\notag\\
 I&=\frac{\Tr_{\langle\sigma\rangle}\bigl(z_y^2(z_x-z_y)\bigr)}{D_y},
 &U_y&=\Nm_{\langle\sigma\rangle}(y)^{-2}.
 \label{eq:trace-invariants}
\end{align}
The origin-preserving error bound gives
\[
 (\sigma^j x)^{-2}=x^{-2}+j+O(\pi^{H-2A}),\quad
 (\sigma^j y)^{-2}=y^{-2}+j+O(\pi^{H-2}).
\]
Consequently
\begin{align}
 D_y&=3y^{-4}+6y^{-2}+5+O(\pi^{H-4})
       =-1+O(\pi^{H-4}),\notag\\
 I&=z_x-z_y+O(\pi^{H-2A-4}),\notag\\
 U_y&=y^{-6}-y^{-2}+O(\pi^{H-6}).
 \label{eq:invariant-expansions}
\end{align}
For the estimate on $I$, first subtract $(z_x-z_y)D_y$ from its numerator; every remaining error has valuation at least $H-2A-4$. The norm estimate follows from
$\prod_{j=0}^2(w+j)=w^3+3w^2+2w$, whose difference from $w^3-w$ is divisible by three. Since $H>4$, $D_y$ is a unit of the chart. These are convergent integral expansions after the indicated normalizations.

Put $a=A/3$ and choose $\lambda\in W^\times$ reducing to $c^{-2}=-1$. Impose the invariant equation $I=\lambda U_y^a$. After multiplication by $x^2$ it is integral and becomes
\begin{equation}\label{eq:simple-branch}
 1-\pi^{2A-2}Z^2-\lambda Z^2(1-\pi^4)^a
                         +O(\pi^{H-4})=0.
\end{equation}
Its reduction is $1-\bar\lambda Z^2$ and has nonzero derivative $-2\bar\lambda c$ at $Z=c$. The complete local implicit-function theorem eliminates $Z$, as well as $V$ in~\eqref{eq:chart-relations}. The remaining ring is a complete DVR with uniformizer $\pi$ and $3=\pi^d$ times a unit. Every primitive coordinate has its assigned weight, so the extrema are $A,1$.

The equation defining the branch is exactly $C_3$-invariant, and the branch is unique over its residue point. It is therefore stable under $\sigma$. The wild action is faithful because
\[
 \sigma\pi=\pi-\tfrac12\pi^3+O(\pi^5).
\]
For the tame generator, conjugation permutes the trace and norm factors, giving $\tau I=-I$ and $\tau U_y=-U_y$. If $A\equiv1\pmod4$, then $a$ is odd, the equation is stable, and $\tau Z=i^{1-A}Z$ preserves the centre. Also $4\mid d$, so the centre for $V$ is stable. The action on $\pi$ is multiplication by $i$. For $\tau^2$, stability holds when $A$ is odd; then $2\mid d$. This gives the asserted groups and their faithful special actions. The universal marking is preserved at every step.
\end{proof}

\subsection{The integer criterion and the examples}
The invariant branch realizes all the integers permitted by the necessary bound.

\begin{proof}[Proof of Theorem~\ref{thm:three-range}]
Necessity is Proposition~\ref{prop:necessary-wild}. Assume~\eqref{eq:wild-criterion} for a chosen $m\in\{3,6,12\}$. Start with $r_0=\max(0,2n-2-s)$, and take the largest $r\in\{r_0,\ldots,n-1\}$ satisfying
\[
 me\ge8\,3^{2n+r-2}.
\]
Define
\begin{equation}\label{eq:construct-A}
 A=\frac{me}{2\,3^{2n-r-2}}-3^{2r+1},\qquad B=1.
\end{equation}
Since $v_3(me/2)=s+1$ and $2n-r-2\le s$, this is an integer divisible by three. The displayed inequality implies $A\ge3^{2r}$. If $r<n-1$, failure of the next inequality gives $A<3^{2r+2}$. Therefore~\eqref{eq:regional-sum} says precisely that the degree $d$ in Lemma~\ref{lem:wild-arc} is $me$. When $m=6$, the first term in~\eqref{eq:construct-A} is even, so $A$ is odd. When $m=12$, that term is divisible by four, and $3^{2r+1}\equiv3\pmod4$, so $A\equiv1\pmod4$.

Apply the lemma and take the fixed field of the indicated subgroup of order $m$. Its absolute index is $d/m=e$. The elliptic curve descends effectively, for example by using the invariant ample divisor $3[0]$. Its marked torsion points descend individually and are rational. Its faithful special inertia action has order $m$, so its exact semistability defect is $m$. It has additive reduction because it is potentially good and its inertia action is nontrivial.

Every additive elliptic curve in the question is covered by Lemma~\ref{lem:additive-potential}. Over $\overline{\F}_3$, the supersingular automorphism group is $C_3\rtimes C_4$: in $Y^2=X^3-X$, its elements have $X\mapsto a^2X+r$, $Y\mapsto a^3Y$, with $a^4=1$ and $r^3-r=0$. Thus the nontrivial possible defects are $2,3,4,6,12$. The tame possibilities satisfy the stronger bound in Proposition~\ref{prop:necessary-wild}; defect twelve gives the weakest bound and is always realized when that bound holds.

Writing $r=n-1$, the defect-twelve criterion is
\[
 r\le s,\qquad 2r\le L+1,\qquad 4r\le L+s+1.
\]
Its largest $r$ is the minimum in~\eqref{eq:maximum-three}. This proves the formula. A potentially good curve with full torsion above this level cannot be additive, and hence must be good. Existence at level $N(e)$ proves minimality.
\end{proof}

\begin{corollary}\label{cor:least-three-index}
The least absolute ramification index admitting an additive elliptic curve with full rational $3^n$-torsion, with arbitrary perfect residue field allowed, is
\[
 e_{\min}(n)=2\,3^{2n-2}.
\]
\end{corollary}

\begin{proof}
If $s\ge2n-2$, the divisibility $2\,3^s\mid e$ gives the bound. If $s\le2n-3$, the exponent $4n-4-s$ in the weakest, defect-twelve inequality is at least $2n-1$, and again $e\ge2\,3^{2n-2}$. At $e=2\,3^{2n-2}$, the defect-twelve criterion holds. Theorem~\ref{thm:three-range} gives an example.
\end{proof}

\begin{example}\label{ex:twenty-seven}
For $e=216$, one has $u=108$, $s=3$ and $L=4$, so $N(e)=3$. At $n=3$ and $m=12$, choose $r=1$ in~\eqref{eq:construct-A}; it gives $A=21$ and $B=1$. Direct substitution yields
\[
 d=18(21+27+6\cdot9+2\cdot21)=2592=12\cdot216.
\]
The invariant equation is $I=-U_y^7$ with $c=i$. The resulting degree-twelve descent gives full rational $27$-torsion with additive reduction. At $e=162$, the same calculation has $A=9$, $r=1$ and $d=1944$. This attains the least index in Corollary~\ref{cor:least-three-index}. For comparison, a tame defect at $e=216$ would require at least $8\,3^5/4=486$, so it cannot give full $27$-torsion there.
\end{example}

\section{Ordinary reduction and explicit limits of the filtration condition}\label{sec:early-reduction}
This section proves the ordinary good-reduction criterion, gives actual elliptic examples with invisible first-level special characters, and explains why good graded pieces do not force potential semistability of the entire representation.

\subsection{Ordinary automorphisms and unramified torsion}
The ordinary criterion rests on two elementary faithfulness arguments.

\begin{lemma}\label{lem:principal-congruence-torsion}
For odd $p$, the kernel of $\GL_g(\Z_p)\to\GL_g(\F_p)$ is torsion-free.
\end{lemma}
\begin{proof}
Write a nonidentity element as $U=1+p^aB$, where $a\ge1$ and $B$ is nonzero modulo $p$. If $r$ is prime to $p$, the binomial expansion gives
\[
 U^r=1+p^a(rB+pC),\qquad
 U^p=1+p^{a+1}(B+pC')
\]
with integral matrices $C,C'$. In the second formula, the middle binomial terms have valuation at least $a+2$, and the last term has valuation $ap\ge a+2$. Thus taking a prime-to-$p$ power preserves the least matrix valuation, while taking a $p$th power raises it by exactly one. An element of finite order $p^br$ would become the identity after these operations, a contradiction.
\end{proof}

\begin{lemma}\label{lem:ordinary-etale-faithful}
Let $B$ be an ordinary abelian variety of dimension $g$ over an algebraically closed field $k$ of characteristic $p$. An endomorphism annihilating $B[p^m](k)$ for every $m$ is zero. Consequently $\operatorname{Aut}(B,0)$ acts faithfully on $T_p^{\mathrm{et}}B$.
\end{lemma}
\begin{proof}
For an abelian variety $C/k$ of dimension $h$, the factorization $[p^m]=V^mF^m$ has $F^m$ radicial and bijective on geometric points, and $V^m$ of degree $p^{mh}$. Hence
\begin{equation}\label{eq:ordinary-point-count}
 |C[p^m](k)|\le p^{mh},\qquad |B[p^m](k)|=p^{mg}.
\end{equation}
The equality is the ordinary condition. The factorization and degree are the standard Frobenius--Verschiebung statements of \citet[Proposition~5.20, paragraph~5.21 and Proposition~5.22, pp.~80--81]{EdixhovenGeerMoonen}.

Suppose $f\ne0$ annihilates all these points. The reduced identity component $H$ of $\ker(f)$ is an abelian subvariety of dimension $h<g$. Indeed the reduced subgroup over the perfect field $k$ is a group scheme; translation carries its nonempty smooth locus to every point, so it is smooth. Its identity component is connected and proper. If it had dimension $g$, then $f$ would vanish on all of $B$. The reduced kernel has finitely many components, say $c$, each a translate of $H$.

Each such translate contains at most $|H[p^m](k)|$ points of $B[p^m](k)$: if it contains one, subtracting that point identifies the intersection with $H[p^m](k)$. Equation~\eqref{eq:ordinary-point-count} would therefore give $p^{mg}\le c p^{mh}$ for all $m$, impossible for $h<g$. Thus $f=0$. If an automorphism is the identity on $T_p^{\mathrm{et}}B$, apply this assertion to its difference from the identity. The finite quotients of that Tate module are precisely $B[p^m](k)$.
\end{proof}

\begin{proof}[Proof of Theorem~\ref{thm:ordinary-unramified-p}]
The dimension-zero case is immediate. Pass to the strict henselian unramified base to describe inertia, and choose a finite Galois extension $L$ over which $A$ has good ordinary reduction, with abelian model $\mathcal A$ and special fibre $B$. Generic descent isomorphisms extend uniquely to the good models by the N\'eron mapping property \citep[Theorem~17.1 and Remark~17.2, p.~70]{Milne2008AV}. Their cocycle identities extend by uniqueness. Since inertia fixes the residue field, reduction gives a finite group of origin-preserving automorphisms of $B$.

The ordinary connected--\'etale sequence of $\mathcal A[p^\infty]$ is functorial under these isomorphisms. Its \'etale quotient has height $g=\dim A$, and its generic $p$-torsion is an equivariant quotient of $A[p]$. The unramified hypothesis makes inertia trivial on that quotient. Finite \'etale specialization over the henselian base identifies it with $B[p](k)=T_p^{\mathrm{et}}B/p$. Thus the finite special inertia image in $\GL_g(\Z_p)$ is congruent to the identity modulo $p$.

Lemma~\ref{lem:principal-congruence-torsion} makes this image trivial; Lemma~\ref{lem:ordinary-etale-faithful} then makes every special inertia automorphism of $B$ trivial. Smooth proper base change identifies the action on $T_\ell B$ with the original inertia action on $T_\ell A$, for $\ell\ne p$. The N\'eron--Ogg--Shafarevich criterion over the original $K$ now gives good reduction \citep[Theorem~1, p.~493]{SerreTate1968}. No cyclotomic determinant or polarization was used, so no roots-of-unity or polarization-degree hypothesis is needed.
\end{proof}

\subsection{A tame defect-two example at every odd prime}
The first family is a direct specialization of the complete tame existence theorem.

\begin{proof}[Proof of Corollary~\ref{cor:defect-two-full-p}]
In Theorem~\ref{thm:tame-level-one} take $m=2$ and $e=(p^2-1)/2$. For odd $p$, one has $2\mid p+1$ and $(p-1)\mid e$, and the lower endpoint is attained. The tame $p=3$ clause gives the same conclusion at $e=4$. The construction in Section~\ref{sec:tame-full-torsion} gives supersingular reduction; alternatively Theorem~\ref{thm:ordinary-unramified-p} excludes the ordinary case for an additive curve with full rational $p$-torsion.

If $P,Q$ form a basis of $E[p](K)$, the perfect alternating Weil pairing $e_p(P,Q)$ has order $p$ and is $K$-rational by equivariance \citep[Section~I.13, p.~57]{Milne2008AV}. Thus $\mu_p\subset K$ without enlarging the field. The nontrivial special inertia automorphism $\alpha$ has order two and fixes the origin. Since $(\alpha-1)(\alpha+1)=0$ and a nonzero elliptic endomorphism is a finite surjective isogeny, $\alpha\ne1$ forces $\alpha=-1$. Its trace on the rank-two rational Dieudonn\'e module is $-2$.
\end{proof}

\subsection{An explicit defect-six model in characteristic three}
The next calculation determines both the full torsion and the exact good-reduction field of one curve.

\begin{proof}[Proof of Proposition~\ref{prop:explicit-cubic-defect-six}]
Put $\zeta=\zeta_3$, $\tau=\zeta-\zeta^2=2\zeta+1$. Then $\tau^2=-3$, $v_K(\tau)=1$ and $v_K(3)=2$. The standard invariants of $y^2+y=x^3$ give $\Delta=-27$ and
\[
                      \psi_3(x)=3x^4+3x.
\]
Its eight nonidentity $3$-torsion points are
\[
 (0,0),\ (0,-1),\quad
 (-\zeta^i,\zeta),\ (-\zeta^i,\zeta^2)\quad(0\le i\le2).
\]
They satisfy the equation and the division polynomial, are distinct, and together with the identity exhaust the nine geometric $3$-torsion points in characteristic zero.

The displayed model is minimal because $v_K(\Delta)=6$ and a change of Weierstrass model changes that valuation by an integer multiple of twelve. A smaller nonnegative value is impossible. The reduced singular point is $(-1,1)$; the substitution $X=x+1$, $Y=y-1$ gives $Y^2=X^3$. This is a cusp and has additive smooth group. Explicitly its smooth locus has parameter $w=X/Y$, with $w=0$ at infinity and $(X,Y)=(w^{-2},w^{-3})$ away from infinity. A line gives an equation $cw^3+aw+b=0$, whose three roots sum to zero, so the chord law is addition of $w$.

Choose $r^3=-1/4$ and $u^2=\tau$, and put $L=K(r,u)$. The unit $r$ reduces to $-1$. If $s=r+1$, then
\[
                     s^3-3s^2+3s-3/4=0.
\]
For $t=v_K(s)>0$, the four term valuations are $3t,2+2t,2+t,2$. A vanishing sum cannot have a unique term of least valuation, so $3t=2$. Thus $K(r)/K$ has ramification index and degree three; it is cyclic because $\mu_3\subset K$. Also $v_K(u)=1/2$, giving a ramified quadratic extension. Their compositum is cyclic of degree six and has ramification index six, since its index is divisible by both two and three.

Under
\[
                    x=r+\tau X,\qquad y=u^3Y-1/2,
\]
the equation becomes
\begin{equation}\label{eq:explicit-cubic-good-model}
                      Y^2=X^3-\tau rX^2-r^2X.
\end{equation}
Its coefficients are integral and its discriminant is
\[
 16r^4(\tau^2r^2+4r^2)=16r^6(\tau^2+4)=16r^6=1.
\]
This is an actual good model. Its special fibre is $E_0:Y^2=X^3-X$. With the standard convention
\[
 \psi_3=3X^4+b_2X^3+3b_4X^2+3b_6X+b_8,
\]
one has $b_8=-1$, and hence $\psi_3=-1$ in characteristic three. There is no nonidentity geometric $3$-torsion point; the elliptic curve has $p$-rank zero and is supersingular. Equivalently, for $Y\ne0$ the tangent slope is $-1/(2Y)$, and the doubling formula gives $X(2P)=X(P)+1/(4Y(P)^2)\ne X(P)$, excluding $2P=-P$; points with $Y=0$ have order two.

The commuting generators $\sigma_3(r)=\zeta r$, $\sigma_3(u)=u$ and $\sigma_2(r)=r$, $\sigma_2(u)=-u$ act in good coordinates by
\[
 \sigma_3(X)=X+\frac{(1-\zeta)r}{\tau}=X+\zeta^{-1}r,
 \quad \sigma_3(Y)=Y,
 \qquad \sigma_2(X)=X,\quad\sigma_2(Y)=-Y.
\]
After reduction these are $(X,Y)\mapsto(X-1,Y)$ and $(X,Y)\mapsto(X,-Y)=[-1]$. The first cyclically permutes the three nonzero $2$-torsion points $(0,0),(1,0),(-1,0)$, so is nontrivial on $T_2E_0$. The second is $-1$ on that Tate module. Any nontrivial subgroup of the cyclic group of order six contains its subgroup of order two or three, so this action is faithful.

Smooth proper base change therefore makes the inertia image on $T_2E$ have exact order six. If $E$ becomes good over $F/K$, NOS makes $I_F$ lie in the kernel of that quotient, so $6\mid[I_K:I_F]=e(F/K)$. The field $L$ attains six. Since $j(E)=0$, multiplicative reduction is impossible after any extension: a nodal minimal cubic has $c_4$ a unit and $v(j)=3v(c_4)-v(\Delta)<0$. Hence good reduction and semistability have the same minimal ramification here.

After unramified base change, $v(\tau)=1$ and $v(3)=2$ remain unchanged. The same two valuation calculations and the same faithful special action still apply. This proves the exact defect also over the maximal unramified base and its completion. The special involution is $[-1]$ and has trace $-2$.
\end{proof}

\subsection{The logarithmic cyclotomic extension}
The final example explains the representation-theoretic scope of the allowed filtration class.

\begin{proof}[Proof of Proposition~\ref{prop:filtered-non-ht}]
Let $\chi_p:G_K\to\Z_p^\times$ be the cyclotomic character. Since $K=\Q_p(\zeta_p)$, its values lie in $1+p\Z_p$. The convergent logarithm gives the additive continuous function $c(g)=\log_p\chi_p(g)\in p\Z_p$. Set
\begin{equation}\label{eq:log-cyclotomic-representation}
 T=\Z_pe_1\oplus\Z_pe_2,\qquad
 \rho(g)=\begin{pmatrix}1&c(g)\\0&1\end{pmatrix}.
\end{equation}
Additivity of $c$ proves the representation identity. The submodule $\Z_pe_1$ and its quotient are trivial Galois modules.

Let $G_m$ be the finite \'etale $K$-group scheme corresponding to $T/p^mT$. The equivariant injections $a\mapsto pa$ from $T/p^mT$ to $T/p^{m+1}T$ give a Barsotti--Tate system of height two. Indeed the kernel of $p^r$ on $T/p^{m+r}T$ is $p^mT/p^{m+r}T$, the image of $T/p^rT$ under the indicated transition. The invariant line and quotient at every level are constant and compatible, giving the exact sequence in the proposition. Its Tate module is $T$.

Put $N(e_2)=e_1$, $N(e_1)=0$. Then $N\ne0$, $N^2=0$ and $\rho(g)=\exp(c(g)N)$. The cyclotomic kernel acts trivially. In the cyclotomic Sen module, choose a sufficiently small $\gamma$ fixing the finite coefficient field of the basis. The normalized logarithm is
\[
               \Theta=\frac{\log\rho(\gamma)}{\log_p\chi_p(\gamma)}=N.
\]
The operator is nonsemisimple, so Sen's criterion excludes Hodge--Tateness. Its compatibility with finite extension leaves the same nonzero nilpotent operator after every finite base extension, excluding potential Hodge--Tateness. Since semistable representations are de Rham and de Rham representations are Hodge--Tate, potential semistability is also excluded.

This is precisely the classical example in \citet[Section~II.1.2, p.~9 of the 2002 author version]{Berger2004}. Each graded piece of the constructed group extends to the constant height-one \'etale group over $\Ocal_K$ and has special inertia character one. Their sum is the constant character two. This graded character does not imply a potentially good realization of the whole group.
\end{proof}

\section{Height-one characters and first-level elliptic torsion}\label{sec:first-level-refinements}
This section proves the all-prime height-one assertion, exact Hesse defects over both a specified field and a variable-ramification family, and the residue-field obstruction over a fixed cyclotomic base.

\subsection{The finite scalar in height one}
The proof identifies the finite descent scalar directly from the first torsion layer.

\begin{proof}[Proof of Proposition~\ref{prop:height-one-index}]
A height-one group in the allowed filtration class has a single nonzero graded piece, so itself acquires good reduction. Its dimension $d$ is zero or one. In dimension zero the good model is \'etale; in dimension one its Cartier dual is \'etale. Over the inertia of a good-reduction field its Tate character is therefore respectively $1$ or $\chi_p$. On the original inertia group it has the form
\[
                         \rho=\chi_p^d\epsilon,
\]
where $\epsilon$ is a finite $\Z_p^\times$-valued character. Indeed $\rho/\chi_p^d$ is trivial on the finite-index inertia subgroup of a good field. This same finite scalar is the special descent character. For the \'etale model it is the scalar on the constant Tate module. For the multiplicative model the assertion follows by Cartier duality; the extra generic cyclotomic factor contributes no additional special descent scalar.

Suppose first that $p$ is odd. Lemma~\ref{lem:principal-congruence-torsion} shows that $1+p\Z_p$ is torsion-free, so reduction embeds every finite subgroup of $\Z_p^\times$ into $\F_p^\times$. Since $\mu_p\subset K$, one has $\chi_p\equiv1\pmod p$. Thus the actual representation $M[p]$ determines $\epsilon\bmod p$, and its unique Teichm\"uller lift recovers every value of $\epsilon$. Recovering the dimension is unnecessary for the special-character conclusion.

For $p=2$, the finite subgroup of $\Z_2^\times$ is $\{1,-1\}$. To see this, reduction modulo four has image of order two, while $1+4\Z_2$ is torsion-free: a nonzero leading term $2^a b$ with $a\ge2$ gains exactly one valuation under squaring and retains its valuation under odd powers. Every odd-order element of the finite inertia image therefore has special scalar one. The requested prime-to-two character is constant one and is determined by $M[2]$. The index is nonnegative by definition, so its least value is zero in every prime.
\end{proof}

\subsection{An explicit Hesse specialization}
The following calculation proves both full torsion for the particular equation and the exact ramification needed for its good reduction.

\begin{proof}[Proof of Theorem~\ref{thm:explicit-hesse-twelve}]
The polynomial $X^2+1$ is separable and irreducible modulo three, so $\Q_3(i)/\Q_3$ is unramified quadratic. Over that field $\pi=\zeta-1$ satisfies the Eisenstein polynomial $X^2+3X+3$. Thus $e_K=2$, the residue field is $\F_9$, and $v_K(\pi)=1$. In particular
\[
 \pi^2=-3(\pi+1),\qquad \pi^3=6\pi+9,\qquad \pi^6=-27.
\]
First allow $t$ to be a variable and put
\[
 A=-27t(t^3+8),\qquad B=54(t^6-20t^3-8).
\]
Direct multiplication gives
\begin{equation}\label{eq:hesse-invariants}
 \Delta=2^{12}3^9(t^3-1)^3,\qquad c_4=2^4 3^4t(t^3+8).
\end{equation}
For the first identity, put $U=t^3$ and use
\[
 (U^2-20U-8)^2-U(U+8)^3=-64(U-1)^3.
\]
The equation is therefore smooth whenever $t^3\ne1$.

Let $\tau=\zeta-\zeta^2$, so $\tau^2=-3$. Its eight nonzero $3$-torsion points are
\begin{align*}
 x_*&=-9t^2,& y_*&=\pm12\tau(t^3-1),\\
 x_j&=3t^2+12(\zeta^jt+\zeta^{-j}),&
 y_j&=\pm36(\zeta^jt^2+\zeta^{-j}t+1),\qquad j=0,1,2.
\end{align*}
Here is a direct verification. The three-division polynomial is
\[
 \psi_3(x)=3x^4+6Ax^2+12Bx-A^2.
\]
It follows from the tangent slope $(3x^2+A)/(2y)$ and the equality $x(2P)=x(P)$; for an affine point with $y\ne0$, that equality gives $2P=-P$. Polynomial division yields
\[
 \frac{\psi_3(x)}{3(x+9t^2)}
 =x^3-9t^2x^2+27t(t^3-16)x-27(t^3+8)^2.
\]
After $x=3t^2+12w$, the cubic on the right becomes
\[
 1728\bigl(w^3-3tw-(t^3+1)\bigr),
\]
whose roots are $w=\zeta^jt+\zeta^{-j}$. At $x_*$ the cubic equation has right side $-432(t^3-1)^2=y_*^2$. At $x_0=3(t+2)^2$ it has right side $1296(t^2+t+1)^2$. For general $j$, replace $t$ by $t'=\zeta^{-j}t$ and multiply $x_0(t')$ by $\zeta^{-j}$; the equation is carried back to the original one and gives the displayed $y_j$.

The four $x$-coordinates are distinct: an equality $x_j=x_k$ for $j\ne k$ forces $t=\zeta^{-j-k}$, while $x_j=x_*$ forces $t^2+\zeta^jt+\zeta^{-j}=0$. Either condition implies $t^3=1$. A displayed $y$ can vanish only under the same condition. Thus the eight points are distinct nonidentity points killed by three, and exhaust the nonzero geometric $3$-torsion in characteristic zero. This proves full rational torsion for the actual curve over any field containing $\mu_3$ with $t^3\ne1$.

Now specialize to $t=1+i\pi$. Exact reduction gives
\begin{align*}
 t^3&=10-9i+(9-3i)\pi,\\
 t^3-1&=9(1-i)+3(3-i)\pi,\\
 t^3+8&=9(2-i)+3(3-i)\pi.
\end{align*}
In each of the last two lines the summands have valuations four and three. Hence $v_K(t^3-1)=v_K(t^3+8)=3$, and \eqref{eq:hesse-invariants} gives $v_K(\Delta)=27$ and $v_K(c_4)=11$.
The change $x=\pi^4X$, $y=\pi^6Y$ gives the exact integral model
\begin{equation}\label{eq:hesse-minimal-model}
 Y^2=X^3+aX+b,\qquad
 a=-3-12i-8i\pi,\quad b=-30+12i-20\pi.
\end{equation}
For example,
\[
 t^6-20t^3-8=-27(15-6i+10\pi)=\pi^6(15-6i+10\pi),
\]
which verifies the constant coefficient after scaling; the linear coefficient follows from $-27=\pi^6$. Both $a$ and $b$ have valuation one. The new discriminant has valuation three, so the model is minimal: an integral model with smaller nonnegative discriminant valuation would differ by a positive multiple of twelve. Its special fibre is the cusp $Y^2=X^3$, hence reduction is additive. Its $c_4$ has valuation three and $j$ has valuation six.

Let $r$ be a root of $X^3+aX+b$. This polynomial is Eisenstein, so $K(r)/K$ is totally ramified of degree three and $v_K(r)=1/3$. The element $c=a+3r^2$ has valuation one, since $v_K(3r^2)=8/3$. Choose $u$ with $u^4=c$. In $K(r)$ the normalized valuation of $c$ is three, so adjoining this fourth root has exact degree and ramification index four. Thus $L=K(r,u)$ has degree and ramification index twelve over $K$, with unchanged residue field.
The substitution $X=r+u^2Z$, $Y=u^3W$ gives
\[
                     W^2=Z^3+(3r/u^2)Z^2+Z.
\]
The middle coefficient has valuation $2+1/3-1/2=11/6>0$. Its special fibre $W^2=Z^3+Z$ has discriminant $-64\ne0$ and three-division polynomial $-1$ in characteristic three. It is smooth and supersingular, proving actual good reduction over $L$.

For the matching lower bound, let $F/K$ give good reduction. The discriminant transformation forces $12\mid3e(F/K)$, hence $4\mid e(F/K)$. Since two is invertible, the $2$-torsion of a good elliptic scheme is finite \'etale, and all its points become rational over an unramified extension of $F$. The point $(r,0)$ in \eqref{eq:hesse-minimal-model} has $v_K(r)=1/3$. Unramified extensions preserve the value group, so $3\mid e(F/K)$. Every good-reduction index is divisible by twelve, and $L$ attains it. These arguments remain valid after unramified base change. Finally $v(j)>0$ excludes multiplicative reduction after extension, since a nodal minimal fibre has $c_4$ a unit and $\Delta$ of positive valuation. Thus the exact semistability defect is twelve.
\end{proof}

\subsection{A variable-ramification Hesse family}
The same Hesse equation realizes exact wild defects over an infinite range of complete bases.

\begin{proof}[Proof of Theorem~\ref{thm:hesse-variable-defect}]
Put $s=\varpi^r$, $t=1+s$, $S=e+r$ and $\delta=t^3-1$. Since $2r<e$,
\[
 v_K(\delta)=v_K(t^3+8)=3r.
\]
Indeed the terms in $\delta=3s+3s^2+s^3$ have values $e+r,e+2r,3r$, and the last is strictly smallest; adding nine does not change that value. With $A=-27t(t^3+8)$ and $B=54(t^6-20t^3-8)$, the identity
\[
 t^6-20t^3-8=\delta^2-18\delta-27
\]
gives $v_K(B)=3e+6r$. The other two terms have larger values because $2r<e$. Thus
\begin{equation}\label{eq:variable-hesse-values}
 \begin{aligned}
 v_K(A)&=3S,&v_K(B)&=3e+6r,\\
 v_K(\Delta)&=9S,&v_K(c_4)&=4e+3r,&v_K(j)&=3e>0.
 \end{aligned}
\end{equation}
Here the invariant identities are~\eqref{eq:hesse-invariants}. In particular $\delta\ne0$. The eight rational points and their universal division-polynomial verification in the proof of Theorem~\ref{thm:explicit-hesse-twelve} apply to every such $t$ over a field containing $\mu_3$. They therefore prove full rational $3$-torsion for this actual curve.

We find the cubic factor in a good-reduction field using the simpler polynomial
\[
                         q(z)=z^3-3tz+2.
\]
The fractional-linear expression
\begin{equation}\label{eq:hesse-two-torsion-map}
              \alpha(z)=\frac{3((t^3-4)z-6t^2)}{2+tz}
\end{equation}
satisfies the exact polynomial identity
\[
 \bigl(\alpha(z)^3+A\alpha(z)+B\bigr)(2+tz)^3
                       =-1728\delta^2q(z).
\]
Its determinant is $24\delta\ne0$, and $q(-2/t)=8(1-t^{-3})\ne0$. It therefore bijects the roots of $q$ with the roots of the two-division cubic and gives $K(\alpha)=K(z)$ for corresponding roots. The polynomial $q$ also results by setting $X=Y$ in the Hesse plane cubic $X^3+Y^3+Z^3-3tXYZ=0$ and putting $z=Z/X$. The displayed fractional-linear identity supplies the field comparison directly.

Set $z=1+w$. Then
\[
                     q(1+w)=w^3+3w^2-3sw-3s.
\]
The coefficient points $(1,S)$ and $(2,e)$ lie strictly above the line joining $(0,S)$ to $(3,0)$: $S>2S/3$ and $e>S/3$, the latter following from $r<2e$. Hence every root has value $v_K(w)=S/3$. This also follows directly because either a larger or a smaller value would give a unique least-valued term. Since $3\nmid S$, the denominator forces the cubic to be irreducible and totally ramified of degree three. Put $K_1=K(z)=K(\alpha)$.

The original cubic $X^3+AX+B$ has a single Newton segment: its coefficient at degree one lies strictly above the segment from degree zero to degree three, since
\[
                    3e+3r>\frac23(3e+6r)=2e+4r.
\]
Thus $v_K(\alpha)=e+2r$. For $c=A+3\alpha^2$ the summands have values $3e+3r$ and $3e+4r$, so $v_K(c)=3S$ and $v_{K_1}(c)=9S$ in the normalized integer valuation. Choose $u$ with $u^4=c$. Every unit of $K_1$ has a fourth root by Hensel's lemma, because its residue field is algebraically closed of characteristic three. Removing that unit reduces the extension to a fourth root of a uniformizer power. Its degree and ramification index are exactly
\[
                     h=\frac4{\gcd(4,9S)}=\frac4{\gcd(4,S)}.
\]
Consequently $L=K_1(u)$ has degree and ramification index $3h$ over $K$.

The substitution $x=\alpha+u^2X$, $y=u^3Y$ gives
\[
                      Y^2=X^3+(3\alpha/u^2)X^2+X.
\]
Its middle coefficient has value
\[
                    e+(e+2r)-\frac{3S}{2}=\frac S2>0.
\]
The reduction is the smooth cubic $Y^2=X^3+X$ in characteristic three. Its three-division polynomial is $-1$, so it has no nonidentity geometric three-torsion and is supersingular. Hence $L$ gives actual good reduction at ramification degree $12/\gcd(4,S)$.

For the matching lower bound let $F/K$ be any finite good-reduction extension and put $d=e(F/K)$. Its valuation ring is strictly henselian, so the finite \'etale group scheme $E[2]$ of the good model is constant. In particular it contains a two-torsion coordinate $\alpha$, and hence $z$ by~\eqref{eq:hesse-two-torsion-map}. The value $v_K(z-1)=S/3$ forces $3\mid d$. A good Weierstrass equation has unit discriminant, and changes of Weierstrass equation multiply discriminants by twelfth powers. Therefore
\[
                         12\mid9dS,
\]
so $4/\gcd(4,S)\mid d$. This divisor is prime to three; their product divides every possible $d$, and $L$ attains it. Finally $v_K(j)>0$ excludes multiplicative reduction after any extension. The curve is potentially good by the displayed model, so the good-reduction and semistability defects coincide. This proves the claimed exact defect.
\end{proof}

\subsection{The first-level defect table}
The complete existence criteria supply actual elliptic curves at every index in the table.

\begin{proof}[Proof of Corollary~\ref{cor:three-first-level-table}]
The Weil pairing forces $\mu_3\subset K$, hence $e$ is even. Theorem~\ref{thm:ordinary-unramified-p} excludes an ordinary good special fibre for an additive curve with full rational $3$-torsion. Over an algebraic closure, a supersingular special fibre has equation $y^2=x^3-x$. Its origin-preserving automorphisms are
\[
 (x,y)\longmapsto(a^2x+r,a^3y),\qquad a^4=1,\quad r^3-r=0.
\]
Comparison of the coefficients in the general Weierstrass coordinate change gives these conditions, so the group has order twelve. Inertia embeds in it by smooth proper base change and faithfulness on a prime-to-three Tate module. The only possible nontrivial orders are therefore $2,3,4,6,12$; trivial inertia would give good reduction.

At $n=1$, Theorem~\ref{thm:tame-full-level} for $m=2,4$ gives exactly $me\ge8$ with positive even $e$. For $m=3,6,12$, Theorem~\ref{thm:three-range} has $s\ge0$ and exponent $\max(0,-s)=0$, giving the same inequality. The five rows follow immediately. Both theorems include actual every-index realizations with algebraically closed residue field, so this argument supplies existence as well as necessity. Proposition~\ref{prop:explicit-cubic-defect-six} gives the asserted example over the particular field $\Q_3(\zeta_3)$.
\end{proof}

\subsection{The residue obstruction at the critical index}
At equality in the basic ramification bound, normalized torsion coordinates must fit into the residue field of the original base.

\begin{lemma}\label{lem:critical-torsion-residue}
Let $E/K$ have full rational $p$-torsion and potentially supersingular reduction. Suppose it has good reduction over a finite Galois extension $L/K$ of ramification degree $m$. Then
\[
                            p^2-1\le me_K.
\]
If equality holds, there is an injective $\F_p$-linear map
\[
                            E[p](K)\longrightarrow k_K.
\]
In particular equality is impossible when $k_K=\F_p$.
\end{lemma}
The inequality is the first-level form of the full-division-field bound; in the number-field setting it is covered by \citet[Theorem~1.1, p.~2]{Smith2023}. The proof below includes the complete-field setting and the equality descent to $k_K$.
\begin{proof}
Choose an integral formal parameter $T$ at the origin of the good model over $\Ocal_L$, with $w(L^\times)=\Z$. Every $p$-torsion point reduces to the origin, because the good special fibre is supersingular. Its $p^2-1$ nonzero coordinates lie in the maximal ideal of $\Ocal_L$. The series $[p](T)/T$ has constant term $p$ and distinguished degree $p^2-1$. Weierstrass preparation gives a unit times a monic polynomial whose roots are exactly these coordinates. Its constant term gives
\[
                    \sum_{P\ne0}w(T(P))=w(p)=me_K.
\]
Each summand is a positive integer, proving the inequality.

At equality every summand equals one. Fix $P_0\ne0$, put $z=T(P_0)$ and define
\[
 \lambda(P)=\overline{T(P)/z},\qquad\lambda(0)=0.
\]
The integral formal law is $S+T$ plus terms of degree at least two, so $\lambda$ is additive. It is injective because every nonzero coordinate has valuation one. Hence it is an injective $\F_p$-linear map into $k_L$.

We must descend the image to $k_K$. Generic descent isomorphisms extend uniquely to the good models by the N\'eron mapping property \citep[Theorem~17.1 and Remark~17.2, p.~70]{Milne2008AV}. For $\sigma\in\Gal(L/K)$ their formal expressions are
\[
 F_\sigma(Z)=a_\sigma Z+\text{terms of degree at least two},
 \qquad a_\sigma\in\Ocal_L^\times.
\]
A rational point satisfies $T(P)=F_\sigma(\sigma T(P))$. Modulo valuation two, numerator and denominator of $T(P)/z$ transform by the same factor $a_\sigma$. Their residue ratio is fixed by every $\sigma$. Since $L/K$ is Galois and the residue field is perfect, the induced residue group is $\Gal(k_L/k_K)$, so $\lambda(P)\in k_K$. This proves the stronger equality statement for the actual finite torsion group.
\end{proof}

\begin{proof}[Proof of Corollary~\ref{cor:cyclotomic-full-p}]
Let $p\ge5$, $K=\Q_p(\zeta_p)$ and suppose such an elliptic curve exists. The shifted cyclotomic polynomial gives $e_K=p-1$ and $k_K=\F_p$. Theorem~\ref{thm:ordinary-unramified-p} excludes ordinary potential good reduction, so the good special fibre is supersingular.

For a prime $\ell\ne p$, potential good reduction makes the inertia image on $T_\ell E$ finite. Smooth proper base change identifies it with a faithful subgroup of the good special fibre's automorphism group. In characteristic at least five, a short equation $y^2=x^3+Ax+B$ shows that origin-preserving automorphisms are $(x,y)\mapsto(u^2x,u^3y)$, with $u^4A=A$ and $u^6B=B$. The translation and extra $y$-terms vanish because three and two are invertible. Thus this automorphism group is cyclic of order dividing four or six, and the finite inertia image has order $m\le6$.

Choose a sufficiently large $r$ that this finite inertia image remains faithful on $E[\ell^r]$. Then $L=K(E[\ell^r])$ is finite Galois with ramification degree exactly $m$. Its inertia acts trivially on the whole $T_\ell E$, so NOS gives good reduction over $L$ \citep[Theorem~1, p.~493]{SerreTate1968}. Lemma~\ref{lem:critical-torsion-residue} yields
\[
                         p^2-1\le m(p-1)\le6(p-1).
\]
This implies $p+1\le6$, excluding $p\ge7$. For $p=5$ it forces $m=6$ and equality. The lemma would then inject the $25$-element group $E[5](K)$ into $\F_5$, a contradiction. Finally Proposition~\ref{prop:explicit-cubic-defect-six} gives the example for $p=3$.
\end{proof}

\section{Rational inverse chains in a height-two formal group}\label{sec:cyclic}
This section proves the cyclic-point restriction by following the Newton slopes at every inverse step.

\begin{lemma}\label{lem:inverse-polygons}
Let $R$ and $F$ satisfy the hypotheses of Theorem~\ref{thm:cyclic-bound}. There is an integral coordinate and a deformation parameter $u\in\mathfrak m_R$ such that, writing $a=v_R(p)$, $b=v_R(u)$, and $q=p^2$, the nonzero $p$-torsion valuations are as follows. If $b\ge ap/(p+1)$, all have valuation $a/(q-1)$. Otherwise they have valuations
\begin{equation}\label{eq:alpha-beta}
 \alpha=\frac{a-b}{p-1}\quad(p-1\text{ points}),\qquad
 \beta=\frac{b}{p(p-1)}\quad(q-p\text{ points}).
\end{equation}
Here $b=+\infty$ is permitted when $u=0$.

If $z$ is a nonzero torsion point of valuation $0<Z<a$, its $p$-preimages have valuations
\begin{equation}\label{eq:inverse-polygons}
 \begin{cases}
 Z/q & b\ge Z(p-1)/p,\\[2pt]
 (Z-b)/p\text{ with multiplicity }p,
 \quad\beta\text{ with multiplicity }q-p
 & b<Z(p-1)/p.
 \end{cases}
\end{equation}
The equality case belongs to the first line. Once an inverse chain takes a step of the first kind, or a step to $\beta$, all its later inverse steps divide the valuation by $q$.
\end{lemma}

\begin{proof}
The integral universal height-two coordinate has
\[
 [p](T)=pT+(1-p^{p-1})uT^p+\sum_{j>p}c_jT^j,
 \qquad c_j\in(p,u)R\ (p<j<q),\quad c_q\in R^\times.
\]
All coefficients are integral and degrees two through $p-1$ vanish. This is the specialization of \citet[Proposition~5.7, equations~(12.3)--(12.5), and Proposition~12.10]{HopkinsGross1994}; the complete local base may have any residue-field extension. The formula for the coefficient of $T^p$ also follows directly by comparing $f([p]T)=pf(T)$ in the normalized logarithm $f(T)=T+(u/p)T^p+O(T^{p^2})$.

The only possible lower Newton vertices of $[p](T)$ at positive roots are $(1,a),(p,b),(q,0)$. Indeed intermediate coefficients have valuation at least $\min(a,b)$, later coefficients are integral, and all the omitted points lie on or above the indicated negative-slope segments. The middle point is below the segment joining the other two exactly when $b<ap/(p+1)$. This proves~\eqref{eq:alpha-beta}.

For $[p](T)-z$, the point $(1,a)$ lies strictly above the lower hull because $Z<a$. The possible vertices are $(0,Z),(p,b),(q,0)$. The middle point is below their single chord exactly when $b<Z(p-1)/p$, giving~\eqref{eq:inverse-polygons}. Every inverse valuation is smaller than its input. If the first inequality holds at $Z$, it continues to hold at $Z/q$. If the next valuation is $\beta$, then
$b=p(p-1)\beta>\beta(p-1)/p$, so the same conclusion applies. Finally every nonzero first-level valuation is at most $a/(p-1)<a$, so the inverse formulas apply at every subsequent step.
\end{proof}

\begin{proof}[Proof of Theorem~\ref{thm:cyclic-bound}]
Let $P$ be rational of order $p^n$. All points in its successive multiplication chain are rational, and hence their positive coordinate valuations are integers. Put $N=p^2-1$ and $q=p^2$.

In the one-slope case the terminal valuation is $a/(Nq^{n-1})$. Integrality implies $a\ge Nq^{n-1}$, which is stronger than the desired inequality. In the two-slope case, if the first-level point $p^{n-1}P$ has valuation $\beta$, then all $n-1$ inverse steps are uniform. Thus $b\ge p(p-1)q^{n-1}$, and the strict two-slope inequality gives
\[
 a>\frac{p+1}{p}b\ge Nq^{n-1}.
\]
We may therefore assume that the first-level point has valuation $Z_0=\alpha$.

Call a step \emph{high} when it uses $(Z-b)/p$ in the second line of~\eqref{eq:inverse-polygons}. After $j$ consecutive high steps, the recurrence gives
\begin{equation}\label{eq:high-recurrence}
 Z_j=\frac{a/p^j-b}{p-1},\qquad
 a=p^j\bigl(b+(p-1)Z_j\bigr).
\end{equation}
Since $Z_j,b$ are integers, $j\le v_p(a)=v$.

Suppose the first non-high step is uniform, following $j\le\min(v,n-2)$ high steps. There remain $r=n-1-j\ge1$ uniform steps. Their terminal valuation is integral, so $Z_j\ge q^r$. The uniform-step inequality says $b\ge Z_j(p-1)/p$. Hence
\[
 a=p^j\bigl(b+(p-1)Z_j\bigr)
   \ge Np^{j-1}Z_j
   \ge Np^{2n-3-j}
   \ge Np^{2n-3-v}.
\]
If $j<v$, the last inequality is strict. If $j=v$ and equality held throughout, then $a=Np^{2n-3-v}$, and therefore $v_p(a)=2n-3-v>v$, since $v=j\le n-2$. This is a contradiction. Thus the final inequality is strict also at the Newton boundary.

Suppose instead that the first non-high step is low, following $j$ high steps. After it, $n-2-j$ uniform steps remain. Consequently
\[
 b\ge p(p-1)q^{n-2-j}.
\]
The strict high-versus-low condition at $Z_j$ is equivalent, using~\eqref{eq:high-recurrence}, to $a/p^j>(p+1)b$. Therefore
\[
 a>Np^{2n-3-j}\ge Np^{2n-3-v}.
\]
Finally, if all $n-1$ inverse steps are high, then $v\ge n-1$, $b\ge1$ and $Z_{n-1}\ge1$. Equation~\eqref{eq:high-recurrence} gives
\[
 a\ge p^n>Np^{n-2}\ge Np^{2n-3-v}.
\]
These cases exhaust every inverse chain and prove the theorem.
\end{proof}

\begin{remark}
For $p\nmid a$ this says $a>(p^2-1)p^{2n-3}$. The dependence on $v_p(a)$ comes entirely from the integrality assertion in~\eqref{eq:high-recurrence}. The Newton slopes are the familiar canonical-subgroup slopes; the inequality is a convenient consequence of their full inverse-chain classification.
\end{remark}

\begin{proposition}\label{prop:mixed-obstruction}
Let $p\ge5$, let $e=e_K$ and $v=v_p(e)$. Let $M/K$ be a height-two Barsotti--Tate group with potentially supersingular good reduction. Let $H/K$ have a two-step filtration by potentially good height-one groups whose dimensions sum to one. If
$M[p^n]\simeq H[p^n]$ for $n\ge2$, then
\begin{equation}\label{eq:mixed-obstruction}
 e>p^{2n-3-v}.
\end{equation}
\end{proposition}

\begin{proof}
Pass to the completed maximal unramified extension, which preserves $e$ and the proposed torsion comparison. Over this base a potentially good height-one character is either $\epsilon$ or $\chi_p\epsilon$, with $\epsilon$ finite of order dividing $p-1$. The determinant of a potentially supersingular height-two group is $\chi_p\epsilon_0$ with the same kind of finite character.

Transport the first height-one piece of $H[p^n]$ to a primitive cyclic subgroup $C$ of $M[p^n]$. If its character is $\epsilon$, twisting by $\epsilon^{-1}$ gives a rational point of exact order $p^n$. If its character is $\chi_p\epsilon$, first quotient $M$ by $C$. The image of $M[p^n]$ in the quotient is cyclic of order $p^n$, with character $\epsilon_0\epsilon^{-1}$; a finite twist again makes it rational. Schematic closure of $C$ in a good model is finite flat, so the quotient still acquires a good formal model of height two. Finite twisting preserves potential good reduction.

The faithful finite special inertia of a height-two formal group embeds in the units of the quaternion division algebra of invariant $1/2$ over $\Q_p$. Every element of this algebra satisfies a quadratic reduced polynomial. An element of order $p$ would generate $\Q_p(\zeta_p)$ of degree $p-1>2$, which is impossible. Its finite subgroups therefore have order prime to $p$ and inject, by reduction of the maximal order, into $\F_{p^2}^\times$. Thus a common faithful good-reduction field has tame degree $m\mid p^2-1$. The passage from faithful special inertia to a good model for the original lattice follows from crystalline descent and the weight-$\{0,1\}$ realization theorem of \citet[Corollary~2.2.6]{Kisin2006}; a finite-flat quotient adjusts the lattice if necessary.

Over that field the rational point is a formal point and the absolute index is $a=me$, with $v_p(a)=v$. Theorem~\ref{thm:cyclic-bound} gives $me>(p^2-1)p^{2n-3-v}$. Since $m\le p^2-1$, this implies~\eqref{eq:mixed-obstruction}.
\end{proof}

\section{Exact cyclotomic base change and the different}\label{sec:different}
This section proves Theorems~\ref{thm:cyclotomic-exact} and~\ref{thm:cyclotomic-bound} for every odd prime.

For a finite totally ramified extension $A/B$, write
$d(A/B)=v_A(\mathfrak D_{A/B})$. We use
\begin{equation}\label{eq:different-tower}
 d(A/C)=d(A/B)+[A:B]d(B/C).
\end{equation}
This is \citet[Chapter~III, \S4, Proposition~8, p.~51]{Serre1979}. When $\Ocal_A=\Ocal_B[\pi]$ and $f$ is the minimal polynomial of $\pi$, the different is generated by $f'(\pi)$ \citep[Chapter~III, \S6, Proposition~11 and Corollary~2, pp.~55--57]{Serre1979}. For a finite Galois extension with lower ramification groups $G_i$,
\begin{equation}\label{eq:different-lower}
 d(A/B)=\sum_{i\ge0}(|G_i|-1),
\end{equation}
as in Chapter~IV, \S1, Proposition~4, p.~64 of that source. These statements require no finite-residue-field hypothesis.

\begin{proof}[Proof of Theorem~\ref{thm:cyclotomic-exact}]
A Cohen coefficient field gives $[K:K_0]=p(p-1)$, hence $[K:K_1]=p$. The Eisenstein derivative calculation given below proves $d\le p^2-1$. Let $T/K_1$ be the normal closure of $K/K_1$. Its Galois group acts faithfully and transitively on the $p$ roots of a minimal polynomial. Since the residue field is algebraically closed, the entire group is inertia. Its wild subgroup is a nontrivial normal $p$-subgroup of the symmetric group on $p$ letters, hence has order $p$ and acts regularly. The normalizer of this subgroup is the affine group $C_p\rtimes\F_p^\times$. Thus
\[
 \Gal(T/K_1)=C_p\rtimes C_m,\qquad m\mid p-1,
\]
with faithful tame action when $m>1$. The extension $T/K$ is tame of degree $m$. If $b$ is the unique positive lower break of $T/K_1$, the different formula in the tower $T/K/K_1$ gives
\[
 pm-1+(p-1)b=(m-1)+md,
 \qquad u_T=\frac bm=\frac{d-(p-1)}{p-1}\le p.
\]
Here $u_T$ is its unique positive upper break.

Set $C_0=K_0(\mu_{p^3})$. The cyclic extension $C_0/K_1$ has upper breaks $p-1$ and $2(p-1)$: its lower breaks, computed on $\zeta_{p^3}-1$, are $p-1$ and $p^2-1$. It is disjoint from $T$ over $K_1$. Indeed, a nontrivial intersection would contain $K_0(\mu_{p^2})$ and give a cyclic quotient of order $p$ of $C_p\rtimes C_m$. If $m>1$, the faithful tame action makes the $C_p$ factor trivial in the abelianization. If $m=1$, the intersection would force $K=K_0(\mu_{p^2})$, contrary to the hypothesis.

Since $2(p-1)>u_T$, the last cyclotomic ramification group persists after the $T$ factor has become trivial. The quotient and subgroup forms of Herbrand's theorem therefore give the last upper break of $TC_0/T$ as
\begin{align*}
 \psi_{T/K_1}\bigl(2(p-1)\bigr)
 &=m u_T+pm\bigl(2(p-1)-u_T\bigr)\\
 &=m\bigl((p-1)(2p+1)-d\bigr).
\end{align*}
The tame extension $T/K$ has degree $m$ and is disjoint from the cyclotomic extension over $K$. Tame base change multiplies its positive upper breaks by $m$. Dividing the displayed expression by $m$ proves the identity; $d\le p^2-1$ gives the lower bound $p(p-1)$.
\end{proof}

\begin{proof}[Proof of Theorem~\ref{thm:cyclotomic-bound}]
Choose a Cohen coefficient ring $W(k)\subset\Ocal_K$ and write $K_0=W(k)[1/p]$. The extension $K/K_0$ is finite of degree $e_K$. For completeness, if $\pi$ is a uniformizer of $K$, then $1,\pi,\ldots,\pi^{e_K-1}$ give a basis of $\Ocal_K/p\Ocal_K$ over $k$. Successive $p$-adic lifting spans $\Ocal_K$ over $W(k)$, and distinct valuations modulo $e_K$ prove independence. This argument works for infinite perfect $k$.

Put
\[
 K_1=K_0(\mu_p),\qquad C_0=K_0(\mu_{p^3}),\qquad C=KC_0.
\]
Cyclotomic Eisenstein polynomials give $[K_1:K_0]=p-1$ and $[K:K_1]=p$. For odd $p$, $C_0/K_1$ is cyclic of degree $p^2$ and its unique degree-$p$ subfield is $K_0(\mu_{p^2})$. The hypothesis $\mu_{p^2}\not\subset K$ therefore implies $K\cap C_0=K_1$. Consequently
\[
 [C:K]=p^2,\qquad[C:C_0]=p.
\]
All fields have the same residue field and all extensions in these towers are totally ramified.

We need two elementary bounds:
\begin{equation}\label{eq:different-bounds}
 d(K/K_1)\le p^2-1,
 \qquad d(C/C_0)\ge p-1.
\end{equation}
A uniformizer of $K$ generates its ring of integers over $\Ocal_{K_1}$. Its minimal polynomial is Eisenstein,
$f(X)=X^p+\sum_{i=0}^{p-1}a_iX^i$.
The leading derivative term $p\pi^{p-1}$ has valuation $p(p-1)+(p-1)=p^2-1$. Each other nonzero term $ia_i\pi^{i-1}$ has valuation congruent to $i-1$ modulo $p$, because $1\le i<p$ and $v_K(a_i)$ is divisible by $p$. These residues are pairwise distinct and exclude the residue $-1$ of the leading term. Thus the least valuation occurs only once and cannot cancel, proving the first inequality. For the degree-$p$ extension $C/C_0$, the Eisenstein derivative formula gives valuation at least $p-1$ for every term, proving the second.

The cyclotomic extension $C_0/K_1$ has lower breaks $p-1$ and $p^2-1$. Indeed, $\zeta_{p^3}-1$ is a uniformizer. An exponent $1+pa$, with $p\nmid a$, displaces it by an element of valuation $p$; an exponent $1+p^2a$, with $p\nmid a$, gives valuation $p^2$. Summing these displacements gives
\[
 d(C_0/K_1)=(p^2-p)p+(p-1)p^2=2p^2(p-1).
\]
Different transitivity in the two towers over $K_1$ yields
\begin{align*}
 d(C/K)+p^2d(K/K_1)
 &=d(C/C_0)+p\,d(C_0/K_1),\\
 d(C/K)&\ge p-1+2p^3(p-1)-p^2(p^2-1)\\
 &=p^2(p-1)^2+p-1.
\end{align*}

A totally ramified cyclic group of order $p^2$ has two positive breaks. In terms of its upper breaks $\delta_1<\delta_2$, the lower breaks are $\delta_1$ and $p\delta_2-(p-1)\delta_1$. The upper-to-lower conversion is the Herbrand function of \citet[Chapter~IV, \S3, pp.~73--76]{Serre1979}. Summing the lower groups gives
\[
 d(C/K)=(p^2-1)(\delta_1+1)
 +(p-1)\bigl(p\delta_2-p\delta_1\bigr)
 =(p-1)(\delta_1+p\delta_2+p+1).
\]
Substitution of the preceding lower bound gives
$\delta_1+p\delta_2\ge p(p^2-p-1)$. If $\delta_1\le p-1$, divide the resulting inequality for $p\delta_2$ by $p$ to obtain~\eqref{eq:second-break}.
\end{proof}

\begin{remark}
No integrality theorem for upper breaks was used. The two-tower estimate alone gives $\delta_2\ge91/5$ at $p=5$ when $\delta_1\le4$. The exact normal-closure argument gives the unconditional stronger bound $\delta_2\ge20$ under the stated hypotheses. Neither argument needs Hasse--Arf.
\end{remark}

\section{Tame weights and integral character comparison}\label{sec:formal-comparison}
We now prove a uniform comparison theorem for two supersingular groups. The point of the argument is to estimate multiplication by $p$ on the quotient appearing in the graph construction. Its coefficients have extra divisibility forced by the common determinant.

\begin{lemma}\label{lem:weighted-kernel}
Let $p\ge5$, $N=p^2-1$, and let $L/K$ be cyclic, totally tamely ramified of degree $N$, with algebraically closed residue field. Write $R=\Ocal_L$, $a=v_L(p)$, and choose $\tau\pi=\zeta\pi$ with $\zeta\in K$ of order $N$. Let $G/R$ be a formal group of finite height with semilinear descent action. Suppose that its special fibre is a product of one-dimensional groups of height two, and that the exponents $c_i$ of the cotangent action satisfy
\[
                         c_i\equiv c_j\pmod{p-1}.
\]
Define $q_0=1$ and
\begin{equation}\label{eq:weighted-recurrence}
                   q_{j+1}=\min\{a+q_j,pq_j+p-1\}.
\end{equation}
If a homomorphism $f:F\to G$ has finite special kernel and no nonzero finite flat subgroup in its kernel, then
\[
 q_u\ge a/(p-1)\quad\Longrightarrow\quad F^u\ker(f_k)=0,
\]
where $F^u$ denotes relative Frobenius iterated $u$ times.
\end{lemma}
\begin{proof}
Average lifts of a cotangent eigenbasis by the tame idempotents. The resulting parameters $X_i$ satisfy $\tau X_i=\zeta^{c_i}X_i$ exactly. If $b_{j,\alpha}X^\alpha$ occurs in the $j$th component of $[p]$, equivariance gives
\[
       \tau b_{j,\alpha}=\zeta^{c_j-\sum_i\alpha_i c_i}b_{j,\alpha}.
\]
The valuation of a nonzero eigenvector of $\tau$ is congruent to its exponent modulo $N$: divide by that power of $\pi$ and reduce the unit. For $|\alpha|=p$ the displayed exponent is divisible by $p-1$. The order of the reduced multiplication map is $p^2$, in these coordinates as in any coordinates with invertible linear part. Every nonzero coefficient of total degree $p$ consequently has valuation at least $p-1$.

Let $U(X)dX$ be a normalized invariant differential matrix. The identity
\[
                 U([p]X)J_{[p]}(X)=pU(X)
\]
shows that every entry of the Jacobian is divisible by $p$. A coefficient of a monomial having some exponent prime to $p$ is therefore divisible by $p$. The linear part is $pX$; all coefficients of degrees less than $p^2$ vanish on reduction. Substitution of a vector in $\pi^q tR[[t]]^d$, for $q\ge1$, now has coefficient valuations at least
\[
 \min\{a+q,pq+p-1,2pq+1,p^2q\}
                      =\min\{a+q,pq+p-1\}.
\]
Indeed the first term handles the linear part and monomials with an exponent prime to $p$, the second handles total degree $p$, the third the remaining degrees below $p^2$, and the fourth all later terms. Iteration proves the bound $q_u$ for $[p]^u(\pi h)$.

We give the integral passage from this bound to the kernel, because a bound for torsion points would not suffice. A normalized logarithm has a coefficient of total degree $r$ of valuation at least $-a\lfloor\log_p r\rfloor$. This follows by integrating the integral invariant differentials and choosing a nonzero exponent in the monomial. Since $\lfloor\log_p r\rfloor\le(r-1)/(p-1)$, substitution into the logarithm at coefficient valuation $q_u\ge a/(p-1)$ yields
\begin{equation}\label{eq:weighted-logarithm}
              p^u\log_G(\pi h)=\log_G([p]^u(\pi h))
                              \in\pi^{q_u}R[[t]]^d.
\end{equation}
This is a formal, coefficientwise assertion, so equality at the bound is allowed.

Here are the Cartier-module facts used to apply it. Put $\mathcal R=R[[\Delta]][1/p]$ and $J=\pi^{\lceil a/(p-1)\rceil}R$. Write $L\{\{\Delta\}\}$ for two-sided series with coefficient valuations bounded below and tending to infinity as the exponent tends to $-\infty$; $J\{\{\Delta\}\}$ imposes the additional condition that every coefficient lies in $J$. The skew series ring $\Omega$ has coefficients in $W(k)$ with this decay condition and product
\[
 (a_i\Delta^i)(b_j\Delta^j)=\sigma^j(a_i)b_j\Delta^{i+j},
\]
where $\sigma$ is Witt Frobenius. Set $G_\pi(X,Y)=\pi^{-1}G(\pi X,\pi Y)$, and let $E(G_k)$ denote the covariant Dieudonn\'e module. The invariant Cartier module $D_G\subset L[[\Delta]]^d$ is obtained from logarithms of integral $p$-typical curves. Its completion $\mathcal M_G=\Omega D_G\subset L\{\{\Delta\}\}^d$ has
\[
 \mathcal M_G=\Omega\pi D_{G_\pi},\qquad
 J[[\Delta]]^d\subset D_G,
 \qquad E(G_k)=D_G/(D_G\cap\mathcal R^d).
\]
For every matrix $A$ over $L$, the completion embedding is the identity on bounded elements of $AD_G$, and a vector maps to a bounded element only if it was bounded. These are \citet[Proposition~9 and Lemma~3, p.~674; Proposition~10(4), p.~675; Proposition~11(1)--(2), pp.~675--676]{Bondarko2006}. The ring $\Omega$ has integral Witt coefficients and preserves coefficient valuation ideals. Equation~\eqref{eq:weighted-logarithm} therefore implies
\[
                         p^u\mathcal M_G\subset J\{\{\Delta\}\}^d.
\]
Consequently, if $A$ is integral, $z\in D_G$, and $Az$ is bounded, then $p^uAz$ has every coefficient in $J$. Since
\[
 \mathbf f^u\left(\sum_{j\ge0}z_j\Delta^j\right)
                  =p^u\sum_{j\ge0}z_{j+u}\Delta^j,
\]
we obtain $\mathbf f^uAz\in J[[\Delta]]$.

Put $I'=df(0)L^{\dim F}$. The absence of a finite flat kernel gives
\[
                  df(0)D_F=D_G\cap I'[[\Delta]].
\]
To see this, the right side corresponds to a formal group by the Cartier classification; $F$ maps to it by an isogeny, which must be an isomorphism because its kernel would be finite flat. Choose an integral complement to $I'\cap R^d$ and let $\operatorname{pr}$ be the resulting projection. If the reduction of $z\in D_G$ lies in the rational saturation of $E(F_k)$ inside $E(G_k)$, then $\operatorname{pr}z$ is bounded. The previous paragraph gives
\[
 c=\mathbf f^u\operatorname{pr}z\in J[[\Delta]]^d\subset D_G.
\]
It reduces to zero, while $\mathbf f^uz-c$ lies in $df(0)D_F$. Thus $\mathbf f^u$ kills the saturation quotient. By the covariant kernel identity, this quotient is the Dieudonn\'e module of $\ker(f_k)$ and $\mathbf f$ is relative Frobenius. This last argument is the kernel argument of \citet[Proposition~16 and proof, p.~681; Theorem~2(8), equation~(6), p.~668]{Bondarko2006}, with the improved coefficient estimate just established. It proves the assertion.
\end{proof}

\begin{theorem}\label{thm:common-determinant}
Let $p\ge5$ and let $K$ have algebraically closed residue field and absolute index $e$. Let $M_1,M_2$ be height-two Barsotti--Tate groups acquiring supersingular good reduction, with common Tate determinant $\chi_p\psi$, where $\psi$ is finite. Set $a=e(p^2-1)$ and define $q_j$ by~\eqref{eq:weighted-recurrence}. If $n\ge2$ and $q_{n-1}\ge a/(p-1)$, an isomorphism $M_1[p^n]\simeq M_2[p^n]$ implies equality of their finite special inertia characters.
\end{theorem}
\begin{proof}
A faithful finite special inertia group embeds in the units of the endomorphism algebra of a height-two formal group. This is the quaternion division algebra over $\Q_p$. It has no element of order $p$ for $p\ge5$, since $\Q_p(\mu_p)$ has degree $p-1>2$, whereas a commutative subfield has degree at most two. Its finite inertia subgroup is therefore tame and cyclic, of order dividing $N=p^2-1$. Both groups have good models over the common tame extension $L/K$ of degree $N$. As explained at the end of Section~\ref{sec:cyclic}, one first obtains good reduction up to isogeny by crystalline descent, and then realizes the given lattice by a finite flat quotient. This justifies using the specified torsion modules on these good models.

Suppose that the two special characters differ. Their eigenvalue multisets are disjoint Frobenius orbits $\{\lambda,\lambda^p\}$ and $\{\mu,\mu^p\}$ in $\mu_N$. Integral polynomials in a tame generator project the product of the two special groups onto its factors: the corresponding characteristic polynomials have relatively prime reductions.

Let $H$ be the schematic closure of the graph of the torsion isomorphism in $\mathcal M_1[p^n]\times\mathcal M_2[p^n]$. Over $L$, the determinant is cyclotomic. The alternating Tate pairings, and their inverse duality maps, extend to the good groups by \citet[Theorem~4, p.~180, and Corollaries~1--2, p.~181]{Tate1967}. Rescaling one pairing by a unit makes $H$ isotropic for their difference. It is then its own orthogonal: both finite flat subgroups have the same rank, and their inclusion becomes equality on coordinate rings, which are free modules of that rank. Cartier duality commutes with reduction, so the same holds for $H_k$.

The tame projectors preserve $H_k$, giving $H_k=J_1\times J_2$. A finite subgroup of rank $p^r$ in a one-dimensional formal group over $k$ has ideal $(T^{p^r})$ and equals $\ker F^r$. Self-orthogonality inside each height-two $p^n$-torsion factor forces $r=n$. Hence
\[
                        H_k=\ker F_1^n\times\ker F_2^n.
\]
Form $Q=(\mathcal M_1\times\mathcal M_2)/H$ and $G=Q\times\mathcal M_1$. The map
\[
 f:\mathcal M_1\times\mathcal M_2\longrightarrow G,
             \qquad z\longmapsto(qz,\operatorname{pr}_1z)
\]
has no finite flat kernel: generically the graph intersects $0\times M_2$ trivially. Its special kernel is $0\times\ker F_2^n$.

Let $c_1,c_2$ be the two cotangent exponents. The common finite determinant is their common reduced norm; hence $c_1\equiv c_2\pmod{p-1}$. Quotient formation and relative Frobenius show that the three cotangent exponents on $G_k$ are
\[
                         p^nc_1,\quad p^nc_2,\quad c_1.
\]
They are congruent modulo $p-1$. Lemma~\ref{lem:weighted-kernel}, applied to this target with $u=n-1$, would make $F^{n-1}\ker F_2^n$ zero. Its rank is $p$, a contradiction. Therefore the characters agree.
\end{proof}

\section{Cyclic points producing character ambiguity}\label{sec:ambiguity}
We construct the comparisons that show the ramification bounds are sharp. The construction has two parts: an invariant Eisenstein equation produces a rational cyclic point, and an extension of rank-one representations reproduces the resulting finite torsion module.

\begin{lemma}\label{lem:extension-lift}
Let $K$ have characteristic zero, and let $\psi:G_K\to\Z_p^\times$ be finite of order prime to $p$. Every continuous extension of free $\Z/p^n$-modules
\[
 0\longrightarrow\mu_{p^n}\longrightarrow V_n
       \longrightarrow(\Z/p^n)(\psi^{-1})\longrightarrow0
\]
lifts to a height-two Barsotti--Tate group over $K$ with graded groups $\mu_{p^\infty}$ and $(\Q_p/\Z_p)(\psi^{-1})$.
\end{lemma}
\begin{proof}
Set $T=\Z_p(1)\otimes\psi$. An underlying module splitting identifies the extension with a class in $H^1(K,T/p^nT)$. Over the finite extension $K'/K$ killing $\psi$, Kummer theory identifies the untwisted group at level $r$ with $K'^\times/(K'^\times)^{p^r}$. Its transition maps are surjective. Inflation--restriction identifies the original cohomology with the invariants under $\Delta=\Gal(K'/K)$ of the suitably twisted Kummer group. Averaging by $|\Delta|^{-1}\sum_{\delta\in\Delta}\delta$ preserves surjectivity. Thus the given class lifts compatibly to every higher level.

Choose continuous cocycles for these classes. At each step the reduction differs from the previous cocycle by a coboundary; lift the correcting zero-cochain and subtract its coboundary. The resulting cocycles are exactly compatible. Their inverse limit is a continuous $T$-valued cocycle, giving an extension of free $\Z_p$-modules with the prescribed reduction. The corresponding system of finite \'etale group schemes over $K$ is a Barsotti--Tate group. Its two graded groups acquire good reduction after killing $\psi$.
\end{proof}

\begin{lemma}\label{lem:point-ambiguity}
Suppose that $K$ has algebraically closed residue field, $p\ge5$, and $G/K$ is height two with supersingular good reduction over a tame extension. Suppose that $G$ has a rational point of exact order $p^n$ and special eigenvalues $a,a^p$ on a tame generator, where $a\in\mu_{p^2-1}$ and $a\ne1$. There exist a potentially supersingular group $M$ and a group $H$ filtered by two potentially good height-one groups, with
\[
                         M[p^n]\simeq H[p^n],
\]
whose finite special characters differ.
\end{lemma}
\begin{proof}
The determinant of $T_pG$ is $\chi_p\psi$, with $\psi$ a finite character valued in $\mu_{p-1}$, and $\psi(\tau)=a^{p+1}$. Indeed its restriction to a good field is crystalline of weight one. After twisting by $\chi_p^{-1}$ it has weight zero; the rank-one case of \citet[Corollary~2.2.6, p.~26 of the author version]{Kisin2006} identifies it with the rational Tate module of a height-one \'etale group. Such a group is constant over an algebraically closed residue field. The determinant is therefore cyclotomic on the good field, and its remaining finite value is the determinant of the special isocrystal action. The same argument shows that a potentially good height-one character is finite or cyclotomic times finite.

Let $C$ be the constant cyclic group generated by the point and put $G'=G/C$. Closure of $C$ in a good model is finite flat, so $G'$ still acquires supersingular good reduction. The multiplication-by-$p^n$ exact sequence gives
\[
 0\longrightarrow C\longrightarrow G[p^n]
       \longrightarrow G'[p^n]\longrightarrow C\longrightarrow0.
\]
The first $C$ is primitive as a free cyclic submodule; its quotient in $G[p^n]$ therefore has character $\chi_p\psi$. Twisting $G'$ by $\psi^{-1}$ gives $M$ and an extension of the form in Lemma~\ref{lem:extension-lift}. Apply that lemma to obtain $H$ with the same truncation. Isogeny preserves the special rational character. The two eigenvalue multisets are
\[
               \{a^{-1},a^{-p}\},\qquad\{1,a^{-(p+1)}\}.
\]
Putting $\lambda=a^{-1}$, the difference of their traces is
\[
                 \lambda+\lambda^p-1-\lambda^{p+1}
                              =-(\lambda-1)(\lambda^p-1)\ne0.
\]
This proves the claim on the finite tame inertia quotient.
\end{proof}

\subsection{Every index above the prime-to-$p$ threshold}
\begin{theorem}\label{thm:general-arc}
Let $p\ge5$, $n\ge2$, and let $e$ be a positive multiple of $p-1$ satisfying $e>p^{2n-3}$. There is a complete field $K$ with residue field $\overline{\F}_p$, absolute index $e$, and $\mu_p\subset K$, and groups $M,H$ as in Lemma~\ref{lem:point-ambiguity} whose $p^n$-torsion is isomorphic but whose finite special characters differ. In particular $\nuMin(p,e,h)\ge n$ for every $h\ge2$.
\end{theorem}
\begin{proof}
Set $q=p^2$, $N=q-1$, $Q=q^{n-1}$ and $x_0=e/Q$. Choose an integer
\[
                          x_0<A<px_0,\qquad N\nmid A.
\]
Such an integer always exists. If $x_0<1$, take $A=1$, using $px_0>1$. If $x_0\ge1$, take $\lfloor x_0\rfloor+1$ and increase it by one when it is divisible by $N$. The resulting integer is at most $x_0+2<px_0$.

Over $W=W(\overline{\F}_p)$ use the integral universal formal group $F_u$ of Section~\ref{sec:cyclic}. Its logarithm satisfies
\begin{equation}\label{eq:homogeneous-logarithm}
 f_u(T)=\sum_{j\ge0}l_j(u)T^{p^j},\quad l_0=1,
 \qquad pl_j=l_{j-1}u^{p^{j-1}}+l_{j-2},\quad l_{-1}=0.
\end{equation}
Put $d=Ne$, $b=d-(p-1)AQ$. The chosen interval gives
\begin{equation}\label{eq:general-arc-weights}
 b>(q-p)Aq^{n-2}>0,\qquad AQ<d<NAQ,
                 \qquad b\equiv-(p-1)A\pmod N.
\end{equation}
For example the first difference is exactly $N(e-Ap^{2n-3})$.
Define
\[
 Y_0=X^A,\qquad Y_{j+1}=[p]_{X^b}(Y_j),\qquad
 P(X)=\Phi_{X^b}(Y_{n-1}),\quad \Phi_u(T)=[p]_u(T)/T.
\]
These are compositions of integral power series in $W[[X]]$.

Modulo $p$, the multiplication series begins with $uT^p+T^q$; intervening coefficients are multiples of $u$, and all later coefficients are integral. The first inequality in~\eqref{eq:general-arc-weights} makes the $T^q$ term strictly earliest at every one of the first $n-1$ compositions. Hence
\[
                        Y_j\equiv X^{Aq^j}+O(X^{Aq^j+1})\pmod p.
\]
In the final series $P$, the $uY_{n-1}^{p-1}$ term has degree $d$. All intervening terms are later, and $Y_{n-1}^N$ has degree $NAQ>d$. Thus
\[
                         P(0)=p,\qquad P\equiv X^d+O(X^{d+1})\pmod p.
\]
Weierstrass preparation makes $R=W[[X]]/(P)$ an Eisenstein discrete valuation ring of absolute index $d$, with uniformizer $x$. The inequality $AQ<d$ prevents coefficients divisible by $p$ from changing the valuations $v_R(Y_j(x))=Aq^j$. Therefore $x^A$ has exact order $p^n$ in $F_{x^b}$.

For a primitive $N$th root $\zeta\in W$, set $a=\zeta^{-A}$. The recursion~\eqref{eq:homogeneous-logarithm} proves
\[
 l_j(a^{p-1}u)=a^{p^j-1}l_j(u),\qquad
 [p]_{a^{p-1}u}(a^{-1}T)=a^{-1}[p]_u(T),\qquad
 \Phi_{a^{p-1}u}(a^{-1}T)=\Phi_u(T).
\]
For the $l_{j-2}$ summand the exponents differ by $p^j-p^{j-2}$, which is divisible by $N$. Thus these are identities of the entire power series. Since $b\equiv-(p-1)A\pmod N$, substitution $X\mapsto\zeta X$ sends $u$ to $a^{p-1}u$, every $Y_j$ to $a^{-1}Y_j$, and $P$ to itself. It induces an automorphism $\sigma(x)=\zeta x$ of $R$, of order $N$.

Let $L=\operatorname{Frac}(R)$ and $K=L^{\langle\sigma\rangle}$. The extension is tame of degree $N$, and $e_K=e$. The integral isomorphism $\sigma^*F_{x^b}\to F_{x^b}$ given by $T\mapsto aT$ satisfies the cocycle identity and fixes the marked point, because $a\sigma(x^A)=x^A$. Descent at every finite generic level produces $G/K$ with that rational point. Its special eigenvalues are $a,a^p$; they are the two embeddings of the quadratic Teichm\"uller unit acting on the Honda isocrystal.

Finally $K$ contains $\mu_p$. Write $-p=\pi_K^e u$ with $u$ a unit. Hensel's lemma supplies a $(p-1)$st root of $u$, and $p-1\mid e$, so $K$ contains a $(p-1)$st root of $-p$. Over $W[1/p]$ this generates the unique tame extension of degree $p-1$, namely $W[1/p](\mu_p)$. Since $N\nmid A$, we have $a\ne1$, and Lemma~\ref{lem:point-ambiguity} applies. Adding the same constant height-$h-2$ group proves the last assertion.
\end{proof}

\begin{corollary}\label{cor:sharp-mixed-prime}
For $p\ge5$, $n\ge2$, and $(p-1)\mid e$ with $p\nmid e$, a supersingular/height-one filtered comparison with isomorphic $p^n$-torsion and different special characters exists over some field of index $e$ if and only if $e>p^{2n-3}$.
\end{corollary}
\begin{proof}
Theorem~\ref{thm:general-arc} proves existence. Proposition~\ref{prop:mixed-obstruction} proves necessity with $v_p(e)=0$.
\end{proof}

\subsection{A third-level construction at the next wild index}
\begin{theorem}\label{thm:two-step-arc}
Let $p\ge5$ and let $e$ be a positive multiple of $p(p-1)$ with $e>p^2$. There are groups $M,H$ over a field with residue field $\overline{\F}_p$, absolute index $e$, and $\mu_p\subset K$, with $M[p^3]\simeq H[p^3]$ and different finite special characters. The group $M$ acquires supersingular good reduction, and $H$ is filtered by two potentially good height-one groups.
\end{theorem}
\begin{proof}
Choose
\[
 \frac{e}{p^3}<A<\frac{e}{p^2},\qquad N\nmid A,
 \qquad d=eN,\quad b=\frac{eN}{p}-p^2(p-1)A,
 \quad J=b+p^3A.
\]
The same integer-choice argument as in the previous theorem applies, because $e/p^2>1$. Direct subtraction gives
\begin{align*}
 b-(p^2-p)A&=N(e-p^2A)/p>0,\\
 (p^2-p)p^2A-b&=N(p^3A-e)/p>0,\\
 d&=b+(p-1)J,\qquad A<p^2A<J<d,\\
 b&\equiv-(p-1)A\pmod N.
\end{align*}
Use $Y_0=X^A$, $Y_1=[p]_{X^b}(Y_0)$, $Y_2=[p]_{X^b}(Y_1)$, and $P=\Phi_{X^b}(Y_2)$. The first displayed inequality selects the degree-$p^2$ term at the first step; the second selects the degree-$p$ term at the next. Therefore
\[
               Y_1\equiv X^{p^2A}+O(X^{p^2A+1}),\qquad
               Y_2\equiv X^J+O(X^{J+1})\pmod p.
\]
The final selected term has degree $b+(p-1)J=d$. Its degree-$p^2$ competitor is later because $b<(p^2-p)J$. Intermediate terms contain $u$ modulo $p$ and are later than the selected terms; terms beyond degree $p^2$ are later by integrality. Hence $P(0)=p$ and $P\equiv X^d+O(X^{d+1})\pmod p$.

The Eisenstein ring $R=W[[X]]/(P)$ has $v_R(p)=d$. Since $J<d$, its marked point $x^A$ has successive valuations $A,p^2A,J$ and exact order $p^3$. The congruence for $b$ permits precisely the descent $x\mapsto\zeta x$, $T\mapsto\zeta^{-A}T$ used above. Its fixed field has absolute index $e$, contains $\mu_p$, and the descended group has special eigenvalues $a,a^p$ with $a=\zeta^{-A}\ne1$. Apply Lemma~\ref{lem:point-ambiguity}.
\end{proof}

\begin{corollary}\label{cor:sharp-mixed-wild}
Among positive multiples $e$ of $p(p-1)$ with $v_p(e)=1$, a comparison of the two classes in Theorem~\ref{thm:two-step-arc} exists with isomorphic $p^3$-torsion and different special characters if and only if $e>p^2$.
\end{corollary}
\begin{proof}
Existence was just proved. Proposition~\ref{prop:mixed-obstruction}, with $n=3$ and $v_p(e)=1$, gives the necessary strict inequality.
\end{proof}

\begin{theorem}\label{thm:second-wild-exact}
For $p\ge5$, in the height-two filtration setting over complete fields with perfect residue field and $\mu_p\subset K$,
\[
                         \nuMin(p,2p(p-1),2)=3.
\]
In particular $\nuMin(5,40,2)=3$.
\end{theorem}
\begin{proof}
Theorem~\ref{thm:two-step-arc} gives failure at $p^3$-torsion. At $e=2p(p-1)$ one may take $A=1$, $b=(p-1)(p^2-2)$ and $J=b+p^3$.

For the upper bound, take a $p^4$-torsion comparison and pass to the completed maximal unramified base. Each determinant has the form $\chi_p^d\epsilon$, with $d\in\{0,1,2\}$ and finite $\epsilon\in\mu_{p-1}$. The wild order of $\chi_p$ modulo $p^4$ is at least $p^2$: the intersection of $K$ with the cyclotomic tower has wild degree at most $p$, since $v_p(e)=1$. The ratio of the two determinants can thus be trivial only when their $d$ agree; their finite factors then agree as well.

For $d=0$, the graded pieces are potentially \'etale, so their special character is the Brauer lift of the semisimplified residual character. Cartier duality gives the case $d=2$. If $d=1$, a height-two potentially good piece is either supersingular or has its connected--\'etale filtration, which descends by uniqueness. Thus the possibilities are a whole supersingular group or a filtration by two height-one groups. Two filtered groups are compared residually, since $\chi_p\equiv1\pmod p$. A mixed comparison would require $e>p^{2\cdot4-3-1}=p^4$ by Proposition~\ref{prop:mixed-obstruction}, which is false.

For the remaining pair of supersingular groups, apply Theorem~\ref{thm:common-determinant} with $a=2p(p-1)(p^2-1)$. The recurrence gives
\[
 q_1=2p-1,\qquad q_2=2p^2-1,\qquad
 q_3=2p^3-1>2p(p^2-1)=a/(p-1).
\]
The $p^4$-torsion therefore determines the special character in every case. Since $p^3$-torsion can fail, the least index is three.
\end{proof}

\section{Quotients and first-level ramification}\label{sec:quotients}
We record the formal-group facts needed to pass between truncation levels. The residue field in this section is algebraically closed and $p\ge5$.

\begin{lemma}\label{lem:tame-good}
A height-two group with potentially supersingular reduction acquires good formal reduction over a tame cyclic extension of degree dividing $N=p^2-1$. Its finite special character has eigenvalues $\zeta^c,\zeta^{pc}$, and its generic determinant is $\chi_p\epsilon$, with $\epsilon$ the finite norm character.
\end{lemma}
\begin{proof}
The endomorphism algebra of the special height-two Honda group is the division algebra
\[
 D=\Q_{p^2}\langle\Pi\rangle,
 \qquad \Pi^2=p,\qquad \Pi a=\sigma(a)\Pi
\]
for $a$ in the unramified coefficient field and its nontrivial Frobenius $\sigma$. Its maximal order has residue field $\F_{p^2}$ and pro-$p$ reduction kernel. An element of order $p$ would generate the field $\Q_p(\zeta_p)$ of degree $p-1>2$, whereas every element of $D$ satisfies a reduced quadratic polynomial. Thus every finite subgroup of its units has order prime to $p$, by Cauchy's theorem, and reduction embeds it into $\F_{p^2}^{\times}$.

The kernel of the special inertia action may be removed from a good extension: the rational representation then has trivial inertial type and is crystalline, and the integral Barsotti--Tate classification applies to the original lattice, not merely to its isogeny class. This is the crystalline descent and lattice statement in \citet[Theorem~2.2.7 and Proposition~A.6]{Kisin2006}; Tate's full faithfulness identifies the resulting model with the original good model after extension \citep[Theorem~4, p.~180, and Corollary~2, p.~181]{Tate1967}. The faithful special quotient therefore supplies the asserted tame extension. The two special eigenvalues are the unramified conjugates in $\F_{p^2}$, and taking reduced norms gives the determinant assertion. Two such groups have a common good extension of degree dividing $N$, since tame inertia is procyclic.
\end{proof}

We may enlarge the good field to degree $N$ and put $s=(p+1)e$. This field also kills the tame quotient of the residual representation. Indeed, a nontrivial residual wild group has a unique invariant line, so its tame quotient acts through the two diagonal $\F_p$-characters; if the wild group is trivial, a tame generator is semisimple with eigenvalues in $\F_{p^2}$. In both cases its order divides $N$. Whenever $p-1\mid e$, the integer $s$ is a multiple of $N$. In a two-slope polygon the parameter $B$ of~\eqref{eq:two-slope-B} satisfies
\[
           0<B<pe,\quad B\in\Z,\quad B\equiv c\pmod{p+1}.
\]
The integrality follows because the invariant high line has trivial character over $L$; the congruence follows because its tame character $\zeta^{c-B}$ is $\F_p$-valued. Over $L$, the first-level image is either trivial or the upper unipotent group of order $p$.

\begin{lemma}\label{lem:actual-quotients}
The quotient by the high line has parameter $pB$ if $B<e$ and $s-B$ if $B>e$; at $B=e$ its first-level polygon has one slope and every nonzero point has valuation $e$. If the first-level module is rational and $J$ is an inertia-stable low line, the quotient by $J$ has parameter $B/p$. A one-slope group has first-level value $e$, and its quotient by any stable line has parameter $e$, high value $pe$, and low value $e/p$.

These are good formal quotients over $L$ with the same finite special character. A comparison at level $p^n$ preserving the chosen lines induces a comparison of their quotients at level $p^{n-1}$.
\end{lemma}
\begin{proof}
The finite flat closure of a stable line $J$ defines an isogeny whose coordinate is
\[
                T\longmapsto U(T)\prod_{Q\in J}(T-T(Q)),
                 \qquad U(T)\in\Ocal_L[[T]]^\times.
\]
For the high line, every argument of value below $\alpha=s-B$ has its value multiplied by $p$. The image of the old first-level quotient has value $B$. If $B>e$, the inverse polygon over a nonzero high point has all $q$ roots of value $\alpha/q$; their quotient images give $q-p$ points of value $\alpha/p$, with $p-1$ old points of value $B$. This is parameter $s-B$. If $B<e$, the inverse roots have values $(s-pB)/p$ with multiplicity $p$, and $B/p$ with multiplicity $q-p$. Quotienting gives $p-1$ high points of value $s-pB$ and $q-p$ low points of value $B$, hence parameter $pB$. At equality the inverse polygon is a single segment and all quotient values are $e$.

For a low line, each nonzero old quotient class has a unique high representative. Its image has value $\alpha+(p-1)\beta=s-B/p$. The remaining images come from low inverse points of value $\beta/q$ and have value $\beta/p=B/p^2$. This gives parameter $B/p$. In the one-slope case every nonzero first-level point has value $e$, by the product formula. Each nonzero old quotient class has $p$ differences of value $e$, giving value $pe$; inverse points above the chosen line have value $e/q$ and quotient value $e/p$.

The special line is the Frobenius kernel, so the special quotient has tangent exponent $pc$ and the same unordered character $\{c,pc\}$. Tate full faithfulness extends the generic semilinear descent maps to the good models and preserves the finite flat closure, hence gives generic descent of the quotient. Finally, as actual finite Galois modules,
\[
 (M/J)[p^{n-1}]
       =\{P\in M[p^n]:p^{n-1}P\in J\}/J.
\]
This identity proves the precise truncation assertion.
\end{proof}

\begin{lemma}\label{lem:prime-residual-comparison}
In a two-slope comparison at level $p^2$ with nontrivial residual wild action, if at least one parameter $B_i$ is prime to $p$, then both parameters are equal and the special characters agree.
\end{lemma}
\begin{proof}
The common wild fixed line is the high line. The common first-level field $R/L$ has degree $p$. A low coordinate has value $B_i$ over $R$ and displacement $p(s-B_i)$. If $p\nmid B_1$, the exact prime-valuation formula gives
\[
                   b(R/L)=ps-(p+1)B_1.
\]
The unconditional inequality in the other coordinate gives $B_1\ge B_2$. If both are prime, symmetry gives equality. If $p\mid B_2$, equality is impossible. A common second-level lift of a low point has value $B_1/p^2$ over $R$; its denominator forces its point field to have degree $p^2$. Every first-level translation occurs, and the two normalized coordinates have values $B_1,B_2$. A low translation has displacement $p^2B_i$. The first model fixes its break as $(p^2-1)B_1$, whereas the second gives $p^2B_2\ge B_2+(p^2-1)B_1$, contradicting $B_1>B_2$. Equal parameters and the identified high-line character identify $c$ modulo $N$.
\end{proof}

\section{Irreducible residual inertia at the first wild index}\label{sec:irreducible}
The comparison with irreducible residual inertia can be settled by the ramification of one point at a time. This keeps the field degrees separate from the dimensions of congruence layers.

\begin{theorem}\label{thm:irreducible-recovery}
Let $p\ge5$, let $K$ have algebraically closed residue field, $e_K=p(p-1)$, and $\mu_p\subset K$. If two height-two groups with potentially supersingular good reduction have isomorphic $p^3$-torsion and their common residual inertia representation is irreducible over $\F_p$, then their finite special inertia characters agree.
\end{theorem}

Put $q=p^2$, $N=q-1$, and $e=p(p-1)$. Take the common tame good field $L/K$ of degree $N$. We argue by contradiction, assuming that the special characters differ.

\subsection{Leading coordinates and common point fields}
The residual wild image is a normal $p$-subgroup of the residual inertia image. Its nonzero fixed space is stable under inertia, so irreducibility makes that subgroup trivial. The residual representation is tame of degree two; its order divides $N$, and its whole first-level module is rational over $L$. A two-slope Newton polygon would give a canonical invariant high line. The polygon therefore has one slope. In the Honda coordinate its nonzero $p^j$-torsion valuations over $L$ are
\begin{equation}\label{eq:irreducible-values}
                              \frac e{q^{j-1}}.
\end{equation}
The first value is $v_L(p)/N=e$; every inverse polygon has a single segment and divides the value by $q$.

Let $c_i$ be the special tangent exponents. The first-level leading-coordinate spaces have tame exponents $c_i+e$. Each is a scalar multiple of $\F_q$: the leading-coordinate map is injective on the two-dimensional residual module, and its irreducible tame scalar generates $\F_q$ over $\F_p$. An equivariant isomorphism between two such spaces has, after scalar normalization, either the form $z\mapsto z$ or the form $z\mapsto z^p$. The scalar form gives equal special exponents. Under our contradiction hypothesis the form must therefore be
\begin{equation}\label{eq:frobenius-orientation}
                                  z_2=z_1^p.
\end{equation}
It remains true after changing the uniformizer of a point field, which only rescales the two leading spaces. At the boundary where the coefficient of $T^p$ could lie on the first Newton segment, the additive initial polynomial would contain a nonzero $T^p$ term. It cannot vanish on a scalar multiple of all of $\F_q$. Thus this boundary is excluded by irreducibility; the later leading multiplication terms are uniquely the degree-$q$ terms.

We use the following elementary ramification calculation. Let $T/S$ be a totally ramified Galois extension in mixed characteristic, with uniformizer $\varpi$. Put
\[
 b(g)=v_T(g\varpi-\varpi)-1.
\]
Expand a coordinate $x=\sum_{j\ge A}[a_j]\varpi^j$ using residue representatives fixed by $g$. Then $v_T(gx-x)\ge A+b(g)$, with equality when $p\nmid A$. Below $v_T(p)$, if $v_p(A)=1$, the first possible terms of the change have degrees
\begin{equation}\label{eq:first-change}
                         A+pb(g),\qquad J+b(g),
\end{equation}
where $J$ is the first nonzero exponent prime to $p$, possibly infinite. To verify this, write $g\varpi=\varpi(1+d\varpi^b+\cdots)$ and expand each monomial. For $j=p^tu$, $p\nmid u$, the first term below $v_T(p)$ is $u[a_j]d^{p^t}\varpi^{j+p^tb}$. Among the exponents divisible by $p$, the leading exponent $A$ is earliest; among the other exponents, $J$ is earliest. If the two displayed degrees differ, the smaller cannot cancel. If they agree, their leading coefficients must be included together. This proves every use of the calculation below.

Write
\[
 F_2=L(M_i[p^2]),\qquad F_3=L(M_i[p^3]).
\]
The torsion isomorphism makes these the same fields for $i=1,2$. Over $L$ the matrices are congruent to $1$ modulo $p$. Let $W_1,W_2\subset M_2(\F_p)$ be the first and terminal congruence layers, and write $r_j=\dim W_j$. Matrix multiplication gives
\begin{equation}\label{eq:matrix-powers}
 (1+pA+p^2B)^p=1+p^2A\pmod{p^3},\qquad W_1\subset W_2.
\end{equation}

\begin{lemma}\label{lem:full-point-orbits}
For every primitive $p^2$-point $P$ and every $p^3$-point $P^+$ with $pP^+=P$, the common point fields $E=L(P)$ and $T=L(P^+)$ satisfy
\[
                [E:L]=q,\qquad[T:E]=q,\qquad[T:L]=q^2.
\]
The two successive extensions are Galois, with the full $p$-torsion translation group. Moreover $W_1v=\F_p^2$ for every $v\ne0$.
\end{lemma}
\begin{proof}
Since $pP$ and every $p$-torsion point are rational over $L$, all conjugates of $P$ lie in $L(P)$. Thus $E/L$ is Galois, of degree at most $q$. The denominator of $v_L(P)=(p-1)/p$ makes its degree at least $p$. Similarly $T/E$ is Galois of degree at most $q$.

Suppose $[E:L]=p$. The denominator of $v_L(P^+)=(p-1)/p^3$ then forces $[T:L]=p^3$ and $[T:E]=q$. In $T$ the two coordinates $x_i$ have valuation $A=p-1$, prime to $p$, and every nonidentity translation displaces them by a nonzero $p$-torsion coordinate of valuation $D=ep^3$. The exact prime-valuation calculation gives one lower break $D-A$ and leading differences
\[
       \overline{(gx_i-x_i)/\varpi^D}=\overline A\,a_i d_g,
       \qquad a_i=\overline{x_i/\varpi^A}\ne0.
\]
These are below $v_T(p)=ND$. The two leading translation maps are scalar multiples of the same $d_g$, although their values run through the whole two-dimensional space and are related by~\eqref{eq:frobenius-orientation}. This would make Frobenius on $\F_q$ scalar, which it is not. Hence $[E:L]=q$. The orbit is $W_1v$, proving its asserted surjectivity. Equation~\eqref{eq:matrix-powers} then makes terminal matrices supply every $p$-torsion translation of $P^+$. Thus $[T:E]=q$ as well.
\end{proof}

\begin{lemma}\label{lem:three-first-dimensions}
One has $r_1\ge3$.
\end{lemma}
\begin{proof}
If $r_1=2$, evaluation at any nonzero vector is an isomorphism by Lemma~\ref{lem:full-point-orbits}. The stabilizer of $P^+$ is then contained in the terminal matrix subgroup, which is central modulo $p^3$. Thus $T/L$ is Galois. Put $G=\Gal(T/L)$ and $U=\Gal(T/E)$. Then
\[
               |G|=p^4,\qquad U\simeq\F_p^2,
                    \qquad G/U\simeq\F_p^2,
\]
and the $p$th-power map $G\to U$ is onto by~\eqref{eq:matrix-powers}.

In $T$ the three relevant coordinate valuations are
\[
 A=e,\qquad D_1=qA,\qquad D_2=q^2A,\qquad v_T(p)=ND_2.
\]
Elements outside $U$ displace the selected point by order-$p^2$ points, of valuation $D_1$; nonzero elements of $U$ displace it by order-$p$ points, of valuation $D_2$. Normalize each leading space using the selected points $pP^+$ and $p^2P^+$ as parameter $1$. If the primitive leading coefficient is $a_i$, the displacement coefficients have the forms
\[
            a_i^q\lambda_i(g),\qquad a_i^{q^2}\lambda_i(h),
                        \qquad\lambda_2=\lambda_1^p.
\]
The parameter maps on $G/U$ and $U$ are additive isomorphisms onto $\F_q$. The relation on the second-level leading spaces follows from the first-level relation by taking the unique $q$th roots.

Let $a$ be the smallest lower break of $G$. Every nonzero element of $U$ has larger break, since it is a $p$th power and the first additive ramification coefficient of a $p$th power is zero. Set
\[
 a_0=(D_1-A)/p=N(p-1),\qquad
 j_0=A+(p-1)a_0=D_1-a_0,\qquad j_0\equiv-1\pmod p.
\]
Consider an element outside $U$ of break $a$ in the two expansions~\eqref{eq:first-change}. If $a<a_0$, cancellation below $D_1$ forces the first ordinary degree to be $J_i=A+(p-1)a$, with nonzero coefficient. If $a>a_0$, achieving degree $D_1$ forces $J_i=D_1-a$, again with nonzero coefficient. An exponent divisible by $p$ in either assertion is already impossible. At every terminal break $b>a$, that fixed ordinary term is uniquely earliest: the differences from the divisible term are respectively $(p-1)(b-a)$ and $(p-1)(b-a)+p(a-a_0)$. Consequently $J_i+b=D_2$, and both terminal translation maps become scalar multiples of the common $d_h$. This contradicts their Frobenius relation. Therefore $a=a_0$.

No ordinary degree below $j_0$ can now occur. Write $f_i$ for its coefficient, allowing zero. At degree $D_1$, for a minimum-break element, the leading equality is
\begin{equation}\label{eq:boundary-leading}
 a_i^q\lambda_i(g)=(A/p)a_i d_g^p+j_0f_i d_g.
\end{equation}
If one $f_i\ne0$, every terminal break equals $b_*=D_2-j_0>a_0$, and this ordinary term is uniquely first. In the other model the same break also forces a nonzero term at exactly $j_0$, because
\[
 b_*-(D_2-A)/p=(p-1)(D_2-D_1)/p>0.
\]
We again obtain scalar proportionality on $U$, a contradiction. Thus both $f_i$ vanish. Every nonterminal element must have break $a_0$, since a larger break would put every coordinate change after $D_1$. The first ramification coefficient identifies $G/U$ with a two-dimensional additive space, and~\eqref{eq:boundary-leading} becomes
\[
                       \lambda_i(g)=(A/p)a_i^{1-q}d_g^p.
\]
The two maps are scalar proportional on all of $\F_q$, again contradicting their Frobenius relation. Hence $r_1=2$ is impossible.
\end{proof}

\subsection{The two point conductors}
\begin{lemma}\label{lem:point-conductors}
The last upper breaks of $E/L$ and $T/E$ are respectively
\[
             c_1=N(p-1),\qquad c_2=(q^2-1)(p-1).
\]
The last upper break of $F_2/L$ is $c_1$, and that of $F_3/F_2$ is at most $p^{r_1-2}c_2$.
\end{lemma}
\begin{proof}
In either successive point extension the group is $\F_p^2$, the coordinate valuation is $A=e$, and every nonidentity element has displacement $D$, with $D=qA$ or $q^2A$. The two leading translation maps fill $\F_q$ and have the Frobenius relation. Put $b_0=(D-A)/p$.

If there is just one lower break $b<b_0$, cancellation in~\eqref{eq:first-change} forces $J=A+(p-1)b$ and $d_g^{p-1}=\kappa\ne0$ on all $q-1$ nonzero leading ramification coefficients. This is impossible for a polynomial of degree $p-1$. If $b>b_0$, the unique ordinary term of degree $D-b$ makes the leading maps scalar proportional, also impossible. Hence a single break equals $b_0$.

If there are two lower breaks $t<u$, then $t<b_0$. For if $t\ge b_0$, no ordinary term can precede degree $D$ at $t$, and every term would follow $D$ at $u$. Cancellation at $t$ therefore forces $J=A+(p-1)t$. At $u$ it is uniquely earliest, so
\[
                               u+(p-1)t=D-A.
\]
Herbrand's formula makes the last upper break $(u+(p-1)t)/p=b_0$. This proves the two stated values.

Choose a basis $P_1,P_2$ of $p^2$-torsion. Then $F_2=L(P_1)L(P_2)$. Both elementary abelian point fields have last break $c_1$, so their compositum has that last break: every cyclic character is a sum of characters coming from the factors, and conductors do not increase under addition.

Lift the basis to $P_i^+$ and set $E_i=L(P_i)$, $T_i=L(P_i^+)$ and $D_i=F_2T_i$. The terminal translations are already the full translations of $T_i/E_i$, giving
\[
 F_2\cap T_i=E_i,\quad[D_i:F_2]=q,\quad
 [F_2:E_i]=p^{r_1-2},\quad D_1D_2=F_3.
\]
For completeness let $H_i$ and $S_i$ be the stabilizers of $P_i$ and $P_i^+$ in $\Gal(F_3/L)$, and let $U$ be the terminal subgroup. Surjectivity of terminal translations says $US_i=H_i$, which proves the intersection. The subgroup $U\cap S_i$ is central. For first-layer matrices $A,C$ of $H_i$, both kill the residual vector $v_i$, so $[A,C]v_i=0$; their $p$th powers also kill it. Hence $D_i/E_i$ is an elementary abelian Galois extension.

For upper index $u>c_2$, its ramification group has trivial image in $\Gal(T_i/E_i)$. Its intersection with $\Gal(D_i/F_2)$ is therefore trivial, since $D_i=F_2T_i$. The subgroup form of Herbrand's theorem bounds the last break of $D_i/F_2$ by $\psi_{F_2/E_i}(c_2)\le p^{r_1-2}c_2$. The two fields $D_i$ generate the elementary abelian terminal extension, proving the same bound for $F_3/F_2$. These are the quotient and subgroup conversions of \citet[Chapter~IV, \S3, Propositions~14--15, pp.~75--76]{Serre1979}.
\end{proof}

\subsection{Contradiction from the determinant}
We finish the proof of Theorem~\ref{thm:irreducible-recovery}. On the good field $L$, the determinant is cyclotomic. First suppose that the trace on $W_1$ is zero. Lemma~\ref{lem:three-first-dimensions} gives $W_1=\mathfrak{sl}_2$ and $r_1=3$. Its determinant is trivial modulo $p^2$, so $\mu_{p^2}\subset L$.

Let $K_0$ be the Cohen field of $K$. The extension $K_0(\mu_{p^2})/K_0(\mu_p)$ has degree $p$. Since $K/K_0(\mu_p)$ also has degree $p$ and $L/K$ is tame, adjoining $\mu_{p^2}$ to $K$ cannot have degree $p$. Thus $K=K_0(\mu_{p^2})$. The cyclotomic extension $C=L(\mu_{p^3})$ has degree $p$ over $L$ and break $N^2$: before the tame base change its break is $p^2-1=N$, and the tame degree is $N$.

Because $N^2>c_1$, $C\cap F_2=L$. In the elementary abelian compositum its cyclotomic character has the unique largest break. After restriction to $F_2$ its conductor is exactly
\[
 \psi_{F_2/L}(N^2)
       =\psi_{F_2/L}(c_1)+p^3(N^2-c_1)>p^3(N^2-c_1).
\]
It is a nonzero terminal character, since $C\subset F_3$. Yet Lemma~\ref{lem:point-conductors} bounds it by $pc_2$, contradicting
\[
 p^3(N^2-c_1)-pc_2=pN(p-1)(p^3-p^2-1)>0.
\]
The trace on $W_1$ is consequently nonzero.

It follows that $\mu_{p^2}\not\subset K$: otherwise the first congruence determinant over $L$ would be trivial. Apply Theorem~\ref{thm:cyclotomic-exact} to $K(\mu_{p^3})/K$, with last upper break $\delta_2$. It gives
\[
 d_2=N\delta_2\ge Np(p-1)>c_1.
\]
The intersection of the cyclic degree-$p^2$ field $L(\mu_{p^3})$ with the elementary abelian field $F_2$ is exactly $L(\mu_{p^2})$. Its nontrivial terminal character consequently has conductor
\[
 \psi_{F_2/L}(d_2)
     =\psi_{F_2/L}(c_1)+p^{r_1}(d_2-c_1)
     >p^{r_1}(d_2-c_1).
\]
One may verify the equality directly from upper ramification groups: above $c_1$ the quotient $F_2/L$ is trivial, while the cyclotomic subgroup remains nontrivial up to $d_2$; the subgroup conversion is therefore exact. But
\[
 q(d_2-c_1)-c_2\ge N(p-1)(p^3-2p^2-1)>0.
\]
This contradicts the terminal bound $p^{r_1-2}c_2$ once more. The assumption of different characters was impossible, proving the theorem.\hfill$\square$

\section{Tame reducible residual inertia at the first wild index}\label{sec:tame-reducible}
This section recovers the character at $p=5$ and bounds the possible wild image of an unequal-character comparison at the general first wild index. The first reduction identifies the canonical lines from a small common point field; two opposite quotient actions, diagonal visibility, and a cyclotomic quotient obstruction force seven matrix dimensions when $p\equiv1\pmod4$.

\begin{proposition}\label{prop:square-determinant}
At $e=p(p-1)$, $p\ge5$, a $p^3$-torsion comparison between two potentially supersingular height-two groups with a common square finite determinant determines their special characters.
\end{proposition}
\begin{proof}
Twist both groups by the inverse of a square root of their finite determinant. The finite special inertia is now norm one in $\F_{p^2}^\times$, so its good degree $m$ divides $p+1$. The good field has absolute index $a=em\le p(p^2-1)$.

The proof of Lemma~\ref{lem:weighted-kernel} without the tame weight congruence gives the recurrence
\[
                      q_0=1,\qquad q_{j+1}=\min\{a+q_j,1+pq_j\}.
\]
Indeed a degree-$p$ coefficient is still divisible by the uniformizer, and the rest of the logarithm and integral Cartier argument is unchanged. Here $q_1=p+1$ and
\[
                q_2=\min\{a+p+1,p^2+p+1\}\ge a/(p-1).
\]
Apply the graph-closure proof of Theorem~\ref{thm:common-determinant}. If the special characters differed, its target would have product-height-two special fibre, and its special kernel would be $\ker F^3$. The unweighted bound would kill that kernel by $F^2$, which is impossible.
\end{proof}

We recall the two-slope notation over the degree-$N$ good field. Put $s=(p+1)e$. The high line $H$ and the low points have valuations
\begin{equation}\label{eq:two-slope-B}
        \alpha=s-B,\qquad\beta=B/p,\qquad0<B<pe.
\end{equation}
The high line is rational over $L$, so $B$ is an integer; its tame character is $\zeta^{c-B}$, where $c$ is the special tangent exponent. These formulas follow from the vertices $(1,eN),(p,(p-1)B),(p^2,0)$ of the Newton polygon. The high-line character follows by multiplying the cotangent factor $\zeta^c$ by the leading coefficient factor $\zeta^\alpha$, noting that $s$ is divisible by $N$ when $e=p(p-1)$.

\begin{lemma}\label{lem:one-prime-parameter}
Suppose that two two-slope groups have rational $p$-torsion over $L$ and isomorphic $p^2$-torsion. Write $B_i=pb_i$, $0<b_i<e$. If $p\nmid b_1$, the high lines agree, $b_1=b_2$, and the special characters agree.
\end{lemma}
\begin{proof}
Take a residual vector outside both high lines and corresponding second-level lifts. Their valuations are $b_i/q$. Their common point field is Galois of degree at most $q$, since its conjugates are rational first-level translates. The denominator of $b_1/q$ forces degree $q$, so every translation occurs.

In this field the coordinate values are $b_i$. For every nonidentity $g$, the first model gives the exact break equality $v(gx_1-x_1)=b_1+b(g)$; the second gives the inequality $v(gx_2-x_2)\ge b_2+b(g)$. If the high lines differed, a translation high in the first and low in the second would imply
\[
                      qb_2\ge b_2+q(s-pb_1)-b_1.
\]
The difference of the left side and right side is increasing in both $b_i$, and at $b_1=b_2=e-1$ it is $-q(p+1)<0$. This is impossible. With the high lines identified, a low translation gives $(q-1)(b_2-b_1)\ge0$, and a high translation gives $(p^3+1)(b_1-b_2)\ge0$. Thus $b_1=b_2$. The equal high-line characters then give equal $c_i$ modulo $N$.
\end{proof}

\begin{corollary}\label{cor:first-wild-tame-pattern}
Let $p\ge5$, $e_K=p(p-1)$ and $\mu_p\subset K$. Suppose two potentially supersingular height-two groups have isomorphic $p^3$-torsion, tame reducible residual inertia and different special characters. Over the common tame good field of degree $p^2-1$, both first-level polygons have two slopes and parameters
\[
                  B_i=p^2k_i,\qquad1\le k_i\le p-2.
\]
Their high lines are the two distinct residual tame eigenspaces and are exchanged under the comparison.
\end{corollary}
\begin{proof}
Proposition~\ref{prop:square-determinant} excludes a square finite determinant. A reducible one-slope leading-coordinate action is scalar, so its determinant is square. Thus both models have two slopes. Their residual modules are rational over $L$, and Lemma~\ref{lem:one-prime-parameter} implies that both parameters $b_i=B_i/p$ are divisible by $p$. The strict range $0<B_i<pe$ gives the displayed bounds on $k_i$.

If the high lines coincided, equality of their characters would give $c_1-c_2\equiv B_1-B_2\pmod{p^2-1}$. Equality of the finite norm characters gives $c_1\equiv c_2\pmod{p-1}$. Hence $k_1\equiv k_2\pmod{p-1}$, and their range forces equality; the high-line character then gives $c_1=c_2\pmod{p^2-1}$, a contradiction. The two lines are distinct. A scalar residual action would again have square determinant, so these are the two distinct tame eigenspaces.
\end{proof}

For the next three lemmas, $K$ is complete discretely valued of mixed characteristic $(0,p)$ with algebraically closed residue field, $p\ge5$, $\mu_p\subset K$, and $e_K=p(p-1)$. The field $L/K$ is the tame extension of degree $N=p^2-1$. For an actual comparison whose $p$-torsion is rational over $L$, write $W\subset1+pM_2(\Z/p^3)$ for its common third-level image. Its first space $W_1$ is the image of $1+pA\mapsto A\pmod p$, and its terminal space is $W_2=\{B:1+p^2B\in W\}$.

\begin{lemma}\label{lem:first-layer-diagonal-visibility}
Let $p\ge5$ and $e_K=p(p-1)$ in the stated base setting. Suppose two whole potentially supersingular height-two groups have isomorphic actual $p^3$-torsion, rational $p$-torsion over $L$, parameters
\[
 B_i=p^2k_i,\qquad1\le k_i\le p-2,
 \qquad H_1=\F_pe_1,\quad H_2=\F_pe_2.
\]
For the common actual first congruence space $W_1$, both functions $A\mapsto A_{11}$ and $A\mapsto A_{22}$ are nonzero. No bound on the full division-field degree is assumed.
\end{lemma}
\begin{proof}
Suppose $A_{11}=0$ for every $A\in W_1$, and choose a second-level point $Q$ with $pQ=e_1$. Its orbit consists of translations by $W_1e_1\subset\F_pe_2$, so its point field has degree at most $p$. In the second model its coordinate has valuation $k_2/p$, whose denominator forces degree exactly $p$. Its stabilizer is the kernel of the column map on the abelian first-layer image; hence the field is Galois and its nonidentity automorphisms translate by nonzero multiples of $e_2$.

Since $B_i=p^2k_i>e$, a division of the high point in the first model has valuation $(N-pk_1)/p$, while the second model gives $k_2/p$. On the common degree-$p$ field both normalized values, $N-pk_1$ and $k_2$, are prime to $p$. The translating point is low in the first model and high in the second. The exact prime-valuation break formula therefore assigns the same automorphism the two breaks
\[
 p(p+1)k_1-N,\qquad p^2N-(p^3+1)k_2.
\]
Equality would require
\[
 p(p+1)k_1+(p^3+1)k_2=p^4-1.
\]
The left side is at most $(p-2)(p^3+p^2+p+1)$, which is smaller than $p^4-1$ by $(p+1)(p^2+1)$. This is impossible. Interchanging the models and axes proves that the $22$-projection is also nonzero.
\end{proof}

\begin{lemma}\label{lem:first-wild-cyclotomic-break}
Let $p\ge5$, $e_K=p(p-1)$, $\mu_p\subset K$, and $[L:K]=N=p^2-1$ as above. If $\mu_{p^2}\subset L$, then $\mu_{p^2}\subset K$, the extension $C=L(\mu_{p^3})/L$ is cyclic of degree $p$, and every nonidentity element has lower break $N^2$.
\end{lemma}
\begin{proof}
The degree of $K(\mu_{p^2})/K$ divides both $p$ and the prime-to-$p$ integer $N$, so this extension is trivial.
Choose a Cohen coefficient ring $\Lambda\subset\mathcal O_K$ with the same perfect residue field, using \citet[Theorem~10.160.8, Tag~032A]{StacksCohen}, and put $K_0=\operatorname{Frac}\Lambda$. The shifted polynomial $\Phi_{p^2}(1+T)$ is Eisenstein of degree $p(p-1)=e_K$. Its root $\zeta_{p^2}-1$ is a uniformizer of $K$, and successive expansion in that uniformizer with coefficients in $\Lambda$ gives $K=K_0(\zeta_{p^2})$. The analogous $p^3$-cyclotomic polynomial gives a cyclic degree-$p$ extension $K_3/K$ with uniformizer $\pi_3=\zeta_{p^3}-1$. Every nonidentity automorphism satisfies
\[
 v_{K_3}(\sigma\pi_3-\pi_3)=v_{K_3}(\zeta_p-1)=p^2,
\]
so its lower break is $N$. The compositum $C=LK_3$ has degree $p$ over $L$. In its normalized valuation, $v_C(\pi_3)=N$ and $v_C(\pi_L)=p$. Choose integers $a,b$ with $aN+bp=1$. Then $\pi_3^a\pi_L^b$ is a uniformizer, $p\nmid a$, and its nonidentity displacement ratio has valuation $N^2$. Thus
\begin{equation}\label{eq:general-cyclotomic-last-break}
                   b(C/L)=N^2.
\end{equation}

\end{proof}

\begin{lemma}\label{lem:traceless-first-wild-exclusion}
Under the hypotheses of Lemma~\ref{lem:first-layer-diagonal-visibility}, assume the residual inertia is tame reducible, the finite determinant is nonsquare, and the ratio of the two residual eigencharacters differs from one. The two actual congruence layers cannot both be $\mathfrak{sl}_2(\F_p)$ in the high-axis basis. This assertion holds for every $p\ge5$.
\end{lemma}
\begin{proof}
Suppose $W_1=W_2=\mathfrak{sl}_2(\F_p)$. On the good field the determinant is cyclotomic. Its first layer is traceless, so $\mu_{p^2}\subset L$. Lemma~\ref{lem:first-wild-cyclotomic-break} gives a nontrivial degree-$p$ cyclotomic quotient $C/L$ with lower break $N^2$.
On the terminal layer, $\det(1+p^2B)=1+p^2\operatorname{tr}B=1$. The nonzero cyclotomic character modulo $p^3$ consequently descends to a nonzero linear functional $\ell$ on $W_1=\mathfrak{sl}_2(\F_p)$. It is invariant under tame conjugation, so it vanishes on the nontrivial off-diagonal weights and is a nonzero multiple of $A_{11}$.

Take a second-level point $Q$ with $pQ=e_1$ and put $D=L(Q)$. Evaluation of $\mathfrak{sl}_2$ on $e_1$ is onto $\F_p^2$, so $D/L$ is Galois of degree $p^2$. Its point stabilizer maps to $\F_pE_{12}$ in $W_1$; the terminal layer and this off-diagonal line are both killed by the determinant. Thus the entire point stabilizer is killed by $\ell$, and $C\subset D$. In the second model $Q$ has normalized $D$-valuation $pk_2$. Every lift of a nonidentity element of $\operatorname{Gal}(C/L)$ translates $Q$ by a point with nonzero $e_1$ coefficient, which is low in that model. Its displacement has valuation $p^3k_2$. The general ramification inequality gives
\[
          b_D(g)\le p^3k_2-pk_2=pNk_2\le pN(p-2)<N^2.
\]
For a nonidentity quotient element, the norm-polynomial formula averages the indices $b_D(g)+1$ over its fibre; see \citet[Proposition~3.7 and its full proof, pp.~14--15]{ColmezLocal}. It therefore gives $b_C<N^2$, contradicting~\eqref{eq:general-cyclotomic-last-break}. This formula assumes characteristic zero and perfect residue, and applies here without a finite-residue hypothesis.
The contradiction proves the lemma.
\end{proof}

\begin{proposition}\label{prop:opposite-quotient-image}
Under the hypotheses of Corollary~\ref{cor:first-wild-tame-pattern}, suppose $p\equiv1\pmod4$. The wild image over $L$ on the $p^3$-torsion has order at least $p^7$.
\end{proposition}
\begin{proof}

First quotient both groups by either of the two common stable residual lines. These quotients have isomorphic $p^2$-torsion: for a line $J$ the relevant group is
\[
 (M/J)[p^2]=\{x\in M[p^3]:p^2x\in J\}/J.
\]
One quotient is canonical and the other noncanonical. Since $B_i=p^2k_i>e$, their parameters are respectively
\[
       p(p^2-1-pk_i),\qquad pk_j.
\]
These are the quotient formulas of Proposition~\ref{rc-prop:quotients}, which do not require $p\nmid e$. After division by $p$, both integers are prime to $p$, and
\[
               p^2-1-pk_i\ge2p-1>p-2\ge k_j.
\]
If the common residual quotient representation were trivial over $L$, Lemma~\ref{lem:one-prime-parameter} would equate these unequal integers. Its image is a $p$-group, being a subquotient of the original third-level image over $L$, so nontriviality means nontrivial wild inertia. The same conclusion holds for the opposite quotient.

Write $P$ for the original wild image, and choose a basis on the two residual tame lines. The first-level module is rational over $L$, so $P\subset1+pM_2(\Z/p^3)$. Its first and terminal matrix spaces $W_1,W_2\subset M_2(\F_p)$ satisfy
\[
 W_1\subset W_2,\qquad[W_1,W_1]\subset W_2,
 \qquad |P|=p^{\dim W_1+\dim W_2}.
\]
Indeed $(1+pA+p^2B)^p\equiv1+p^2A\pmod{p^3}$, and a group commutator reduces to $1+p^2[A,B]$.

The quotient by the first residual axis replaces a Tate lattice $T$ by $T+\Z_p(e_1/p)$. In the basis $(e_1/p,e_2)$ a matrix $\rho$ becomes
\[
 \begin{pmatrix}\rho_{11}&p\rho_{12}\\
                 \rho_{21}/p&\rho_{22}\end{pmatrix}.
\]
For $\rho=1+pA+p^2B$, its residual lower coefficient is $A_{21}$. The nontrivial residual wild action of this quotient makes the $21$-projection of $W_1$ nonzero. The opposite quotient makes its $12$-projection nonzero.

Let $a,d\in\F_p^*$ be the residual tame eigenvalues. Their product is the nonsquare finite determinant, because $\chi_p\equiv1\pmod p$. The ratio $r=a/d=a^2/(ad)$ is nonsquare. Since $p\equiv1\pmod4$, both $1$ and $-1$ are squares, and the three adjoint weights $1,r,r^{-1}$ are distinct. Tame spectral projectors preserve $W_1$ and therefore put both $E_{12}$ and $E_{21}$ in it. They also belong to $W_2$, as does their commutator $E_{11}-E_{22}$. These three matrices are independent. Consequently $\dim W_1\ge2$ and $\dim W_2\ge3$. This is the first, five-dimensional bound.

Lemma~\ref{lem:first-layer-diagonal-visibility} makes both diagonal projections nonzero. Tame projectors isolate a nonzero diagonal direction, independent of $E_{12},E_{21}$. Thus $\dim W_1\ge3$, and the power inclusion gives $\dim W_2\ge3$. If their total dimension were six, they would coincide. The two off-diagonal units and their commutator would then give
\[
                    W_1=W_2=\mathfrak{sl}_2(\F_p),
\]
contrary to Lemma~\ref{lem:traceless-first-wild-exclusion}. The total dimension is at least seven, proving the proposition.
\end{proof}

The congruence condition separates the opposite tame weights. When $p\equiv3\pmod4$, a nonsquare ratio can be $-1$; two nonzero off-diagonal projections then need not supply the two individual matrix units. The proposition makes no assertion in that case and constructs no unequal-character pair.

\begin{lemma}\label{lem:opposite-line-conductor}
Let $p$ be a prime, and let $T/S$ be a totally ramified Galois extension of complete discretely valued mixed-characteristic fields with the same perfect residue field, group $\F_p^2$, and two lower breaks $t<u$, with last line $V$. Suppose a coordinate $x$ has $v_p(A)=1$, where $A=v_T(x)$, and displacement $D_1$ on a line $H\setminus\{0\}$ and $D_0$ off it, with $A<D_0<D_1<v_T(p)$. If $H\ne V$, the last upper break is $(D_0-A)/p$.
\end{lemma}
\begin{proof}
Use the expansion~\eqref{eq:first-change}. A nonzero element of $H$ has break $t$ and displacement $D_1$, while a nonzero element of $V$ has break $u$ and displacement $D_0$. Both possible first-change degrees increase strictly with the break. Thus the first change at $t$ would occur before $D_0$ unless it cancels. Cancellation forces a nonzero ordinary term of degree $J=A+(p-1)t$. At $u$ that term is uniquely first, giving $J+u=D_0$. The last upper break is therefore $(u+(p-1)t)/p=(D_0-A)/p$.
\end{proof}

\begin{theorem}\label{thm:tame-five-recovery}
Let $K$ have algebraically closed residue field of characteristic five, $e_K=20$, and $\mu_5\subset K$. Two potentially supersingular height-two groups with isomorphic $125$-torsion and tame reducible residual inertia have equal finite special characters.
\end{theorem}
\begin{proof}
Assume otherwise. The determinant comparison fixes the common finite determinant; a square is excluded by Proposition~\ref{prop:square-determinant}. A common finite twist makes its value a specified primitive fourth root. Thus the special tangent exponents satisfy $c_i\equiv1\pmod4$ in $\Z/24$.

A one-slope polygon has a scalar leading-coordinate action on every invariant residual line. If it is reducible, the tame scalar belongs to $\F_5$, and its determinant is square. Thus both polygons have two slopes. Their first-level modules are rational over the common tame good field $L/K$ of degree 24. Lemma~\ref{lem:one-prime-parameter} forces $B_i=25k_i$, with $1\le k_i\le3$. If the high lines agreed, their common characters and determinant would give $k_1\equiv k_2\pmod4$, hence $k_1=k_2$ and equal special characters. The high lines are therefore distinct.

A high-line character takes values in $\F_5^\times$, so $c_i-B_i$ is divisible by six. Since $c_i$ is odd, $k_i$ is odd. Thus the parameters must be 25 and 75. For $B=25$ the possible exponents are $c=1,13$; for $B=75$ they are $c=9,21$. Equality of the exchanged residual character pairs leaves the two rows
\[
 (25,1;75,9),\qquad(25,13;75,21).
\]
A common quadratic twist exchanges them. Put the high lines at $\F_5e_1$ and $\F_5e_2$, in the parameter order 25,75. The two tame residual exponents are $0,6$, up to adding twelve to both.

Write $F_2=L(M_i[25])$, $F_3=L(M_i[125])$ and let $W_1,W_2$ be their congruence layers. We first prove that evaluation $W_1e_i\to\F_5^2$ is onto. A second-level lift of $e_i$ has, in the two models, valuations
\[
          (19/5,3/5)\quad(i=1),\qquad(1/5,9/5)\quad(i=2)
\]
over $L$. Its common field has degree five or 25. At degree five, both normalized values are prime to five. A translation line can be the first high line, the second high line, or a line different from both. The two break values for these three possibilities would be
\[
 \begin{array}{c|ccc}
 i=1&(456,72)&(6,222)&(6,72)\\
 i=2&(474,66)&(24,216)&(24,66).
 \end{array}
\]
For example, in the first row the leading coordinate values are 19 and three; the high/low translating values are 475,25 and 225,75. Subtraction gives the displayed pairs. No pair agrees, contradicting the exact prime-valuation break equality. Thus both point fields $E_i$ have degree 25, and both evaluations are onto.

Tame projectors separate the diagonal plane and the two off-diagonal axes, whose weights are $0,6,18$. The two evaluations consequently put both off-diagonal matrix units in $W_1$ and make both diagonal projections nonzero. In particular $r=\dim W_1\ge3$. Since $W_1\subset W_2$, each lift $Q_i$ of order 125 gives a common field $T_i=L(Q_i)$ with $[T_i:E_i]=25$. A prime-to-five tame lift diagonalizes on the torsion module by the coprime factors of $X^{24}-1$; choose $Q_i$ in its two direct summands. This makes all four point fields stable under the tame lift.

Here are the normalized coordinate and displacement rows. The designation in the last column is the line on which the larger displacement occurs.
\[
\begin{array}{c|c|r|r|r|c}
\text{extension}&\text{model}&A&D_0&D_1&H\\\hline
E_1/L&1&95&125&2375&\text{diagonal}\\
E_1/L&2&15&375&1125&\text{off-diagonal}\\
E_2/L&1&5&125&2375&\text{off-diagonal}\\
E_2/L&2&45&375&1125&\text{diagonal}\\\hline
T_1/E_1&1&95&3125&59375&\text{diagonal}\\
T_1/E_1&2&15&9375&28125&\text{off-diagonal}\\
T_2/E_2&1&5&3125&59375&\text{off-diagonal}\\
T_2/E_2&2&45&9375&28125&\text{diagonal}.
\end{array}
\]
The values come directly from the first-level low/high values $(5,95)$ and $(15,45)$: every inverse step divides its value by 25, while the field degree multiplies by 25. The displacements are the translating first-level values multiplied by the total field degree. The absolute indices are 12000 and 300000, respectively, strictly above every listed displacement.

In each degree 25 extension the translation group has two different tame weights, diagonal and off-diagonal. Its diagonal cyclic quotient break is a positive multiple of 24, and its other quotient break is six or eighteen modulo 24. To see this, on a uniformizer in the cyclic degree-five quotient the tame leading multiplier is $\eta$ with $\eta^5=\zeta$, or $\eta^{125}=\zeta$ at the higher point level. In both cases $\eta=\zeta^5$. Conjugating a wild expansion gives weight $\eta^b=\zeta^{5b}$. Thus the two quotient breaks differ. The group has two lower breaks and its last line is one of the two tame eigenaxes.

The two coordinate rows of each extension have opposite high-displacement lines. At least one of them is not the last ramification line, so Lemma~\ref{lem:opposite-line-conductor} bounds the last upper break by the maximum of $(D_0-A)/5$ for the two rows. We obtain
\[
 c(E_1/L)\le72,\quad c(E_2/L)\le66,\quad
 c(T_1/E_1)\le1872,\quad c(T_2/E_2)\le1866.
\]
Since the two $E_i$ generate $F_2$, its last break is at most 72. The same stabilizer and elementary abelian compositum argument as in Lemma~\ref{lem:point-conductors} gives
\begin{equation}\label{eq:tame-terminal-bound}
                       c(F_3/F_2)\le5^{r-2}\,1872.
\end{equation}
Indeed terminal translations are onto both columns, giving $F_2\cap T_i=E_i$, and $[F_2:E_i]=5^{r-2}$. Upper subgroup conversion bounds the two terminal point quotients by that degree times their original conductors; they generate $F_3/F_2$.

If the first-layer trace vanishes, $\mu_{25}\subset K$ and $L(\mu_{125})/L$ is cyclic of degree five with break 576, by the Cohen-field calculation in Section~\ref{sec:irreducible}. Its break exceeds 72, so its terminal restriction has conductor greater than $5^r(576-72)$, contradicting~\eqref{eq:tame-terminal-bound}, since $25(576-72)=12600>1872$.

Otherwise $\mu_{25}\not\subset K$. Theorem~\ref{thm:cyclotomic-exact} gives $\delta_2\ge20$, so the second break over $L$ is at least $480$. It is greater than 72. The cyclotomic field meets the elementary abelian $F_2$ in its first layer, and its nontrivial terminal restriction has conductor greater than $5^r(480-72)$. This again contradicts~\eqref{eq:tame-terminal-bound}, since
\[
                         25(480-72)=10200>1872.
\]
Both determinant cases are impossible, proving the theorem.
\end{proof}

\section{Wild reducible residual inertia at $p=5$}\label{sec:wild-reducible}
We now treat nontrivial residual wild inertia. The high line is then intrinsic: it is the unique fixed line of the residual wild group. Throughout this section $e=20$, $N=24$, $s=120$, and $v_L(5)=480$. A common finite twist normalizes a nonsquare finite determinant so that $c_i\equiv1\pmod4$. A square determinant has already been treated in Proposition~\ref{prop:square-determinant}.

\subsection{Parameters and the two possible matrix groups}
By Lemma~\ref{lem:prime-residual-comparison}, different special characters force $B_i=5b_i$. Equality of the high-line characters and finite determinants gives
\[
 b_i\in\{1,3,5,7,9,11,13,15,17,19\},\qquad
 b_1\equiv b_2\pmod4.
\]
Indeed $c_i-B_i$ is divisible by six, so $b_i$ is odd, and equality of parameters would recover $c_i$. Both parameters cannot be below 20, since the only possibilities are 5 and 15. The boundary 20 is impossible.

If the parameters straddle 20, their high quotients have rational first-level torsion: in the smaller-parameter quotient the old image is a stable low line distinct from its high line. A $5$-group acting on both lines is trivial. Lemma~\ref{lem:one-prime-parameter}, applied to the quotient $25$-comparison, forces the larger $b$ to satisfy $5\mid24-b$. The congruence modulo four leaves exactly
\begin{equation}\label{eq:crossing-pairs}
                         (B_1,B_2)=(5,45),(15,95).
\end{equation}
Their quotient parameters are 25 and 75, with exchanged high lines. If both $B_i>20$, their high quotient parameters are $5(24-b_i)$. Their old first-level images are the common quotient high line. A rational quotient first level would force both $24-b_i$ divisible by five by Lemma~\ref{lem:one-prime-parameter}; the possibilities 9 and 19 for $b_i$ are incongruent modulo four. Thus both original and quotient residual actions are wild.

In the both-large case choose an integral tame eigenbasis with high line $\F_5e_1$. Write the wild image at level 125 as $W$ and its congruence layers as $W_1,W_2$. The functions
\[
                 v_0=\rho_{12}\bmod5,
          \qquad w_0=\rho_{21}/5\bmod5
\]
are additive and nonzero; the second is the quotient residual character. Their tame weights are opposite primitive fourth roots. Thus $w_0$ cannot factor through $v_0$, and the first layer contains $E_{21}$ by a tame projector. Conjugating by $1+tE_{12}$ for $t=1,-1$ gives $E_{12}$ and $\operatorname{diag}(1,-1)$. Fifth powers give $W_1\subset W_2$. Consequently
\[
                  \mathfrak{sl}_2(\F_5)\subset W_1\subset W_2.
\]
The possibility $W_1=W_2=\mathfrak{sl}_2$ is excluded as follows. On the first congruence subgroup the determinant modulo 125 factors through an invariant linear functional on $\mathfrak{sl}_2$. Tame conjugation kills its two off-diagonal values; conjugating $E_{21}$ by $1+E_{12}$ kills the diagonal value. Hence the determinant is trivial there. Its total image has order at most five, so $\mu_{25}\subset K$ and its remaining cyclotomic field equals the residual field $R/L$. The Cohen-field calculation gives break 576. But a low first-level coordinate gives $b(R/L)\le600-6B_i<576$, a contradiction.

There remain two actual groups. Put
\[
 \mathcal I=\left\{
 \begin{pmatrix}1+5a&v\\5w&1+5d\end{pmatrix}:
 a,d,w\in\Z/25,\ v\in\Z/125\right\},
\]
and write each entry in base-five digits, for example $a=a_0+5a_1$ and $v=v_0+5v_1+25v_2$. Then
\begin{equation}\label{eq:two-iwahori-groups}
 W=\mathcal I\quad\hbox{or}\quad
 W=\{\rho\in\mathcal I:a_0+d_0=v_0w_0\}.
\end{equation}
For the first case this follows from the order $5^{1+4+4}=5^9$. In the second, the determinant modulo 25 is trivial on the first layer; its remaining character factors through the residual $v_0$, whose nontrivial tame weight excludes an invariant determinant. Thus the displayed determinant relation holds, and the orders $5^{1+3+4}=5^8$ agree. Notice that both groups retain independent coordinates $v_0,a_0,w_1$ over the quotient residual field, and independent coordinates $w_0,d_0,v_1$ over the original residual field. This fact allows a single point argument to treat both groups.

\subsection{Two fixed-coefficient lemmas}
\begin{lemma}\label{lem:common-high-numbers}
Suppose an elementary abelian degree 25 extension has two coordinate rows
\[
 A_i=5m_i,\qquad D_{0,i}=125m_i,\qquad
 D_{1,i}=15000-625m_i,\qquad v(5)=60000,
\]
with a common high-displacement line, where $1\le m_1<m_2\le19$ and both $m_i$ are prime to five. Then that line is last, and the two lower breaks and the larger upper break are
\[
 t=24m_1,\qquad u=15000-630m_2-96m_1,
 \qquad C=3000-126m_2.
\]
\end{lemma}
\begin{proof}
A single break $h$ is impossible. If the first coordinate change were below $D_{0,i}$, a nonzero additive polynomial of degree at most five would vanish on all 25 injective leading coefficients. At $D_{0,i}$ a unique monomial could not vanish precisely on the high line. A tie is therefore necessary, forcing $h=(D_{0,i}-A_i)/5=24m_i$ for both models. These values differ.

For two breaks $t<u$, a last line different from the high line would give larger upper break $24m_i$ in both models by Lemma~\ref{lem:opposite-line-conductor}. Hence the high line is last. Put $r_i=24m_i$. The expansion~\eqref{eq:first-change}, with first ordinary exponent $J_i$, gives
\[
\begin{array}{c|c|c}
 &J_i&\text{last-break equality}\\\hline
 t<r_i&5m_i+4t&u+4t=15000-630m_i\\
 t>r_i&125m_i-t&u-t=15000-750m_i.
\end{array}
\]
In the first row the two earlier terms must cancel; in the second the ordinary term is uniquely first. At the later break it is uniquely first in both cases. The same strict row in both models forces equal parameters. The mixed strict case gives $t=150m_1-126m_2<24m_1$, contrary to its assumed position. At $t=r_2$, the first model gives $u=15000-750m_1+24m_2$. For the second model, $A_2+5u-D_{1,2}=750(80-5m_1+m_2)>0$, because $m_1<m_2\le19$. Thus the ordinary term controls the later displacement. The boundary requires $J_2\ge101m_2$, but that later displacement forces $J_2=750m_1-649m_2$, implying $m_1\ge m_2$. Thus $t=r_1$, and the larger parameter supplies the asserted formulas.
\end{proof}

\begin{lemma}\label{lem:large-gap}
Let $T/S$ be totally ramified elementary abelian of degree 25, with a line $H$ in its group. Suppose coordinates have values $A_i>0$, $5\mid A_i$, and displacements $D_{1,i}$ on $H\setminus\{0\}$ and $D_{0,i}$ off $H$, with
\[
                    A_i<D_{0,i}<D_{1,i}<v_T(5).
\]
Put $\Delta_i=D_{0,i}-A_i$. If the cyclic quotient killing $H$ has conductor $z$ satisfying
\[
 z>\frac9{25}\min_i\Delta_i,\qquad
 z>\frac15\max_i\Delta_i,
\]
then every gap $D_{1,i}-D_{0,i}$ is the same.
\end{lemma}
\begin{proof}
In the expansion in a uniformizer with fixed coefficient representatives, every five-divisible exponent has its first possible change at degree at least $A_i+5h$ at break $h$. The first nonzero ordinary exponent $J_i$ changes at $J_i+h$ with nonzero coefficient. This holds also when $25\mid A_i$; terms with five-divisible binomial coefficients are beyond the stated mixed-characteristic cutoff.

If there is a single break $z$, the second inequality places all divisible terms after $D_{0,i}$. An ordinary term must give $J_i+z=D_{0,i}$, and it is nonzero for every nonidentity element, including those of $H$. This contradicts $D_{1,i}>D_{0,i}$.

Suppose the lower breaks are $t<u$, with last line $V\ne H$. Choose $i$ minimizing $\Delta_i$. The displacement on $V$ gives $u\le\Delta_i$. If $t>\Delta_i/5$, the ordinary term forced at $u$ would already give $J_i+t<D_{0,i}<D_{1,i}$ on $H$, a contradiction. Thus $t\le\Delta_i/5$, and the quotient conductor is $(u+4t)/5\le9\Delta_i/25$, again impossible. Hence $V=H$ and $t=z$. At $t$, every divisible term is later than $D_{0,i}$, so $J_i+t=D_{0,i}$. At $u$ these terms increase at least five times as fast as the ordinary term; therefore $D_{1,i}=J_i+u$. The common gap is $u-t$.
\end{proof}

\subsection{A selected point excludes both large parameters}
\begin{proposition}\label{prop:both-large-excluded}
A different-special-character comparison cannot have $B_1,B_2>20$.
\end{proposition}
\begin{proof}
The possible different-character parameter pairs, after ordering, are
\[
 (5,9),(5,13),(5,17),(7,15),(7,19),
 (9,13),(9,17),(11,15),(11,19),(15,19).
\]
Here the entries are $b_i=B_i/5$. This list follows by combining $b_1\equiv b_2\pmod4$, $c_i\equiv1\pmod4$, $c_i\equiv5b_i\pmod6$, and equality of $c_i-5b_i$ modulo 24, then removing the two pairs $(7,11),(13,17)$ giving equal unordered orbits $\{c_i,5c_i\}$. The involution $(b_1,b_2)\mapsto(24-b_2,24-b_1)$ pairs the ten rows. We may choose the original or high-quotient orientation to obtain
\begin{equation}\label{eq:large-five-pairs}
                    (5,9),(5,13),(5,17),(7,15),(11,15).
\end{equation}
This change of orientation will use a selected quotient point only, not an extra level of quotient torsion.

For the original orientation set $P^+=e_1/125$, $P=5P^+$, $E_h=L(P)$, $S=R_oE_h$, and $T=R_oL(P^+)$, where $R_o$ is the field of $v_0$. In either group~\eqref{eq:two-iwahori-groups}, their stabilizers respectively impose
\[
 a_0=w_0=0;\qquad v_0=a_0=w_0=0;\qquad v_0=a=w=0.
\]
Their degrees over $L$ are 25,125,3125. On the subgroup for $S$, $(a_1,w_1)$ are independent additive coordinates and their joint kernel fixes $T$. Thus $T/S$ is elementary abelian of degree 25, with high line $w_1=0$.

In the quotient orientation use the actual image $M[125]/H$, of type $\Z/125\oplus\Z/25$, and choose $P^+=q(e_2/125)$. Its order is 125 and its first-level image is the quotient high point. In the reordered quotient lattice $(e_2,e_1/5)$ the matrix is
\[
                 \begin{pmatrix}1+5d&w\\5v&1+5a\end{pmatrix}.
\]
Its first column is determined modulo 125 by the original module. The stabilizers of $P$, $S$, and $T$ impose respectively
\[
 v_0=d_0=0;\qquad w_0=v_0=d_0=0;\qquad
 w_0=v_0=v_1=d_0=d_1=0.
\]
The same degrees result in both groups, and $T/S$ has additive coordinates $(d_1,v_1)$ and high line $v_1=0$. The trace relation in~\eqref{eq:two-iwahori-groups} only determines the unused diagonal digit in each of these calculations.

For either orientation the two successive high inverse polygons are single segments. A high first-level value $5(24-b_i)$ becomes $(24-b_i)/5$ and then $(24-b_i)/125$. Hence in $T$ the coordinates have
\begin{equation}\label{eq:large-point-rows}
 A_i=25(24-b_i),\quad D_{0,i}=3125b_i,\quad
 D_{1,i}=15625(24-b_i),\quad v_T(5)=1500000.
\end{equation}
These displacements are the translating first-level values multiplied by 3125.

We compute a lower bound for the quotient conductor killing the high line. In the opposite orientation take a primitive low 25 point. Its field $Z$ has degree 125 over $L$ and degree 25 over its residual field $R_*$. Its two rows are those of Lemma~\ref{lem:common-high-numbers}, with parameter pairs
\[
                    (15,19),(11,19),(7,19),(9,17),(9,13).
\]
The diagonal and off-diagonal tame weights differ, so there are two ramification breaks. For pairs prime to five, the common high line is last by Lemma~\ref{lem:opposite-line-conductor}. For $(15,19)$ this conclusion still holds, as follows. If the high line were not last, the prime row would give $u+4t=120\cdot19$. In the other row $A=75$, and the first possible changes have degrees
\[
                \min\{75+25h,N_1+5h,J+h\},
\]
where $N_1>75$ is the first exponent of five-adic order one and $J>75$ the first ordinary exponent. At $t$ at least two terms must tie, because the required high displacement is larger than the low displacement at $u$. A tie of the last two terms gives $J=N_1+4t$ and low displacement $N_1+120\cdot19>1875$. A tie of the first and last gives $J=75+24t$ and an even larger displacement. A tie of the first two with the last later gives $N_1=75+20t$, while at $u$ both divisible terms and the ordinary term are later than 1875. All cases contradict the prescribed low displacement. Thus the high line is last.

Let $t'$ and $C'$ be the diagonal and off-diagonal conductors of $Z/R_*$. The resulting values are
\[
\begin{array}{c|c|r}
 (m_1,m_2)&t'&C'\\\hline
 (15,19)&24\le t'\le336,\ 24\mid t'&606\\
 (11,19)&264&606\\
 (7,19)&168&606\\
 (9,17)&216&858\\
 (9,13)&216&1362.
\end{array}
\]
The four prime rows follow from Lemma~\ref{lem:common-high-numbers}. For the mixed row, if $t'>24\cdot19$, the first ordinary exponents in both rows force incompatible values $u-t'=15000-750m_i$; equality at $24\cdot19$ gives the same boundary contradiction as in that lemma. Thus $t'<456$, and the prime row gives $u=3030-4t'$. In the row $m_1=15$, one has $5u+75>D_1$ because $t'<456$, and its forced ordinary exponent is $J=5625-u$. Its change at $t'$ is later than 1875, since $J+t'-1875=720+5t'>0$. If $t'\ge360$, every divisible term is also later than 1875: the leading degree 75 has order two, and all order-one degrees exceed 75. Therefore $t'<360$. The diagonal tame weight gives $24\mid t'$, and the upper conversion gives $C'=606$.

In the original orientation, on the subgroup over $R_*$ the three coordinates $(v_0,a_0,w_1)$ are independent additive characters in both groups~\eqref{eq:two-iwahori-groups}. The field $Z$ has characters $(a_0,w_1)$, and $S$ has $(v_0,a_0)$. If $\sigma$ is the conductor of $v_0$, conjugation by an opposite residual unipotent sends $a_0$ to $a_0\pm t v_0$. Invariance of conductors gives $\sigma\le t'$, and the distinct tame weights exclude equality through the injective leading-coefficient map. Thus $0<\sigma<t'<C'$. The actual compositum $SZ/R_*$ has these three independent characters, so the conductor of $w_1$ restricted to $S$ is
\[
 z=\psi_{S/R_*}(C')
   =\sigma+5(t'-\sigma)+25(C'-t')>25C'-24t'.
\]
It is the quotient of $T/S$ killing its high line. For the quotient orientation the same calculation uses $(w_0,d_0,v_1)$, with $Z$ defined by $(d_0,v_1)$ and $S$ by $(w_0,d_0)$. Thus it proves the same bound without any hypothetical quotient 125 comparison.

Substitution gives
\[
\begin{array}{c|r|r|r}
 (b_1,b_2)&z>&9\min_i(D_{0,i}-A_i)/25&\max_i(D_{0,i}-A_i)/5\\\hline
 (5,9)&7086&5454&5550\\
 (5,13)&8814&5454&8070\\
 (5,17)&11118&5454&10590\\
 (7,15)&16266&7722&9330\\
 (11,15)&28866&12258&9330.
\end{array}
\]
All hypotheses of Lemma~\ref{lem:large-gap} hold. It forces equality of the gaps in~\eqref{eq:large-point-rows}, but these are $375000-18750b_i$, different for $b_1<b_2$. This excludes both actual groups and every both-large pair.
\end{proof}

\subsection{The two crossing parameters}
\begin{lemma}\label{lem:opposite-twenty-five}
Suppose a comparison of $25$-torsion has rational first-level modules over $L$, parameters 25 and 75, and exchanged high lines $\F_5e_1,\F_5e_2$. Its image is $1+5M_2(\F_5)$, and its coordinate characters have breaks
\[
             c(A_{21})=6,\quad c(A_{22})=24,
             \quad c(A_{12})=66,\quad c(A_{11})=72.
\]
\end{lemma}
\begin{proof}
The degree-five point-field exclusions in the proof of Theorem~\ref{thm:tame-five-recovery} use only the $25$-comparison: the coordinate values are $(19/5,3/5)$ and $(1/5,9/5)$, and the six pairs of exact degree-five breaks displayed there all disagree. Thus both point fields $E_i$ have degree 25. Tame projection puts both off-diagonal axes in the image and gives nonzero projections to both diagonal coordinates. Its dimension is at least three.

The four coordinate rows are the first four rows in that theorem's table. In $E_1$, an off-diagonal last line would give last upper break 6 by Lemma~\ref{lem:opposite-line-conductor}, impossible since the diagonal cyclic break is a positive multiple of 24. Thus the diagonal is last, with upper break 72. In $E_2$, a diagonal last line would have break 24; an off-diagonal last line would have break 66. If the matrix image had dimension three, its two diagonal projections would be the same nonzero cyclic character up to scalar, and hence have the same conductor. Neither alternative in $E_2$ can share the diagonal break 72 of $E_1$. Therefore the image is all $M_2(\F_5)$.

We determine the remaining breaks without a computation of any higher torsion. In $E_1$, the smaller lower break satisfies $t\equiv6\pmod{24}$ and $0<t<72$, and the larger is $u=360-4t$. For $t=30$, the second row forces the ordinary exponent $15+4t=135$, divisible by five. For $t=54$, the first row forces the exponent $125-t=71<95$. Both are impossible, so $(t,u)=(6,336)$.

In $E_2$, a diagonal last line would give $(t,u)=(18,48)$. The fourth coordinate row forces the first ordinary exponent 117, which would change at 165 at the later break instead of 1125. If the off-diagonal is last, its upper break is 66 and the diagonal smaller break is 24 or 48. The latter forces the third row's ordinary exponent 77 and a later change 215 instead of 2375. Hence $(t,u)=(24,234)$. Upper conversion gives precisely the four stated cyclic conductors.
\end{proof}

In either crossing~\eqref{eq:crossing-pairs}, the high quotients satisfy this lemma. Their lattice $(e_1/5,e_2)$ has matrix
\[
 \begin{pmatrix}\rho_{11}&5\rho_{12}\\\rho_{21}/5&\rho_{22}\end{pmatrix}.
\]
Rationality of its first level forces $25\mid\rho_{21}$. Its full 25 image makes the four functions $v_0=\rho_{12}\bmod5$, $a_0=(\rho_{11}-1)/5\bmod5$, $d_0=(\rho_{22}-1)/5\bmod5$, and $c=\rho_{21}/25\bmod5$ jointly surjective. The original first layer is upper triangular, with the full diagonal plane by tame projection and $E_{12}$ from fifth powers of residual unipotents. The terminal layer contains this plane and $E_{12}$. If its lower coordinate vanished, the nonzero character $c$ would factor successively through the first layer and residual upper character, both of different tame weights. This is impossible. Hence the terminal layer is all $M_2(\F_5)$, and the full image is
\begin{equation}\label{eq:crossing-matrix}
 W=\left\{\begin{pmatrix}1+5a&v\\25c&1+5d\end{pmatrix}:
 a,d\in\Z/25,\ v\in\Z/125,\ c\in\F_5\right\},
 \qquad |W|=5^8.
\end{equation}
The order proves equality with the displayed containing group.

Let $R$ be the original residual field and $D$ the cyclic field of $d_0$. In the first crossing the original residual upper entry is the quotient character $A_{12}$ and $d_0$ is $A_{22}$; their breaks are 66 and 24. In the second crossing the parameter 25 quotient is the other model, so reordering its high basis identifies these with $A_{21}$ and $A_{11}$; their breaks are 6 and 72.

A low 25 point has field $E_\ell/L$ of degree 125, with $E_\ell/R$ elementary abelian of degree 25 and translation coordinates $(v_1,d_0)$. Its rows are those of Lemma~\ref{lem:common-high-numbers}, with $(m_1,m_2)=(1,9)$ or $(3,19)$. The quotient killing the high line is the actual field $DR/R$. In the second crossing that lemma gives its break 72, whereas the actual distinct fields $R,D$ over $L$ have breaks 6 and 72, giving restricted break $6+5(72-6)=336$. This contradiction excludes $(15,95)$.

\subsection{The terminal determinant excludes the last crossing}
Theorem~\ref{thm:cyclotomic-exact} gives the last upper break of $K(\mu_{125})/K$ as $44-d\ge20$, where $d=v_K(\mathfrak D_{K/K_0(\mu_5)})\le24$. Its last break after the tame extension $L/K$ is therefore at least $480$.

Assume now the remaining crossing $(5,45)$. Let $F$ and $F_2$ be the original 125 and 25 division fields. In~\eqref{eq:crossing-matrix}, the four cyclic fields defined by $(v_0,a_0,d_0,c)$ have breaks $(66,72,24,6)$ over $L$, by Lemma~\ref{lem:opposite-twenty-five}. Put $S=R A$, where $A$ is the field of $a_0$. The common-high lemma gives lower breaks 24 and 9234 for $E_\ell/R$, hence upper breaks 24 and 1866 for its characters $(d_0,v_1)$. The character $a_0$ restricted to $R$ has conductor $66+5(72-66)=96$. Thus $F_2/R$ has upper breaks 24,96,1866 and lower breaks 24,384,44634.

For precision we also compute the passage to $L$. The nonidentity elements of $\Gal(F_2/R)$ have these three lower breaks with multiplicities 100,20,4. If $M$ is the maximal lower break in a nontrivial residual coset, the quotient-average formula gives
\[
 \frac{100\min(24,M)+20\min(384,M)+4\min(44634,M)+M}{125}=66.
\]
Its unique solution is 234. Indeed the left side is strictly increasing and equals $24+(M-24)/5$ between 24 and 384. The other elements of the coset have the minimum of $M$ and the corresponding subgroup break, by the ultrametric inequality and maximality. Therefore
\begin{equation}\label{eq:crossing-filtrations}
\begin{array}{c|c|c}
 \text{extension}&\text{lower breaks}&\text{upper breaks}\\\hline
 F_2/L&24,234,384,44634&24,66,72,426\\
 F_2/S&24,44634&24,8946\\
 F_2/E_\ell&384&384.
\end{array}
\end{equation}
The last two rows also follow directly from the actual stabilizers: over $S$ one has $v_0=a_0=0$ and free $(d_0,v_1)$; over $E_\ell$ one has $v_0=v_1=d_0=0$ and free $a_0$.

We bound the two diagonal characters on the terminal group $\Gal(F/F_2)=1+25M_2(\F_5)$. Choose a high 25 inverse in the first model's high inverse coset. Its value is 19; the second model's value is 3. Tame projection to the first eigenline preserves that coset, so the point and a $125$-torsion lift may be chosen in the same integral tame eigenline. Both next inverse polygons are single segments, giving values $19/25$ and $3/25$.

For the 125 point $P_h^+$ put $T_h=R L(P_h^+)$. Its stabilizer is $v_0=a=c=0$, so $[T_h:L]=625$ and $T_h/S$ is elementary abelian with independent additive coordinates $(a_1,c)$. These assertions follow by multiplying the matrices on the subgroup $v_0=a_0=0$. The character $c$ has conductor 6 over $S$, since both breaks of $S/L$ are larger. In $T_h$ the second coordinate has value 75, while every nonzero high translation has displacement $625\cdot75=46875$. Hence its lower break is at most 46800. If the high line is last, the first break is 6 and the last upper break is at most $(46800+24)/5$; if it is not last, the $c$ conductor 6 is already the largest. Thus in every case
\[
                         c_S(a_1)\le46824/5.
\]
Restriction to $F_2$, using~\eqref{eq:crossing-filtrations}, gives
\begin{equation}\label{eq:crossing-diagonal-bounds}
 c_{F_2}(a_1)\le44634+25(46824/5-8946)=55104.
\end{equation}
No disjointness assertion is needed for this upper restriction bound.

For a low 125 lift $P_\ell^+$, its field has degree 3125 over $L$ and is elementary abelian of degree 25 over $E_\ell$, with additive coordinates $(v_2,d_1)$. The first-model coordinate has value 5. Every element with $d_1\ne0$ displaces it by 3125, so every such lower break is at most 3120. Averaging over the high line gives $c_{E_\ell}(d_1)\le3120$. Restricting through the cyclic last row of~\eqref{eq:crossing-filtrations} gives
\[
                    c_{F_2}(d_1)\le384+5(3120-384)=14064.
\]
Both characters are the actual terminal diagonal coordinates. The first is already additive on the subgroup for $S$, and the second on the subgroup for $E_\ell$; this verifies their identities before applying any restriction formula.

The actual determinant modulo 25 is $a_0+d_0$, whose conductor over $L$ is 72. Hence $\mu_{25}\not\subset K$. By Theorem~\ref{thm:cyclotomic-exact}, the last cyclotomic break over $L$ is at least 480, greater than the last break 426 of $F_2/L$. Its terminal restriction therefore has conductor at least
\[
                  \psi_{F_2/L}(480)=44634+625(480-426)=78384.
\]
On the terminal group, however, the same determinant is $1+25(a_1+d_1)$, whose conductor is at most 55104 by the two proved bounds. This contradiction excludes the last crossing.

\begin{theorem}\label{thm:wild-five-recovery}
Two potentially supersingular height-two groups over a field with algebraically closed residue field of characteristic five, $e_K=20$, and $\mu_5\subset K$, having isomorphic 125-torsion and nontrivial residual wild inertia, have equal finite special characters.
\end{theorem}
\begin{proof}
A square finite determinant is covered by Proposition~\ref{prop:square-determinant}. Otherwise the parameter reduction leaves precisely the crossings~\eqref{eq:crossing-pairs} or both-large parameters. The preceding arguments exclude the crossings, and Proposition~\ref{prop:both-large-excluded} excludes both-large parameters for both possible congruence images. These cases exhaust the different-character assumption.
\end{proof}

\section{The first wild index at five by residual comparison}\label{sec:exact-five}
This section gives a proof of the first wild index at five by separate residual comparisons; Section~\ref{sec:first-wild-general} proves the general-prime theorem from the graph valuation bound.
\begin{theorem}\label{thm:exact-five}
For the height-two filtration class over complete fields with perfect residue field, $e_K=20$, and $\mu_5\subset K$, one has
\[
                         \nuMin(5,20,2)=2.
\]
Thus $125$-torsion determines the finite special character uniformly, and $25$-torsion does not.
\end{theorem}
\begin{proof}
For the lower bound apply Theorem~\ref{thm:general-arc} with $p=5$, $n=2$, $e=20$, $A=1$. Its weights are $d=480$ and $b=380$, and its tangent multiplier is a primitive 24th root. The construction gives a supersingular group and a group filtered by height-one groups with the same 25-torsion but different special characters. This is an actual example over a base of the specified index, so the uniform index is at least two.

For the upper bound, take a comparison of 125-torsion and pass to the completed maximal unramified base. The determinant on inertia has the form $\chi_5^d\epsilon$, with $0\le d\le2$ and finite $\epsilon\in\mu_4$. Its wild cyclotomic order modulo 125 is at least five, because $v_5(e_K)=1$. Equality of the determinants therefore gives the same $d$ and the same finite factor. As in the proof of Theorem~\ref{thm:second-wild-exact}, dimensions zero and two are determined residually, and the dimension-one case consists of filtered height-one groups or whole potentially supersingular groups. Two filtered groups are determined residually. A mixed comparison would require $e_K>5^{2\cdot3-3-1}=25$ by Proposition~\ref{prop:mixed-obstruction}, contrary to $e_K=20$.

For two supersingular groups, their common residual inertia is either irreducible, tame reducible, or wild reducible. Theorems~\ref{thm:irreducible-recovery}, \ref{thm:tame-five-recovery}, and \ref{thm:wild-five-recovery} settle these cases. Passing to the completed unramified base does not change the inertia character being compared, so the conclusion holds for the original perfect residue field. This proves the upper bound two.
\end{proof}

\section{Intervals with absolute index prime to $p$}\label{sec:prime-to-p}
This section proves the uniform prime-to-$p$ character bound by recovering canonical lines and derives its first exact intervals.
Assume throughout this section that $p\ge5$, $p\nmid e$, and $\mu_p\subset K$. After completed unramified base change, take the common tame good field $L/K$ of degree $N=p^2-1$ and put $s=(p+1)e$.

\begin{lemma}\label{lem:prime-index-irreducible}
A $p^2$-torsion comparison with irreducible residual inertia determines the finite special characters when $p\nmid e$.
\end{lemma}
\begin{proof}
The normal residual wild subgroup is trivial by irreducibility. Over $L$ the first-level module is rational, and irreducibility forces a one-slope polygon. Its first-level coordinate value is $e$ and every primitive second-level value is $e/q$, where $q=p^2$. As in Section~\ref{sec:irreducible}, the first-level leading spaces are scalar multiples of $\F_q$. If the special characters differed, their identification would be Frobenius-semilinear after scalar normalization.

Choose corresponding primitive second-level points. Their common point field $E/L$ is Galois with group a subgroup of the rational first-level translation group. Its degree is at most $q$, and the denominator of $e/q$ forces degree exactly $q$. Thus every first-level translation occurs. In $E$ the two coordinates have value $e$, prime to $p$, and every nonidentity displacement has value $qe$. With a fixed uniformizer and fixed coefficient representatives, write the leading coordinates as $a_i\varpi^e$. The exact prime-valuation formula gives common break $Ne$ and leading displacement $\bar e a_i d_g$ for every $g\ne1$. Consequently the two first-level leading maps on the entire translation group are scalar multiples of the same map $g\mapsto d_g$.

Their prescribed Frobenius-semilinear identification would therefore make Frobenius a fixed scalar multiplication on a scalar multiple of all $\F_q$. This is impossible: after rescaling it would assert $z^p=c z$ for all $z\in\F_q$, although that nonzero polynomial has degree $p<q$. Hence the special characters agree. The argument uses a degree-$q$ point field, independently of the degree of the full second division field.
\end{proof}

\begin{lemma}\label{lem:rational-high-recovery}
For two two-slope groups with first-level torsion rational over $L$, every $p^2$-torsion isomorphism identifies their high lines. There is no upper bound on $e$ in this assertion.
\end{lemma}
\begin{proof}
Write their parameters as $D_i$. Rationality gives $p\mid D_i$, so $D_i\ne e$. Suppose the high lines are not identified.

If $D_1<e$, choose a high inverse of a nonzero high point in the first model. Its value is $(s-pD_1)/p$. The $p$ roots of this value form one coset of the high line: translation by a high point preserves the value, whereas translation by a low point lowers it to $D_1/p$. The numerator $s-pD_1$ is prime to $p$, so the point field has degree $p$ and translation group exactly the high line. Its coordinate has normalized value $s-pD_1$, and its nonzero displacements have value $p(s-D_1)$. Its exact break is therefore $eN$.

The corresponding second-model point has a low target and value $D_2/p^3$ over $L$. In the same degree-$p$ field its value is $D_2/q$; nonintegrality is already impossible. Otherwise every nonidentity translation is low in the second model and has displacement $D_2$. The coordinate inequality would give
\[
                  D_2\ge D_2/q+eN,
                  \qquad D_2\ge qe,
\]
contrary to $D_2<pe$.

It remains that both $D_i>e$; reorder so $D_1\le D_2$. A high-target inverse in the first model has value $(s-D_1)/q$, with prime-to-$p$ numerator. Its common field has degree $q$ and every first-level translation occurs. The two normalized values are $s-D_1$ and $D_2/p$. A translation high in the first model and low in the second has exact first-model break $N(s-D_1)$. The second-model inequality implies $D_2/p\ge s-D_1$. But
\[
                    s-D_1\ge s-D_2>D_2/p
\]
because $D_2<pe$. This is the final contradiction. No choice of a complementary tame eigenline is involved, so scalar residual action causes no exception.
\end{proof}

\begin{lemma}\label{lem:one-slope-recovery}
A $p^3$-torsion comparison cannot have one one-slope polygon and one two-slope polygon. If both polygons have one slope, their special characters agree.
\end{lemma}
\begin{proof}
For a one-slope group the first-level leading map, obtained by dividing a coordinate by $\pi_L^e$, is injective and linear. Wild inertia over $L$ fixes it, so first-level torsion is $L$-rational. Tame inertia acts on its two-dimensional image by scalar multiplication by $\lambda=\zeta^{c+e}$. If the residual representation is reducible, an invariant line makes $\lambda\in\F_p$, and the entire residual action is scalar. If both polygons are one-slope, the irreducible case is Lemma~\ref{lem:prime-index-irreducible}; in the reducible case equality of these actual scalars gives equal $c$.

Suppose one polygon is two-slope with parameter $B$. Its high line forces the common residual representation to be reducible, hence scalar. In particular $p\mid B$. Transport that high line to the one-slope group and take both actual quotients; a $p^2$-comparison remains. The one-slope quotient has parameter $e$, prime to $p$. The other quotient has parameter $pB$ if $B<e$, or $s-B$ if $B>e$. In the first case opposite divisibility contradicts Lemma~\ref{lem:prime-residual-comparison}. In the second both parameters are prime, and the first-level break formula forces $s-B=e$, hence $B=pe$, contrary to the strict two-slope range. The equality $B=e$ is excluded by divisibility. These cases exhaust the possibilities.
\end{proof}

\begin{theorem}\label{thm:prime-index-upper}
Let $n\ge3$ and $e<p^n$. In the full height-two filtration class over a complete field with perfect residue field, $p\ge5$, $p\nmid e$, and $\mu_p\subset K$, a $p^n$-torsion comparison determines the finite special character. Equivalently,
\[
                         \nuMin(p,e,2)\le n-1.
\]
\end{theorem}
\begin{proof}
We first compare two whole supersingular groups. The one-slope cases and mixed polygons are covered by Lemma~\ref{lem:one-slope-recovery}. For two two-slope groups, nontrivial residual wild inertia identifies the high line as its unique fixed line. If the residual action is tame, its stable high line has a stable complement, both characters take values in $\F_p^\times$, and $L$ kills them. Thus their first-level modules are rational, and Lemma~\ref{lem:rational-high-recovery} identifies the high lines from the available $p^2$-comparison.

Write the original parameters and tangent exponents as $B_i,c_i$. Equality of the determinant and high-line characters gives
\begin{equation}\label{eq:prime-index-congruences}
 c_1\equiv c_2\pmod{p-1},\qquad
 c_1-B_1\equiv c_2-B_2\pmod N.
\end{equation}
If at least one parameter is prime to $p$, the residual action is wild and Lemma~\ref{lem:prime-residual-comparison} recovers the character. Assume both are divisible by $p$. They cannot equal $e$. If both exceed $e$, their high quotients have parameters $s-B_i$ prime to $p$. Their first-level comparison suffices: a low point field has degree $p$ and exact break $ps-(p+1)(s-B_i)$, so the parameters and high characters agree. If the original parameters straddle $e$, the next first-level modules have trivial and nontrivial wild action respectively. Indeed the small high quotient has two distinct stable lines, its old image and its high line, whereas the larger quotient has nonintegral low value $(s-B_i)/p$. This is impossible.

The remaining case is $B_i=pb_i<e$. We iterate actual high quotients. After $j$ quotients the available level is exactly $p^{n-j}$ and, while all preceding parameters are below $e$, the current parameter is
\[
                         D_i^{(j)}=p^{j+1}b_i.
\]
At every stage $1\le j\le n-2$, the preceding small quotient supplies a stable low line distinct from the current high line. Thus the current first-level module is $L$-rational. Its available level is at least $p^2$, so Lemma~\ref{lem:rational-high-recovery} identifies the high lines and permits the next quotient. This next comparison retains at least $p$-torsion, even at $j=n-2$.

At any such stage the two divisible parameters cannot equal $e$. If both exceed $e$, the next prime parameters are recovered from their first-level breaks. If they straddle $e$, the next residual wild actions disagree. The only remaining possibility is that both stay below $e$ through $j=n-2$. Then
\[
                 0<b_i<e/p^{n-1}<p,
                 \qquad b_i\in\{1,\ldots,p-1\}.
\]
Equation~\eqref{eq:prime-index-congruences} gives $b_1\equiv b_2\pmod{p-1}$. Reduction modulo $p-1$ is injective on that precise interval, so $b_1=b_2$ and the common original high character gives $c_1=c_2$ modulo $N$. This completes the whole supersingular comparison, with every loss of a torsion level accounted for.

For the full filtration class, pass to the completed maximal unramified base. Since $p\nmid e$, the cyclotomic character modulo $p^n$ has wild order $p^{n-1}$. The determinant identifies total dimensions and finite determinant factors. Dimensions zero and two, and two groups filtered by height-one pieces, are determined residually as before. A mixed filtered/supersingular comparison would require $e>p^{2n-3}$ by Proposition~\ref{prop:mixed-obstruction}, whereas $e<p^n\le p^{2n-3}$ for $n\ge3$. Thus it is excluded. The whole-group argument just proved settles the remaining case.
\end{proof}

\begin{corollary}\label{cor:prime-index-intervals}
Assume $p\ge5$, $p\nmid e$, and $\mu_p\subset K$. Then
\[
\begin{array}{c|c}
\text{absolute ramification interval}&\nuMin(p,e,2)\\\hline
 p<e<p^3&2\\
 p^3<e<p^4&3.
\end{array}
\]
More generally, if $e>p$ and $p^k<e<p^{k+1}$, then
\[
 \left\lfloor\frac{k+3}{2}\right\rfloor
       \le\nuMin(p,e,2)\le\max\{2,k\}.
\]
\end{corollary}
\begin{proof}
Take $n=3$ or $4$ in Theorem~\ref{thm:prime-index-upper}. Theorem~\ref{thm:general-arc} gives failures at $p^2$ when $e>p$, and at $p^3$ when $e>p^3$, proving both exact values. For the general upper bound use $n=\max\{3,k+1\}$. Put $r=\lfloor(k+3)/2\rfloor$; then $2r-3\le k$ and the same construction gives a failure at $p^r$, hence $\nuMin\ge r$. No equality is asserted when these bounds differ.
\end{proof}

\section{Restricted supersingular comparisons}\label{sec:restricted}
This section proves additional comparisons for quadratic Lubin--Tate twists and for two-slope models, and determines the obstruction imposed by a stable cyclic subgroup.
\subsection{Conventions and formal-group input}
This subsection fixes the descent character and the valuation conventions used in the comparisons.


Let $K$ be a complete discretely valued field of characteristic zero with
perfect residue field of characteristic $p\geq5$. Write $e_K=v_K(p)$, where
$v_K(K^\times)=\mathbf Z$. Unramified base extension does not change this
index. When considering inertia, we pass to the completion of the maximal
unramified extension and hence assume that the residue field $k$ is
algebraically closed. Put $K_0=\operatorname{Frac} W(k)$.

A potentially supersingular Barsotti--Tate group has height two and acquires
a good dimension-one formal model with special slope $1/2$. Its special
inertia trace is the trace of finite descent on its rational covariant
Dieudonn\'e module, evaluated on prime-to-$p$ inertia. A finite torsion
isomorphism always means an isomorphism of actual generic-fibre finite
group schemes, equivalently of the corresponding finite Galois modules.

We use the height-two Lubin--Tate deformation theorem
\citep{LT1966}. For the Honda formal group, an integral universal deformation
over $W(k)[[u]]$ can be chosen in the form
\begin{equation}\label{rc-eq:ptypical}
 [p]_{F_u}(T)=pT+_{F_u}uT^p+_{F_u}T^{p^2}.
\end{equation}
The Teichm\"uller scalar action is
\begin{equation}\label{rc-eq:scalar}
 F_{a^{1-p}u}(aX,aY)=aF_u(X,Y),\qquad
 a\in\mu_{p^2-1}.
\end{equation}
This normalization and action are recorded explicitly in
\citet[\S2.1]{BDMMRR2019}. No orbit computation from that source is used.

For completeness, the tame good-reduction extension used below can be
bounded uniformly. The endomorphism algebra of the height-two simple
special isocrystal is the quaternion division algebra of invariant $1/2$.
A finite subgroup of its units has no element of order $p$: such an element
would generate a field containing a primitive $p$th root, of degree
$p-1>2$, whereas commutative subfields have degree at most two. A finite
subgroup is conjugate into the units of a maximal order. Reduction to
$\mathbf F_{p^2}^{\times}$ has pro-$p$ kernel, so the finite subgroup injects
into this cyclic group. Its order therefore divides $p^2-1$.
Killing this finite special descent action makes the potentially crystalline
representation crystalline, and the crystalline/good-reduction equivalence
for integral Barsotti--Tate representations over a complete DVR with
perfect residue field gives a good model; see
\citet[Theorem 2.2.1]{Liu2013}. Thus the tame extension of degree
$m=p^2-1$ suffices. Its special trace has the form
\begin{equation}\label{rc-eq:type}
 \theta=\alpha+\alpha^p,\qquad \alpha^m=1.
\end{equation}

\subsection{The deformation weight}

Take the tame extension $L/K$ of degree $m=p^2-1$ and write
$E=v_L(p)=me_K$. Choose a uniformizer $\pi$ of $L$ and a generator $\tau$
of $\Gal(L/K)$ with $\tau(\pi)=\zeta\pi$, where $\zeta$ is a primitive
$m$th root in $W(k)$. A special-fibre framing conjugates the finite stabilizer
into the Teichm\"uller subgroup. The representing property of the framed
deformation functor, together with \eqref{rc-eq:scalar}, then gives the
parameter action $u\mapsto\alpha^{1-p}u$.
If $\alpha(\tau)=\zeta^a$ and $u\ne0$ has valuation $h<E$, comparison of
leading residues gives
\begin{equation}\label{rc-eq:hasseweight}
 h\equiv(1-p)a\pmod{p^2-1}.
\end{equation}
In particular $p-1\mid h$. We will use this statement only below the
two-slope boundary, where the uncapped parameter valuation is well defined.

\subsection{Newton polygons and division points}

\begin{lemma}\label{rc-lem:newton}
Let $h=v_L(u)>0$, allowing $h=+\infty$. If
$h\geq pE/(p+1)$, all nonzero $p$-torsion points have valuation
$E/(p^2-1)$. If $0<h<pE/(p+1)$, put
\[
 A=\frac{E-h}{p-1},\qquad B=\frac{h}{p(p-1)}.
\]
There are $p-1$ nonzero $p$-torsion points of valuation $A$ and
$p^2-p$ of valuation $B$, with $A>B$. The points of valuation at least
$A$, together with zero, form a subgroup $C$ of order $p$.

For a nonzero torsion target of valuation $0<q\leq A<E$ in the two-slope
case, its $p$-division points have valuations
\[
 \begin{cases}
 q/p^2 & q<p^2B,\quad\text{all }p^2\text{ points};\\
 B & q=p^2B,\quad\text{all }p^2\text{ points};\\
 (q-h)/p\ \text{and}\ B & q>p^2B,
 \quad\text{respectively }p\text{ and }p^2-p\text{ points}.
 \end{cases}
\]
\end{lemma}

\begin{proof}
The lower Newton polygon of $[p](T)/T$ is determined by
$(1,E),(p,h),(p^2,0)$ before division by $T$.
Terms arising from formal addition in \eqref{rc-eq:ptypical} lie strictly
above the relevant supporting lines: a mixed monomial is a product of at
least two positive-valuation summands. The two stated polygons and root
counts follow.
The high-valuation points form a subgroup because an integral formal group
law preserves the valuation filtration. Their count is $p$ including zero.

For division of the given target, the relevant polygon has points
$(0,q),(p,h),(p^2,0)$. The linear term is above it because $q<E$.
The middle point lies above, on, or below the segment from $(0,q)$ to
$(p^2,0)$ precisely as $q<p^2B$, $q=p^2B$, or $q>p^2B$.
Reading slopes proves the three cases.
\end{proof}

The restriction $q\leq A<E$ in this lemma is essential. It is preserved
along a chain of nonzero torsion division points: for every positive
valuation $r$, the formal multiplication law gives $v([p]x)>v(x)$, so
successive divisions have strictly decreasing positive valuation.

\subsection{Unramified quadratic Lubin--Tate twists}

Let $F/\Q_p$ be unramified quadratic, let $\Ocal=\mathbf Z_{p^2}\subset W(k)$,
and let $\sigma$ be its nontrivial automorphism. Let $G_0$ be the formal
$\Ocal$-module with
\[
 [p](T)=pT+T^{p^2}.
\]
Its existence and $\Ocal$-module identities are given by
\citet[Lemma 1 and Theorem 1, pp.~381--382]{LT1965}. It is a height-two,
dimension-one supersingular Barsotti--Tate group as a $\Z_p$-group. Denote its
$\Ocal$-linear Tate character by $\psi:G_{K_0}\to\Ocal^\times$.
For a finite character $\eta:G_K\to\mu_{p^2-1}$, let $T_\eta=\Ocal$ with
character $\eta\psi|_{G_K}$. This is the standard lattice of the finite twist
of $G_0$ by $\eta$, with special trace $\eta+\eta^p$.

\begin{theorem}\label{rc-thm:LT}
Suppose
\[
 T_\eta/p^n\simeq T_{\eta'}/p^n
\]
as $\mathbf Z/p^n[G_K]$-modules. If $p^{n-1}\nmid e_K$, then
$\eta=\eta'$ with the displayed coefficient conventions, and in particular
their special inertia traces agree. This also applies to any $\Ocal$-stable
rank-one lattices in these characters.
\end{theorem}

\begin{proof}
We first derive the division-field degree needed in the argument, including
over the possibly infinite residue field $k$. Put $q=p^2$,
$f(T)=pT+T^q$, and let $f_j$ be its $j$-fold iterate. Then
\[
 Q_j(T)=\frac{f_j(T)}{f_{j-1}(T)}
       =p+f_{j-1}(T)^{q-1}
\]
is a monic Eisenstein polynomial over $W(k)$ of degree
$(q-1)q^{j-1}$. A root $P_j$ has exact order $p^j$, since
$f_{j-1}(P_j)=0$ would force $p=0$. Its annihilator in $\Ocal$ is $(p^j)$,
so its $\Ocal$-span contains $q^j$ points, all the roots of $f_j$.
The endomorphism power series $[a]$ converges in $K_0(P_j)$, so this field
contains every division point and is Galois. Its automorphisms act through
$(\Ocal/p^j)^\times$. The Eisenstein degree is the order of this unit group,
so the action is surjective. The division extension is totally ramified.

Consider the character
\[
 c_n(u)=\frac{\sigma(u)}u\quad\text{on }(\Ocal/p^n)^\times.
\]
Its image on Teichm\"uller units has order $p+1$. On principal units,
the logarithm identifies it with $\sigma-1$ on $p\Ocal/p^n\Ocal$.
Choose a unit $\delta\in\Ocal$ of trace zero; then
$\Ocal=\Z_p\oplus\Z_p\delta$ and $\sigma(\delta)=-\delta$.
Since $2$ is a unit, the image is
$p\Z_p\delta/p^n\Z_p\delta$, of order $p^{n-1}$.
Consequently the pro-$p$ part of $c_n\circ\psi$ cuts out a totally ramified
cyclic extension of $K_0$ of degree $p^{n-1}$. If that character vanishes on
$G_K$, the extension lies in $K$, forcing $p^{n-1}\mid e_K$.

Write $\Ocal_n=\Ocal/p^n$. Every $\mathbf Z/p^n$-linear map
$f:\Ocal_n\to\Ocal_n$ has a unique expression
\begin{equation}\label{rc-eq:intertwiner}
 f(x)=ax+b\sigma(x),\qquad a,b\in\Ocal_n.
\end{equation}
Indeed, evaluation at $1$ and $\delta$ gives
$a+b=f(1)$ and $a-b=f(\delta)/\delta$, determining $a,b$ uniquely.
If $f$ is invertible, at least one of $a,b$ is a unit, since otherwise
its reduction modulo $p$ is zero.

Set $\lambda=\eta\psi$ and $\lambda'=\eta'\psi$.
Intertwining and uniqueness in \eqref{rc-eq:intertwiner} give, for every $g$,
\[
 a(\lambda(g)-\lambda'(g))=0,
 \qquad b(\sigma(\lambda(g))-\lambda'(g))=0.
\]
If $a$ is a unit, $\eta\equiv\eta'\pmod p$, hence $\eta=\eta'$ because
reduction is injective on $\mu_{p^2-1}$. If $b$ is a unit, then
\[
 \frac{\sigma(\psi)}\psi=\frac{\eta'}{\sigma(\eta)}
 \pmod{p^n}\quad\text{on }G_K.
\]
The right side has prime-to-$p$ order. The pro-$p$ quotient just computed
therefore vanishes on $G_K$, contradicting $p^{n-1}\nmid e_K$.
Finally, every $\Ocal$-stable rank-one lattice is an $\Ocal$-fractional ideal
and hence isomorphic to the standard lattice.
\end{proof}

\begin{remark}
The theorem uses the full quadratic coefficient action on the lattice.
It does not assert that an arbitrary potentially supersingular group, or an
arbitrary stable $\Z_p$-lattice in such a representation, belongs to this
family. For example, it applies to the $125$-torsion comparison at
$(p,e_K)=(5,20)$ within this explicitly specified family.
\end{remark}

\subsection{A residual wild break and a second comparison theorem}\label{rc-sec:residual}
This subsection recovers a prime-to-$p$ Hasse valuation from the actual residual splitting field.


Return to the common tame field $L$ and the two-slope notation $A,B,h$
of Lemma~\ref{rc-lem:newton}.

\begin{proposition}\label{rc-prop:break}
If $p\nmid h$, the full $p$-torsion field over $L$ is cyclic of degree $p$,
with lower ramification break
\begin{equation}\label{rc-eq:break}
 b=\frac{pE-(p+1)h}{p-1}.
\end{equation}
If $p\mid h$ and the residual representation over $L$ is nontrivial, its
full splitting extension is again cyclic of degree $p$, and its break
satisfies the strict inequality
\[
 b<\frac{pE-(p+1)h}{p-1}.
\]
\end{proposition}

\begin{proof}
The canonical $p$-torsion points lie in $L$. Their valuation $A$ is integral
by \eqref{rc-eq:hasseweight}; scaling their factor by $\pi^A$ gives a nonzero
constant plus a unit times $X^{p-1}$, with simple roots in $k$.
Hensel's lemma gives all of them over $L$.
The quotient line is also constant over $L$: the determinant of a good
dimension-one Barsotti--Tate group on inertia is cyclotomic, and
$\mu_p\subset L$ because $p-1\mid E$ and $k$ is algebraically closed.
Indeed choose $b\in L$ with $b^{p-1}=-p$; after $X=bY$, the polynomial
$((1+X)^p-1)/(pX)$ reduces to $1-Y^{p-1}$, whose simple nonzero roots lift.
For the determinant background in this generality, see
\citet[proof of Proposition 5.2.2]{Liu2007}.

Choose a noncanonical point $Q$ with coordinate $x$. Galois acts by
$Q\mapsto Q+jP$ for a fixed generator $P$ of the canonical line. Thus
$L(x)$ is the full $p$-torsion field and has degree at most $p$.
If $p\nmid h$, its valuation $B=h/[p(p-1)]$ has denominator $p$, so the
degree is exactly $p$. In the second assertion nontriviality gives the same
conclusion.

Normalize $w=pv_L$ on $L(x)$. For every nontrivial $\gamma$ in its cyclic
Galois group, formal addition gives
\[
 w(\gamma(x)-x)=pA,\qquad w(x)=\frac h{p-1}.
\]
If a cyclic order-$p$ extension has lower break $b$, expansion in a
uniformizer with fixed Teichm\"uller coefficients gives
\[
 w(\gamma(y)-y)\geq w(y)+b.
\]
The inequality is an equality when $p\nmid w(y)$: the first exponent then
has nonzero derivative in the residue field. It is strict when
$p\mid w(y)$, because the first exponent's leading change vanishes and
all subsequent terms have larger valuation. Applying these statements to
$x$ proves both assertions.
\end{proof}

\begin{theorem}\label{rc-thm:primehasse}
Let $M_1,M_2$ be potentially supersingular groups over the same $K$.
Suppose $M_1[p]\simeq M_2[p]$, and over the common tame field $L$ both
models are two-slope with $p\nmid h_1h_2$. Then their special inertia traces
agree.
\end{theorem}

\begin{proof}
The actual residual representations determine the same splitting field over
$L$, hence the same lower break. Formula~\eqref{rc-eq:break} gives
$h_1=h_2=h$, and therefore the same high valuation $A$.
The canonical line is the unique fixed line of nontrivial wild inertia,
so the torsion isomorphism identifies its character $\kappa$.

For a canonical point, its coordinate is $\pi^A$ times a unit.
Field application of $\tau$ multiplies its leading residue by $\zeta^A$.
The good-model descent map has tangent multiplier $\alpha(\tau)$.
Thus the descended Galois action on the canonical line satisfies
\begin{equation}\label{rc-eq:kappa}
 \kappa(\tau)=\zeta^A\alpha(\tau)^{-1}.
\end{equation}
The same $A$ and $\kappa$ give the same $\alpha$, and
\eqref{rc-eq:type} gives equality of special traces.
\end{proof}

\begin{corollary}\label{rc-cor:strictordering}
If two compatible residual representations have two-slope parameters with
$p\nmid h_1$ and $p\mid h_2$, then $h_2<h_1$.
\end{corollary}
\begin{proof}
The common residual representation is nontrivial over $L$ by
Proposition~\ref{rc-prop:break}. Its exact formula for $h_1$ and strict upper
bound for $h_2$ imply the claimed ordering.
\end{proof}

\subsection{Stable cyclic subgroups and the threshold \texorpdfstring{$p^n$}{p\^n}}
This subsection obtains a ramification obstruction by comparing a primitive point valuation with the degree allowed by cyclic Galois action.


\begin{definition}
In the two-slope case, a cyclic subgroup $D\subset M[p^n]$ of order $p^n$
is called noncanonical at level one if its order-$p$ subgroup is one of the
low-valuation lines, rather than the canonical subgroup of
Lemma~\ref{rc-lem:newton}.
The outer canonical range is $0<h<E/(p+1)$.
\end{definition}

\begin{theorem}\label{rc-thm:cyclic}
If $D$ is $G_K$-stable and noncanonical at level one, then
$v_p(h)\geq n$. If the model is in the outer canonical range, this implies
$e_K>p^n$. Consequently no such subgroup exists there when $e_K\leq p^n$.
\end{theorem}

\begin{proof}
Let $Q$ generate $D$. Its order-$p$ multiple has valuation $B$.
Each subsequent division remains on the low branch of
Lemma~\ref{rc-lem:newton}, so
\begin{equation}\label{rc-eq:lowprimitive}
 v_L(T(Q))=\frac h{(p-1)p^{2n-1}}.
\end{equation}
The Galois action on $Q$ is faithful through a subgroup of
$(\mathbf Z/p^n)^\times$, so the $p$-part of $[L(Q):L]$ is at most
$p^{n-1}$. The $p$-part of the denominator of
\eqref{rc-eq:lowprimitive} is $p^{\max(0,2n-1-v_p(h))}$ and must divide its
ramification index. Hence $v_p(h)\geq n$.
By \eqref{rc-eq:hasseweight}, $h/(p-1)$ is a positive integer, so it is at
least $p^n$. In the outer range,
\[
 \frac h{p-1}<\frac E{(p-1)(p+1)}=e_K,
\]
which proves the strict bound.
\end{proof}

\begin{proposition}\label{rc-prop:oneslopecyclic}
If a one-slope supersingular group has a $G_K$-stable cyclic subgroup of
order $p^n$, then $v_p(e_K)\geq n-1$.
\end{proposition}
\begin{proof}
Every primitive point in this range has valuation
\[
 \frac E{(p^2-1)p^{2n-2}}=\frac{e_K}{p^{2n-2}},
\]
by repeated division along the one-slope polygon. Its denominator has $p$-part
\[p^{\max(0,2n-2-v_p(e_K))}.\] The field of a generator of a stable
cyclic subgroup has $p$-degree at most $p^{n-1}$. Comparing the two proves
the claim.
\end{proof}

These arguments do not identify the character of the stable cyclic
subgroup with a finite twist of the trivial or cyclotomic character. Such
an identification would be an additional assertion, not a consequence of
cyclicity or of the denominator estimate.

\subsection{Degree-\texorpdfstring{$p$}{p} quotient valuations}
This subsection computes how the Hasse valuation changes under each degree-$p$ quotient.


\begin{proposition}\label{rc-prop:quotients}
Let $C$ be a subgroup of order $p$ of a two-slope formal group. The Hasse
valuation $h'$ of its quotient, whenever below the one-slope boundary, is
\[
\begin{array}{c|c}
\text{case}&h'\\\hline
C\text{ canonical},\ h<E/(p+1)&ph\\
C\text{ canonical},\ h>E/(p+1)&E-h\\
C\text{ noncanonical}&h/p.
\end{array}
\]
At $h=E/(p+1)$, the canonical quotient is one-slope. A quotient of a
one-slope group by any subgroup of order $p$ has $h'=E/(p+1)$.
\end{proposition}

\begin{proof}
The quotient formal isogeny $q_C$ can be normalized, up to a unit, as the
product of $T-_{F}P$ for $P\in C$. The existence of these quotient
isogenies is classical; see \citet{Lubin1967}. For a point away from
equal-valuation cancellation this gives
\[
 v(q_C(x))=\sum_{P\in C}v(x-_{F}P).
\]

If $C$ is canonical, a low $p$-torsion point has image valuation
$pB=h/(p-1)$. A division point of a high point has valuation
$(A-h)/p$ when $A>p^2B$, and $A/p^2$ when $A<p^2B$.
Its image has valuation $p$ times that value. Since
$A>p^2B$ is equivalent to $h<E/(p+1)$, the resulting high and low
$p$-torsion valuations identify $h'$ as $ph$ in the first regime and
$E-h$ in the second. At equality all nonzero quotient $p$-torsion points
have valuation $E/(p^2-1)$.

If $C$ is noncanonical, a high point has image valuation
\[
 A+(p-1)B=\frac{E-h/p}{p-1},
\]
while division of a low point gives image valuation
$B/p=h/[p^2(p-1)]$. These identify $h'=h/p$.

Finally, in the one-slope case write $r=E/(p^2-1)$.
The distinct leading residues of the $p$-torsion form a two-dimensional
$\F_p$-space. The leading quotient polynomial has as kernel precisely the
one-dimensional residue space of $C$. Thus nonzero points of the dual
kernel have image valuation $pr$, while division points of $C$ have image
valuation $r/p$. The two quotient valuations give $h'=E/(p+1)$.
\end{proof}

In particular, a noncanonical $K$-defined subgroup of order $p$ forces
$p\mid h$, because its quotient over the same good-reduction field has an
integral Hasse valuation.


\section{Canonical quotients and an elementary valuation bound}
We refine the residual comparisons by retaining quotient markings, then obtain an elementary proof of the supersingular upper step at \(p=5\) and \(e=20\).

\subsection{Rational points and quotient markings}
This subsection records the rationality and marking facts needed for the two comparisons.
Keep the notation of Lemma~\ref{rc-lem:newton}: \(p\geq5\), \(e=v_K(p)\), \(L/K\) is the common tame extension of degree \(m=p^2-1\), \(R=\Ocal_L\), and \(E=me\). The residue field is algebraically closed. In the two-slope range put \(h=v_L(u)\), \(A=(E-h)/(p-1)\), and \(B=h/[p(p-1)]\). Choose \(\pi,\tau,\zeta\) with \(\tau(\pi)=\zeta\pi\) as before.

\begin{lemma}\label{fc-lem:rational}
The canonical \(p\)-torsion points in the two-slope case are rational over
\(L\). All \(p\)-torsion points in the one-slope case are rational over
\(L\). In the two-slope case, if a noncanonical \(p\)-line is also
\(L\)-stable, then all \(p\)-torsion is rational over \(L\).
\end{lemma}
\begin{proof}
The high valuation \(A=(E-h)/(p-1)\) is integral by \eqref{rc-eq:hasseweight}.
Substituting \(T=\pi^A X\) in \([p](T)/T\) and dividing by its smallest
coefficient valuation gives a nonzero constant plus a unit times
\(X^{p-1}\) on reduction. Its nonzero roots are simple and lie in \(k\),
so Hensel's lemma gives the canonical points over \(L\).

In the one-slope case substitute \(T=\pi^eX\) in \([p](T)\) and divide
by \(\pi^{p^2e}\). The reduction is an additive polynomial of degree
\(p^2\), with nonzero linear coefficient given by the residue of
\(p/\pi^E\). The coefficient of \(X^p\) may vanish, and that of
\(X^{p^2}\) is a unit. The terms contributed by formal addition have
larger valuation. Thus this polynomial has \(p^2\) distinct roots in
\(k\), which all lift over \(L\). This includes equality in the
one-slope boundary.

Finally, the determinant of the good height-two, dimension-one Tate
module over \(L\) is cyclotomic. Since \(\mu_p\subset L\), its reduction
modulo \(p\) is trivial. The canonical line is already constant, so its
quotient character is trivial. A second stable line splits these two
trivial characters, proving the last assertion.
\end{proof}


For a canonical quotient \(q:G\to G/C\), the dual kernel \(q(G[p])\) is noncanonical when \(h<E/(p+1)\), canonical when \(h>E/(p+1)\), and the quotient is one-slope at equality. Indeed a low \(p\)-point has quotient valuation \(h/(p-1)\); comparison with the quotient slopes in Proposition~\ref{rc-prop:quotients} gives exactly these alternatives. That proposition also proves that every \(p\)-quotient of a one-slope group has Hasse valuation \(E/(p+1)\).

An isomorphism of \(p^2\)-torsion carrying \(C\) to \(C'\) induces an isomorphism of the quotient \(p\)-torsion including these marked dual lines. The quotient groups are the subquotients
\[
\{x\in G[p^2]:px\in C\}/C,\qquad
\{y\in H[p^2]:py\in C'\}/C',
\]
and the markings are the images of \(G[p]\) and \(H[p]\).

\subsection{Two canonical-neighbor comparisons}
These comparisons retain the dual-kernel marking when passing through a canonical quotient.


\begin{proposition}\label{fc-prop:central}
Suppose \(p\mid E\) and \(G[p^2]\simeq H[p^2]\) over \(K\).
If \(G\) has a two-slope Hasse valuation \(h_G\) prime to \(p\) with
\(h_G>E/(p+1)\), then the special inertia traces agree.
\end{proposition}
\begin{proof}
The residual splitting field of \(G\) is nontrivial. Therefore \(H\) is
two-slope. If \(p\nmid h_H\), the residual comparison of Section~\ref{rc-sec:residual} gives the conclusion.
Otherwise \(p\mid h_H\) and \(h_H<h_G\). Wild inertia identifies the
canonical lines. Quotient by them to obtain \(G_1,H_1\), whose residual
representations and marked dual lines are identified.

We have \(h_{G_1}=E-h_G\), prime to \(p\), and its dual marking is
canonical. If \(h_H>E/(p+1)\), then \(h_{H_1}=E-h_H\) is divisible by
\(p\) and greater than \(h_{G_1}\), contradicting
the residual comparison of Section~\ref{rc-sec:residual}. If \(h_H<E/(p+1)\), the marking on \(H_1\) is
noncanonical, whereas it must correspond to the unique wild-fixed
canonical line of \(G_1\). Equivalently, Lemma~\ref{fc-lem:rational} makes
\(H_1[p]\) trivial over \(L\), again a contradiction. At
\(h_H=E/(p+1)\), the quotient \(H_1\) is one-slope and has trivial
residual representation over \(L\), with the same contradiction.
\end{proof}

\begin{proposition}\label{fc-prop:oneslope}
Suppose \(p\nmid E\) and \(G[p^2]\simeq H[p^2]\) over \(L\). If \(G\)
is one-slope, then \(H\) is one-slope.
\end{proposition}
\begin{proof}
Both groups have all \(p\)-torsion rational over \(L\). If \(H\) were
two-slope, its low-point valuation would force \(p\mid h_H\).
Transport its canonical line to a line of \(G[p]\) and quotient both.
Every quotient on the \(G\) side has Hasse valuation \(E/(p+1)\), prime
to \(p\), and residual lower break exactly \(E\).

For \(h_H<E/(p+1)\), the marked dual line of the canonical quotient of
\(H\) is noncanonical and stable. Its residual representation is trivial,
which is impossible. For \(h_H>E/(p+1)\), the quotient Hasse valuation
\(E-h_H\) is prime to \(p\); its break is
\[
 \frac{(p+1)h_H-E}{p-1}<E,
\]
also impossible. Equality \(h_H=E/(p+1)\) is incompatible with
\(p\mid h_H\). Thus \(H\) is one-slope.
\end{proof}

\begin{remark}
Proposition~\ref{fc-prop:oneslope} recognizes a valuation regime over the
common good field. It does not, by itself, recover descent characters
over \(K\).
\end{remark}

\begin{corollary}\label{fc-cor:six}
Let \(p=5\) and \(e=20\). Consider a hypothetical unequal-type pair with
isomorphic \(p^2\)-torsion, one two-slope Hasse valuation prime to \(5\) and
the other divisible by \(5\). Write \(h_i=4z_i\), with \(z_G\) prime to \(5\).
Then its possible ordered normalized valuations are confined to
\[
 (z_G,z_H)\in
 \{(9,5),(13,5),(17,5),(14,10),(18,10),(19,15)\}.
\]
\end{corollary}
\begin{proof}
Here \(E=480\), the two-slope range is \(1\leq z\leq99\), and the
outer range is \(z<20\). the residual comparison of Section~\ref{rc-sec:residual} and
Proposition~\ref{fc-prop:central} leave only
\(0<z_H<z_G<20\), with \(5\mid z_H\) and \(5\nmid z_G\).
The common canonical character in \eqref{rc-eq:kappa} gives
\(a_H-a_G\equiv z_G-z_H\pmod{24}\).
The finite determinant characters agree, so
\(\alpha_G^6=\alpha_H^6\); hence \(4\mid z_G-z_H\).
For \(z_H=5,10,15\), the listed possibilities follow immediately.
No occurrence of a listed pair at any level is asserted.
\end{proof}

\subsection{The graph input and its compatible flag}

We use the graph-specialization statement proved in the proof of Theorem~\ref{thm:common-determinant}.
If an actual isomorphism \(G[p^n]_L\simeq H[p^n]_L\) descends to \(K\)
and the special inertia traces differ, its finite flat graph closure
\(X_n\subset G[p^n]\times H[p^n]\) satisfies
\begin{equation}\label{fc-eq:graphinput}
 (X_n)_k=\ker(F_G^n)\times\ker(F_H^n).
\end{equation}
Here \(F^n\) is the iterated relative Frobenius of the special formal
group. The proof of that statement, based on disjoint descent-character
orbits and maximal isotropy of the graph, is not repeated here.

\begin{proposition}\label{fc-prop:flag}
For \(0\leq i\leq n\), let \(X_i\) be the closure of the restricted
graph at level \(p^i\), with \(X_0=0\). Then these form a finite flat
flag, with \((X_i)_k=\ker(F_G^i)\times\ker(F_H^i)\).
The quotients \(Q_i=X_i/X_{i-1}\) have rank \(p^2\), are killed by \(p\),
and specialize to \(\alpha_p^2\). For \(i\geq2\), multiplication by \(p\)
induces \(Q_i\to Q_{i-1}\) which is a generic-fiber isomorphism and is
zero on the special fiber.
\end{proposition}
\begin{proof}
The generic graph inclusions extend to their flat closures. Multiplication
by \(p\) maps \(X_i\) to \(X_{i-1}\), since the factorization holds
generically and the source is flat. The quotients are finite flat of rank
\(p^2\) and are killed by \(p\), again by generic-fiber faithfulness.
Apply \eqref{fc-eq:graphinput} at each level. Its successive special
quotients are the first Frobenius kernels of the corresponding Frobenius
twists of the two supersingular groups. Each is \(\alpha_p\): it and its
dual are connected, excluding the etale and multiplicative order-\(p\)
groups. The generic multiplication maps on successive quotients of
\((\mathbf Z/p^i)^2\) are isomorphisms. Specially, \([p]\) has initial
degree \(p^2\), so it sends \(\ker(F^i)\) into \(\ker(F^{i-2})\) and
therefore induces zero on the asserted successive quotients.
\end{proof}

\subsection{An invariant augmentation-ideal bound}
The graph algebra gives an elementary numerical obstruction without choosing a canonical subgroup.


\begin{proposition}\label{fc-thm:valuation}
Let \(M,H/K\) be potentially supersingular height-two Barsotti--Tate
groups with isomorphic \(p^n\)-torsion and unequal special inertia traces.
Over the common tame good field $L$ of degree $p^2-1$, every nonzero $p^n$-torsion coordinate in either good integral formal model satisfies
\[
 v_L(T(P))\ge\frac{p-1}{p^n-1}.
\]
Consequently,
\begin{equation}\label{fc-eq:mainbound}
 e(p^n-1)\geq(p-1)p^{2n-2}.
\end{equation}
No assumption on \(v_p(e)\) is required.
\end{proposition}
\begin{proof}
Use the common good field \(L\) and the notation above. The determinant
comparison gives
\(\alpha^{p+1}=\beta^{p+1}\) for the two tangent descent characters.
Indeed both determinant Tate characters have the same cyclotomic factor,
and their remaining finite inertia factors have order prime to \(p\).
The torsion isomorphism makes these factors congruent modulo \(p^n\),
hence equal. The good dimension-one determinant identity is recorded in
the proof of \citet[Proposition~5.2.2]{Liu2007}. Inverting both character
conventions leaves the displayed equality unchanged.

Let \(\Delta\subset\Gal(L/K)\) have order \(p-1\), generated by
\(\delta=\tau^{p+1}\). Its two special tangent scalars are the same
\(c\in\mu_{p-1}\subset\Z_p^*\). Put \(R_0=R^\Delta\). Every nonzero
element of its maximal ideal has \(v_L\) at least \(p-1\).

We first choose exact scalar coordinates. For a formal parameter \(T\),
the projector
\[
 \frac1{p-1}\sum_{j=0}^{p-2}c^{-j}\delta^j(T)
\]
has unit linear coefficient and is a \(c\)-eigenparameter. The canonical
scalar multiplications \([a]\), for \(a\in\mu_{p-1}\), commute with
descent. Applying also
\[
 \frac1{p-1}\sum_{a\in\mu_{p-1}}a^{-1}[a]^*(T)
\]
preserves that property and makes \([a]^*T=aT\). Both projectors preserve
a unit tangent term. Choose such parameters \(T,S\) for the two formal
groups. Then
\[
 \delta(T)=cT,\quad\delta(S)=cS,\qquad
 [a]^*T=aT,\quad[a]^*S=aS.
\]
If pullback conventions give the simultaneous scalar \(c^{-1}\), the same
construction uses that scalar instead. Integral parameter changes with
unit linear term preserve all positive point valuations.

Let \(A_X\) be the algebra of the graph closure, with augmentation ideal
\(I_X\), and write \(x,y\) for the coordinate images. By
\eqref{fc-eq:graphinput},
\[
 A_X/\pi A_X=k[x,y]/(x^{p^n},y^{p^n}).
\]
Thus the monomials \(x^iy^j\), for \(0\leq i,j<p^n\), form an
\(R\)-basis: they generate by Nakayama and their number equals the
finite free rank \(p^{2n}\). Express uniquely
\[
 x^{p^n}=\sum c_{ij}x^iy^j,
 \qquad y^{p^n}=\sum d_{ij}x^iy^j.
\]
Every coefficient is in the maximal ideal of \(R\), and
\(c_{00}=d_{00}=0\) by augmentation.

The graph is stable under descent and simultaneous scalar multiplication.
The semilinear operator \([c]^{*\,-1}\circ\delta\) therefore preserves
\(A_X\), acts as \(\delta\) on \(R\), and fixes \(x,y\). Uniqueness of
the two basis expansions forces \(c_{ij},d_{ij}\in R_0\). Consequently
\begin{equation}\label{fc-eq:augmentation}
 x^{p^n},y^{p^n}\in\pi^{p-1}I_X.
\end{equation}

For a nonzero torsion point \(P\) and its nonzero graph image, put
\[
 r=\min\{v_L(T(P)),v_L(S(f(P)))\}>0.
\]
Every nonconstant monomial has
valuation at least \(r\). Evaluating the two expansions in
\eqref{fc-eq:augmentation} gives
\[
 p^n r\geq p-1+r,\qquad
 r\geq\frac{p-1}{p^n-1}.
\]
Cancellation can only increase the valuation of either sum. Thus every
nonzero \(p^n\)-torsion coordinate of both groups satisfies this bound.

Finally consider
\[
 Q_n(T)=\left.\frac{[p](U)}U\right|_{U=[p^{n-1}](T)}.
\]
Its constant term is \(p\), and it is distinguished of degree
\(D=(p^2-1)p^{2n-2}\), because \([p]\) has special initial degree
\(p^2\). Its Weierstrass roots are precisely the primitive \(p^n\)-torsion
coordinates. The constant term of its Weierstrass polynomial has
valuation \(E\), so the sum of their valuations is \(E\). Their average
is therefore
\[
 \frac ED=\frac e{p^{2n-2}}.
\]
Comparing this average with the lower bound for every root proves
\eqref{fc-eq:mainbound}.
\end{proof}

\begin{corollary}\label{fc-cor:Aupper}
For \(p=5\) and \(e_K=20\), two potentially supersingular height-two
groups with isomorphic \(125\)-torsion have equal special inertia traces.
\end{corollary}
\begin{proof}
Otherwise Proposition~\ref{fc-thm:valuation}, with \(n=3\), would give
\[
 20(125-1)\geq4\cdot5^4,
 \qquad\text{that is,}\qquad2480\geq2500.
\]
This is false. Equivalently, the primitive-torsion average is \(4/125\),
whereas every root would have to have valuation at least \(4/124\).
\end{proof}

Together with the already-proved structural and mixed comparisons and the
unequal-trace \(25\)-torsion construction, Corollary~\ref{fc-cor:Aupper} gives
\(\nu_{\min}(5,20,2)=2\) in the convention where index \(\nu\) compares
level \(p^{\nu+1}\). The pointwise estimate in Proposition~\ref{fc-thm:valuation} also combines with the central point-field criterion in Section~\ref{sec:central-ramification} to prove the general prime-to-$p$ recovery theorem and its stronger supersingular range. That argument uses actual canonical quotients and transported kernels; it does not require lifting the special graph splitting in \eqref{fc-eq:graphinput} to generic cyclic subgroups.

\section{The first wild index for every prime at least five}\label{sec:first-wild-general}
This section proves the exact index at $e=p(p-1)$ by combining the graph estimate with an actual cyclic-point construction.

\begin{proposition}\label{prop:first-wild-ss}
Let $p\ge5$, $\mu_p\subset K$ and $e_K=p(p-1)$, with $K$ complete and its residue field perfect. Two whole potentially supersingular height-two groups with isomorphic $p^3$-torsion have equal special inertia characters, regardless of their residual inertia type.
\end{proposition}
\begin{proof}
If the characters differed, Proposition~\ref{fc-thm:valuation}, which has no condition on $v_p(e_K)$, would imply
\[
                    p(p-1)(p^3-1)\ge(p-1)p^4.
\]
The left side minus the right side equals $-p(p-1)<0$. Equivalently, the primitive-coordinate average $(p-1)/p^3$ is strictly below the required minimum $(p-1)/(p^3-1)$. This excludes every residual branch.
\end{proof}

\begin{proof}[Uniform upper bound in Theorem~\ref{thm:first-wild-general}]
Pass to the completed maximal unramified base, preserving $e$ and inertia. For an allowed height-two group the determinant on inertia is $\chi_p^d\epsilon$, where $0\le d\le2$ is its total dimension and $\epsilon$ is a finite $\Z_p$-valued character of order dividing $p-1$, as in Theorem~\ref{thm:second-wild-exact}.

The cyclotomic image modulo $p^3$ has nontrivial wild part. To see this without assuming $\mu_{p^2}\not\subset K$, take a Cohen coefficient field $K_0$. The polynomial $\Phi_{p^3}(1+X)$ is monic, has constant term $p$ and reduction $X^{p^2(p-1)}$, so is Eisenstein over $\Ocal_{K_0}$. The extension $K_0(\mu_{p^3})/K_0$ is totally ramified of index $p^2(p-1)$, which cannot divide $e_K=p(p-1)$. Thus $\mu_{p^3}\not\subset K$. Since $\mu_p\subset K$, the remaining cyclotomic image is a nontrivial $p$-group.

An isomorphism of $p^3$-torsion identifies determinants modulo $p^3$. On wild inertia the finite factors are trivial, so $p\mid d-d'$. As $|d-d'|\le2<p$, one has $d=d'$. Reduction modulo $p$ then identifies $\epsilon,\epsilon'$, because reduction is injective on roots of unity of order dividing $p-1$.

For $d=0$, all potentially good graded pieces are potentially \'etale; their special character is the Brauer lift of the residual semisimplification. Hence it is determined by the common $p$-torsion. Cartier duality gives $d=2$. If $d=1$, the group either has two height-one graded pieces or is a whole potentially good height-two group. In the latter case ordinary reduction gives the connected--\'etale filtration, which descends by uniqueness. The remaining whole possibility is supersingular.

Two height-one-filtered groups are compared residually, since $\chi_p\equiv1\pmod p$ over $K$. A mixed comparison is impossible: Proposition~\ref{prop:mixed-obstruction}, with $n=3$ and $v_p(e)=1$, would require $e>p^2$, contrary to $e=p^2-p$. Two whole supersingular groups are covered by Proposition~\ref{prop:first-wild-ss}. All dimensions and filtration types are accounted for, proving the uniform upper bound two.
\end{proof}

\begin{proof}[Actual lower bound in Theorem~\ref{thm:first-wild-general}]
We specialize Theorem~\ref{thm:general-arc}, whose hypotheses do not require $p\nmid e$, and give the parameters and exact point construction. Put
\[
 k=\overline{\F}_p,\quad W=W(k),\quad q=p^2,\quad N=q-1,
 \qquad d=Np(p-1),\qquad b=p(p-1)(p^2-p-1).
\]
Use the integral universal group $F_u$ and its homogeneous logarithm from \eqref{eq:homogeneous-logarithm}. Define
\[
 Y_0=X,\qquad Y_1=[p]_{X^b}(X),\qquad
 P(X)=\Phi_{X^b}(Y_1),\qquad \Phi_u(T)=[p]_u(T)/T.
\]
The two strict leading-term comparisons are
\[
 b-(q-p)=Np(p-2)>0,\qquad Nq-d=Np>0.
\]
Modulo $p$, the first selects $T^q$ in $Y_1$, so $Y_1=X^q+O(X^{q+1})$. In $P$ the term $uY_1^{p-1}$ has degree $b+(p-1)q=d$, strictly before $Y_1^{q-1}$. Intervening coefficients are divisible by $u$ modulo $p$ and are later; higher coefficients are integral. Consequently
\[
                     P(0)=p,\qquad P(X)=X^d+O(X^{d+1})\pmod p.
\]
Weierstrass preparation gives an Eisenstein polynomial of degree $d$, so $R=W[[X]]/(P)$ is a complete discrete valuation ring with uniformizer $x$, residue field $k$ and $v_R(p)=d$. Since $q<d$, the value of $Y_1(x)$ is exactly $q$, hence is nonzero. The relation $\Phi(Y_1(x))=0$ makes $[p]Y_1(x)=0$. Thus $x$ has exact order $p^2$ in $F_{x^b}$.

Choose $\zeta\in W$ of order $N$ and let $a=\zeta^{-1}$. The congruence $b\equiv-(p-1)\pmod N$ and the exact homogeneous identities give
\[
 [p]_{a^{p-1}u}(a^{-1}T)=a^{-1}[p]_u(T),\qquad
 \Phi_{a^{p-1}u}(a^{-1}T)=\Phi_u(T).
\]
Therefore $X\mapsto\zeta X$ fixes $P$ and defines an automorphism $\sigma$ of $R$ of exact order $N$. With $L=\operatorname{Frac}R$ and $K=L^{\langle\sigma\rangle}$, the extension $L/K$ is totally tame of degree $N$, so $e_K=d/N=p(p-1)$ and the residue field is $k$. The full formal-group isomorphism $\sigma^*F\to F$, $T\mapsto aT$, follows from the homogeneous logarithm, satisfies the cocycle condition and fixes the point because $a\sigma(x)=x$. Effective descent of the generic finite levels gives an actual height-two group $G/K$ with a rational point of exact order $p^2$ and special eigenvalues $a,a^p$.

The field contains $\mu_p$. Indeed $x^N$ is a uniformizer of $K$; write $-p=(x^N)^eU$ with $U$ a unit. Hensel's lemma provides a $(p-1)$st root of $U$, and $(p-1)\mid e$. Thus $K$ contains a $(p-1)$st root of $-p$, generating the unique tame degree-$(p-1)$ extension of its Cohen field, namely the cyclotomic extension.

Finally apply the actual point-to-ambiguity construction of Lemmas~\ref{lem:point-ambiguity} and~\ref{lem:extension-lift}. More explicitly, write $\det G=\chi_p\psi$, let $C$ be the constant cyclic subgroup of order $p^2$ generated by the point, and form $G'=G/C$. The isogeny sequence gives
\[
0\longrightarrow(\Z/p^2)(\chi_p\psi)
 \longrightarrow G'[p^2]\longrightarrow\Z/p^2\longrightarrow0.
\]
Twisting by $\psi^{-1}$ gives a whole potentially supersingular group $M$ and an extension of $(\Z/p^2)(\psi^{-1})$ by $\mu_{p^2}$. The extension-lifting lemma realizes this finite extension as $H[p^2]$ for an actual height-two group $H$ with graded pieces $\mu_{p^\infty}$ and $(\Q_p/\Z_p)(\psi^{-1})$.

On a tame generator, the special eigenvalues of $M$ are $a^{-1},a^{-p}$. Those of $H$ are $1,a^{-(p+1)}$. Their trace difference is
\[
                -(a^{-1}-1)(a^{-p}-1)\ne0.
\]
Hence $M[p^2]\simeq H[p^2]$ with unequal special characters at the required absolute index. This proves the lower bound two and completes Theorem~\ref{thm:first-wild-general}.
\end{proof}

\section{Uniform recovery and exact bands at wild indices}\label{sec:all-wild-indices}
This section combines the pointwise graph bound with an equivariant cyclic-point construction whose number of degree-$p$ stages is prescribed by the divisibility of the base index.

\subsection{A graph bound for the full filtration class}
The determinant first fixes the total dimension; the mixed and supersingular bounds then exhaust the remaining comparisons.

\begin{proposition}\label{prop:graph-uniform-full}
Let $p\ge5$ and $n\ge2$. In the standing height-two filtration class, $e_K\le(p-1)p^{n-2}$ implies recovery of the special character from actual $p^n$-torsion. Consequently
\[
 \nuMin(p,e,2)\le1+\left\lceil\log_p\frac{e}{p-1}\right\rceil.
\]
This is a consequence of Theorem~\ref{thm:all-index-upper}; the following direct proof uses the pointwise graph valuation bound.
\end{proposition}
\begin{proof}
Pass to the completed maximal unramified base. This preserves $e=e_K$, inertia and the actual torsion comparison. For an allowed height-two group the determinant on inertia is $\chi_p^d\epsilon$, where $0\le d\le2$ is the total dimension of its potentially good graded pieces and $\epsilon$ is finite of order dividing $p-1$, as in the proof of Theorem~\ref{thm:first-wild-general}.

The cyclotomic image modulo $p^n$ is nontrivial. Indeed a primitive $p^n$th root satisfies
\[
 v_K(\zeta_{p^n}-1)=\frac{e}{(p-1)p^{n-1}}\le\frac1p<1,
\]
where the value follows from the Eisenstein polynomial $\Phi_{p^n}(1+X)$ over a Cohen coefficient ring. It cannot belong to the integer-valued field after unramified passage. Since $\mu_p\subset K$, that nontrivial image is a $p$-group, and restriction to wild inertia has the same image because the tame quotient has no nontrivial finite $p$-quotient. The finite determinant factors become trivial on wild inertia. Equality of determinants modulo $p^n$ therefore implies $p\mid d-d'$, and $|d-d'|\le2<p$ gives $d=d'$. Reduction modulo $p$ then identifies $\epsilon$ and $\epsilon'$ by injectivity on $\mu_{p-1}$.

If $d=0$, the special sum character is the Brauer lift of the residual semisimplification; the common $p$-torsion determines it. Cartier duality gives $d=2$. If $d=1$, each group is either whole potentially supersingular or has a two-step potentially good height-one filtration. In the ordinary case the unique connected--\'etale filtration descends, giving the second possibility. Two height-one-filtered groups compare residually: their special characters are sums of the finite factors, and $\chi_p\equiv1\pmod p$.

Put $v=v_p(e)$. The assumed bound gives
\[
                    e\le(p-1)p^{n-2}<p^{n-1},\qquad v\le n-2.
\]
A mixed comparison would require, by Proposition~\ref{prop:mixed-obstruction},
\[
                        e>p^{2n-3-v}\ge p^{n-1},
\]
which is impossible. Two whole supersingular groups with unequal special characters would require, by Proposition~\ref{fc-thm:valuation},
\[
                         e(p^n-1)\ge(p-1)p^{2n-2}.
\]
But
\[
 \begin{aligned}
 e(p^n-1)&\le(p-1)p^{n-2}(p^n-1)\\
         &=(p-1)\bigl(p^{2n-2}-p^{n-2}\bigr)
          <(p-1)p^{2n-2}.
 \end{aligned}
\]
This excludes the last case, including the endpoint of the asserted range.

For the logarithmic formulation let $t$ be the least nonnegative integer with $e/(p-1)\le p^t$. Applying the result with $n=t+2$ gives $\nuMin\le t+1$. This is an integer-threshold statement; no approximation to a logarithm is needed.
\end{proof}

\subsection{Prescribing the degree-$p$ stages}
The next construction produces actual groups at each index, including those with an arbitrarily large $p$-part.

\begin{proposition}\label{prop:mixed-stage-construction}
Let $p\ge5$, $n\ge2$ and $0\le t\le n-2$. Suppose $e$ is a positive multiple of $(p-1)p^t$ and
\[
                             e>p^{2n-3-t}.
\]
There are actual groups $M,H$ over a complete field of residue $\overline{\F}_p$, exact absolute index $e$ and containing $\mu_p$, with $M[p^n]\simeq H[p^n]$ and different special characters. The group $M$ is whole potentially supersingular, while $H$ has a two-step potentially good height-one filtration.
\end{proposition}
\begin{proof}
Set
\[
 q=p^2,\qquad N=q-1,\qquad j=n-1-t\ge1,\qquad d=Ne,
 \qquad x_0=\frac e{p^{2n-2-t}}.
\]
Choose an integer $A$ with $x_0<A<px_0$ and $N\nmid A$. Such a choice exists because $px_0>1$. If $x_0<1$, take $A=1$; otherwise take $\lfloor x_0\rfloor+1$ and increase it by one if divisible by $N$. The result is at most $x_0+2<px_0$, since $p\ge5$.

Put
\[
                       b=\frac{Ne}{p^t}-(p-1)Ap^{2j}.
\]
It is integral, and the two strict comparisons are
\begin{align}
 b-(q-p)Aq^{j-1}
   &=\frac N{p^t}\bigl(e-Ap^{2n-3-t}\bigr)>0,
       \label{eq:stage-first}\\
 (q-p)Aq^j-b
   &=\frac N{p^t}\bigl(Ap^{2n-2-t}-e\bigr)>0.
       \label{eq:stage-second}
\end{align}
In particular $b>0$. Also
\begin{equation}\label{eq:stage-descent-congruence}
                         b\equiv-(p-1)A\pmod N,
\end{equation}
because $q^j\equiv1\pmod N$ and $p^t\mid e$.

Over $W=W(\overline{\F}_p)$ use the integral height-two formal group $F_u$ and the full homogeneous identities of~\eqref{eq:homogeneous-logarithm}. Define
\[
 Y_0=X^A,\qquad Y_{i+1}=[p]_{X^b}(Y_i),\qquad
 P(X)=\Phi_{X^b}(Y_{n-1}),\qquad \Phi_u(T)=[p]_u(T)/T.
\]
These are integral power series. In the first $j$ compositions the selected degree is
\[
                             h_i=Aq^i\quad(0\le i\le j).
\]
Indeed~\eqref{eq:stage-first} says $b>(q-p)h_{j-1}$, hence the same inequality holds for all earlier inputs. It makes $T^q$ strictly earlier than $uT^p$. Intermediate coefficients contain $u$ modulo $p$, so are later; degrees above $q$ are later by integrality.

For the remaining $t$ compositions, put $h_{i+1}=b+ph_i$, starting from $h_j=Aq^j$. Inequality~\eqref{eq:stage-second} gives $b<(q-p)h_j$, and the $h_i$ strictly increase. Thus $uT^p$ is strictly first at every remaining step. All selected coefficients are one modulo $p$, so
\[
 Y_i\equiv X^{h_i}+O(X^{h_i+1})\pmod p,
 \qquad H:=h_{n-1}=p^tAq^j+b\frac{p^t-1}{p-1}.
\]
For $t=0$ the last formula is $H=Aq^j$ and there are no degree-$p$ compositions before the final divided series.

In $P$ the term $uY_{n-1}^{p-1}$ has degree
\begin{equation}\label{eq:stage-eisenstein-degree}
 b+(p-1)H=p^tb+(p-1)p^tAq^j=Ne=d.
\end{equation}
Its competitor $Y_{n-1}^{q-1}$ is later because $b<(q-p)H$, by~\eqref{eq:stage-second} and $H\ge Aq^j$. The intermediate terms have a factor $u$ modulo $p$ and are later than the selected term. Consequently
\[
                    P(0)=p,\qquad P(X)\equiv X^d+O(X^{d+1})\pmod p.
\]
Weierstrass preparation gives an Eisenstein ring $R=W[[X]]/(P)$ with uniformizer $x$ and $v_R(p)=d$. Equation~\eqref{eq:stage-eisenstein-degree} implies $H<d$, and all earlier $h_i$ are smaller. Coefficients divisible by $p$ therefore do not affect the values $v_R(Y_i(x))=h_i$. In particular $Y_{n-1}(x)\ne0$, while $\Phi(Y_{n-1}(x))=0$ makes its $p$-multiple zero. Thus $x^A$ has exact order $p^n$.

Choose $\zeta\in W$ of order $N$ and set $a=\zeta^{-A}$. Congruence~\eqref{eq:stage-descent-congruence} and the exact homogeneous identities show that $X\mapsto\zeta X$ carries $u$ to $a^{p-1}u$, every $Y_i$ to $a^{-1}Y_i$, and $P$ to itself. Hence $\sigma(x)=\zeta x$ defines an automorphism of $R$ of exact order $N$. For $L=\operatorname{Frac}R$ and $K=L^{\langle\sigma\rangle}$, the extension $L/K$ is totally tame of degree $N$. Thus $e_K=d/N=e$ and the residue field is $\overline{\F}_p$.

The full formal-group isomorphism $\sigma^*F\to F$, $T\mapsto aT$, satisfies the cocycle identity and fixes the marked point because $a\sigma(x^A)=x^A$. Descent of the compatible finite generic levels gives an actual height-two group $G/K$ with that rational point. Its special eigenvalues are $a,a^p$, and $a\ne1$ because $N\nmid A$. This suffices; $a$ need not have exact order $N$.

The field $K$ contains $\mu_p$: write $-p=\pi_K^eU$, extract a $(p-1)$st root of the unit $U$ by Hensel's lemma, and use $(p-1)\mid e$. The resulting root of $-p$ generates the unique tame degree-$(p-1)$ extension of the Cohen field, namely its cyclotomic extension.

Finally Lemmas~\ref{lem:point-ambiguity} and~\ref{lem:extension-lift} supply actual $M,H$. Concretely, quotient $G$ by the marked constant cyclic subgroup and twist by the inverse finite determinant factor $\psi^{-1}$. The finite extension is then an extension of $(\Z/p^n)(\psi^{-1})$ by $\mu_{p^n}$. Compatible Kummer lifts and prime-to-$p$ averaging give the height-one filtered group $H$ with that same truncation. Its special eigenvalues are $1,a^{-(p+1)}$, while those of $M$ are $a^{-1},a^{-p}$. Their trace difference
\[
                         -(a^{-1}-1)(a^{-p}-1)
\]
is nonzero. This proves the exact field and actual-group assertions.
\end{proof}

\subsection{The sharp mixed criterion and full-class bands}
The number of degree-$p$ stages can be chosen from $v_p(e)$, and the uniform recovery range makes the resulting lower examples sharp on explicit infinite bands.

\begin{proof}[Proof of Theorem~\ref{thm:mixed-all-indices}]
Necessity is Proposition~\ref{prop:mixed-obstruction}. For sufficiency put $v=v_p(e)$ and $t=\min(v,n-2)$. If $v\le n-2$, the desired strict inequality is exactly the hypothesis of Proposition~\ref{prop:mixed-stage-construction}. If $v\ge n-1$, then
\[
                 e\ge(p-1)p^v\ge(p-1)p^{n-1}>p^{n-1}
                    =p^{2n-3-t},
\]
so that proposition again applies. In both cases $p^t(p-1)\mid e$, and its construction has all the required properties.
\end{proof}

\begin{proof}[Theorem~\ref{thm:wild-bands-exact} for $2\le k\le p+1$]
In this paragraph assume $2\le k\le p+1$. Put $e=kp^r(p-1)$ and $n=r+2$. Use Proposition~\ref{prop:mixed-stage-construction} with $t=r=n-2$. Its strict threshold is $e>p^{r+1}$, which follows from $k(p-1)>p$. It supplies an actual mixed pair with isomorphic $p^{r+2}$-torsion and unequal special characters, so $\nuMin\ge r+2$.

The construction permits $A=1$ uniformly, since
\[
                \frac{k(p-1)}{p^2}<1<\frac{k(p-1)}p.
\]
Its deformation exponent is $b=(p-1)((k-1)p^2-k)>0$, independent of $r$. There is one degree-$p^2$ stage, then $r$ degree-$p$ stages, before the final divided series. Thus each index is attained directly, without requiring that a lower example survive a prescribed base change.

For $k\le p$, Proposition~\ref{prop:graph-uniform-full} at level $r+3$ gives the upper bound, since $e\le(p-1)p^{r+1}$. For the entire range $k\le p+1$, Theorem~\ref{thm:all-index-upper} applies because
\[
 e\le(p^2-1)p^r<p^{r+2}.
\]
Thus $\nuMin\le r+2$ in this range, proving equality. In particular this endpoint is $\nuMin(p,(p^2-1)p^r,2)=r+2$. The corresponding finite band follows by $r=1$, and the diagonal formula follows by $k=p$, $r=s-1$. At $k=p$, the valuation of $e$ is $r+1$ rather than $r$; the direct construction above remains valid because it requires only $(p-1)p^r\mid e$. The larger coefficient interval, with its distinction between coefficients prime to $p$ and divisible by $p$, is treated in Section~\ref{sec:extended-wild-bands}.
\end{proof}

\section{Primitive coordinates and uniform recovery}\label{sec:one-slope-all-indices}
This section compares the reduced supports of scaled torsion coordinates, proves the uniform bound and one-slope results of Section~1, and isolates the remaining central comparison at wild indices.

Throughout, $p\ge5$. After completing the maximal unramified base, take the common tame good field $L/K$ of degree $m=p^2-1$. Write $R=\Ocal_L$, choose $\tau(\pi)=\zeta\pi$ with $\zeta$ of order $m$, and normalize $v_L(\pi)=1$. Put $e=e_K$ and $v=v_p(e)$, so $v_L(p)=me$. We use the good formal models and exact tame eigenparameters of Sections~\ref{sec:quotients} and~\ref{sec:restricted}. There is no assumption that $p\nmid e$.

\subsection{Minimum values and the scaled graph}
The torsion valuations determine the multiplicities of zero in the special coordinate algebras.

\begin{lemma}\label{os-lem:minima}
For a one-slope model, every primitive $p^n$-point has value $e/p^{2n-2}$; exactly $p^{2n-2}$ points of the full $p^n$-torsion, including zero, have larger value. For a two-slope model of parameter $z$, the minimum nonzero value is $z/p^{2n-1}$; exactly $p^{2n-1}$ points, including zero, have larger value.
\end{lemma}
\begin{proof}
In the one-slope case the first-level values are all $e$. For a nonzero torsion target of value at most $e$, the degree-$p$ and degree-one coefficients of its division equation lie above the segment ending in degree $p^2$. All $p^2$ divisions therefore have the target value divided by $p^2$. Induction gives the primitive values, and the higher-valued points are exactly the $p^{n-1}$-torsion.

In the two-slope case the first-level minimum is $z/p$, attained at $p^2-p$ points. Every minimum target divides on the degree-$p^2$ branch. A larger target below the middle Newton break has strictly larger division values, and above that break all division values are at least $z/p$. Divisions of zero cannot lower the minimum. Thus the minimum occurs at $(p^2-p)p^{2n-2}$ points, leaving $p^{2n-1}$ larger-valued points. The Newton range of Lemma~\ref{rc-lem:newton} is preserved because multiplication by $p$ strictly raises a positive nonzero torsion value. These formulas agree with the elliptic valuation and counting formulas of \citet[Theorem~4.4 and Corollaries~4.5--4.6, pp.~10--11]{Smith2023}; the argument here applies to the formal models under consideration.
\end{proof}


Let $f:G[p^n]\simeq H[p^n]$ and put $D=p^{2n-1}$. Make the Eisenstein extension
\[
 R'=R[\varpi]/(\varpi^D-\pi).
\]
For a one-slope coordinate assign exponent $a=pe$, and for a two-slope coordinate assign $a=z$. If $x,y$ are the coordinates in the generic graph algebra, define
\[
 X=x/\varpi^a,\qquad Y=y/\varpi^b,\qquad B'=R'[X,Y].
\]
All geometric values are integral by Lemma~\ref{os-lem:minima}. Their characteristic polynomials have coefficients in the integrally closed DVR $R'$, so $B'$ is finite and torsion-free, hence finite free. Scaling is invertible generically; its rank is therefore $p^{2n}$.

Every nonlinear term of total degree $r$ in the scaled first formal law has the extra factor $\varpi^{a(r-1)}$, and similarly for the second law. The series converge in the finite complete algebra $B'\otimes_{R'}B'$. They give integral coproducts, inverses and counits. Thus
\[
 B_0=B'/\varpi B'
\]
is a Hopf algebra generated by primitive $X,Y$. The classical primitive-presentation lemma, proved in full as Lemma~\ref{pg-lem:primitives}, gives
\begin{equation}\label{os-eq:primitive-dimension}
 V=\operatorname{span}_k\{X,X^p,\ldots,Y,Y^p,\ldots\},
 \qquad \dim_k V=2n.
\end{equation}

The group scheme $\mu_D$ acts by $\varpi\mapsto t\varpi$, fixing the original graph coordinates. It preserves $B'$ and gives a $\Z/D\Z$ grading on $B_0$, including in characteristic $p$. The weights of the Frobenius powers are
\[
 \operatorname{wt}(X^{p^j})=-ap^j,\qquad
 \operatorname{wt}(Y^{p^j})=-bp^j\pmod D.
\]
For an exponent of $p$-adic valuation $s$, there are exactly $\max(0,2n-1-s)$ distinct nonzero weights before the chain reaches zero.

\subsection{Reduced supports and stable Frobenius spans}
We show that passage to the graph preserves every reduced coordinate value; the resulting support sizes distinguish the two polygon types.

\begin{lemma}\label{os-lem:stable-spans}
In the scaled graph, the stable cyclic Frobenius span of a one-slope coordinate has dimension two; that of a two-slope coordinate has dimension one. Every Frobenius power of either coordinate is nonzero.
\end{lemma}
\begin{proof}
For one coordinate alone, $C'=R'[X]$ is monogenic of rank $p^{2n}$. Its generic coordinate has that many distinct values. The monic polynomial with those roots has integral coefficients; monic division proves that it is the exact presentation of $C'$. The scaled group law makes the special coordinate primitive.

Its monic special polynomial $P(U)$ is additive. Indeed $P(U+W)-P(U)-P(W)$ vanishes in the tensor square of the quotient and has degree below $\deg P=p^{2n}$ in each variable. The normal monomial basis makes it identically zero. The binomial-coefficient criterion then leaves only $p$-power exponents.

For a one-slope coordinate, Lemma~\ref{os-lem:minima} makes the zero-root multiplicity exactly $p^{2n-2}$. Hence
\[
 P(U)=U^{p^{2n}}+a_1U^{p^{2n-1}}+a_0U^{p^{2n-2}},
 \qquad a_0\ne0.
\]
Over the perfect field $k$, this is a $p^{2n-2}$th power of a separable additive polynomial of degree $p^2$. Its distinct roots form a two-dimensional $\F_p$-space. For a two-slope coordinate the zero multiplicity is $p^{2n-1}$, so
\[
 P(U)=U^{p^{2n}}+b_0U^{p^{2n-1}},\qquad b_0\ne0,
\]
and its distinct roots form an $\F_p$-line.

The extension $C'\subset B'$ is finite integral. Lying over makes its map on spectra surjective, including on the closed fibre. Thus all these distinct coordinate values survive in the graph, even though the map of special coordinate algebras need not be injective.

Let $F$ be $p$th power on $V$. On the two-dimensional reduced coordinate support, the functions $U,U^p$ are linearly independent: a relation $U^p=cU$ could have at most $p$ roots. The separable degree-$p^2$ relation expresses $U^{p^2}$ in their span with nonzero linear coefficient, making that span $F$-bijective. Some Frobenius power kills the nilradical of $B_0$. Applying it to the support relation gives a two-dimensional stable cyclic span in $B_0$ itself; evaluation proves that its dimension cannot drop. On the one-dimensional coordinate support the same argument gives stable dimension one. Nonzero support values show that every Frobenius power is nonzero. The stable dimensions remain the same for every tail of a chain.
\end{proof}

\subsection{Comparing different polygons}
The different stable dimensions prevent a common nonzero weight from identifying the two coordinate chains.

\begin{proof}[Proof of Theorem~\ref{thm:mixed-polygon-obstruction}]
Let $X$ be the one-slope coordinate and $Y$ the two-slope coordinate, with $w=v_p(z)$. The respective numbers of distinct nonzero weights are
\[
 d_X=\max(0,2n-2-v),\qquad
 d_Y=\max(0,2n-1-w).
\]
At an overlapping nonzero weight the two vectors are independent. Otherwise $X^{p^i}=cY^{p^j}$ with $c\ne0$, and successive Frobenius powers would identify stable cyclic spans of dimensions two and one, contrary to Lemma~\ref{os-lem:stable-spans}. Every remaining nonzero weight contributes a nonzero vector, and different weights are independent. Weight zero contains the two-dimensional stable span of $X$. Consequently
\begin{equation}\label{os-eq:dimension-inequality}
 2n\ge d_X+d_Y+2.
\end{equation}
If $w=0$, then $d_Y=2n-1$, an impossibility. Thus $w\ge1$. If $v\le2n-2$ and $w<2n-1$, substitution gives $v+w\ge2n-1$. If $w\ge2n-1$, that inequality is automatic. When $v>2n-2$, the condition $w\ge2n-1-v$ is already weaker than $w\ge1$. Therefore
\[
 w\ge\max\{1,2n-1-v\}.
\]
Finally $z\ge p^w$ and $z<pe$ give
\[
 e>z/p\ge p^{w-1}\ge p^{2n-2-v}.
\]
This includes negative exponents on the right and proves both assertions without a hypothesis on the special characters.
\end{proof}

\subsection{Comparing two one-slope models}
Equal nonzero weights now distinguish different tangent characters, while a single coordinate supplies two dimensions at weight zero.

\begin{proof}[Proof of Theorem~\ref{thm:one-slope-all-indices}]
Choose exact tame eigenparameters with tangent scalars $\alpha_1,\alpha_2\in\mu_m$. Choose $rD\equiv1\pmod m$ and extend the tame generator by $\tau(\varpi)=\zeta^r\varpi$. Both scaling exponents are $pe$, so the scaled tangent eigenvalues are
\[
 \lambda_i=\alpha_i\zeta^{-rpe}.
\]
Suppose $\alpha_1\ne\alpha_2$. Their $j$th Frobenius powers have distinct tame eigenvalues $\lambda_1^{p^j},\lambda_2^{p^j}$, since Frobenius on $k$ is injective. Put $d=2n-2-v$. The hypothesis $n\ge v+2$ gives $d\ge n$. For $j=0,\ldots,d-1$, each of the $d$ distinct common nonzero weights therefore contributes two independent vectors. Weight zero contains a two-dimensional stable span by Lemma~\ref{os-lem:stable-spans}. Equation~\eqref{os-eq:primitive-dimension} would imply
\[
 2n\ge2d+2\ge2n+2,
\]
a contradiction. Thus $\alpha_1=\alpha_2$, which gives equality of their unordered Frobenius orbits and hence of the special characters. No numerical upper bound on $e$ was used.
\end{proof}

\subsection{Matching canonical lines}
For two two-slope models, different kernels of the leading coordinate maps provide the independence that the polygon comparison supplied above.

\begin{lemma}[A necessary condition for unmatched canonical lines]\label{os-lem:canonical-matching}
Let $n\ge1$ and let $f:G[p^n]\simeq H[p^n]$ compare two two-slope models of parameters $z_1,z_2$, with canonical $p$-lines $C_1,C_2$. Put $s_i=v_p(z_i)$. If $f(C_1)\ne C_2$, then
\[
 \min(s_1,2n-1)+\min(s_2,2n-1)\ge2n.
\]
In particular, $s_1,s_2\ge1$ and $s_1+s_2\ge2n$. Thus $s_1+s_2<2n$ forces the given isomorphism to identify the canonical lines. There is no bound on $e$ and no hypothesis that the special characters differ.
\end{lemma}
\begin{proof}
Use the scaled graph with exponents $z_1,z_2$ and $D=p^{2n-1}$. A geometric $p^n$-point $P$ has first scaled coordinate zero on specialization precisely when $p^{n-1}P\in C_1$. Indeed the minimum value $z_1/D$ is attained exactly by successive low divisions of the noncanonical first-level points. Every division of such a point remains on the low branch. There are $(p^2-p)p^{2n-2}$ such points, and the remaining points have strictly larger value, as in Lemma~\ref{os-lem:minima}. The corresponding assertion for the second coordinate is $p^{n-1}f(P)\in C_2$.

After extension to a splitting valuation ring, the scaled coordinate specialization maps are additive. Each image is a nonzero $\F_p$-line by Lemma~\ref{os-lem:stable-spans}. Choose a nonzero basis in each line. The resulting maps to $\F_p$ have kernels
\[
 \{P:p^{n-1}P\in C_1\},\qquad
 \{P:p^{n-1}P\in f^{-1}(C_2)\}.
\]
Multiplication by $p^{n-1}$ maps the common torsion group onto its two-dimensional $\F_p$-space of $p$-points. If the canonical lines do not match, these kernels are distinct. The two maps are therefore independent, and their joint image is all of $\F_p^2$.

Consequently no Frobenius power of $X$ is proportional to a Frobenius power of $Y$ in $B_0$. Evaluation on this full product would identify two independent coordinate functions: a $p$th power multiplies each chosen coordinate basis by a nonzero scalar and fixes its $\F_p$ argument. This also proves that sufficiently high powers of $X$ and $Y$, whose weights are both zero, are independent.

The two chains have respectively
\[
 d_i=\max(0,2n-1-s_i)
\]
distinct nonzero weights. Every overlapping weight contributes two independent vectors, every other weight contributes one, and weight zero contributes two. Equation~\eqref{os-eq:primitive-dimension} gives
\[
 2n\ge d_1+d_2+2
   =4n-\min(s_1,2n-1)-\min(s_2,2n-1).
\]
This is the asserted inequality. If either $s_i$ were zero, its left-hand truncated sum would be at most $2n-1$; hence both are positive. Removing the truncations gives the remaining assertions.
\end{proof}

\subsection{The full filtration class in the remaining range}
The cyclotomic determinant fixes dimensions before the mixed and polygon obstructions are applied.

\begin{proof}[Proof of Corollary~\ref{cor:all-index-two-slope-reduction}]
The assumptions and $e\ge(p-1)p^v$ give
\[
 p-1\le p^{2n-3-2v},\qquad v\le n-2.
\]
A primitive $p^n$th root satisfies
\[
 v_K(\zeta_{p^n}-1)=\frac{e}{(p-1)p^{n-1}}.
\]
This follows from the Eisenstein polynomial $\Phi_{p^n}(1+T)$ over a Cohen coefficient ring: its constant term is $p$ and its reduction is $T^{(p-1)p^{n-1}}$. The displayed value is not an integer, since its $p$-adic valuation is $v-n+1<0$. Thus the completed maximal unramified base does not contain $\mu_{p^n}$. Because it contains $\mu_p$, the cyclotomic image modulo $p^n$ is a nontrivial $p$-group, already attained by wild inertia.

The determinants have the form $\chi_p^d\epsilon$, with $0\le d\le2$ and $\epsilon$ of order dividing $p-1$. Equality of truncated determinants on wild inertia forces $p\mid d-d'$, hence $d=d'$. Reduction modulo $p$ then identifies the finite factors. Dimension zero is recovered by the residual Brauer character, and Cartier duality treats dimension two. In dimension one, uniqueness descends an ordinary connected--\'etale filtration, leaving the alternatives of two potentially good height-one graded pieces or a whole potentially supersingular group. Two filtered groups compare residually. A mixed unequal-character pair is excluded by Theorem~\ref{thm:mixed-all-indices}.

For two whole supersingular groups, a comparison involving different polygons contradicts Theorem~\ref{thm:mixed-polygon-obstruction}, because its necessary bound is larger by a factor of $p$. Two one-slope groups have equal characters by Theorem~\ref{thm:one-slope-all-indices}, since $n\ge v+2$. Only two two-slope models remain, as asserted.
\end{proof}

\subsection{Two Frobenius chains at arbitrary ramification}

We first isolate the dimension calculation that applies without a
prime-to-$p$ hypothesis on the absolute ramification index.  Retain the
common tame good field, the exact tame eigenparameters, and the scaled
actual graph constructed above.  Thus
\[
 m=p^2-1,\qquad D=p^{2n-1},\qquad
 \tau(x_i)=\alpha_i x_i,
 \qquad \rho D\equiv1\pmod m,
\]
and the scaled coordinates have tame eigenvalues
\[
 \lambda_i=\alpha_i\zeta^{-\rho z_i}.
\]
We use the same notation for prime-to-$p$ roots of unity and their
reductions.  The inverse $\rho$ is unrelated to a remaining torsion
level.

\begin{lemma}\label{os-lem:two-chains}
Suppose that two whole potentially supersingular groups have two-slope
parameters $z_1,z_2$, unequal special characters, and isomorphic actual
$p^n$-torsion.  Put $M=2n-1$, and write, after relabeling,
\[
 z_1=p^s a,\qquad z_2=p^t b,
 \qquad p\nmid ab,\qquad 0\le s\le t<M.
\]
Then $\lambda_1,\lambda_2\in\mathbf F_p^{\times}$.  If
$\lambda_1\ne\lambda_2$, then $s+t\ge2n$.  If
$\lambda_1=\lambda_2$, then $a\ne b$ and
\begin{equation}\label{os-eq:two-chain-divisibility}
 (p-1)p^{\max\{0,\,2n-1-s-t\}}\mid a-b.
\end{equation}
These conclusions impose no condition that $p\nmid e$.
\end{lemma}

\begin{proof}
The rescaled graph is a finite free order of rank $p^{2n}$, with
primitive special coordinates $X,Y$.  Its primitive space $V$ satisfies
\eqref{os-eq:primitive-dimension}, so $\dim V=2n$.  The diagonalizable
group scheme $\mu_D$ gives the weights
\[
 \operatorname{wt}(X^{p^j})=-z_1p^j\pmod D,
 \qquad
 \operatorname{wt}(Y^{p^j})=-z_2p^j\pmod D.
\]
This is a group-scheme grading, including in characteristic $p$; it
does not assume that $\mu_D$ has $D$ distinct geometric points.
Every displayed Frobenius iterate is nonzero, as in
Lemma~\ref{os-lem:stable-spans}: a vanishing iterate would make a
positive power of a rescaled coordinate divisible by the uniformizer,
contradicting evaluation at a point where that coordinate is a unit.

We give the tame leading-coefficient calculation explicitly.  By
Lemma~\ref{os-lem:minima}, a two-slope coordinate has exactly
$(p-1)D$ points of minimum value $z_i/D$, and exactly $D$ points,
including zero, of larger value.  In the extension with uniformizer
$u$, $u^D=\pi_L$, consider the monogenic order generated by
$X_i=u^{-z_i}x_i$.  It is finite free of rank $p^{2n}$ and has a
monic defining polynomial $P_i$ of that degree.  Its rescaled formal
addition has nonlinear coefficients divisible by $u$, so the
reduction of $X_i$ is primitive.  Hence $\overline P_i$ is additive.
Indeed, the mixed terms of
$\overline P_i(U+V)-\overline P_i(U)-\overline P_i(V)$ have degree
less than $\deg P_i$ in each variable; their vanishing modulo
$(\overline P_i(U),\overline P_i(V))$ forces their coefficients to
vanish.  The constant coefficient is zero because the identity is a
root.

Exactly $D$ rescaled roots have positive valuation, while the other
$(p-1)D$ roots are units.  Thus zero has multiplicity exactly $D$ in
$\overline P_i$.  Additivity, and the fact that $D$ and $pD$ are
consecutive powers of $p$, give
\[
 \overline P_i(T)=T^{pD}+b_iT^D,\qquad b_i\ne0.
\]
Tame inertia acts trivially on the residue field and sends $X_i$ to
$\lambda_iX_i$.  Comparing the two nonzero coefficients therefore
gives $\lambda_i^{(p-1)D}=1$.  Injectivity of Frobenius in the residue
field implies $\lambda_i^{p-1}=1$, proving
$\lambda_i\in\mathbf F_p^{\times}$.

The first chain has $M-s$ distinct nonzero weights, followed by
nonzero vectors of weight zero.  Its span consequently has dimension
at least $M-s+1$; the other chain similarly has dimension at least
$M-t+1$.  Since each $\lambda_i$ lies in $\mathbf F_p$, its entire
Frobenius chain belongs to the $\lambda_i$ tame eigenspace.  If the
eigenvalues differ, these spans are disjoint, and
\[
 (M-s+1)+(M-t+1)\le2n.
\]
This is $s+t\ge2n$.

Suppose henceforth that the leading eigenvalues agree.  If $a=b$,
then
\[
 \alpha_2
 =\lambda_1\zeta^{\rho p^t a}
 =\bigl(\lambda_1\zeta^{\rho p^s a}\bigr)^{p^{t-s}}
 =\alpha_1^{p^{t-s}}.
\]
The special Frobenius orbits would agree, contradicting the assumed
inequality of characters.  Hence $a\ne b$; put
$\ell=v_p(a-b)$.

A nonzero weight of the two chains can coincide only at equal
$p$-adic order $w$, where $t\le w<M$.  At that order the equality
is equivalent to
\[
 a\equiv b\pmod {p^{M-w}}.
\]
Thus at most $\min\{\ell,M-t\}$ nonzero weights are shared.  The
number of distinct represented nonzero weights, together with at
least one weight-zero vector, gives
\[
 (M-s)+(M-t)-\min\{\ell,M-t\}+1\le\dim V=M+1.
\]
It follows that $\ell\ge M-s-t$ whenever the latter integer is
positive.

It remains to retain the factor $p-1$.  The actual residual
isomorphism identifies determinant characters.  Since $\mu_p\subset
K$, the residual cyclotomic factor is trivial; reduction is injective
on prime-to-$p$ roots of unity.  The two finite determinant factors
are therefore equal.  For the special tangent characters this says
\[
 \alpha_1^{p+1}=\alpha_2^{p+1}.
\]
Equality of $\lambda_1,\lambda_2$ now implies
$\zeta^{\rho(p+1)(z_1-z_2)}=1$.  Since $\rho$ is invertible modulo
$m=(p-1)(p+1)$, we obtain $(p-1)\mid z_1-z_2$.  Reducing
$p^s a-p^t b$ modulo $p-1$ yields $(p-1)\mid a-b$.
Combining this coprime divisibility with the preceding power of $p$
proves \eqref{os-eq:two-chain-divisibility}.
\end{proof}


\subsection{The uniform bound at every index}
The two-chain divisibility forces a contradiction in the larger uniform range, without using a bound on the ramification of a chosen point field.

\begin{proof}[Proof of Theorem~\ref{thm:all-index-upper}]
Let $e\le p^{n-1}$ and put $v=v_p(e)$. Since $e\ge(p-1)p^v>p^v$, one has $v\le n-2$. Hence
\[
 n\ge v+2,\qquad e\le p^{n-1}\le p^{2n-3-v}.
\]
Corollary~\ref{cor:all-index-two-slope-reduction} treats the determinant, dimension, filtered, mixed and one-slope cases. An unequal-character comparison would therefore consist of two whole supersingular two-slope models. Their parameters satisfy
\[
 0<z_i<pe\le p^n.
\]
After relabeling write $z_1=p^s a$, $z_2=p^t b$, with $p\nmid ab$ and $0\le s\le t\le n-1<2n-1$. Lemma~\ref{os-lem:two-chains} applies at this original torsion level. Different leading eigenvalues would force $s+t\ge2n$, contrary to $s+t\le2n-2$.

For equal leading eigenvalues, the lemma gives $a\ne b$ and
\[
 (p-1)p^{2n-1-s-t}\mid a-b.
\]
The exponent is at least $n-s$, whereas positivity and the strict upper bounds on $z_i$ give
\[
 0<|a-b|<p^{n-s}.
\]
A nonzero integer of that size cannot be divisible by $(p-1)p^{n-s}$. This contradiction excludes the last case and proves the uniform recovery assertion for every $n\ge2$.

For its index formulation put $N=\lceil\log_p e\rceil$. Since $e\ge p-1>1$, one has $N\ge1$. Apply the assertion with $n=N+1$ to obtain $\nuMin(p,e,2)\le N$. The graph range in Proposition~\ref{prop:graph-uniform-full} follows because $(p-1)p^{n-2}<p^{n-1}$, and its direct proof gives an additional route in that subrange.
\end{proof}

\subsection{High Hasse valuation in the proposed range}
A sharper proposed recovery range forces at least one Hasse parameter to have two more factors of $p$ than the absolute index.

\begin{lemma}\label{os-lem:high-hasse}
Let $v=v_p(e)$, $n\ge2$, and $e\le p^{2n-3-v}$.  In an
unequal-character actual $p^n$-torsion comparison of two two-slope
whole potentially supersingular groups,
\[
 \max\{v_p(z_1),v_p(z_2)\}\ge v+2.
\]
\end{lemma}

\begin{proof}
Admissibility gives $e\ge(p-1)p^v>p^v$, so the assumed bound implies
$n\ge v+2$.  Also
\[
 0<z_i<pe\le p^{2n-2-v}<p^{2n-1}.
\]
Thus Lemma~\ref{os-lem:two-chains} applies.  Suppose, to the contrary,
that its ordered valuations satisfy $s\le t\le v+1$.
Different leading eigenvalues would require $s+t\ge2n$, whereas
$s+t\le2v+2<2n$.

For equal leading eigenvalues, write $z_1=p^sa$ and $z_2=p^tb$ as
in that lemma.  Its nonzero difference satisfies
\[
 0<|a-b|<\frac{pe}{p^s}\le p^{2n-2-v-s}.
\]
But \eqref{os-eq:two-chain-divisibility} makes that integer divisible
by $(p-1)p^{2n-1-s-t}$, and $t\le v+1$ gives
\[
 (p-1)p^{2n-1-s-t}\ge(p-1)p^{2n-2-v-s}
 >p^{2n-2-v-s}.
\]
This contradiction proves the assertion.
\end{proof}

\subsection{The central terminal comparison}

The next argument specifies exactly the comparison that remains
when the proposed sharper numerical range does not already force
recovery.  It does not retain that range after a loss of torsion
levels.

\begin{proof}[Proof of Theorem~\ref{thm:all-index-central-reduction}]
Put $v=v_p(e)$ and suppose $e\le p^{2n-3-v}$.  An
unequal-character comparison reduces by
Corollary~\ref{cor:all-index-two-slope-reduction} to two whole
potentially supersingular groups with two slopes.  By
Lemma~\ref{os-lem:high-hasse}, choose one side with parameter $z$
satisfying $v_p(z)\ge v+2$.

The minimum-value formula and Proposition~\ref{fc-thm:valuation} give
\[
 z\ge\frac{(p-1)p^{2n-1}}{p^n-1}>(p-1)p^{n-1}.
\]
Let $c$ be the least nonnegative integer with $p^cz>e$.  Such an
integer exists, and equality $p^jz=e$ never occurs: its two sides
would have different $p$-adic valuations.  Moreover
\[
 p^{n-v-2}z>(p-1)p^{2n-v-3}\ge e,
\]
so $c\le n-v-2$.

At each stage before $c$, the chosen parameter is strictly below
$e$.  The canonical quotient formula in
Proposition~\ref{rc-prop:quotients} therefore multiplies it by $p$.
Quotient the partner by the transported actual $K$-stable kernel.
The actual finite-level quotient construction gives a comparison
one level lower, and the $K$-isogenies preserve the special rational
inertia characters.  The transported kernel is not presumed to be
canonical in the partner.

After exactly $c$ such steps, the remaining level is $p^r$, where
$r=n-c\ge v+2$.  The chosen parameter is $z_c=p^cz$.  If $c=0$,
its original two-slope range gives $z_c<pe$; if $c>0$, the preceding
parameter was strictly below $e$, again giving $z_c<pe$.  Hence
\[
 e<z_c<pe.
\]
Finally $v_p(z_c)\ge v+2+c$, whereas
$v_p((p+1)e)=v$.  Therefore
\[
 v_p\bigl((p+1)e-z_c\bigr)=v.
\]
The partner remains whole potentially supersingular under these
isogenies, but its terminal polygon is unrestricted.  The original
inequality involving $n$ has not been asserted with the smaller
integer $r$.  These are exactly the stated terminal conclusions.
\end{proof}

\subsection{The full-class index and the supersingular index}

We now separate the mixed obstruction from the remaining whole
supersingular comparison.  Define $\nu_{\mathrm{SS}}(p,e)$ to be the
least integer $a\ge0$ for which actual $p^{a+1}$-torsion determines
the special character uniformly among two whole potentially
supersingular groups over all standing fields of absolute index $e$.
Its finiteness follows from the uniform upper theorem.  Both this
index and $\nu_{\min}(p,e,2)$ use universal upper quantifiers; an
existential lower example at the specified index is therefore enough
to disprove recovery at a given level.

\begin{proof}[Proof of Theorem~\ref{thm:full-supersingular-maximum}]
Write $v=v_p(e)$ and
\[
 R=\left\lceil\frac{1+\log_p e+v}{2}\right\rceil,
 \qquad T=\max\{\nu_{\mathrm{SS}}(p,e),R\}.
\]
Since $e\ge(p-1)p^v$, one has $R\ge v+1$.  For $n=T+1$ it
follows that $n\ge v+2$ and
\[
 e\le p^{2R-1-v}\le p^{2T-1-v}=p^{2n-3-v}.
\]
An unequal-character full-class comparison at level $p^n$ would
therefore reduce, by
Corollary~\ref{cor:all-index-two-slope-reduction}, to a comparison
of two whole potentially supersingular groups.  This contradicts
$T\ge\nu_{\mathrm{SS}}(p,e)$, by restriction of torsion if necessary.
Consequently $\nu_{\min}(p,e,2)\le T$.

Conversely, inclusion of the supersingular comparison class gives
$\nu_{\min}(p,e,2)\ge\nu_{\mathrm{SS}}(p,e)$.  If $R\ge2$, the
strict inequality defining the ceiling gives
\[
 2R-2<1+\log_p e+v,
 \qquad e>p^{2R-3-v}.
\]
The actual mixed existence assertion of
Theorem~\ref{thm:mixed-all-indices}, applied at level $p^R$, supplies
unequal-character groups with isomorphic actual $p^R$-torsion over
some standing field of that exact absolute index.  Hence
$\nu_{\min}(p,e,2)\ge R$.

If $R=1$, the already proved inequality $R\ge v+1$ gives $v=0$,
and the defining bound gives $e\le p$.  A positive multiple of
$p-1$ in this range must be $e=p-1$ because $p\ge5$.
Theorem~\ref{thm:minimal-index} supplies the required lower bound
$\nu_{\min}(p,p-1,2)=1$.  Thus the same lower estimate holds in
this sole case not covered by the mixed theorem's $n\ge2$ range.
Combining both lower bounds with the upper bound proves
\[
 \nu_{\min}(p,e,2)
 =\max\left\{\nu_{\mathrm{SS}}(p,e),
 \left\lceil\frac{1+\log_p e+v_p(e)}2\right\rceil\right\}.
\]
\end{proof}

\subsection{Extending the exact wild bands}\label{sec:extended-wild-bands}

We now enlarge the coefficient interval by matching the canonical
residual characters before taking one actual quotient.  The argument
first treats coefficients prime to $p$; the different values at
multiples of $p$ then follow from the smaller coefficient bands.

\begin{proof}[Completion of Theorem~\ref{thm:wild-bands-exact}]
Assume first that
\[
 2\le k\le\frac{p^2-1}{2},\qquad p\nmid k,
 \qquad e=k(p-1)p^r,
 \qquad n=r+3,\qquad N=n-1=r+2.
\]
In particular $v_p(e)=r$.  Since $k<p^2$, we have
$e<p^{r+3}=p^{2n-3-r}$.  An unequal-character actual $p^n$-torsion
comparison in the full filtration class therefore reduces, by
Corollary~\ref{cor:all-index-two-slope-reduction}, to two whole
potentially supersingular groups with two slopes.
Lemma~\ref{os-lem:minima} and Proposition~\ref{fc-thm:valuation} give
\begin{equation}\label{os-eq:half-square-range}
 (p-1)p^{n-1}<z_i<pe=k(p-1)p^{n-2}.
\end{equation}
Both models are already central, because the lower bound divided by
$e$ is $p^2/k>1$.  The upper bound is less than $p^{n+1}$, so
$v_p(z_i)\le n$.

\smallskip\noindent
\emph{A parameter with valuation $n$.}
Suppose both valuations were at most $n-1$, ordered $s\le t$.
Different leading tame eigenvalues would contradict the inequality
$s+t\ge2n$ in Lemma~\ref{os-lem:two-chains}.  For equal eigenvalues,
write $z_1=p^sa$ and $z_2=p^tb$, with $p\nmid ab$.  Unequal special
characters make $a-b$ nonzero, and the same lemma makes it divisible
by at least $(p-1)p^{n-s}$.  But
\[
 |a-b|<k(p-1)p^{n-2-s}<(p-1)p^{n-s},
\]
since $k<p^2$.  Thus one valuation is exactly $n$.  Relabel that
parameter as
\begin{equation}\label{os-eq:half-square-high}
 z_1=bp^n,
 \qquad 1\le b<\frac{k(p-1)}{p^2}<\frac{p-1}{2}.
\end{equation}
In particular $b$ is prime to $p$.  Every parameter of valuation $n$
has its coefficient in this same strict interval.

\smallskip\noindent
\emph{Matching the canonical lines.}
Let $m=p^2-1$, let $\zeta$ have order $m$, and write
$\alpha_i=\zeta^{a_i}$ for the special tangent scalars.  Put
$\delta=\zeta^{p+1}$, of order $p-1$.  The tame inverse of
$p^{2n-1}$ modulo $m$ is $p$, so the leading scalars of
Lemma~\ref{os-lem:two-chains} are
\[
 \lambda_i=\alpha_i\zeta^{-pz_i}\in\mathbf F_p^{\times}.
\]
The canonical residual character $\kappa_i$ and its quotient
character $\eta_i$ satisfy
\begin{equation}\label{os-eq:half-square-residual}
 \kappa_i=\zeta^{(p+1)e-z_i-a_i},
 \qquad
 \eta_i=\zeta^{pz_i-a_i}=\lambda_i^{-1},
 \qquad
 \frac{\kappa_i}{\eta_i}=\delta^{-z_i}.
\end{equation}
For the first formula, a canonical point has high value
$(p+1)e-z_i$, and its leading field multiplier is divided by
$\alpha_i$ under descent, as in \eqref{rc-eq:kappa}.
For the second, the minimum-point description in
Lemma~\ref{os-lem:minima} identifies the leading noncanonical
quotient multiplier with $\zeta^{pz_i}$, again divided by
$\alpha_i$.  Equivalently this is the leading-coordinate computation
of \eqref{pg-eq:residual-characters}.  These leading-term calculations
use no prime-to-$p$ condition on $e$.  The last equality uses
$(p-1)\mid e$.

Suppose $\lambda_1\ne\lambda_2$.  The valuation-sum bound in
Lemma~\ref{os-lem:two-chains} forces $v_p(z_2)=n$, so
$z_2=b'p^n$ with $0<b'<(p-1)/2$.  The two quotient characters
$\eta_i$ differ.  Equality of the unordered residual composition
characters must therefore interchange them:
\[
 \eta_2=\kappa_1,\qquad \kappa_2=\eta_1.
\]
Equation~\eqref{os-eq:half-square-residual} then gives
$(p-1)\mid z_1+z_2$, hence $(p-1)\mid b+b'$.  This contradicts
$0<b+b'<p-1$.  The leading eigenvalues consequently agree.

If the second valuation were also $n$, the two-chain congruence
would give $(p-1)\mid b-b'$.  The strict interval in
\eqref{os-eq:half-square-high} would force $b=b'$, and equal
parameters and leading scalars would give equal tangent scalars,
contrary to unequal special characters.  Thus, writing
$s=v_p(z_2)$, we have $s\le n-1$.  The two-chain congruence gives
an integer $j$ such that
\begin{equation}\label{os-eq:half-square-partner}
 z_2=bp^s+j(p-1)p^{n-1}.
\end{equation}
The lower bound in \eqref{os-eq:half-square-range} implies $j\ge1$.
No upper bound on $j$ will be needed.

Equal leading scalars give equal quotient characters $\eta_i$.
The residual determinant then gives equal canonical characters
$\kappa_i$.  Their ratio in the first model is $\delta^{-b}\ne1$,
by \eqref{os-eq:half-square-high}.  Thus both representations have
distinct canonical and quotient characters.  A two-dimensional
representation with these composition characters has at most one
stable line of character $\kappa_i$: two distinct such lines would
span the space and make both composition factors equal to
$\kappa_i$.  The actual residual isomorphism therefore identifies
the canonical lines.

\smallskip\noindent
\emph{The quotient parameters.}
Take these two matched canonical quotients.
Lemma~\ref{lem:actual-quotients} gives an actual $p^N$ comparison,
and the $K$-isogenies preserve the unequal special characters.
Both original models were proved central.  Hence
Proposition~\ref{rc-prop:quotients} gives two central, two-slope
quotient parameters
\begin{equation}\label{os-eq:half-square-quotients}
 \begin{aligned}
 A&=p^rU, & U&=k(p^2-1)-bp^3,\\
 B&=p^rV-bp^s, & V&=k(p^2-1)-j(p-1)p^2.
 \end{aligned}
\end{equation}
In particular $e<A,B<pe$.  Since $U\equiv V\equiv-k\pmod p$,
one has $v_p(A)=r$.  If $s<r$, then $v_p(B)=s$; if $s>r$, then
$v_p(B)=r$.  If $s=r$ and $p\nmid k+b$, again $v_p(B)=r$.

In each of these cases the ordered quotient valuations $u\le t$
satisfy $u,t\le r=N-2$.  They are below $2N-1$ and their sum is
less than $2N$, so Lemma~\ref{os-lem:two-chains} forces equal
leading eigenvalues.  The nonzero difference of their unit parts
is then divisible by
\[
 (p-1)p^{2N-1-u-t}\ge(p-1)p^{r+3-u}.
\]
Its absolute value is instead less than
\[
 \frac{pe}{p^u}=k(p-1)p^{r+1-u},
\]
which is strictly smaller because $k<p^2$.  This contradiction
excludes every such case.  Only the explicit valuation hypotheses
of the two-chain lemma have been used at level $N$; the original
numerical candidate range has not been carried to that lower level.

\smallskip\noindent
\emph{The cancellation case.}
It remains to treat $s=r$ and $p\mid k+b$.  Write
\[
 k=pq+c,\qquad 1\le c\le p-1.
\]
The coefficient bounds give
\[
 q\le\frac{p-1}{2},\qquad b\le\frac{p-3}{2}.
\]
Since $0<c+b<2p$, the divisibility forces $c+b=p$.  Thus
$k+b=p(q+1)$, with $1\le q+1<p$, and
\eqref{os-eq:half-square-quotients} becomes
\[
 B=p^{r+1}W,
 \qquad W=p\bigl[k-j(p-1)\bigr]-(q+1).
\]
Consequently $v_p(B)=r+1$ exactly.  The quotient valuations $r,r+1$
are below $2N-1$ and have sum less than $2N$.  The two-chain lemma
again forces equal leading eigenvalues and requires $p^2$ to divide
$U-W$.  But
\[
 U-W\equiv-k+q+1\equiv b+q+1\pmod p,
 \qquad 1\le b+q+1\le p-1.
\]
This is the final contradiction.  The uniform upper bound
$\nuMin(p,e,2)\le r+2$ has been proved for the prime-to-$p$
coefficients in the stated interval.

\smallskip\noindent
\emph{Lower examples and coefficients divisible by $p$.}
For every such prime-to-$p$ coefficient, the strict mixed threshold
at level $p^{r+2}$ is satisfied:
\[
 e=k(p-1)p^r>p^{r+1}=p^{2(r+2)-3-r},
\]
because $k\ge2$ and $p\ge5$.  Theorem~\ref{thm:mixed-all-indices}
therefore supplies actual unequal-character groups with isomorphic
$p^{r+2}$-torsion over some standing field of the exact index $e$,
which may have residue field $\overline{\mathbf F}_p$.  It follows
that $\nuMin(p,e,2)\ge r+2$, proving equality.

Finally suppose $p\mid k$ in the same coefficient interval.  Write
$k=jp$.  Since $k\le(p^2-1)/2$, one has
$1\le j\le(p-1)/2$.  If $j=1$, the already established coefficient
$k=p$ case in Section~\ref{sec:all-wild-indices} gives
\[
 \nuMin\bigl(p,p^{r+1}(p-1),2\bigr)=r+2.
\]
If $j\ge2$, apply the coefficient-$j$ case proved in that section,
with exponent $r+1$, to obtain
\[
 \nuMin\bigl(p,jp^{r+1}(p-1),2\bigr)=r+3.
\]
These earlier small-coefficient proofs include both their uniform
upper statements and their actual existential lower constructions.
They do not use the prime-to-$p$ valuation calculation of the present
argument.  This proves every case of the theorem, with the stated
universal upper and existential lower quantifiers.
\end{proof}

\subsection{The first coefficient beyond the exact bands}
\label{sec:first-boundary-normal-form}

At the next coefficient, the strict residual-character separation
used for the exact bands becomes an equality. We determine the
resulting necessary form of an unequal-character comparison.

\begin{proof}[Proof of Theorem~\ref{thm:first-boundary-normal-form}]
Write $k_0=(p^2+1)/2$, $b_0=(p-1)/2$, $e=k_0(p-1)p^r$ and
$n=r+3$. Since $p\nmid k_0$, one has $v_p(e)=r$. The inequality
$k_0<p^2$ gives
\[
 e<p^{r+3}=p^{2n-3-r}.
\]
Corollary~\ref{cor:all-index-two-slope-reduction} therefore reduces
an unequal-character comparison to two whole potentially
supersingular groups, both with two-slope good models.
Lemma~\ref{os-lem:minima} and Proposition~\ref{fc-thm:valuation}
give
\begin{equation}\label{os-eq:first-boundary-range}
 (p-1)p^{n-1}<z_i<pe=k_0(p-1)p^{n-2}.
\end{equation}
Both models are central: the lower bound divided by $e$ is
$p^2/k_0>1$. The upper bound is less than $p^{n+1}$, so
$v_p(z_i)\le n$.

\smallskip\noindent
\emph{A parameter of valuation $n$.}
Suppose both valuations were at most $n-1$, ordered $s\le t$.
Different leading tame eigenvalues would contradict
$s+t\ge2n$ in Lemma~\ref{os-lem:two-chains}. For equal leading
eigenvalues, write $z_1=p^sa$ and $z_2=p^tb$, with $p\nmid ab$.
The same lemma makes the nonzero integer $a-b$ divisible by
$(p-1)p^{2n-1-s-t}$, hence by $(p-1)p^{n-s}$. But
\[
 |a-b|<k_0(p-1)p^{n-2-s}<(p-1)p^{n-s}.
\]
This is impossible. Relabel the model of valuation $n$ as
\begin{equation}\label{os-eq:first-boundary-high}
 z_1=bp^n,\qquad
 1\le b<\frac{k_0(p-1)}{p^2}
       =b_0+\frac{b_0}{p^2}.
\end{equation}
Since $0<b_0/p^2<1$, it follows that $1\le b\le b_0$.
Every parameter of valuation $n$ has its coefficient in this
same interval. All valuations just used are less than $2n-1$,
as required in Lemma~\ref{os-lem:two-chains}.

\smallskip\noindent
\emph{Residual inertia characters.}
Let $m=p^2-1$, let $\zeta$ have order $m$, write
$\alpha_i=\zeta^{a_i}$, and put $\delta=\zeta^{p+1}$, of order
$p-1$. The leading eigenvalues and the canonical and quotient
residual inertia characters are
\begin{equation}\label{os-eq:first-boundary-residual}
 \begin{aligned}
 \lambda_i&=\alpha_i\zeta^{-pz_i}\in\mathbf F_p^{\times},\\
 \kappa_i&=\zeta^{(p+1)e-z_i-a_i},\qquad
 \eta_i=\zeta^{pz_i-a_i}=\lambda_i^{-1},\qquad
 \kappa_i/\eta_i=\delta^{-z_i}.
 \end{aligned}
\end{equation}
These are the arbitrary-index leading-term formulas explained in
\eqref{os-eq:half-square-residual}; the last equality uses
$(p-1)\mid e$. We use the formulas on inertia, without asserting
equality of any unramified character factors. The canonical lines,
and their images under the given torsion isomorphism, are actual
$K$-stable lines. The isomorphism identifies the unordered residual
composition-character pairs on inertia.

\smallskip\noindent
\emph{Equal leading eigenvalues cannot occur.}
Suppose $\lambda_1=\lambda_2$. If $v_p(z_2)=n$, write
$z_2=b'p^n$ with $1\le b'\le b_0$. Lemma~\ref{os-lem:two-chains}
gives $(p-1)\mid b-b'$. This interval forces $b=b'$, and equal
parameters and leading eigenvalues give equal tangent scalars and
special characters, a contradiction. Thus $t=v_p(z_2)\le n-1$.
The same lemma gives an integer $j$ such that
\begin{equation}\label{os-eq:first-boundary-partner}
 z_2=bp^t+j(p-1)p^{n-1}.
\end{equation}
The lower bound in \eqref{os-eq:first-boundary-range} implies
$j\ge1$, because $bp^t<(p-1)p^{n-1}$. Its upper bound and the
positivity of $bp^t$ imply
\[
 j<\frac{k_0}{p}<b_0+1,
 \qquad\text{and hence}\qquad 1\le j\le b_0<p.
\]

Equal leading eigenvalues give equal quotient inertia characters
$\eta_i$. Equality of the residual determinants gives equal
canonical inertia characters $\kappa_i$. Their ratio in the first
model is $\delta^{-b}\ne1$, since $1\le b\le b_0<p-1$.
There is at most one stable line with inertia character $\kappa_i$:
two distinct such lines would span the representation and make
both composition characters equal to $\kappa_i$.
Consequently the actual isomorphism identifies the canonical lines.

Take these two matched canonical quotients.
Lemma~\ref{lem:actual-quotients} gives an actual comparison at level
$p^N$, where $N=r+2$, and the $K$-isogenies preserve the unequal
special characters. Since the original models are central,
Proposition~\ref{rc-prop:quotients} gives central, two-slope
quotient parameters
\begin{equation}\label{os-eq:first-boundary-quotients}
 \begin{aligned}
 A&=p^rU,& U&=k_0(p^2-1)-bp^3,\\
 B&=p^rV-bp^t,& V&=k_0(p^2-1)-j(p-1)p^2.
 \end{aligned}
\end{equation}
In particular $e<A,B<pe$. Since $U\equiv V\equiv-k_0\pmod p$,
we have $v_p(A)=r$. If $t<r$, then $v_p(B)=t$; if $t>r$, then
$v_p(B)=r$; and if $t=r$ but $p\nmid k_0+b$, again $v_p(B)=r$.

In each of these cases the ordered quotient valuations $q\le s$
satisfy $q,s\le r$. Their sum is less than $2N$ and both are less
than $2N-1$, so Lemma~\ref{os-lem:two-chains} forces equal leading
eigenvalues. The nonzero difference of their unit parts is then
divisible by at least
\[
 (p-1)p^{2N-1-q-s}\ge(p-1)p^{r+3-q}.
\]
Its absolute value is smaller than
\[
 pe/p^q=k_0(p-1)p^{r+1-q},
\]
which is strictly less than that divisor because $k_0<p^2$.
This is impossible. The lower-level argument uses the explicit
valuation hypotheses of the two-chain lemma, without transferring
the original numerical bound to level $N$.

The only case left is $t=r$ and $p\mid k_0+b$.
Put $c_0=(p+1)/2$. Since $k_0\equiv c_0\pmod p$ and
$1\le b\le b_0$, the positive integer $c_0+b$ is at most $p$.
Therefore $b=b_0$ and $k_0+b_0=pc_0$. The first unit part in
\eqref{os-eq:first-boundary-quotients} is now
\[
 U=(p^3-1)/2,
\]
and the second parameter is
\[
 B=p^{r+1}W,\qquad W=p\bigl[k_0-j(p-1)\bigr]-c_0.
\]
Since $p\nmid c_0$, one has $v_p(B)=r+1$ exactly.
The quotient valuations $r,r+1$ are less than $2N-1$ and their
sum is less than $2N$. The two-chain lemma therefore forces equal
leading eigenvalues and requires $p^2\mid U-W$. But the exact
identity is
\[
 U-W=p(p-1)j.
\]
Its $p$-adic valuation is one, since $1\le j\le b_0<p$.
This contradiction excludes equal leading eigenvalues.

\smallskip\noindent
\emph{The remaining form and the actual residual splitting.}
Now $\lambda_1\ne\lambda_2$. Since $v_p(z_1)=n$ and
$v_p(z_2)\le n$, the two-chain inequality forces $v_p(z_2)=n$.
Write $z_2=b'p^n$ with $1\le b'\le b_0$. The quotient inertia
characters $\eta_1,\eta_2$ differ, so equality of the unordered
composition-character pairs exchanges their entries:
\[
 \eta_2=\kappa_1,\qquad \kappa_2=\eta_1.
\]
Equation~\eqref{os-eq:first-boundary-residual} gives
$(p-1)\mid z_1+z_2$, hence $(p-1)\mid b+b'$. Since
$0<b+b'\le2b_0=p-1$, this forces $b=b'=b_0$ and
$z_1=z_2=b_0p^n$.

The equal parameters give
$\alpha_2/\alpha_1=\lambda_2/\lambda_1\in\mathbf F_p^\times$.
Equality of the finite determinant factors also gives
$(\alpha_2/\alpha_1)^{p+1}=1$. Its order thus divides both $p-1$
and $p+1$, and hence divides two. The ratio is not one, because
the leading eigenvalues differ. Therefore $\alpha_2/\alpha_1=-1$.

Finally $\kappa_i/\eta_i=\delta^{-b_0}=-1$, so the two inertia
characters in each model are distinct. The actual isomorphism
sends the canonical line of the first representation to a
$K$-stable line in the second of inertia character $\eta_2$.
This line differs from the second canonical line and is therefore
a stable complement. Applying the inverse isomorphism to the
second canonical line supplies the complementary stable line in
the first representation. Both residual representations split
over $K$. The two distinct-inertia summands are unique, so the
given isomorphism exchanges canonical and complementary lines
in both directions. This proves all the asserted necessary
conditions.
\end{proof}

\subsection{Proper canonical point fields}
\label{sec:proper-canonical-point-fields}
An exchanged canonical line strengthens the point-field degree
constraints at proper torsion levels. The argument uses an exact lower
break only when a coordinate has valuation prime to $p$ in its
actual point field.

Pass to the completed maximal unramified base and take the common
tame good field $L/K$ of degree $p^2-1$, with algebraically closed
residue field and $v_L(L^\times)=\Z$. Write $e=e_K$.

\begin{proposition}\label{prop:proper-canonical-point-fields}
Let $p\ge5$, and suppose two good models have the same central
parameter $e<z<pe$. Put
\[
 A=(p+1)e-z,\qquad B=z/p,\qquad a=v_p(A)\ge1.
\]
Let $f:G[p^n]\simeq H[p^n]$, $n\ge3$, be an actual torsion
isomorphism carrying the canonical line $C_G$ to a noncanonical
line of $H$. For an exact $p^j$-point $Q$ with
$p^{j-1}Q\in C_G\setminus\{0\}$, its point-field degree is a
power $[L(Q):L]=p^d$. If $2\le j<n$, then
\[
 d\ge\max\{0,2j-1-a\}.
\]
At $j=n$ one always has $d\ge\max\{0,2n-2-a\}$.
If $d=2n-2-a\ge0$, then $L(Q)/L(pQ)$ has degree at most $p$,
and its translation subgroup has trivial intersection with $C_G$.
In particular a nonzero such subgroup is a noncanonical line.
\end{proposition}

\begin{proof}
The canonical points of each group are $L$-rational by
Lemma~\ref{fc-lem:rational}. The image $f(C_G)$ is also pointwise
$L$-rational and differs from $C_H$. These two lines span $H[p]$,
so $H[p]$, and then $G[p]$, are pointwise $L$-rational. In
particular $B$ is an integer. No higher cyclic subgroup has been
inferred from this argument.

Choose any compatible tower $Q_1,\ldots,Q_n$ with
$pQ_i=Q_{i-1}$ and $Q_1\in C_G\setminus\{0\}$, and put
$R_i=f(Q_i)$. Since $A<p^2B$, the Newton division formula gives
\begin{equation}\label{eq:proper-point-values}
 v_L(T(Q_i))=\frac{A}{p^{2i-2}},\qquad
 v_L(S(R_i))=\frac{B}{p^{2i-2}}.
\end{equation}
For the first formula the initial target is canonical and lies
strictly below the middle division breakpoint; later target
values decrease. The second begins at a noncanonical point and
always remains on the degree-$p^2$ branch. Thus these values hold
for every choice of lifts, by Lemma~\ref{rc-lem:newton}.

Let $F_i=L(Q_i)$. The actual isomorphism identifies point
stabilizers, so $F_i=L(R_i)$. Each extension $F_i/F_{i-1}$ is
Galois: every conjugate of $Q_i$ over $F_{i-1}$ is $Q_i+P$ for
some $P\in G[p](L)$, and its coordinate belongs to $F_i$ by
formal addition. The map
\[
 \sigma\longmapsto\sigma(Q_i)-Q_i
\]
identifies its Galois group with an $\F_p$-subspace
$V_i\subset G[p](L)$. It is a homomorphism because all the
translating points are fixed over $L$. Consequently
\[
 [F_i:L]=p^{d_i},\quad d_1=0,\quad
 d_i-d_{i-1}=\dim_{\F_p}V_i\in\{0,1,2\},\quad d_i\le2i-2.
\]
The extensions are totally ramified because the residue field
is algebraically closed. No assertion that $F_i/L$ is Galois
has been used.

We first prove that
\begin{equation}\label{eq:prime-coordinate-no-canonical-translation}
 p\nmid w_i(T(Q_i))\quad\Longrightarrow\quad
 V_i\cap C_G=0\qquad(i\ge2),
\end{equation}
where $w_i=p^{d_i}v_L$ is the integer valuation of $F_i$.
Suppose instead that an automorphism $\sigma$ translates $Q_i$
by a nonzero point $P_0\in C_G$. It has order $p$. Apply
Lemma~\ref{cr-lem:valuation} to the cyclic extension of $F_i$
over its fixed field under $\sigma$.

For a good formal law, $F(T,U)-T$ is $U$ times a power series
with constant term one. Evaluated at points of positive value,
the remaining factor is a unit. Thus
\[
 w_i\bigl(\sigma(T(Q_i))-T(Q_i)\bigr)=p^{d_i}A.
\]
The prime-valuation clause of Lemma~\ref{cr-lem:valuation}
determines the lower break in this particular cyclic extension:
\[
 b_\sigma=p^{d_i}A\bigl(1-p^{-2i+2}\bigr).
\]
For the transported point the same automorphism gives
$\sigma(R_i)=R_i+f(P_0)$, and $f(P_0)$ is noncanonical, of
value $B$. The unconditional inequality of the same lemma and
\eqref{eq:proper-point-values} yield
\[
 b_\sigma\le p^{d_i}B\bigl(1-p^{-2i+2}\bigr).
\]
This contradicts $A>B$ and $i\ge2$. It proves
\eqref{eq:prime-coordinate-no-canonical-translation}, and hence
$\dim V_i\le1$ whenever that hypothesis holds. Equality in the
ramification estimate has not been asserted for a coordinate of
$p$-divisible valuation.

Write $A=p^aA_0$, with $p\nmid A_0$, and set
\[
 \delta_i=a+d_i-2i+2.
\]
The coordinate in \eqref{eq:proper-point-values} has positive
integer valuation $A_0p^{\delta_i}$ in $F_i$, so $\delta_i\ge0$.
Moreover
\[
 \delta_{i+1}-\delta_i=(d_{i+1}-d_i)-2\in\{-2,-1,0\}.
\]
If $2\le i<n$ and $\delta_i=0$, then integrality at the next
level forces $\delta_{i+1}=0$ and $d_{i+1}-d_i=2$.
Therefore $V_{i+1}=G[p]$ contains the canonical line, while
the coordinate of $Q_{i+1}$ has valuation prime to $p$.
This contradicts
\eqref{eq:prime-coordinate-no-canonical-translation}.
We have proved $\delta_i\ge1$ for every proper level
$2\le i<n$, which gives the asserted stronger lower bound.

At $i=n$, integrality gives $d_n\ge\max\{0,2n-2-a\}$.
If $d_n=2n-2-a\ge0$, then $\delta_n=0$, so
\eqref{eq:prime-coordinate-no-canonical-translation} gives the
claimed degree and translation restriction. Every point in the
statement extends to a compatible tower through level $n$.
The proof therefore covers all the stated point fields.
\end{proof}

Rational first-level torsion imposes a further constraint on these
point-field towers: their actual translation spaces can only increase.
This observation strengthens the bound at the top available level.

\begin{lemma}\label{lem:point-field-translation-persistence}
Let $p$ be odd, let $G$ be a height-two $p$-divisible group over a
field $L$ of characteristic zero, and suppose that $G[p]$ is
pointwise $L$-rational. Let $Q_1,\ldots,Q_n$ be compatible points
of exact orders $p,\ldots,p^n$, with $pQ_j=Q_{j-1}$, and put
$E_j=L(Q_j)$. Identify the relative Galois group
$\operatorname{Gal}(E_j/E_{j-1})$ with its translation subspace
$V_j\subset G[p]$. Then
\[
 V_2\subseteq V_3\subseteq\cdots\subseteq V_n.
\]
Only the relative extensions $E_j/E_{j-1}$ are asserted to be
Galois; the fields $E_j/L$ need not be Galois.
\end{lemma}

\begin{proof}
Every conjugate of $Q_j$ over $E_{j-1}$ differs from it by an
$L$-rational $p$-point. Thus $E_j/E_{j-1}$ is Galois, and the
translation map identifies its group with an $\mathbf F_p$-subspace
of $G[p]$.

Let $T=T_p(G)$ and choose a primitive $u\in T$ whose reduction
modulo $p^n$ represents $Q_n$. Its reduction modulo $p^j$ represents
$Q_j$. Write $\rho$ for the Galois action on $T$. Given
$v\in V_j$, choose a lift $\gamma$ of its relative automorphism
to the absolute Galois group of $E_{j-1}$, and put
$M=\rho(\gamma)-1$. Rationality of $G[p]$ and the fact that
$\gamma$ fixes $Q_{j-1}$ give
\[
 M(T)\subseteq pT,\qquad Mu\in p^{j-1}T,
 \qquad v=\frac{Mu}{p^{j-1}}\pmod{pT}.
\]
The binomial expansion of the single endomorphism $M$ gives
\[
 (\rho(\gamma)^p-1)u
 =pMu+\sum_{k=2}^{p-1}\binom pk M^ku+M^pu.
\]
For $2\le k\le p-1$, the corresponding summand belongs to
$p^{j+k-1}T\subseteq p^{j+1}T$. The last term belongs to
$p^{j+p-2}T\subseteq p^{j+1}T$, where oddness of $p$ is used.
Consequently
\[
 (\rho(\gamma^p)-1)u\equiv pMu\pmod{p^{j+1}T}.
\]
In particular $\gamma^p$ fixes $Q_j$, and its translation on
$Q_{j+1}$ is
\[
 \frac{(\rho(\gamma^p)-1)u}{p^j}
 \equiv\frac{Mu}{p^{j-1}}=v\pmod{pT}.
\]
Hence $v\in V_{j+1}$. This proves each asserted inclusion.
The lift $\gamma$ need not have finite order or fix all the
$(j-1)$st-level torsion; only the displayed hypotheses were used.
\end{proof}

\begin{corollary}\label{cor:proper-point-field-full-steps}
In the configuration of Proposition~\ref{prop:proper-canonical-point-fields},
suppose $n\ge a+2$, where $a=v_p(A)\ge1$.
For every compatible canonical point tower, each relative extension
$L(Q_j)/L(Q_{j-1})$ with $a+1\le j\le n$ has degree $p^2$.
In particular every canonical exact $p^n$-point $Q$ satisfies
\[
 [L(Q):L]=p^d,\qquad d\ge2n-1-a.
\]
\end{corollary}

\begin{proof}
Put $E_j=L(Q_j)$, $d_j=\log_p[E_j:L]$ and $i_j=\dim_{\mathbf F_p}V_j$.
The preceding proposition gives $E_1=L$ and
$d_j=\sum_{l=2}^j i_l$. The lemma makes $i_j$ nondecreasing.
The index $a+1$ is a proper level and is at least two, so the
proper-level bound gives
\[
 d_{a+1}\ge2(a+1)-1-a=a+1.
\]
If $i_{a+1}\le1$, all the $a$ terms in this sum would be at most
one, giving $d_{a+1}\le a$, a contradiction. Thus
$V_{a+1}=G[p]$, and every subsequent translation space is also
$G[p]$. All the stated relative extensions have degree $p^2$.
Finally the proper-level bound at $n-1$ gives
\[
 d_n=d_{n-1}+2\ge2n-3-a+2=2n-1-a.
\]
Every canonical exact $p^n$-point supplies such a tower by
successive multiplication, which proves the all-point assertion.
\end{proof}

\begin{proof}[Proof of Corollary~\ref{cor:boundary-point-field-degrees}]
Theorem~\ref{thm:first-boundary-normal-form} supplies the exchanged
central configuration, in both directions. Its first-level values are
\[
 A=\frac{p^3-1}{2}p^r,\qquad
 B=\frac{p-1}{2}p^{r+2},\qquad v_p(A)=r.
\]
Proposition~\ref{prop:proper-canonical-point-fields} gives the
stated proper-level bounds. Here $n=r+3$ and $a=r$, so
Corollary~\ref{cor:proper-point-field-full-steps} applies and yields
\[
 d_n\ge2(r+3)-1-r=r+5.
\]
It also shows that the relative steps at $j=r+1,r+2,r+3$ all
have degree $p^2$.

For $r=1$ and $n=4$, the earlier proper-level bounds and the
fibre-cardinality upper bounds already give $d_2=2$ and $d_3=4$;
the denominator bound alone leaves $d_4\in\{5,6\}$. If $d_4=5$,
the earlier top equality clause would give a cyclic final extension
of degree $p$ with noncanonical translation line. The new
translation-space inclusion makes the final relative step have
degree $p^2$, and therefore $d_4=6$. Equivalently the three total
degrees are exactly $p^2,p^4,p^6$. This excludes the previous
degree-$p^5$ possibility without using an exact break formula
at a coordinate of $p$-divisible valuation.
\end{proof}

These are necessary point-field restrictions. The remaining
degree-$p^6$ alternative when $r=1$ is neither constructed nor
excluded, and the boundary character index is not evaluated here.

\subsection{A divisorial calculation for the last Cartan layer}\label{sec:divisorial-cartan}
This subsection computes the ramification and inseparable residue extension of a concrete quotient of the universal deformation surface. Its complete discrete valuation fields have imperfect residue field, so the calculation is a supporting result rather than an example in the standing perfect-residue-field class.

Let $k=\overline{\F}_p$, and let $R_m$ be the complete local ring of the universal deformation of a one-dimensional height-two formal group over $k$, with full Drinfeld level-$p^m$ structure. Write $G$ for the universal formal group over $R_0$. For $m\ge1$, write $x_m,y_m$ for the coordinates of a chosen level basis. These are regular parameters of the two-dimensional ring $R_m$, and, for $m\ge2$, the transition is
\[
 x_{m-1}=[p](x_m),\qquad y_{m-1}=[p](y_m).
\]
The level action of $\GL_2(\Z/p^m)$ fixes the deformation base $R_0$ and $W(k)$. These are Drinfeld's regularity and level-parameter results, in the form of \citet[Theorem~2.1.2 and its proof, pp.~9--11; \S2.2, pp.~12--14]{Strauch2006}; see also \citet[Theorem~2.1.2, p.~437; Proposition~2.2.3(i), p.~439]{CarterStrauch2021}. The height assumption gives a nonzero degree-$p^2$ leading term of the special $p$-series. Scaling the coordinate over $k$ normalizes its coefficient to one. No exact polynomial expression for a mixed-characteristic lift is required.

For comparison, \citet[Section~2.2.8 and Corollary~2.2.9, p.~443]{CarterStrauch2021} obtain ramification index $(p-1)p^{n-1}$, when $\mathcal O=\Z_p$, for a tower formed by localizing at the height-one prime generated by one level coordinate. The valuation used here assigns positive weight to both level parameters, so its scalar fixed fields require a separate calculation.

\begin{proposition}\label{prop:divisorial-last-cartan}
Let $p$ be odd and $n\ge2$. Put
\[
 A=p^2+p+1,\qquad B=p^2,\qquad N=p^{2n-2}.
\]
Let $w$ be the divisorial monomial valuation on $\operatorname{Frac}(R_n)$ with $w(x_n)=A$ and $w(y_n)=B$, normalized to value group $\Z$, and let $F$ be its completion. Then its residue field is $k(t)$, where
\[
 t=\operatorname{res}(x_n^B/y_n^A),
\]
and
\begin{equation}\label{eq:divisorial-level-values}
 w(x_j)=A p^{2n-2j},\qquad w(y_j)=B p^{2n-2j}
 \quad(1\le j\le n),\qquad
 w(p)=(p^4-1)N.
\end{equation}
Fix a nonsquare $\epsilon\in\Z_p^\times$ and put
\[
 J=\begin{pmatrix}0&\epsilon\\1&0\end{pmatrix},\qquad
 U=I+p^{n-1}I,\qquad V=I+p^{n-1}J,\qquad H=\langle U,V\rangle .
\]
The level actions preserve $w$ and extend faithfully to $F$. The elements $U,V$ have order $p$, and $H$ has order $p^2$. The exact relative invariants are
\[
\begin{array}{c|cc}
 \text{extension}&\text{ramification index}&\text{residue degree}\\ \hline
 F/F^{\langle U\rangle}&1&p\\
 F/F^{\langle V\rangle}&p&1\\
 F/F^H&p&p .
\end{array}
\]
Each residue extension is purely inseparable, with the trivial extension included in this terminology. The full principal Cartan subgroup
\[
 C=(1+p\Z_p[J])/(1+p^n\Z_p[J])
\]
also preserves $w$, has order $p^{2n-2}$, and satisfies
\begin{equation}\label{eq:divisorial-full-cartan-bound}
 f(F/F^C)\ge p,\qquad e(F/F^C)\le p^{2n-3}.
\end{equation}
No equality for the full subgroup $C$ is asserted.
\end{proposition}

\begin{proof}
\emph{The valuation and the torsion values.}
The monomial valuation exists over the mixed-characteristic regular surface without choosing a coefficient field $k$ inside $R_n$. Blow up the closed point repeatedly. When the positive weights are $a>b$, the chart with parameters $x/y,y$ has weights $a-b,b$; use the other chart when $b>a$. The Euclidean algorithm reaches weights $1,1$, and the generic point of the next exceptional divisor is the required discrete valuation ring. Reversing these unimodular monomial changes gives residue field $k(t)$ and value group $\Z$. If $X,Y$ denote the initial forms of $x=x_n,y=y_n$, the associated graded fraction field is
\begin{equation}\label{eq:divisorial-graded-field}
 \operatorname{Frac}(\operatorname{gr}_w F)=k(X,Y).
\end{equation}
Indeed, choose integers $a,b$ with $aA+bB=1$. The monomials $x^ay^b$ and $x^B/y^A$ give a uniformizer and a transcendental residue coordinate, and their exponent matrix has determinant $-1$.

Let $m$ be the minimum value of a nonzero $p$-torsion point. The full-level product gives
\[
 [p](T)=q(T)\prod_{P\in G[p]}(T-T(P)),
\]
where $q$ is a unit whose constant residue is one. If $0<w(z)<m$, every factor has initial form $z$, including the factor for the zero point. Thus
\begin{equation}\label{eq:divisorial-small-root}
 w([p]z)=p^2w(z),\qquad
 \operatorname{in}_w([p]z)=\operatorname{in}_w(z)^{p^2}.
\end{equation}
One of $x_1,y_1$ attains $m$: a formal $\F_p$-linear combination cannot have value smaller than both basis points.

Multiplication by $p$ strictly raises the value of every positive-valued point with nonzero image, because its linear term is $pT$ and every other term has degree at least two. All ancestors of a minimum point therefore have value below $m$. Formula~\eqref{eq:divisorial-small-root} shows that either $m=NA$ or $m=NB$. The first possibility is impossible: starting from $y$ of value $B<A$, iteration would give $w(y_1)=NB<NA=m$. Hence $m=NB$ and the entire $y$-chain has the displayed values. In the $x$-chain the last input before reaching $x_1$ has value $AN/p^2<BN$. The same formula proves all its values and, for both chains,
\[
 \operatorname{in}_w(x_j)=X^{p^{2n-2j}},
 \qquad \operatorname{in}_w(y_j)=Y^{p^{2n-2j}}.
\]
At the first level, $p-1$ nonzero points have value $AN$, and the other $p^2-p$ have value $BN$. The linear coefficient of the product is $p$ up to a unit. Therefore
\[
 w(p)=\bigl((p-1)A+(p^2-p)B\bigr)N=(p^4-1)N,
\]
which proves~\eqref{eq:divisorial-level-values}.

\emph{The action and its leading derivations.}
An element of $C$ has the form $aI+bJ$, with $a\equiv1\pmod p$ and $b\equiv0\pmod p$. Its off-diagonal contribution to $x$ has value at least $p^2B>A$, and the correction to its diagonal contribution has value at least $p^2A>A$. The corresponding corrections to $y$ also have greater value than $y$. Thus the initial forms of $x,y$ are fixed. The action fixes $W(k)$, preserves $w$, and extends to $F$. It remains faithful because the original fraction field embeds in its completion. The two last-layer matrices have order $p$ and commute. Their coefficients $I,J$ in the last congruence layer are linearly independent over $\F_p$, giving $|H|=p^2$.

Use the basis-column convention $\phi\mapsto\phi\circ g$. Formal addition and the torsion initial forms give
\[
\begin{aligned}
 \operatorname{in}_w(Ux-x)&=X^N,&
 \operatorname{in}_w(Uy-y)&=Y^N,\\
 \operatorname{in}_w(Vx-x)&=Y^N,&
 \operatorname{in}_w(Vy-y)&=\overline\epsilon X^N .
\end{aligned}
\]
A transposed level convention moves the nonzero scalar $\overline\epsilon$ between the last two formulas and has no effect on the argument.
The smaller displacement excess for $U$ is $B(N-1)$, attained on $y$; for $V$ it is $BN-A$, attained on $x$. Both are positive. The first nonzero graded differences therefore define derivations of~\eqref{eq:divisorial-graded-field}:
\begin{equation}\label{eq:divisorial-derivations}
 D_U=Y^N\frac{\partial}{\partial Y},\qquad
 D_V=Y^N\frac{\partial}{\partial X}.
\end{equation}

Here the derivation assertion applies to the complete field. To justify it, first work on $R_n$: lift a homogeneous initial form by a polynomial in $x,y$ with Teichm\"uller coefficients from $W(k)$, subtract, and repeat. The positive monomial weights give convergent approximations; this filtration is comparable with the maximal-ideal filtration. The action fixes the coefficient lifts. If $d>0$ is the smaller displacement excess, the identity
\[
 g(fh)-fh=(g(f)-f)h+f(g(h)-h)+(g(f)-f)(g(h)-h)
\]
shows that the degree-$d$ term is a derivation, since the last summand has larger degree. The bound and its initial term pass to the convergent approximations, to fractions by the quotient rule, and then to the completion. In particular a higher-weight tail cannot cancel a nonzero value of the derivation in~\eqref{eq:divisorial-derivations}.

\emph{Fixed graded fields.}
For a finite group $\Gamma$ of these automorphisms, $E=F^\Gamma$ is complete and $F/E$ is Galois of degree $|\Gamma|$. A finite extension of complete discrete valuation fields satisfies $[F:E]=ef$, with the total residue degree, even if that residue extension is inseparable. One way to see this is to use the finite integral closure $\Ocal_F$ over $\Ocal_E$: it is free of rank $[F:E]$, and reduction modulo a uniformizer of $E$ has length $ef$. Finiteness follows by placing integral elements in the trace-dual of a lattice generated by an integral field basis; separability makes that trace pairing nondegenerate. The same valuation classes and a residue-field basis give
\begin{equation}\label{eq:divisorial-graded-degree}
 [k(X,Y):\operatorname{Frac}(\operatorname{gr}_w E)]=ef=|\Gamma|.
\end{equation}
Independence is checked first in distinct valuation classes modulo $e$, then in the residue field; successive leading-term subtraction proves spanning.

An element fixed by $U$ has initial form killed by $D_U$, so
\[
 \operatorname{Frac}(\operatorname{gr}_w F^{\langle U\rangle})
 \subset k(X,Y^p).
\]
Both fields on the right and left have degree $p$ in $k(X,Y)$, by~\eqref{eq:divisorial-graded-degree}; hence equality holds. The weight group of $k(X,Y^p)$ is generated by $A,pB$ and is $\Z$. Its degree-zero field is $k(t^p)$. This gives $e=1,f=p$ for the scalar subgroup.

For $V$ the same argument gives
\[
 \operatorname{Frac}(\operatorname{gr}_w F^{\langle V\rangle})
 =k(X^p,Y).
\]
Its weight group is $p\Z$, since $\gcd(pA,B)=p$, and its degree-zero field is $k(t)$, since $(X^p)^p/Y^A=X^B/Y^A$. This gives $e=p,f=1$.

For $H$, both derivations kill every invariant initial form. Their common kernel is $k(X^p,Y^p)$, because $k$ is perfect. Its degree is $p^2$, so~\eqref{eq:divisorial-graded-degree} again gives equality with the fixed graded fraction field. Its weight group is $p\Z$ and its degree-zero field is $k(t^p)$. Thus $e=f=p$ for $H$, proving the table and its purely inseparable residue assertions.

\emph{The full principal subgroup.}
The principal-unit filtration of the unramified quadratic ring $\Z_p[J]$ gives $|C|=p^{2n-2}$. Since $F^C\subset F^H$, residue degrees in the tower multiply, giving $f(F/F^C)\ge p$. The degree formula gives~\eqref{eq:divisorial-full-cartan-bound}. This does not compute that full residue degree.
\end{proof}

The last-layer bound alone gives a weaker numerical observation at $p=5,n=4$. If $F^C/K$ is a tame extension of ramification index $12$ with unchanged residue field, then
\[
 e_K=\frac{624\cdot5^6}{12\,e(F/F^C)}
 \ge\frac{624\cdot5^6}{12\cdot5^5}=260.
\]
Equality in this estimate would require $e(F/F^C)=5^5$, or equivalently $f(F/F^C)=5$. The next proposition excludes that equality for this valuation by computing the full scalar subgroup. The tame extension remains an additional hypothesis, and no perfect-residue specialization follows from the divisorial calculation.

\subsection{The full scalar subgroup at level four}\label{sec:divisorial-scalar}
This subsection computes the full scalar fixed field at level four and strengthens the ramification obstruction for the associated principal Cartan quotient.

\begin{proposition}\label{prop:divisorial-scalar-level-four}
Use the notation of Proposition~\ref{prop:divisorial-last-cartan}, and take $n=4$. Thus
\[
 A=p^2+p+1,\qquad B=p^2,\qquad N=p^6,
 \qquad w(x)=A,\quad w(y)=B,
\]
where $x=x_4$, $y=y_4$, and $p$ is odd. Write $X,Y$ for their initial forms and put
\[
 S=(1+p\Z_p)/(1+p^4\Z_p)\subseteq\GL_2(\Z/p^4),
 \qquad t=\operatorname{res}(x^B/y^A).
\]
Then
\begin{equation}\label{eq:divisorial-scalar-graded}
 \operatorname{Frac}(\operatorname{gr}_w F^S)
       =k(X^p,Y^{p^2}),
 \qquad e(F/F^S)=p,
 \qquad \kappa(F^S)=k(t^{p^2}).
\end{equation}
In particular $f(F/F^S)=p^2$, and the residue extension is purely inseparable. For the full principal nonsplit Cartan subgroup
\[
 C=(1+p\Z_p[J])/(1+p^4\Z_p[J])
\]
one has
\begin{equation}\label{eq:divisorial-cartan-scalar-bound}
 f(F/F^C)\ge p^2,\qquad e(F/F^C)\le p^4.
\end{equation}
\end{proposition}

We first give the approximation argument needed to pass between the actual fixed fields. For a valued field $K$, write
\[
 \mathcal G(K)=\operatorname{Frac}(\operatorname{gr}_w K).
\]
The degree-zero subfield of its homogeneous elements is its residue field. Every homogeneous element of $\mathcal G(K)$ is the initial form of an element of $K$: a homogeneous quotient can be represented using homogeneous numerator and denominator and then lifted.

\begin{lemma}\label{lem:divisorial-controlled-lift}
Let $K$ be a complete discretely valued field and let $\sigma$ be a finite-order valuation-preserving automorphism. Put $L=K^{\langle\sigma\rangle}$. Suppose that $\sigma-1$ raises every value by at least $d>0$, its degree-$d$ initial operator is a derivation $D$ of $\mathcal G(K)$, and
\[
 \mathcal G(L)=\ker D.
\]
If $c=w(\sigma z-z)>w(z)+d$, there is $a\in L$ such that
\[
 w(z-a)\ge c-d.
\]
The conclusion also holds with $a=z$ when $z$ is fixed.
\end{lemma}

\begin{proof}
The strict inequality implies $D(\operatorname{in}(z))=0$. Choose $a\in L$ with the same initial form as $z$. If the error $b=z-a$ has value below $c-d$, then
\[
 \sigma b-b=\sigma z-z,
 \qquad w(\sigma b-b)>w(b)+d.
\]
Hence $D(\operatorname{in}(b))=0$, and $\operatorname{in}(b)$ lifts to an element of $L$. Adding that element to $a$ strictly increases the error value. Discreteness makes this a finite procedure before the value reaches $c-d$.
\end{proof}

We shall use two elementary facts about graded fields. If $L\subset K$ is finite, then
\begin{equation}\label{eq:divisorial-graded-inequality}
 [\mathcal G(K):\mathcal G(L)]
      =e(K/L)f(K/L)\le [K:L].
\end{equation}
To see this, choose representatives of the value-group cosets and a basis of the residue-field extension. Their homogeneous products form a basis of the graded field extension. Any nontrivial linear relation among their lifts would give a nontrivial relation among the terms of least value, so those lifts are linearly independent over $L$.

Second, a displacement bound on lifts of homogeneous rational generators extends to the complete field. Products, inverses and sums first give the bound on their rational expressions. Every homogeneous initial form has such a lift. Successively subtracting these lifts approximates any element to arbitrarily high value, so continuity gives the bound on every element. The same argument determines the initial operator. Here residues in $k$ have fixed Teichm\"uller representatives in $W(k)$.

\begin{proof}[Proof of Proposition~\ref{prop:divisorial-scalar-level-four}]
Put
\[
 E=w(p)=(p^4-1)p^6,\qquad Q_j=p^{2j},
 \qquad d_j=(Q_j-1)B,\quad d'_j=(Q_j-1)A
 \quad(1\le j\le3).
\]
Let $\sigma_j$ denote scalar multiplication by $1+p^j$. The coordinate calculation in Proposition~\ref{prop:divisorial-last-cartan}, with the height-two leading coefficient normalized to one, gives
\begin{equation}\label{eq:divisorial-scalar-displacements}
 \operatorname{in}([p^j]x)=X^{Q_j},\qquad
 \operatorname{in}([p^j]y)=Y^{Q_j}.
\end{equation}
The formal group law satisfies $G(T,Z)-T=Z(1+\text{terms in }(T,Z))$. Consequently
\[
 \operatorname{in}(\sigma_jx-x)=X^{Q_j},\qquad
 \operatorname{in}(\sigma_jy-y)=Y^{Q_j}.
\]
These substitutions preserve the weighted filtration and induce the identity on its associated graded ring; the inverse substitutions do likewise. They therefore preserve $w$ and extend to $F$.

The least displacement excess on $F$ is $d_j$, since $B<A$. Its degree-$d_j$ operator is
\begin{equation}\label{eq:divisorial-scalar-first-derivation}
 D_j=Y^{Q_j}\frac{\partial}{\partial Y}
       \quad\text{on }\mathcal G(F)=k(X,Y).
\end{equation}
It is a derivation because the product of two positive displacements has strictly higher degree. The preceding continuation argument extends the generator calculation to the complete field. Its kernel is $k(X,Y^p)$.

For odd $p$, the scalar group $S$ is cyclic of order $p^3$, generated by $\sigma_1$. Its subgroups of orders $p^2$ and $p$ are generated by $\sigma_2$ and $\sigma_3$, respectively: the binomial formula gives
\[
 v_p\bigl((1+p)^{p^r}-1\bigr)=r+1.
\]
Write
\[
 L=F^{\langle\sigma_3\rangle},\qquad
 K=F^{\langle\sigma_2\rangle},\qquad M=F^S.
\]
The level action is faithful, so each step in $F/L/K/M$ has degree $p$. These fixed fields are complete for the restricted valuation. The last scalar calculation in Proposition~\ref{prop:divisorial-last-cartan} gives
\begin{equation}\label{eq:divisorial-scalar-last-input}
 \mathcal G(L)=k(X,Y^p)=\ker D_3.
\end{equation}

We will take ordinary $p$th powers as well as formal-group multiples. If $w(z)=r$ and $w(\sigma z-z)=r+s$, then
\begin{equation}\label{eq:divisorial-ordinary-power}
 \operatorname{in}\bigl(\sigma(z^p)-z^p\bigr)
       =\operatorname{in}(\sigma z-z)^p
       \quad\text{provided }E>(p-1)s.
\end{equation}
Indeed the mixed binomial terms have value at least $E+pr+s$, whereas the pure last term has value $pr+ps$. All applications below satisfy this condition because
\[
 A<p^3,\qquad
 E>(p-1)(p^6-1)A.
\]
For the second inequality, its right side is less than
$(p-1)p^9=p^{10}-p^9$, while $E=p^{10}-p^6$. We also have
\begin{equation}\label{eq:divisorial-weight-order}
 B<A<pB.
\end{equation}

Apply Lemma~\ref{lem:divisorial-controlled-lift} to $\sigma_3$, first with $x$ and then with the ordinary power $y^p$. It gives $a,b\in L$ with
\begin{align}
 \operatorname{in}(a)&=X,&
 w(a-x)&\ge C_x:=p^6A-(p^6-1)B,\label{eq:divisorial-first-lift-x}\\
 \operatorname{in}(b)&=Y^p,&
 w(b-y^p)&\ge C_y:=p^7B-(p^6-1)B.\label{eq:divisorial-first-lift-y}
\end{align}
The error bounds exceed $A$ and $pB$, respectively. For $j=1,2$ one has the exact strict inequalities
\begin{align*}
 C_x+d_j-Q_jA&=(p^6-Q_j)(A-B)>0,\\
 C_y+d_j-pQ_jB&=(p-1)(p^6-Q_j)B>0.
\end{align*}
Thus the errors in these lifts do not affect the leading displacements:
\begin{equation}\label{eq:divisorial-middle-displacements}
 \operatorname{in}(\sigma_ja-a)=X^{Q_j},\qquad
 \operatorname{in}(\sigma_jb-b)=Y^{pQ_j}
 \quad(j=1,2).
\end{equation}
In the independent coordinates $X,Z=Y^p$ of $\mathcal G(L)$, inequality~\eqref{eq:divisorial-weight-order} makes the least displacement excess $d'_j$. Its initial derivation is
\[
 D'_j=X^{Q_j}\frac{\partial}{\partial X},\qquad
 \ker D'_j=k(X^p,Y^p).
\]
The continuation argument gives the same bound on all of $L$. Every initial form from $K$ lies in $\ker D'_2$. This kernel has degree $p$ in $\mathcal G(L)$, while $[L:K]=p$. Equation~\eqref{eq:divisorial-graded-inequality} therefore forces
\begin{equation}\label{eq:divisorial-middle-fixed}
 \mathcal G(K)=k(X^p,Y^p).
\end{equation}

There are now controlled representatives for the final scalar action. Apply Lemma~\ref{lem:divisorial-controlled-lift} on $L$, using $\sigma_2$, its shift $d'_2$, and~\eqref{eq:divisorial-middle-fixed}. Starting with $a^p$ and $b$ gives $u,v\in K$ such that
\begin{align}
 \operatorname{in}(u)&=X^p,&
 w(u-a^p)&\ge C_u:=(p^5-p^4+1)A,\label{eq:divisorial-second-lift-x}\\
 \operatorname{in}(v)&=Y^p,&
 w(v-b)&\ge C_v:=p^5B-(p^4-1)A.\label{eq:divisorial-second-lift-y}
\end{align}
Here~\eqref{eq:divisorial-ordinary-power} gives initial displacement $X^{p^5}$ for $a^p$, and~\eqref{eq:divisorial-middle-displacements} gives $Y^{p^5}$ for $b$. The second error exceeds $pB$ because
\[
 C_v-pB=(p^4-1)(pB-A)>0.
\]
To compute $\sigma_1$ on these lifts, use its shift $d'_1$ on $L$. The two error estimates needed are
\begin{align*}
 C_u+d'_1-p^3A&=p^2(p-1)(p^2-1)A>0,\\
 C_v+d'_1-p^3B&=p^2(p^2-1)(pB-A)>0.
\end{align*}
It follows that
\[
 \operatorname{in}(\sigma_1u-u)=X^{p^3},\qquad
 \operatorname{in}(\sigma_1v-v)=Y^{p^3}.
\]
The independent coordinates $U_0=X^p,V_0=Y^p$ of $\mathcal G(K)$ have weights $pA,pB$. The smaller displacement excess is now $p(p^2-1)B$, and the initial derivation is
\[
 D''=V_0^{p^2}\frac{\partial}{\partial V_0},\qquad
 \ker D''=k(U_0,V_0^p)=k(X^p,Y^{p^2}).
\]
The action of $S/\langle\sigma_2\rangle$ on $K$ has order $p$ and fixed field $M$. Its invariant initial forms lie in this kernel, which has degree $p$ in $\mathcal G(K)$. Since $[K:M]=p$, equation~\eqref{eq:divisorial-graded-inequality} forces
\[
 \mathcal G(M)=k(X^p,Y^{p^2}).
\]

The two final homogeneous generators have weights $pA,p^2B$, whose greatest common divisor is $p$. Thus the value group of $M$ is $p\Z$, and $e(F/M)=p$. Since $\gcd(A,pB)=1$, its degree-zero subfield is generated by
\[
 \frac{(X^p)^{pB}}{(Y^{p^2})^A}
      =\left(\frac{X^B}{Y^A}\right)^{p^2}=t^{p^2}.
\]
Therefore $\kappa(M)=k(t^{p^2})$, and the residue extension from $F$ has degree $p^2$, purely inseparable. This proves~\eqref{eq:divisorial-scalar-graded}.

Finally $C$ preserves $w$, has order $p^6$ and contains $S$, by Proposition~\ref{prop:divisorial-last-cartan}. The inclusion $F^C\subseteq F^S$ gives
\[
 f(F/F^C)\ge [k(t):k(t^{p^2})]=p^2.
\]
Since $[F:F^C]=p^6$, the fundamental inequality gives $e(F/F^C)\le p^4$, proving~\eqref{eq:divisorial-cartan-scalar-bound}.
\end{proof}

\begin{corollary}\label{cor:divisorial-1300-bound}
In Proposition~\ref{prop:divisorial-scalar-level-four}, let $p=5$. If $F^C/K$ is a tame extension of ramification index $12$ with unchanged residue field, then
\[
 e_K\ge1300.
\]
Thus this level-four divisorial Cartan construction cannot yield absolute ramification index $260$.
\end{corollary}

\begin{proof}
Ramification indices multiply, so
\[
 e_K=\frac{w(p)}{12\,e(F/F^C)}
     \ge\frac{(5^4-1)5^6}{12\cdot5^4}
     =1300.
\]
\end{proof}

This obstruction concerns the specified valuation and Cartan quotient. The existence of other unequal-character boundary comparisons at index $260$ remains open. The argument does not construct the tame quotient, the required descent datum, or a perfect-residue specialization.

\subsection{The full principal Cartan at every level}\label{sec:divisorial-full-cartan}
This subsection computes the fixed field of the full principal nonsplit Cartan at every level by constructing invariant coordinates with controlled errors.

\begin{proposition}\label{prop:divisorial-full-cartan}
Use the notation and valuation of Proposition~\ref{prop:divisorial-last-cartan}, with $p$ odd and $n\ge2$. Thus
\[
 A=p^2+p+1,\qquad B=p^2,\qquad
 w(x)=A,\quad w(y)=B,\quad
 w(p)=(p^4-1)p^{2n-2},
\]
where $x=x_n$, $y=y_n$, and $F$ is the completed valued field. Let $X,Y$ be the initial forms of $x,y$, put $t=\operatorname{res}(x^B/y^A)$, and write
\[
 \mathcal O=\Z_p[J],\qquad
 C=(1+p\mathcal O)/(1+p^n\mathcal O),\qquad r=p^{n-1}.
\]
Then
\begin{equation}\label{eq:full-cartan-field}
 \mathcal G(F^C)=k(X^r,Y^r),\qquad
 e(F/F^C)=r,\qquad \kappa(F^C)=k(t^r).
\end{equation}
In particular $f(F/F^C)=r$, and the residue extension is purely inseparable.
At $n=4$, let $S=(1+p\Z_p)/(1+p^4\Z_p)$ be the scalar subgroup. Then
\begin{equation}\label{eq:full-cartan-scalar-relative}
 e(F^S/F^C)=p^2,\qquad f(F^S/F^C)=p,
\end{equation}
and the latter residue extension is purely inseparable.
\end{proposition}

\begin{proof}
We use $\mathcal G(K)=\operatorname{Frac}(\operatorname{gr}_wK)$ with the restriction of the same valuation $w$ on every intermediate field.

\paragraph{Numerical and fixed-field preliminaries.}
We first record the bounds needed to lift initial forms inside actual fixed fields. Set
\[
 D=A-B=p+1,\qquad E=w(p)=(p^4-1)p^{2n-2},
 \qquad C_j=(1+p^j\mathcal O)/(1+p^n\mathcal O).
\]
Here $C_n=1$, $|C_j|=p^{2(n-j)}$, and $C_j/C_{j+1}$ is the additive group $\mathcal O/p\mathcal O$. Its independent generators are represented by
\[
 \sigma_j=1+p^j,\qquad \tau_j=1+p^jJ.
\]
All these elements commute. Their induced actions on $F^{C_{j+1}}$ have order $p$, and their joint fixed field is $F^{C_j}$. These degrees follow from the faithful finite group action on $F$. Faithfulness persists in the completion because it holds on the dense fraction field. The fixed subfields are complete.

The following inequalities hold for every $p\ge3$:
\begin{equation}\label{eq:full-cartan-elementary}
 p^2D>A,\qquad
 (p-1)B-pD=p(p^2-2p-1)>0,\qquad p^4-1>pA.
\end{equation}
If $0\le m\le n-2$, $r_m=p^m$, $j=n-m-1$, and $Q=p^{2j}$, they imply
\begin{equation}\label{eq:full-cartan-power-bound}
 E>p r_m Q A,
\end{equation}
since $p r_m Q=p^{2n-m-1}\le p^{2n-1}$.

For positive-valued $z,h$ and a power $s$ of $p$, all intermediate binomial coefficients in $(z+h)^s-z^s$ are divisible by $p$. If $E>s\,w(h)$, their terms have value greater than $s\,w(h)$, so
\begin{equation}\label{eq:full-cartan-ordinary-power}
 \operatorname{in}\bigl((z+h)^s-z^s\bigr)
       =\operatorname{in}(h)^s.
\end{equation}
For a lower bound, if $w(h)\ge c>0$ and $E>sc$, the same expansion gives $w((z+h)^s-z^s)\ge sc$. These statements concern ordinary powers in the characteristic-zero field, not formal multiplication by $p$.

We shall repeatedly use Lemma~\ref{lem:divisorial-controlled-lift}. If an order-$p$ automorphism $g$ has positive displacement shift $d$ and first derivation $D_g$ on $\mathcal G(K)$, then
\[
 \mathcal G(K^{\langle g\rangle})\subseteq\ker D_g.
\]
When $[\mathcal G(K):\ker D_g]=p$, equation~\eqref{eq:divisorial-graded-inequality} and $[K:K^{\langle g\rangle}]=p$ force equality. Only after establishing that equality do we choose an invariant lift $z_0$ with
\begin{equation}\label{eq:full-cartan-control}
 w(z_0-z)\ge w(gz-z)-d
\end{equation}
for an initial form in $\ker D_g$. A fixed $z$ needs no correction. Displacement bounds on homogeneous rational generators extend to the complete field by successive leading-term approximation, as proved above. The coefficient field $k$ has fixed Teichm\"uller representatives in $W(k)$.

\paragraph{The invariant-coordinate induction.}
We keep a quantitative comparison with the original coordinate powers. For $0\le m\le n-1$, put
\[
 K_m=F^{C_{n-m}},\qquad r_m=p^m.
\]
We prove
\begin{equation}\label{eq:full-cartan-induction-field}
 \mathcal G(K_m)=k(X^{r_m},Y^{r_m}),
\end{equation}
and choose $u_m,v_m\in K_m$ with those initial forms such that, for $m\ge1$,
\begin{equation}\label{eq:full-cartan-induction-error}
 \begin{split}
 w(u_m-x^{r_m})&\ge H_m,\\
 w(v_m-y^{r_m})&\ge H_m,
 \end{split}
 \qquad H_m=r_m p^{2(n-m)}D.
\end{equation}
At $m=0$, take $u_0=x$, $v_0=y$. The assertion about the graded field is already known. The errors are zero. For $m\ge1$, $H_m>r_mA>r_mB$, so the stated error bound preserves the initial forms.

Suppose the assertions hold at $m\le n-2$. Write
\[
 \rho=r_m,\qquad j=n-m-1,\qquad Q=p^{2j},
 \qquad u=u_m,\quad v=v_m,\quad U=X^{\rho},\quad V=Y^{\rho}.
\]
When $m\ge1$, the previous error bound satisfies
\begin{equation}\label{eq:full-cartan-old-error}
 H_m=\rho p^2QD>\rho QA.
\end{equation}
Thus valuation preservation alone makes the displacement of each old error too large to affect any of the coordinate initial forms below.

\paragraph{The induced coordinate displacements.}
The torsion initial forms from Proposition~\ref{prop:divisorial-last-cartan} and the identity
\[
 G(T,Z)-T=Z\bigl(1+\text{terms in }(T,Z)\bigr)
\]
give initial differences $X^Q,Y^Q$ for $\sigma_j$ on $x,y$, and $Y^Q,\bar\epsilon X^Q$ for $\tau_j$, with $\bar\epsilon\in\F_p^*$. Equations~\eqref{eq:full-cartan-power-bound} and \eqref{eq:full-cartan-ordinary-power} give their ordinary $\rho$-th powers. Equation~\eqref{eq:full-cartan-old-error} then permits replacement by the actual invariant coordinates $u,v$. The table of initial differences is
\begin{equation}\label{eq:full-cartan-displacements}
 \begin{array}{c|cc}
       &u&v\\ \hline
 \sigma_j-1&U^Q&V^Q\\
 \tau_j-1&V^Q&\bar\epsilon U^Q .
 \end{array}
\end{equation}
The weights of the four entries are $\rho QA,\rho QB,\rho QB,\rho QA$, respectively.

For $\tau=\tau_j$ on $K_m$, the smallest relative shift is
\[
 d_\tau=\rho(QB-A)>0.
\]
The shift on $v$ is larger by $\rho(Q+1)D$. The first derivation is therefore
\begin{equation}\label{eq:full-cartan-cross-derivation}
 D_\tau=V^Q\frac{\partial}{\partial U}
       \quad\text{on }k(U,V).
\end{equation}
For $\sigma=\sigma_j$ on $K_m$, the smallest relative shift and derivation are
\begin{equation}\label{eq:full-cartan-scalar-shift}
 d_\sigma=\rho(Q-1)B>0,\qquad
 D_\sigma=V^Q\frac{\partial}{\partial V}.
\end{equation}
The bounds hold on the entire complete field by the approximation argument stated above.

\paragraph{First fix the cross generator.}
We construct coordinates in $L=K_m^{\langle\tau\rangle}$ before treating the scalar action there. The kernel of \eqref{eq:full-cartan-cross-derivation} is $k(U^p,V)$, of index $p$ in $k(U,V)$. Since $[K_m:L]=p$, the graded-degree argument proves
\begin{equation}\label{eq:full-cartan-cross-field}
 \mathcal G(L)=k(U^p,V).
\end{equation}
The ordinary-power estimate gives $w(\tau(u^p)-u^p)=p\rho QB$, while the table gives $w(\tau v-v)=\rho QA$. Formula~\eqref{eq:full-cartan-control} produces $a,b\in L$ satisfying
\begin{equation}\label{eq:full-cartan-first-errors}
 \begin{split}
 w(a-u^p)&\ge E_a:=\rho((p-1)QB+A),\\
 w(b-v)&\ge E_b:=\rho(QD+A).
 \end{split}
\end{equation}
Here $E_a>p\rho A$ and $E_b>\rho B$. These strict inequalities verify the hypotheses of the controlled-lifting lemma and give initial forms $U^p,V$ and weights $p\rho A,\rho B$.

\paragraph{Then fix the scalar generator.}
We determine its action on the actual coordinates $a,b$. The Cartan is abelian, so $\sigma$ preserves $L$, acts there with order $p$, and has fixed field $K_{m+1}$. The bound on $K_m$ gives
\[
 w\bigl(\sigma(b-v)-(b-v)\bigr)\ge E_b+d_\sigma>\rho QB,
\]
where the strict margin is $E_b-\rho B=\rho(Q+1)D$. Hence
\begin{equation}\label{eq:full-cartan-scalar-b}
 \operatorname{in}(\sigma b-b)=V^Q.
\end{equation}
For the other coordinate, the ordinary-power estimate and the same displacement bound give
\begin{equation}\label{eq:full-cartan-scalar-a}
 w(\sigma a-a)\ge\min(p\rho QA,E_a+d_\sigma).
\end{equation}
Both terms on the right exceed $p\rho A+d_\sigma$, because
\[
 \begin{split}
 p\rho QA-p\rho A-d_\sigma&=\rho(Q-1)(pA-B)>0,\\
 E_a-p\rho A&=\rho(p-1)(QB-A)>0.
 \end{split}
\]
Thus on $\mathcal G(L)=k(U^p,V)$ the first scalar derivation still has shift $d_\sigma$, vanishes on $U^p$, and sends $V$ to $V^Q$. Its kernel is $k(U^p,V^p)$, of index $p$. The second actual degree-$p$ fixed-field extension therefore gives
\begin{equation}\label{eq:full-cartan-next-field}
 \mathcal G(K_{m+1})=k(U^p,V^p)=k(X^{p\rho},Y^{p\rho}).
\end{equation}

\paragraph{Errors for the next congruence layer.}
We complete the induction by constructing sufficiently accurate coordinates in $K_{m+1}$. Put $H=p\rho QD=H_{m+1}$. Applying the controlled-lifting lemma to $a$ inside $L$ gives an invariant $u'$ with
\begin{equation}\label{eq:full-cartan-second-u}
 w(u'-a)\ge\min(p\rho QA-d_\sigma,E_a).
\end{equation}
For $b^p$, equations~\eqref{eq:full-cartan-power-bound}, \eqref{eq:full-cartan-ordinary-power} and \eqref{eq:full-cartan-scalar-b} give $w(\sigma(b^p)-b^p)=p\rho QB$. An invariant $v'$ satisfies
\begin{equation}\label{eq:full-cartan-second-v}
 w(v'-b^p)\ge E_v:=p\rho QB-d_\sigma
                  =\rho((p-1)Q+1)B.
\end{equation}
The required error margins are
\begin{equation}\label{eq:full-cartan-error-margins}
 \begin{split}
 E_a-H&=\rho\{Q((p-1)B-pD)+A\}>0,\\
 p\rho QA-d_\sigma-H&=\rho\{(p-1)Q+1\}B>0,\\
 E_v-H&=\rho\{Q((p-1)B-pD)+B\}>0,\\
 pE_b-H&=p\rho A>0.
 \end{split}
\end{equation}
The first and third inequalities use \eqref{eq:full-cartan-elementary}. Since $E>p\rho QA>H$, the binomial expansion of $b^p-v^p$ gives
\[
 w(b^p-v^p)\ge\min(pE_b,E)>H.
\]
Equations~\eqref{eq:full-cartan-first-errors}, \eqref{eq:full-cartan-second-u} and \eqref{eq:full-cartan-error-margins} give $w(u'-u^p)>H$. Equations~\eqref{eq:full-cartan-second-v} and \eqref{eq:full-cartan-error-margins} similarly give $w(v'-v^p)>H$.

For $m\ge1$, taking ordinary $p$-th powers of the old approximation errors gives valuation at least $\min(pH_m,E)>H$, because $pH_m=p^2H$. At $m=0$ those errors vanish. Consequently
\[
 w(u'-x^{p\rho})\ge H,\qquad w(v'-y^{p\rho})\ge H.
\]
Taking $u_{m+1}=u'$ and $v_{m+1}=v'$ completes \eqref{eq:full-cartan-induction-field}--\eqref{eq:full-cartan-induction-error}.

\paragraph{Ramification and residue fields.}
The final induction step has $m=n-1$, so $K_m=F^C$ and $r_m=r$. Its graded field is $k(X^r,Y^r)$. The value group is generated by $rA,rB$ and is therefore $r\Z$, since $\gcd(A,B)=1$. Its homogeneous degree-zero field is
\[
 k\bigl((X^r)^B/(Y^r)^A\bigr)=k(t^r).
\]
This proves \eqref{eq:full-cartan-field}: the residue extension $k(t)/k(t^r)$ is purely inseparable of degree $r$. The product $ef=r^2$ agrees with the group order after the fixed field has been computed.

At $n=4$, Proposition~\ref{prop:divisorial-scalar-level-four} gives $e(F/F^S)=p$ and $f(F/F^S)=p^2$. Multiplicativity in the actual tower $F/F^S/F^C$ now yields \eqref{eq:full-cartan-scalar-relative}. The intermediate residue extension is $k(t^{p^2})/k(t^{p^3})$, hence purely inseparable.
\end{proof}

The exact Cartan ramification degree gives a sharper conditional index for the further tame quotient.
\begin{corollary}\label{cor:divisorial-full-cartan-6500}
Take $p=5$ and $n=4$ in Proposition~\ref{prop:divisorial-full-cartan}. If $F^C/K$ is a further tame extension of ramification index $12$ with unchanged residue field, then the absolute ramification index of $K$ is exactly $6500$.
\end{corollary}
\begin{proof}
The valuation on $F$ gives $w(5)=624\cdot5^6$, and the proposition gives $e(F/F^C)=5^3$. Thus
\[
 e_K=\frac{624\cdot5^6}{12\cdot5^3}=6500.
\]
The existence of the further tame extension is a hypothesis. This calculation concerns the specified divisorial valuation and supplies no perfect-residue specialization or actual unequal-character boundary pair.
\end{proof}

\subsection{Regularity on the exceptional torus}\label{sec:cartan-torus-regularity}
The field calculation extends to the local rings on the exceptional torus and gives an integral obstruction to specializing this Cartan construction.

\begin{proposition}\label{prop:cartan-torus-regularity}
Use the notation of Proposition~\ref{prop:divisorial-full-cartan}. In the fraction field of the regular full-level ring $R_n$, put
\[
 \pi=x^{1-p}y^p,\qquad t=x^B/y^A,\qquad
 S_0=R_n[\pi,t,t^{-1}].
\]
For $c\in k^*$, let $[c]\in W(k)$ be its Teichm\"uller representative and set
\[
 B_c=(S_0)_{(\pi,t-[c])},\qquad A_c=B_c^C,
 \qquad P=p^{n-1}.
\]
Then $B_c$ is a regular local ring of dimension two, the principal Cartan acts on $B_c$, and $A_c$ is regular. The unique reduced vertical divisor of $A_c$ is regular and Cartier. It has a local equation $\rho$ for which
\[
 \rho/\pi^P\in B_c^*,\qquad
 A_c/(\rho)=k[t^P,t^{-P}]_{(t^P-c^P)}
\]
as subrings of the divisorial residue field. After multiplication by a unit of $A_c$, one may arrange $\operatorname{in}(\rho)=\operatorname{in}(\pi)^P$. Moreover
\begin{equation}\label{eq:cartan-torus-p-factorization}
 p=u\rho^{(p^4-1)p^{n-1}},\qquad u\in A_c^*.
\end{equation}
\end{proposition}

\begin{proof}
\par\medskip\noindent\textbf{The regular chart and the action.}\enspace
The identity $(1-p)A+pB=1$ gives the inverse monomial formulas
\begin{equation}\label{eq:cartan-torus-coordinates}
 x=\pi^A t^p,\qquad y=\pi^B t^{p-1}.
\end{equation}
The quotient $S_0/(\pi)$ is exactly $k[t,t^{-1}]$. Indeed $x,y$ vanish in this quotient by~\eqref{eq:cartan-torus-coordinates}, and they generate the maximal ideal of $R_n$. Thus the quotient is generated by $t,t^{-1}$ over $k$. The monomial valuation sends $t$ to a transcendental residue in $k(t)$, which excludes any further Laurent-polynomial relation. Consequently $(\pi,t-[c])$ is a maximal ideal, and the prime chain
\[
 (0)\subsetneq(\pi)\subsetneq(\pi,t-[c])
\]
shows that its localization has dimension at least two. Its maximal ideal has two generators, so Krull's height theorem gives the reverse inequality. This proves regularity of $B_c$, with parameters $\pi,t-[c]$, and
\begin{equation}\label{eq:cartan-torus-original-divisor}
 B_c/(\pi)=k[t,t^{-1}]_{(t-c)}.
\end{equation}
These are uncompleted local rings; their fraction field remains $\operatorname{Frac}(R_n)$.

For every $g\in C$, the principal level action and the torsion initial forms proved above give
\[
 w(gx-x)>A,\qquad w(gy-y)>B.
\]
In the regular local ring $B_c$, divisibility by $\pi^a$ is equivalent to divisorial value at least $a$. One can verify this by intersecting the height-one valuation rings: $\pi$ is a unit at every other height-one prime. Since $gx-x$ and $gy-y$ belong to $R_n\subset B_c$, formula~\eqref{eq:cartan-torus-coordinates} implies
\[
 gx/x,\ gy/y\in1+\pi B_c.
\]
The monomial definitions then give $g\pi/\pi,gt/t\in1+\pi B_c$. Thus $g$ sends the generators of $S_0$ into $B_c$, fixes its reduction modulo $\pi$, and preserves the closed point $t=c$. Denominators outside that point map to units. Applying the same argument to $g^{-1}$ proves that $g$ is an automorphism of $B_c$.

All invariant rings below are finite Noetherian local quotients. Here is a direct verification. The full-level ring is finite over the fixed Noetherian deformation base $R_0$. Form the $R_0$-algebra $S$ generated by all translates of a finite set of $R_0$-algebra generators of $S_0$. Then $S$ is a finite-type $C$-stable subalgebra of $B_c$. If $\mathfrak m$ is the contraction of the maximal ideal of $B_c$, the inclusions $S_0\subset S\subset B_c$ give $B_c=S_{\mathfrak m}$. For any subgroup $H$, the coefficients of the finitely many orbit polynomials generate a Noetherian subring of $S^H$ over which $S$ is finite. Its submodule $S^H$ is therefore finite and Noetherian. Since $\mathfrak m$ is $H$-stable, the norm of any denominator outside $\mathfrak m$ is invariant and remains outside it. Hence
\[
 (B_c)^H=(S^H)_{\mathfrak m\cap S^H},
\]
and $B_c$ is finite over this ring. Normality follows from normality of $B_c$. All closed residue fields are $k$, because the action fixes $W(k)$.

\par\medskip\noindent\textbf{The unique vertical divisor.}\enspace
Write $N=p^{2n-2}$ and $E=(p^4-1)N$. The full level structure gives a Weierstrass factorization of $[p](T)$. Differentiating it at the zero point expresses $p$ as a unit times the product of all nonzero first-level torsion coordinates. There are $p-1$ points with initial form a nonzero multiple of $X^N$, and $p(p-1)$ with initial form a nonzero multiple of $Y^N$. Thus
\[
 \operatorname{in}(p)=a_0X^{(p-1)N}Y^{p(p-1)N},
 \qquad a_0\in k^*.
\]
By~\eqref{eq:cartan-torus-coordinates}, this is
\[
 a_0\operatorname{in}(\pi)^E t^{p^2(p-1)N}.
\]
The Laurent coefficient is a unit at every $c\in k^*$. Divisibility in $B_c$ therefore proves
\begin{equation}\label{eq:cartan-torus-upstairs-unit}
 p=\pi^E h,\qquad h\in B_c^*.
\end{equation}
There is exactly one vertical prime through this chart, and the same holds in every intermediate finite invariant quotient.

\par\medskip\noindent\textbf{Field information and actual local rings.}\enspace
We use the complete induction of Proposition~\ref{prop:divisorial-full-cartan}, including its first derivations and actual fixed-field degrees. Its order is cross first, then scalar. The completed fixed fields also describe the divisorial completions of the invariant fraction fields of the local rings. To verify this point, approximate an invariant element of $F$ by an element of $\operatorname{Frac}(R_n)$ and average over the finite subgroup. Division by its order loses only a fixed finite amount of valuation. Thus invariant elements of the dense fraction field are dense in the completed fixed field. This argument takes place in fields and makes no claim about integral averaging.

At the beginning of a complete congruence layer write
\[
 r=p^m,\quad j=n-m-1\ge1,\quad Q=p^{2j},\quad
 U=X^r,\quad V=Y^r,
\]
and introduce the homogeneous quantities
\begin{equation}\label{eq:cartan-torus-layer-coordinates}
 P_0=U^{1-p}V^p=\operatorname{in}(\pi)^r,
 \qquad T=U^B/V^A=t^r.
\end{equation}
They satisfy $U=P_0^A T^p$ and $V=P_0^B T^{p-1}$. The current graded field is $k(U,V)$, its value group is $r\Z$, and its residue field is $k(T)$.

Inductively, let $\mathcal B$ be the current actual local invariant ring. We prove simultaneously that $\mathcal B$ and its reduced vertical divisor are regular. If $f$ generates the vertical prime, its pullback to $B_c$ is $\pi^r$ times a unit: the unique vertical divisor has that multiplicity by the field ramification calculation. Thus
\[
 \operatorname{in}(f)=P_0 a(T),
\]
where $a$ is regular and nonzero at $T=c^r$. The divisor ring is the local ring of $k(T)$ at $c^r$. Indeed it is a discrete valuation ring with that fraction field, dominated by the finite local ring~\eqref{eq:cartan-torus-original-divisor}; the restricted valuation is precisely the valuation at $T=c^r$. A parameter on this curve therefore has residue $g(T)$ with a simple zero at $c^r$, and $g'(c^r)\ne0$. It lifts to an element of $\mathcal B$. These local induction assertions hold at $m=0$ by the explicit chart construction.

\par\medskip\noindent\textbf{The cross step and its integral parameter.}\enspace
By~\eqref{eq:full-cartan-cross-derivation}, the first derivation of $\tau_j$ on $k(U,V)$ is $V^Q\partial_U$, with shift $r(QB-A)$. As $B=p^2$ is zero in characteristic $p$, it satisfies
\begin{equation}\label{eq:cartan-torus-cross-local}
 D_\tau(T)=0,\qquad
 D_\tau(P_0)=(1-p)P_0^{QB-A+1}T^{(p-1)Q-p}.
\end{equation}
Put $d=QB-A>0$. The second expression has a nonzero Laurent-monomial coefficient at the closed point. Since $D_\tau$ kills $a(T)$, the difference $\tau_j(f)-f$ is $f^{d+1}$ times a unit in $\mathcal B$.

For any element of positive divisorial value in $\mathcal B$, the global shift bound gives displacement at least $r(d+1)$. For an element of value zero, the initial form belongs to $k(T)$ and is killed by~\eqref{eq:cartan-torus-cross-local}. Its displacement has value strictly greater than $rd$, hence at least $r(d+1)$ by discreteness. Divisibility along the prime parameter $f$ now shows that the entire augmentation ideal is
\[
 I_{\tau_j}=(\tau_j(b)-b:b\in\mathcal B)=(f^{d+1}).
\]
The principal-augmentation theorem of \citet[Theorem~2 and Proposition~5, pp.~64--69]{KiralyLutkebohmert2013} proves that $\mathcal A=\mathcal B^{\langle\tau_j\rangle}$ is regular, that $\mathcal B$ is finite free over $\mathcal A$, and, since the augmentation ideal is generated by $\tau_j(f)-f$, that
\[
 \mathcal B=\mathcal A[f],\qquad \mathcal B\simeq\mathcal A[Z]/(P_f(Z)),
\]
where $P_f$ is the monic degree-$p$ orbit polynomial and $1,f,\ldots,f^{p-1}$ is an $\mathcal A$-basis. All the hypotheses of that theorem hold here: the rings are Noetherian normal local domains, $\mathcal B$ is regular, and the action has order $p$.

The actual fixed graded field is $k(U^p,V)$ by~\eqref{eq:full-cartan-cross-field}. Its value group is $rp\Z$, and its residue field remains $k(T)$. Therefore $q=\operatorname{Norm}(f)$ has multiplicity one at the unique vertical prime of $\mathcal A$. It has no other prime divisor, since each conjugate of $f$ is an equation of the vertical divisor. Regularity of $\mathcal A$ implies that $q$ generates that prime up to a unit. The monogenic presentation gives the exact isomorphism
\begin{equation}\label{eq:cartan-torus-norm-divisor}
 \mathcal B/(f)\simeq\mathcal A/(q).
\end{equation}
The reduced divisor of $\mathcal A$ is consequently regular. Moreover
\[
 \operatorname{in}(q)=P_0^p a(T)^p,
\]
because all conjugates of $f$ have the same initial form. This constructs the needed parameter inside the actual local invariant ring, without extending a fractional controlled field lift across the closed point.

\par\medskip\noindent\textbf{The scalar step and descent of the reduced divisor.}\enspace
Equations~\eqref{eq:full-cartan-scalar-b}--\eqref{eq:full-cartan-scalar-a} in the field induction show that on $k(U^p,V)$ the first scalar derivation is $D_\sigma=V^Q\partial_V$, with shift $rB(Q-1)$. It kills $P_0^p$ and satisfies
\begin{equation}\label{eq:cartan-torus-scalar-local}
 D_\sigma(T)=-A P_0^{B(Q-1)}T^{1+(p-1)(Q-1)}.
\end{equation}
Relative to the current value group $rp\Z$, its shift is the positive integer
\[
 d_s=B(Q-1)/p=p(Q-1).
\]
Choose a lift $z\in\mathcal A$ of a parameter $g(T)$ of the regular divisor in~\eqref{eq:cartan-torus-norm-divisor}. Its simple zero gives $g'(c^r)\ne0$. Formula~\eqref{eq:cartan-torus-scalar-local} shows that $\sigma_j(z)-z$ is $q^{d_s}$ times a unit: its normalized leading coefficient is
\[
 \frac{-A g'(T)T^{1+(p-1)(Q-1)}}{a(T)^{p d_s}},
\]
which is nonzero at $c^r$. The global shift bound puts every local difference in $(q^{d_s})$. Hence the augmentation ideal is exactly this principal ideal.

The same principal-augmentation theorem proves that $\mathcal A'=\mathcal A^{\langle\sigma_j\rangle}$ is regular and $\mathcal A$ is finite free over $\mathcal A'$. The actual scalar field step has ramification index one and residue field $k(T^p)$, by~\eqref{eq:full-cartan-next-field}. Let $q'$ generate the unique vertical prime of the regular ring $\mathcal A'$. Its pullback is $q$ times a unit, because that prime has ramification index one and no other vertical prime lies above it. Thus $\mathcal A/(q')=\mathcal A/(q)$ is regular, and
\[
 \mathcal A'/(q')\longrightarrow\mathcal A/(q')
\]
is faithfully flat. Regularity descends under this map, so the reduced divisor of $\mathcal A'$ is regular as well as Cartier. Its local ring is
\[
 k[T^p,T^{-p}]_{(T^p-c^{rp})}.
\]
Its vertical parameter pulls back to $\pi^{rp}$ times a unit. In the graded field $k(U^p,V^p)$ its initial coefficient is a rational function of $T^p$ that is regular and nonzero at the indicated point. We have obtained every local induction assertion with $r$ replaced by $rp$.

\par\medskip\noindent\textbf{The final unit factorization.}\enspace
Iteration through all $n-1$ congruence layers proves regularity of $A_c$ and of its reduced vertical divisor. With $P=p^{n-1}$, a local equation $\rho$ has pullback $\pi^P$ times a unit and the asserted divisor ring. Its initial coefficient belongs to $k(t^P)$ and is a unit in that divisor ring. Lifting its inverse to a unit of $A_c$ normalizes $\operatorname{in}(\rho)=\operatorname{in}(\pi)^P$.

Equation~\eqref{eq:cartan-torus-upstairs-unit} and the pullback-unit statement show that $p/\rho^{E/P}$ and its inverse are units in $B_c$. They are invariant, so they belong to $A_c^*$. Since $E/P=(p^4-1)p^{n-1}$, this proves~\eqref{eq:cartan-torus-p-factorization}.
\end{proof}

\begin{corollary}\label{cor:cartan-torus-specialization}
Every local homomorphism $A_c\to\mathcal O_L$ to a mixed-characteristic discrete valuation ring, with $p$ nonzero, satisfies
\[
 (p^4-1)p^{n-1}\mid e_L.
\]
For $p=5$ and $n=4$, this gives $78000\mid e_L$. If in addition $L/K$ is a finite extension with $e(L/K)\le12$, then $e_K\ge6500$.
\end{corollary}
\begin{proof}
Under the local map, $\rho$ has positive integral valuation and the unit in~\eqref{eq:cartan-torus-p-factorization} remains a unit. Therefore
\[
 e_L=(p^4-1)p^{n-1}v_L(\rho).
\]
At $p=5,n=4$, the coefficient is $624\cdot125=78000$. Multiplicativity $e_L=e(L/K)e_K$ gives the last inequality. Existence of such an extension is not asserted.
\end{proof}

The proposition concerns the points $t\ne0,\infty$ of this particular Cartan chart. It supplies no reduction of every surviving torsion comparison to these local rings. The general first-boundary occurrence question is therefore unchanged. The known non-Cartier quotient below explains why a field ramification calculation, without such a local regularity proof, does not imply the same specialization divisibility for arbitrary normal models.

\subsection{Contact orders in a perfect-residue specialization}
\label{subsec:boundary-specialization-contact}

This subsection gives a necessary contact condition for a specialization
that realizes an unequal-character comparison at the first unresolved
index.  A separate elementary example explains how an inseparable
residue extension can become ramification along a transverse curve.

\begin{proposition}[A necessary upstairs contact order]
\label{prop:boundary-specialization-contact}
Let $K$ be a complete discretely valued field of mixed characteristic
$(0,5)$ with algebraically closed residue field, $\mu _5\subset K$, and
$e_K=260$.  Suppose that $G,H$ belong to the standing height-two
filtration class, that there is an actual isomorphism
$G[5^4]\simeq H[5^4]$ over $K$, and that their special characters differ.
Let $F/K$ be a finite extension over which all the geometric
$5^4$-torsion points of $G$ are rational.

Put $E_*=624\cdot5^6$.  Suppose a local ring $S$ admits a local
homomorphism $S\to\mathcal O_F$ and elements $\pi,h\in S$ satisfying
\[
  5=\pi^{E_*}h,\qquad h\in S^\times,
  \qquad q:=v_F(\pi)>0,
\]
where $v_F$ is normalized to have value group $\mathbf Z$.
Then $5\mid q$.  In particular, the displayed chart cannot realize such
a comparison by an upstairs transverse curve with $q=1$.
\end{proposition}

\begin{proof}
The hypotheses are the case $p=5$, $r=1$ of
Theorem~\ref{thm:first-boundary-normal-form}.  Since the residue field is
already algebraically closed, Corollary~\ref{cor:boundary-point-field-degrees}
applies after a common tame good extension $L/K$, which may be chosen
Galois of degree dividing $24$.  If $Q$ is an exact $5^4$-torsion point
whose $5^3$-multiple is a nonzero canonical point, that corollary gives
\[
  [L(Q):L]=5^6.
\]
The full-level hypothesis on $F$ implies $Q\in G(FL)$.  Consequently
$5^6\mid[FL:L]$.  The two tower formulas give
\[
  [FL:L][L:K]=[FL:F][F:K].
\]
Both $[L:K]$ and $[FL:F]$ are prime to $5$, because $L/K$ is Galois
and tame.  It follows that
\[
  5^6\mid[F:K].
\]
This argument requires neither $F/K$ nor a whole point-field tower
to be Galois.

Every finite extension of $K$ has residue field equal to the given
algebraically closed field.  The degree formula for finite extensions
of complete discrete valuation fields therefore gives
$[F:K]=e(F/K)$; see
\citet[Section~15.113, Remark~15.113.6]{StacksDVR}.
The local homomorphism carries $h$ to a unit, so
\[
  v_F(5)=E_*q,
  \qquad
  [F:K]=\frac{v_F(5)}{e_K}
       =\frac{624\cdot5^6q}{260}
       =12\cdot5^5q.
\]
The divisibility $5^6\mid[F:K]$ now implies $5\mid q$.
\end{proof}

Here $\pi$ is a local equation of the upstairs divisor, and the unit
hypothesis requires that $5/\pi^{E_*}$ remain invertible at the chosen
closed point.  The proposition does not apply in this form if that
factor vanishes.  It also does not exclude a downstairs transverse
curve whose lift has contact $5$ or a larger multiple of $5$ upstairs.
For example, if a finite extension of degree $12\cdot5^6$ were retained
under these hypotheses, the displayed degree identity would force
$q=5$.  No existence of such an extension or of the required comparison
is asserted.

The specialization must be made on a model.  There is no $k$-algebra
homomorphism $k(t)\to k$: if $t$ mapped to $c\in k$, the invertible
element $t-c$ would map to zero.  Thus one must choose the closed point
and transverse curve before inverting all nonzero residue functions.

\begin{example}[An inseparable residue extension becoming ramification]
\label{ex:fierce-transverse-specialization}
Let $k=\overline{\mathbf F}_5$.  Choose a complete discrete valuation
field $K_0$ with residue field $k$, $\mu _5\subset K_0$, and
$e_{K_0}=260$.  Such a field is obtained from $\operatorname{Frac}W(k)$
by adjoining a primitive fifth root of unity $\zeta$, then a root
$\varpi$ of
\[
  \varpi^{65}=\zeta-1.
\]
The translated cyclotomic polynomial is Eisenstein of degree $4$ over
$W(k)$, and the second polynomial is Eisenstein of degree $65$ over
the resulting valuation ring.  Hence $\varpi$ is a uniformizer of
$K_0$ and $v_{K_0}(5)=4\cdot65=260$.

Consider the finite flat degree-five cover
\[
 R=\mathcal O_{K_0}[s]_{(\varpi,s-1)},
 \qquad B=R[a]/(a^5-s).
\]
At the generic point of the divisor $\varpi=0$, complete the downstairs
discrete valuation ring and denote its fraction field by $E$.
Its residue field is $k(s)$.  The polynomial $T^5-s$ is irreducible
over $k(s)$, so the corresponding integral algebra over the complete
valuation ring of $E$ has residue field $k(s^{1/5})$.  It is local,
with maximal ideal generated by $\varpi$, and is therefore a discrete
valuation ring.  It is free of rank $5$ over the downstairs ring.
Consequently the extension $E(a)/E$ has
\[
  e(E(a)/E)=1,\qquad
  f(E(a)/E)=5,
\]
with purely inseparable residue extension.  It is a cyclic extension
of fields because $\mu _5\subset E$.

Now restrict the cover to the downstairs curve $s=1+\varpi$.
This curve is transverse to $\varpi=0$, since its coordinate ring is
$\mathcal O_{K_0}$ and its intersection with the divisor has length
one.  Putting $z=a-1$ gives
\[
  z^5+5z^4+10z^3+10z^2+5z-\varpi=0.
\]
This polynomial is Eisenstein at $\varpi$.  The specialized fraction
field extension therefore has degree $5$, ramification index $5$,
and residue degree $1$.  The downstairs absolute index remains
$260$, while the upstairs contact order with $\varpi=0$ is $5$.

By comparison, the downstairs transverse curve $s=1$ gives $a^5=1$;
its generic fibre splits into five copies of $K_0$.
Thus the generic inseparable residue extension determines neither
retention nor loss of degree under transverse specialization.
This example concerns a Kummer cover only.  It supplies neither a
Drinfeld-level specialization nor a pair of groups with unequal
special characters and isomorphic $5^4$ torsion.
\end{example}

\subsection{Cartier indices and specialization through a singular quotient}
\label{subsec:noncartier-specialization}

This subsection applies the cyclic quotient of
\citet[Example~5.10, pp.~24--25]{LiedtkeMatsumoto2018} to specialization through a
non-Cartier vertical divisor. Their construction has a regular covering
surface and a normal singular quotient with nonreduced special fibre.
We use the same construction over a complete discrete valuation ring
containing $\mu_p$, and compute its Cartier index and an explicit
perfect-residue arc. The purpose is to identify the extra local
information needed to pass from a divisorial ramification bound to a
bound on every specialization.
For broader context on wild surface quotients and their divisor and
discriminant groups, see \citet{Lorenzini2013WildQuotients} and
\citet{LorenziniSchroer2023WildQuotients}.

\begin{lemma}[The local Cartier-index constraint]
\label{lem:cartier-index-contact}
Let $R$ be a normal Noetherian local domain of dimension two and mixed
characteristic $(0,p)$. Suppose that
\[
                   \operatorname{div}_R(p)=N D
\]
for a single height-one prime divisor $D$ and an integer $N>0$.
Let $i$ be the least positive integer for which $iD$ is Cartier at the
closed point. Then $i$ exists and divides $N$. If
$\phi:R\longrightarrow\mathcal O_K$ is a local map to a complete
mixed-characteristic discrete valuation ring, with $p\ne0$ and
$e_K=v_K(p)$, then
\begin{equation}
               \frac{N}{\gcd(N,e_K)}\mid i.
\label{eq:cartier-specialization-divisibility}
\end{equation}
In particular, $N=1300$ and $e_K=260$ require $5\mid i$.
More generally, if $N=52\cdot5^j$ with $j\ge2$, an index-$260$
specialization requires $5^{j-1}\mid i$.
\end{lemma}

\begin{proof}
The integers $a$ for which $aD$ is locally principal form a subgroup of
$\mathbf Z$ containing $N$. Since every invertible module over a local
ring is free, its positive generator is the Cartier index $i$, and
$i\mid N$.
Choose $f\in\operatorname{Frac}(R)$ with
$\operatorname{div}_R(f)=iD$. Effectivity and the height-one intersection
property for a normal Noetherian domain give $f\in R$; see
\cite[Lemma~10.157.6]{StacksHeightOneIntersection}.
The rational function $p/f^{N/i}$ and its inverse have nonnegative
valuation at every height-one prime. The same intersection property
therefore gives
\[
                      p=u f^{N/i},\qquad u\in R^\times.
\]
The image of $u$ under $\phi$ is a unit, and the image of $f$ is nonzero.
Taking valuations yields
\[
                      e_K=(N/i)v_K(\phi(f)).
\]
Thus $N\mid i e_K$, which is equivalent to
\eqref{eq:cartier-specialization-divisibility}.
The two numerical consequences follow by taking the indicated greatest
common divisors.
\end{proof}

If the reduced vertical divisor is Cartier, the same proof gives
$e_K\ge N$ for a local specialization through it. In general it gives
only $e_K\ge N/i$. The following example attains this latter bound.

\begin{example}[The Liedtke--Matsumoto quotient at index $260$]
\label{ex:noncartier-quotient-specialization}
There is a complete field $K$ of mixed characteristic $(0,5)$ with
algebraically closed residue field $k$, $\mu_5\subset K$, and $e_K=260$,
and a finite quotient $A=B^{C_5}$ of a regular local surface $B$, such
that $A$ is a normal local surface and the following hold.
The generic-fibre cover is finite \'etale of degree five.
The unique reduced vertical prime $D$ has residue field $k(Y)$,
normalized value $v_D(5)=1300$, and Cartier index five.
The cover of divisorial valuation rings has ramification index one and
purely inseparable residue degree five.
Nevertheless there is a local map $A\to\mathcal O_K$, not contained in
$D$, for which
\[
            \operatorname{length}_{\mathcal O_K}
            (\mathcal O_K/I_D\mathcal O_K)=1.
\]
Here $I_D$ denotes the prime ideal of the reduced divisor.
The generic pullback of the cover is a connected degree-five field
extension of $K$. Its normalization has absolute index $1300$ and
residue field $k$.
\end{example}

\begin{proof}
We first give the same cyclic construction for any prime $p$ and any complete
mixed-characteristic discrete valuation ring $\mathcal O$ with
uniformizer $\pi$, residue field $k$ of characteristic $p$, and a
primitive $p$th root of unity $\zeta\in\mathcal O$.
Put $e_0=v_{\mathcal O}(p)$ and
\[
 \mathcal O'=\mathcal O[t]/(t^p-\pi),\qquad
 B_0=\mathcal O'[U],\qquad B=(B_0)_{(t,U)}.
\]
Eisenstein's criterion shows that $\mathcal O'$ is a complete discrete
valuation ring with uniformizer $t$ and residue field $k$.
Thus $B$ is a regular local ring of dimension two. The formulas
\[
                 \sigma(t)=\zeta t,\qquad
                 \sigma(U)=\zeta^{-1}U
\]
define a faithful action of the constant cyclic group $C_p$, fixing
$\mathcal O$. Set $X=tU$ and $Y=U^p$. We claim that
\begin{equation}
 A_0:=B_0^{C_p}=\mathcal O[X,Y]/(X^p-\pi Y),\qquad
 A:=B^{C_p}=(A_0)_{(\pi,X,Y)}.
\label{eq:noncartier-cyclic-ring}
\end{equation}

Indeed, every element of $B_0$ has a unique finite expansion in monomials
$t^aU^b$, with $0\le a<p$, $b\ge0$, and coefficients in $\mathcal O$.
The monomial has character $\zeta^{a-b}$. Since the coefficient domain
has characteristic zero, invariance forces $b=a+pj$ for an integer
$j\ge0$, giving $t^aU^b=X^aY^j$.
Conversely these monomials are invariant. They show that $A_0$ is free
on $1,X,\ldots,X^{p-1}$ over $\mathcal O[Y]$, with the single relation
$X^p=\pi Y$. The ring $B_0$ is finite over $A_0$, generated as a module
by $t^aU^b$ with $0\le a,b<p$.

There is only one prime of $B_0$ above $(\pi,X,Y)$, because its closed
fibre has ring
\[
                         k[t,U]/(t^p,tU,U^p).
\]
Localization over $A_0$ therefore gives $B$.
To verify invariants after localization directly, write a fixed
element as $f/s$ with $s\notin(t,U)$. The norm
$\prod_{g\in C_p}g(s)$ is invariant and is outside $(\pi,X,Y)$.
Multiplication by this norm gives an invariant numerator in $B_0$,
hence one in $A_0$. This proves~\eqref{eq:noncartier-cyclic-ring}.

The ring $A$ is a Noetherian local domain of dimension two. It is normal:
an element of $\operatorname{Frac}(A)$ integral over $A$ is integral over
$B$, lies in the normal ring $B$, and is fixed, so belongs to $A$.
Here normality of regular rings is the standard result
\cite[Lemma~10.157.5]{StacksRegularNormal}.
After inverting $p$, both $\pi$ and $t$ are units and $U=X/t$, so
\[
                    B[1/p]=A[1/p][t]/(t^p-\pi).
\]
The degree is $p$: the $p$ distinct elements $\zeta^a t$ are conjugate,
and $t^p=\pi$. The derivative $pt^{p-1}$ is a unit, proving finite
\'etaleness of the generic-fibre cover.

Modulo $\pi$, the equation of $A$ is $X^p=0$. Hence its unique reduced
vertical prime is $D=(\pi,X)$, with residue field $k(Y)$.
At $D$, the element $Y$ is a unit and $\pi=X^p/Y$.
The discrete valuation ring $A_D$ has uniformizer $X$, and
\begin{equation}
       v_D(\pi)=p,\qquad v_D(p)=p e_0,
       \qquad\operatorname{div}_A(\pi)=pD.
\label{eq:noncartier-vertical-values}
\end{equation}
The corresponding prime of $B$ is $E=(t)$, where $U$ is a unit.
Thus $v_E(X)=v_E(tU)=1=v_D(X)$, giving relative ramification index one.
The residue extension is $k(Y)\subset k(U)$ with $Y=U^p$, purely
inseparable of degree $p$.

The Cartier index of $D$ divides $p$ by
\eqref{eq:noncartier-vertical-values}. It is not one.
The classes of $\pi,X,Y$ are linearly independent in the cotangent
space of $A$, because $X^p-\pi Y$ belongs to the square of the maximal
ideal of the ambient regular local ring
$\mathcal O[X,Y]_{(\pi,X,Y)}$.
Consequently $I_D=(\pi,X)$ has two independent classes modulo
$\mathfrak m_A I_D$: this latter ideal is contained in $\mathfrak m_A^2$.
The ideal $I_D$ is therefore not principal. In a normal domain it is
the divisorial ideal $\mathcal O(-D)$, so a Cartier divisor $D$ would
make it invertible, hence principal over the local ring. This is a
contradiction. Since $p$ is prime, the Cartier index is exactly $p$.

Now define
\begin{equation}
                  A\longrightarrow\mathcal O,
                  \qquad X\longmapsto\pi,\quad
                         Y\longmapsto\pi^{p-1}.
\label{eq:noncartier-specialization-map}
\end{equation}
The relation is satisfied, and every localized denominator maps to a
unit. The map is local, is not contained in the vertical divisor,
and takes $I_D$ to $(\pi)$, proving the stated intersection length one.
Its target has absolute index $e_0$, while the vertical valuation has
index $pe_0$.
On generic fibres the pullback is
$\operatorname{Frac}(\mathcal O)[t]/(t^p-\pi)$, with
$U=\pi/t=t^{p-1}$. It is an Eisenstein degree-$p$ field extension.
The normalization of its dominating integral component is
$\mathcal O'$, of absolute index $pe_0$ and residue field $k$.
This assertion does not require the whole closed-point tensor product
to be flat over $\mathcal O$.

For the numerical example take $p=5$ and $k=\overline{\mathbf F}_5$.
Start with $\operatorname{Frac}W(k)$, adjoin a primitive fifth root
$\zeta$, and then adjoin $\pi$ with $\pi^{65}=\zeta-1$.
The polynomial
\[
                   \Phi_5(1+Z)=Z^4+5Z^3+10Z^2+10Z+5
\]
is Eisenstein at $5$, so the first extension has degree four and
uniformizer $\zeta-1$. The second polynomial is Eisenstein of degree
$65$. The resulting complete field $K$ has residue $k$, contains
$\mu_5$, and has $e_K=4\cdot65=260$.
Applying the construction with $\mathcal O=\mathcal O_K$ gives
$v_D(5)=5\cdot260=1300$ and every asserted property.
\end{proof}

The example shows that normality, a finite constant-group quotient,
and generic \'etale structure do not by themselves
promote the specified divisorial bound to all perfect-residue curves.
Its reduced vertical prime has precisely the Cartier-index denominator
allowed by Lemma~\ref{lem:cartier-index-contact}.
This ring and curve are auxiliary examples: they have not been identified
with the full Drinfeld-level or nonsplit Cartan quotient, and they supply
no pair of height-two groups with unequal special characters and
isomorphic $5^4$-torsion.
For that actual comparison, it remains to decide whether its finite-level
marking factors through the exceptional-torus chart of
Proposition~\ref{prop:cartan-torus-regularity}, or through another quotient.
That proposition resolves the local divisor question on its stated charts;
it does not identify all curves compatible with the torsion isomorphism.

\subsection{Point orbits and quadratic self-twists}
\label{sec:cartan-image-constraints}

This subsection distinguishes the information carried by a primitive-point
orbit from the stronger existence of a quadratic self-twist of the entire
finite representation.

\begin{proposition}[Maximal point orbits with noncommutative image]
\label{prop:cartan-orbit-data}
Let $p$ be odd, $n\ge2$, $\Lambda_n=\mathbf Z/p^n\mathbf Z$, and let
$D$ be any subgroup of $1+p\Lambda_n$. Put
\[
 U_n(D)=\{g\in\mathrm{GL}_2(\Lambda_n):g\equiv I\pmod p,
                                  \ \det(g)\in D\}.
\]
At each level $1\le j\le n$, the orbit of every primitive vector
$v\in \Lambda_j^2$ is exactly $v+p\Lambda_j^2$, of size $p^{2j-2}$.
Moreover
\[
                  |U_n(D)|=p^{3n-3}|D|,
\]
and $U_n(D)$ is noncommutative when $n\ge3$.
In particular, at $n=4$, the orbit sizes $p^2,p^4,p^6$, even together
with a prescribed determinant image, do not force containment in a
Cartan subgroup.
\end{proposition}

\begin{proof}
The residue of every vector is fixed. To prove transitivity on its
residue fibre, it suffices to use the determinant-one congruence
subgroup. First take $v=e_1$. Any target in this fibre has coordinates
$(a,c)^t$, with $a\equiv1\pmod p$ and $c\equiv0\pmod p$.
Lift $a,c$ to $\Lambda_n$. The matrix
\[
                    \begin{pmatrix}a&0\\c&a^{-1}\end{pmatrix}
\]
has determinant one, is the identity modulo $p$, and carries $e_1$
to the target at level $j$. For arbitrary primitive $v$, lift it
primitively to $\Lambda_n$ and complete it to a basis. Conjugating the
preceding construction by this basis change proves transitivity.
There are $p^{2j-2}$ vectors in the fibre.

The full first congruence subgroup has $p^{4n-4}$ elements: each of
the four entries of $g-I$ is an arbitrary element of $p\Lambda_n$, and
such a matrix is invertible. Its determinant maps onto $1+p\Lambda_n$,
since $\operatorname{diag}(d,1)$ realizes every element $d$ of that
group. Each determinant fibre therefore has $p^{3n-3}$ elements,
which gives the displayed order. Finally, the two elements
$I+pE_{12}$ and $I+pE_{21}$ have determinant one and their products
in the two orders differ by
\[
                         p^2(E_{11}-E_{22}).
\]
This is nonzero modulo $p^n$ when $n\ge3$. Every Cartan arising from
a quadratic commutative algebra is commutative, proving the last
assertion.
\end{proof}

One can also impose the residual markings of the boundary form
without changing this abstract obstruction. The diagonal matrix
$T=\operatorname{diag}(1,-1)$ normalizes $U_n(D)$ and has order two.
The residual representation of
$U_n(D)\rtimes\langle T\rangle$ is split with character ratio $-1$.
Give two copies opposite choices of their distinguished coordinate
line; their identity isomorphism exchanges the distinguished and
complementary markings. Under $U_n(D)$, every lift of either nonzero marked point
still has the stated orbit sizes; the full semidirect product may
have larger orbits because of its tame factor. If $p\equiv1\pmod4$, choose a
root $i\in \Lambda_n$ of $i^2=-1$ and instead take
$T=\operatorname{diag}(i,-i)$. Its determinant is one and its order
is four. The root lifts from $\mathbf F_p$ by successively solving
a linear equation with unit derivative $2i$. Thus even the tame
determinant may be fixed. These are abstract marked representations;
no good formal model or actual unequal-character pair is supplied.

\begin{lemma}[A quadratic self-twist]
\label{lem:quadratic-self-twist-cartan}
Let $p$ be odd, $R=\mathbf Z/p^n\mathbf Z$,
$\rho:\Gamma\to\mathrm{GL}_2(R)$, and let
$\chi:\Gamma\to\{1,-1\}$ be nontrivial. Suppose that an invertible
matrix $J$ satisfies
\[
              J\rho(g)=\chi(g)\rho(g)J\qquad(g\in\Gamma).
\]
Then $\operatorname{tr}(J)=0$ and $J^2=\epsilon I$ for a unit
$\epsilon\in R$. The algebra $R[J]$ is quadratic \'etale, and
\[
                   \rho(\ker\chi)\subset R[J]^\times.
\]
If a subgroup of $\ker\chi$ acts trivially modulo $p$, its image
lies in $1+pR[J]$. This Cartan is nonsplit precisely when the
residue of $\epsilon$ is a nonsquare. A self-twist alone permits
the split case.
\end{lemma}

\begin{proof}
Choose $g$ with $\chi(g)=-1$. Conjugation takes $J$ to $-J$,
so $2\operatorname{tr}(J)=0$. Since $2$ is invertible, the trace
vanishes. Cayley--Hamilton gives
$J^2=-\det(J)I$; take $\epsilon=-\det(J)\in R^\times$.
The reduction of $J$ is not scalar, since an invertible scalar
matrix in odd characteristic has nonzero trace. A nonscalar
$2\times2$ matrix over a field has a cyclic vector: if every
vector spanned a stable line, its action on two basis vectors
and their sum would force it to be scalar. Lift a cyclic vector
$v$ to $R^2$. Then $(v,Jv)$ is a basis.

In this basis $J$ has matrix
$\left(\begin{smallmatrix}0&\epsilon\\1&0\end{smallmatrix}\right)$,
so $R[J]\simeq R[T]/(T^2-\epsilon)$. This is quadratic \'etale,
because $2T$ is a unit in that quotient. If $X$ commutes with $J$
and $Xv=av+bJv$, then $X(Jv)=aJv+b\epsilon v$. Thus
$X=aI+bJ$, proving that the full matrix centralizer is exactly
$R[J]$. Its invertible part is $R[J]^\times$, since the inverse
of a commuting invertible matrix still commutes with $J$.
This proves the asserted containment.

The matrices $I,J$ are linearly independent modulo $p$.
Consequently $aI+bJ\equiv I\pmod p$ forces
$a\equiv1$ and $b\equiv0\pmod p$, giving the principal-unit
containment. If the residue of $\epsilon$ is a square, its two
distinct roots lift and the Chinese remainder theorem gives
$R[J]\simeq R\times R$. If it is a nonsquare, the residue algebra
is $\mathbf F_{p^2}$ and the quadratic \'etale algebra is nonsplit.
Both types are possible: $J=\operatorname{diag}(1,-1)$
anticommutes with the coordinate-swap matrix, giving a split
self-twist example.
\end{proof}

A finite tame extension retains the wild inertia image: its
absolute Galois subgroup contains the subgroup fixing the maximal
tame extension. Thus a tame quadratic self-twist, if actually
identified, places wild inertia in the Cartan of the lemma.
The opposite tangent scalars and exchanged residual lines of
Theorem~\ref{thm:first-boundary-normal-form} do not by themselves
identify the two full deformations as quadratic twists. The
point-orbit restrictions alone are likewise insufficient by
Proposition~\ref{prop:cartan-orbit-data}. Neither elementary result
asserts an actual boundary realization or a boundary character
index.

\subsection{The Cartan chart under an actual image hypothesis}
\label{sec:boundary-nonsplit-image-proof}

This subsection constructs the relevant invariant-ring map from
an actual nonsplit-Cartan image and derives its incompatible
ramification divisibility. We first make the coefficient-field
scope of the universal frame map explicit.

\begin{lemma}[Scalar-compatible universal frames]
\label{lem:scalar-compatible-frame}
Put $k_0=\overline{\mathbf F}_p$, and let $R_n(k_0)$ be the universal
full-level deformation ring of a fixed one-dimensional height-two
framing group over $k_0$. Let $A$ be a complete mixed-characteristic
discrete valuation ring with algebraically closed residue
$\kappa\supset k_0$ and a fixed compatible coefficient map
$W(k_0)\subset W(\kappa)\to A$.
A good formal group over $A$ with the scalar-extended special
framing, together with a full $p^n$-frame over a finite extension
valuation ring $B$, gives a continuous homomorphism
\[
                           R_n(k_0)\longrightarrow B.
\]
The level-zero deformation and its identity-reducing formal
isomorphism can be chosen over $A$. If the fraction field of $B$
is the Galois splitting field, the displayed map is equivariant
for the actual Galois image acting on the frame. The integral
formal coordinate used in this evaluation preserves the valuation
of every positive-valuation point.
\end{lemma}

\begin{proof}
The printed deformation category of
\citet[Propositions~4.1--4.3, pp.~570--573]{Drinfeld1974} has residue
field $k_0$. The coefficient extension needed here follows from
the coefficientwise proof; we do not attribute a broader literal
statement to that source.

In the level-zero proof, the difference between successive
deformations over a square-zero layer
$M=\mathfrak m_A^a/\mathfrak m_A^{a+1}$ is a formal-module
$2$-cocycle with coefficients in $M$. Such a layer may be viewed
as an arbitrary $k_0$-vector space. The equations for the cocycle
and its coboundaries have coefficients in the original framing
group over $k_0$. The degree-elimination assertions~(a) and~(b)
in that proof, pp.~571--572, only invert nonzero coefficients of
$k_0$. They eliminate coefficients one degree at a time; each
fixed degree is modified only finitely often. They therefore
apply to $M$-valued coefficients as well as scalar coefficients.
There is no interchange of an infinite product of power-series
coefficients with a tensor product. At height two, one parameter of deformation classes remains,
represented by a change of the single universal parameter.

Starting from the given special framing, choose that parameter
modulo successive powers of $\mathfrak m_A$ to correct the
remaining cocycle class. The same degree elimination gives the
formal coordinate isomorphism at each step. The rigidity argument
of Proposition~4.1 is valid with these coefficients: applying
separating $k_0$-linear functionals to an infinitesimal
homomorphism reduces it to the scalar argument, and finite height
excludes a nonzero homomorphism to the additive formal group.
Thus an isomorphism reducing to the identity is unique. The
successive choices are compatible, and completeness gives the
parameter and formal isomorphism over $A$, hence the map
$R_0(k_0)\to A$. Its linear coordinate coefficient is a unit.
For a positive-valuation argument, the linear term then has
strictly smaller valuation than every higher term, proving
coordinate invariance of the valuation.

The higher-level construction in Proposition~4.3 adjoins roots of
successive primitive division factors at level one and then roots
of $[p](X_i)-x_i$ at later levels. The chosen actual frame supplies
these roots. At a primitive step, its next point is outside the
span already chosen; in the characteristic-zero fraction field,
all corresponding linear factors are nonzero at that point.
Since the full division series vanishes there, the remaining
primitive factor vanishes. The later division relations hold by
the definition of a compatible full frame. The root-adjunction
construction therefore gives the asserted map $R_n(k_0)\to B$.
The level-zero data were constructed over $A$, and their
identity-reducing isomorphism is unique. Change of frame is
functorial. These facts give the stated Galois equivariance.
\end{proof}

\begin{proof}[Proof of Corollary~\ref{cor:boundary-nonsplit-cartan-image}]
Work over the completed maximal unramified base, with algebraically
closed residue field $\kappa$, and let $L$ be the common tame good
extension of degree $m=p^2-1$. The residual representations split
into $\mathbf F_p$-valued characters by
Theorem~\ref{thm:first-boundary-normal-form}. Their wild inertia
is trivial and their tame orders divide $p-1$; hence first torsion
is rational over $L$. We keep the universal deformation and its
Cartan charts over $k_0=\overline{\mathbf F}_p\subset\kappa$.

Put
\[
 c_0=\frac{p-1}{2},\qquad A_0=p^2+p+1,\qquad B_0=p^2.
\]
The normal form gives $z=c_0p^{r+3}$ and
$e=c_0(p^2+1)p^r$. The canonical and noncanonical first-torsion
valuations are therefore
\begin{align*}
 A&=(p+1)e-z=c_0p^r A_0,\\
 B&=z/p=c_0p^r B_0.
\end{align*}
Both models are central. Lemma~\ref{rc-lem:newton} shows that every
exact $p^n$-point above a nonzero canonical, respectively
noncanonical, first-torsion point has valuation $A/p^{2n-2}$,
respectively $B/p^{2n-2}$. Here $v_L(L^\times)=\mathbf Z$ and
$v_L(p)=me$.

Suppose that the finite image for the first group $G$ is contained
in a conjugate of the principal nonsplit Cartan $C$. Write
$R=\mathbf Z/p^n\mathbf Z$ and realize its quadratic algebra as
$R[J]$, where $J^2=\epsilon I$ and the residue of $\epsilon$ is a
nonsquare. Choose a primitive point $v$ whose $p^{n-1}$-multiple
is a nonzero canonical point. The reduction of $J$ has no
$\mathbf F_p$-eigenline. Thus $(v,Jv)$ is an $R$-basis and
$p^{n-1}Jv$ is noncanonical. In this basis
\[
                          J=\begin{pmatrix}0&\epsilon\\1&0\end{pmatrix},
\]
so the Cartan has the standard form of
Proposition~\ref{prop:cartan-torus-regularity}. This basis makes
the first point canonical without changing the Cartan type.

Let $F=L(G[p^n])$ be the full splitting field, with Galois group
$\Gamma\subset C$, and let $x_F,y_F$ be the coordinates of
$v,Jv$. After choosing the special framing and a coefficient
map, Lemma~\ref{lem:scalar-compatible-frame} gives the equivariant
homomorphism $R_n(k_0)\to\mathcal O_F$ carrying $x,y$ to these
coordinates. All torsion points lie in the formal group because
the good special fibre is supersingular.

Extend $v_L$ to $F$ and put $q=c_0p^r/p^{2n-2}>0$.
Then $v_L(x_F)=qA_0$ and $v_L(y_F)=qB_0$. Since
$(1-p)A_0+pB_0=1$,
\[
 v_L(x_F^{1-p}y_F^p)=q>0,\qquad
 v_L(x_F^{B_0}/y_F^{A_0})=0.
\]
Thus $\pi_F=x_F^{1-p}y_F^p$ is integral of positive valuation
and $t_F=x_F^{B_0}/y_F^{A_0}$ is a unit. Substitution defines an
actual map
\[
 S_0=R_n(k_0)[\pi,t,t^{-1}]\longrightarrow\mathcal O_F.
\]
The images of $x,y$ are nonzero, so every rational expression used
here is defined. Its integrality follows from the displayed
valuations, without an assertion about fractional invariant lifts.
The residue of $t_F$ is some $c\in\kappa^\times$, and
$S_0/(\pi)=k_0[t,t^{-1}]$.

If $c$ is algebraic over $k_0$, then $c\in k_0$. Every denominator
outside $(\pi,t-[c])$ has nonzero residue, giving the actual local
map $B_c\to\mathcal O_F$ of the original closed chart.
If $c$ is transcendental over $k_0$, choose any $d\in k_0^\times$.
A denominator outside $(\pi,t-[d])$ reduces to a Laurent polynomial
$f(t)$ with $f(d)\ne0$. It is nonzero, so $f(c)\ne0$ as well.
Consequently there is still a ring map
$B_d\to\mathcal O_F$. In this second case the map is generally
not local at the closed point $d$; we do not apply a closed-point
locality assertion to it.

In either case, write $d=c$ in the first case and retain the chosen
$d$ in the second. The map on $B_d$ is $\Gamma$-equivariant.
Every $C$-invariant is $\Gamma$-invariant, even if $\Gamma$ is a
proper subgroup. Hence restriction gives a ring map
\[
                   A_d=B_d^C\longrightarrow
                   \mathcal O_F^\Gamma=\mathcal O_L.
\]
Use the exact identity inside the original invariant ring from
Proposition~\ref{prop:cartan-torus-regularity}:
\[
              p=u_d\rho_d^N,\qquad u_d\in A_d^\times,
              \qquad N=(p^4-1)p^{n-1}.
\]
A unit maps to a unit because its inverse maps into $\mathcal O_L$
as well. The image of $\rho_d$ is nonzero, since $p\ne0$, and has
positive integral valuation by the identity. Thus $N\mid e_L$
in both residue cases. This use of the unit identity requires no
locality at a transcendental closed point and no extension of
the chart theorem beyond $k_0$.

Finally,
\[
 e_L=(p^2-1)e=c_0(p^4-1)p^r,\qquad
 \frac{e_L}{N}=\frac{c_0}{p^2}
              =\frac{p-1}{2p^2}\in(0,1),
\]
which is impossible. The same argument applies to the second group.
\end{proof}

No equality between a point-orbit size and the full image order
was used. At $p=5,r=1$, the contradiction is
$e_L=6240<78000=N$. The proof excludes the nonsplit-Cartan image
hypothesis for an actual comparison; it does not force that
hypothesis for a general comparison.

\begin{proof}[Proof of Corollary~\ref{cor:boundary-nonsplit-self-twist}]
The completed unramified passage kills the unramified part of the
quadratic character. Its remaining tame part is killed by the
common extension of degree $p^2-1$. The intertwiner therefore
commutes with the image over $L$. By
Lemma~\ref{lem:quadratic-self-twist-cartan}, irreducible quadratic
reduction makes its centralizer the nonsplit quadratic \'etale
algebra. First torsion is rational over $L$, so the image is in
its principal-unit subgroup. This contradicts
Corollary~\ref{cor:boundary-nonsplit-cartan-image}.
\end{proof}

The split-intertwiner case is outside this deduction. The general
non-Cartan boundary configuration and its character index remain
open.

\section{The least index at minimal cyclotomic ramification}\label{sec:minimal-index}
This section proves Theorem~\ref{thm:minimal-index} and the exact characterization in Corollary~\ref{cor:exact-index-two}. The upper bound compares the graph's pointwise valuation constraint with the average valuation of primitive torsion. The lower bound is a finite twist of a Lubin--Tate group; its residual representations are identified by the Frobenius of a quadratic finite field.

\begin{proof}[The upper bound in Theorem~\ref{thm:minimal-index}]
Suppose $G[p^2]\simeq H[p^2]$ and $e=p-1$. Pass to the completed maximal unramified base. A determinant in the filtration class is $\chi_p^d\epsilon$, where $d\in\{0,1,2\}$ and $\epsilon\in\mu_{p-1}$. The wild cyclotomic image modulo $p^2$ has order $p$: otherwise $K$ would contain $\mu_{p^2}$, contrary to its absolute index. Equality of the determinant truncations forces the two dimensions $d$ to agree, because a nonzero difference of absolute value at most two cannot be divisible by $p\ge5$. Reduction is injective on $\mu_{p-1}$, so the finite determinant factors also agree.

For dimension zero, the special character is the Brauer lift of the residual semisimplification; duality treats dimension two. In dimension one, a potentially ordinary group has a unique connected--\'etale filtration, which descends. Thus each group either is whole potentially supersingular or is filtered by two potentially good height-one groups. Two filtered groups are compared residually, since $\chi_p\equiv1\pmod p$. A mixed comparison would require $e>p$ by Proposition~\ref{prop:mixed-obstruction} with $n=2$ and $v_p(e)=0$, which is impossible.

If both groups are potentially supersingular and their special characters differ, Proposition~\ref{fc-thm:valuation} gives
\[
 (p-1)(p^2-1)\ge(p-1)p^2.
\]
This contradiction proves recovery from $p^2$-torsion throughout the filtration class. The strict inequality comes from the graph bound on every point, not from a choice of Newton slope.
\end{proof}

\subsection{A pair with identical first-level torsion}
Put $k=\overline{\F}_p$, $R_0=W(k)$, $K_0=\operatorname{Frac}R_0$ and
\[
 \mathcal O=W(\F_{p^2}),\qquad q=p^2,
 \qquad K=K_0(\zeta_p).
\]
The cyclotomic polynomial $\Phi_p(1+X)$ is Eisenstein over $R_0$, so $K$ is complete, has residue field $k$, contains $\mu_p$ and has absolute index $p-1$.

Apply \citet[Lemma~1 and Theorem~1, pp.~381--382]{LT1965} over the unramified quadratic local field $\operatorname{Frac}\mathcal O$, with uniformizer $p$ and the series
\[
                         [p](T)=pT+T^q.
\]
This gives a one-dimensional formal $\mathcal O$-module, which we then base change to $R_0$. Its kernels of $[p^n]$ are finite flat of rank $q^n$, and their transition sequences define a height-two Barsotti--Tate group $G_0$. Its reduction is supersingular because $[p](T)$ reduces to $T^{p^2}$. The theorem is applied over a field with finite residue field before this base change.

The Tate module is free of rank one over $\mathcal O$ \citep[Theorem~2(b)--(d), pp.~383--384]{LT1965}. One can also check this directly: a primitive $p^n$-point has annihilator $p^n\mathcal O$, so the $q^n$ elements of its cyclic $\mathcal O$-module exhaust the $p^n$-torsion. Compatible primitive points give the assertion on inverse limits. All $\mathcal O$-endomorphisms are defined over $K_0$, so Galois acts through a character
\[
                       \psi:G_{K_0}\longrightarrow\mathcal O^*.
\]

Choose $t$ with $t^{q-1}=-p$. The Eisenstein polynomial $T^{q-1}+p$ gives a cyclic tame extension $E=K_0(t)$ of degree $q-1$, since $K_0$ contains $\mu_{q-1}$. Its nonzero roots are exactly the first-level torsion coordinates. We verify the subfield $K\subset E$ explicitly. If $\lambda=\zeta_p-1$, the cyclotomic equation gives
\[
 \lambda^{p-1}=-pu,\qquad
 u=1+\sum_{i=1}^{p-2}\frac{\binom p{i+1}}p\lambda^i
                      \in1+\lambda\mathcal O_K.
\]
Hensel's lemma supplies $v\in1+\lambda\mathcal O_K$ with $v^{p-1}=u$. Hence $(\lambda/v)^{p-1}=-p$. This root generates $K$ over $K_0$; all roots differ by elements of $\mu_{p-1}\subset K_0$. Since $t^{p+1}$ is another such root,
\[
                    K=K_0(t^{p+1}),\qquad [E:K]=p+1.
\]
For $a\in\mu_{q-1}$, the series $aT$ commutes with $pT+T^q$. Uniqueness in the Lubin--Tate construction gives $[a](T)=aT$. In the $\mathcal O/p$-basis $t$ of $G_0[p]$, the character $\overline\psi$ is therefore the reduction of $g(t)/t$. It follows that
\[
                 \overline\psi(G_K)=\mu_{p+1}\subset\F_{p^2}^*.
\]

Let $a(g)$ be the Teichm\"uller lift of $\overline\psi(g)$, and define
\[
                   \eta(g)=a(g)^{p-1}.
\]
The character $\eta$ has order $(p+1)/2\ge3$ and satisfies $\eta^p=\eta^{-1}$. Let $K'/K$ be the cyclic tame extension killing $\eta$. Twist the generic fibre of $G_0$ by the automorphisms $[\eta]$, obtaining $G_1$ with Tate character $\eta\psi$.

Here finite-level descent makes the construction explicit. For $M_n=\mathcal O/p^n\mathcal O$ with its $G_K$-action through $\eta\psi$, choose a finite Galois extension $D_n/K$ through which that action factors. The algebra
\[
 B_n=\operatorname{Map}(M_n,D_n)^{\operatorname{Gal}(D_n/K)},
 \qquad (gf)(x)=g\bigl(f(g^{-1}x)\bigr)
\]
is a product of fixed fields of point stabilizers and is finite \'etale over $K$. Addition, zero and inverse on $M_n$ give its Hopf operations. The equivariant maps $M_n\to M_{n+1}$, $x\mapsto px$, and $M_{n+1}\to M_n$, $x\mapsto x\bmod p^n$, give the closed embeddings and multiplication maps. Their kernel sequences are exact after separable base change and hence over $K$. Thus $\operatorname{Spec}B_n$ form a height-two Barsotti--Tate group. Over $K'$, the twist is trivial at every level, so $G_1$ has the same good supersingular model as $G_0$ after that extension.

Modulo $p$ one has $\overline\eta\,\overline\psi=\overline\psi^{p}$. The invertible $\F_p$-linear map
\[
                  \F_{p^2}\longrightarrow\F_{p^2},
                         \qquad x\longmapsto x^p
\]
intertwines the two Galois actions. It therefore gives $G_0[p]\simeq G_1[p]$ as finite group schemes over $K$.

The special character of $G_0$ is $2$, because it has good reduction over $K$. On the special isocrystal of $G_1$, the two embeddings of $\mathcal O$ give one-dimensional eigenspaces: Frobenius interchanges them, and their dimensions sum to two. The finite descent trace is consequently $\eta+\eta^p=\eta+\eta^{-1}$. A nontrivial value $a$ cannot satisfy $a+a^{-1}=2$, since that would imply $(a-1)^2=0$ in characteristic zero. The special characters differ. This proves the lower bound in Theorem~\ref{thm:minimal-index}, even within the class of whole potentially supersingular groups.

\begin{proof}[Proof of Corollary~\ref{cor:exact-index-two}]
For $p<e<p^3$, Theorem~\ref{thm:prime-index-upper} gives recovery from $p^3$-torsion, while Corollary~\ref{cor:sharp-mixed-prime} with $n=2$ gives unequal-character pairs with isomorphic $p^2$-torsion. Thus the least index is two.

Among positive multiples of $p-1$ with $e\le p$, the only possibility is $e=p-1$, where Theorem~\ref{thm:minimal-index} gives index one. Equality $e=p^3$ is excluded by $p\nmid e$. For $e>p^3$, the same mixed construction with $n=3$ gives failure of recovery from $p^3$-torsion, hence index at least three. These cases exhaust the complement of the displayed interval. The index is finite for every permitted $e$, since one can choose $n\ge3$ with $e<p^n$ and apply Theorem~\ref{thm:prime-index-upper}.
\end{proof}

\section{Central rigidity from relative point fields}\label{sec:central-ramification}
This section proves Theorems~\ref{thm:central-p2-rigidity} and \ref{thm:ss-prime-range}. The first argument compares a cyclic extension of point fields with the residual ramification of the quotients by a canonical kernel and its transported image. The second combines that criterion with the pointwise graph valuation bound, losing at most one torsion level.

Pass to the completed maximal unramified base and write $e=v_K(p)$. Let $L/K$ be the common tame good-reduction field of degree $m=p^2-1$, with algebraically closed residue field $k$ and $E=v_L(p)=me$. Lemma~\ref{lem:tame-good} supplies this field. For a two-slope model write $h=(p-1)z$, so
\begin{equation}\label{cr-eq:AB}
 A=(p+1)e-z,\qquad B=z/p,\qquad 0<z<pe.
\end{equation}
The deformation weight makes $z$ integral. Lemmas~\ref{rc-lem:newton} and \ref{fc-lem:rational} give the division valuations and the rationality of the canonical points. The canonical line and its quotient character are trivial over $L$, since the determinant over the good field is cyclotomic and $\mu_p\subset L$.

We use the following exact quotient data, proved in Proposition~\ref{rc-prop:quotients} and the marking calculation preceding Proposition~\ref{fc-prop:central}. A canonical quotient has normalized Hasse value $pz$ for $z<e$, value $(p+1)e-z$ for $z>e$, and one slope for $z=e$. Its dual kernel is noncanonical in the first case and canonical in the second. A quotient by a noncanonical stable line has normalized Hasse value $z/p$. A one-slope model has all nonzero $p$-point values equal to $e$, all primitive $p^2$-point values equal to $e/p^2$, and every $p$-quotient has normalized Hasse value $e$.

If a $p^2$-torsion isomorphism carries $C$ to $C'$, it induces
\begin{equation}\label{cr-eq:quotient}
 (G/C)[p]\simeq(H/C')[p],\qquad
 (G/C)[p]=\{x\in G[p^2]:px\in C\}/C.
\end{equation}
The quotients exist over $K$ for $K$-stable kernels and preserve special inertia traces; see Lemma~\ref{lem:actual-quotients}. These facts require no extension of the given torsion isomorphism to good integral models.

\subsection{The cyclic relative extension and its break}
The elementary valuation estimate below fixes the lower-numbering convention; compare \citet[Chapter~IV, \S1]{Serre1979}.

\begin{lemma}\label{cr-lem:valuation}
Let $F/F_0$ be a totally ramified cyclic extension of degree $p$ of
complete discretely valued fields of characteristic zero, with common
perfect residue field of characteristic $p$. Let $w$ be the integer
valuation of $F$, let $\sigma$ be nonidentity, and let $b$ be its lower
ramification break. For every $x$ in the maximal ideal of $\mathcal O_F$,
\[
 w(\sigma(x)-x)\geq w(x)+b.
\]
Equality holds when $p\nmid w(x)$.
\end{lemma}

\begin{proof}
Choose a uniformizer $\varpi$ of $F$. The residue of
$\sigma(\varpi)/\varpi$ has order dividing $p$ in the multiplicative
group of a field of characteristic $p$, so it is~$1$. Thus $b\geq1$, and
\[
 \sigma(\varpi)=\varpi+u\varpi^{b+1}
                  +\text{terms of larger valuation},\qquad u\in\mathcal O_F^*.
\]
Use Teichm\"uller representatives from the coefficient ring of $F_0$ for
the common residue field; they are fixed by $\sigma$. Every element of
$\mathcal O_F$ has a convergent expansion in powers of $\varpi$ with
these coefficients. The binomial expansion gives
\[
 w\bigl(\sigma(\varpi^j)-\varpi^j\bigr)\geq j+b\qquad(j\geq1).
\]
When $p\nmid j$, the first binomial term has exactly this valuation, and
all later terms have larger valuation because $b\geq1$. Apply these
estimates to the expansion of $x$. If its first exponent $w(x)$ is prime
to $p$, its change has strictly smaller valuation than the changes of all
subsequent terms. This proves the equality assertion as well.
\end{proof}

\begin{proof}[Proof of Theorem~\ref{thm:central-p2-rigidity}]
Let $C$ be the canonical $p$-subgroup of $G$, choose
$0\ne P\in C(L)$, and choose $Q$ with $[p]Q=P$. The point $Q$ has exact
order $p^2$. If $T$ is the good formal coordinate, centrality gives
\begin{equation}\label{cr-eq:qval}
 v_L(T(Q))=\frac{A}{p^2}.
\end{equation}
The denominator is $p^2$, because $p\nmid A$. Thus the ramification
degree of $L(Q)/L$ is at least $p^2$. Its field degree is at most $p^2$,
since the orbit of $Q$ lies in the set of $p^2$ preimages of the fixed
point $P$. Therefore
\begin{equation}\label{cr-eq:field1}
 F=L(Q),\qquad [F:L]=p^2.
\end{equation}
The extension is totally ramified, since $k$ is algebraically closed.

Form $q:G\to G_1=G/C$. Its normalized coordinate is, up to a unit, the
product of $T-_G R$ over $R\in C$. At $Q$, all $p$ factors have valuation
$A/p^2$, since the nonzero kernel points have the larger valuation $A$.
Thus
\[
 v_L\bigl(T_1(q(Q))\bigr)=\frac Ap.
\]
The point $q(Q)$ is nonzero $p$-torsion. The normalized Hasse value of
$G_1$ is $A$, and this is a noncanonical point of $G_1[p]$. Its valuation
forces field degree at least $p$. Its orbit has at most $p$ elements,
because the canonical line of $G_1[p]$ and its quotient character are
trivial over $L$. Hence
\begin{equation}\label{cr-eq:field0}
 F_0=L(q(Q)),\qquad [F_0:L]=p.
\end{equation}

The extension $F/F_0$ has degree $p$ and is cyclic Galois. Indeed the
$p$ points $Q+R$, $R\in C(L)$, are the full fiber of $q$ above $q(Q)$ and
all belong to $F$. The formal addition series converges in the complete
field $F$. The minimal polynomial of $T(Q)$ over $F_0$ has degree $p$,
and all its roots, precisely the coordinates of this fiber, lie in $F$.
Thus a nonidentity $\sigma\in\Gal(F/F_0)$ sends $Q$ to $Q+P_0$ for some
$0\ne P_0\in C(L)$.

Normalize $w=p^2v_L$ on $F$. A formal coordinate difference under
translation by $P_0$ is $T(P_0)$ times a unit. Consequently
\[
 w(T(Q))=A,\qquad
 w\bigl(\sigma(T(Q))-T(Q)\bigr)=p^2A.
\]
Lemma~\ref{cr-lem:valuation}, with its equality clause, determines the
relative lower break:
\begin{equation}\label{cr-eq:break}
 b=(p^2-1)A.
\end{equation}

Now let $f:G[p^2]\to H[p^2]$ be the given isomorphism. Put
$C'=f(C)$, $R_Q=f(Q)$, and $q':H\to H_1=H/C'$. The line $C'$ is
$K$-stable and all its points are $L$-rational. The point stabilizers of
$Q$ and $R_Q$ in the absolute Galois group of $L$ are identical. The
isomorphism~\eqref{cr-eq:quotient} also identifies the stabilizers of
$q(Q)$ and $q'(R_Q)$. Thus, as subfields of the fixed separable closure,
\begin{equation}\label{cr-eq:transport}
 L(R_Q)=F,\qquad L(q'(R_Q))=F_0.
\end{equation}
For the same automorphism $\sigma$,
\begin{equation}\label{cr-eq:translation}
 \sigma(R_Q)=R_Q+f(P_0).
\end{equation}
These assertions use only the finite torsion isomorphism; they do not
require its extension to integral good models.

The group $G_1[p]$ has nontrivial wild action over $L$. Therefore $H$ cannot
be two-slope with $C'$ canonical in the outer range: its quotient would
have a stable noncanonical dual line and hence trivial residual action.
The canonical boundary $z_H=e$ is excluded for the same reason, since its
quotient is one-slope.

In every remaining case, $H_1$ is two-slope. Denote its normalized Hasse
value by $t$. The exact values are as follows:
\begin{enumerate}
\item If $C'$ is canonical and $H$ is central, then
      $t=(p+1)e-z_H$. The point $[p]R_Q$ is canonical, central division
      gives $v_L(S(R_Q))=t/p^2$, and $f(P_0)$ has valuation $t$.
\item If $C'$ is noncanonical, then $t=z_H/p$.
      The point $[p]R_Q$ is noncanonical, so again
      $v_L(S(R_Q))=t/p^2$, while $f(P_0)$ has valuation $t$.
      In particular $t$ is integral, since the nonzero points of $C'$
      are $L$-rational.
\item If $H$ is one-slope, then $t=e$. All nonzero $p$-points have
      valuation $e$ and every primitive $p^2$-point has valuation
      $e/p^2$, including the equality case of the Newton boundary.
\end{enumerate}
In all three cases, \eqref{cr-eq:translation} therefore gives
\[
 w(S(R_Q))=t,\qquad
 w\bigl(\sigma(S(R_Q))-S(R_Q)\bigr)=p^2t.
\]
Applying Lemma~\ref{cr-lem:valuation} to $S(R_Q)$ and using
\eqref{cr-eq:break}, we obtain
\begin{equation}\label{cr-eq:opposite}
 (p^2-1)A\leq(p^2-1)t,\qquad A\leq t.
\end{equation}

The quotient $p$-torsion representations are isomorphic, and the
normalized Hasse value $A$ is prime to $p$. If $t$ is also prime to $p$,
Theorem~\ref{rc-thm:primehasse} gives $t=A$ and equal special
trace characters of $G_1,H_1$. Their $K$-isogenous sources then have equal
special traces. If $p\mid t$, Corollary~\ref{rc-cor:strictordering} instead
gives $t<A$, contradicting~\eqref{cr-eq:opposite}. This proves the
theorem.
\end{proof}

\begin{remark}
The two ramification comparisons retain different information. The
residual representation of the quotients makes a divisible quotient
Hasse value strictly smaller than the prime-to-$p$ one. The cyclic
extension of point fields above those quotients makes it at least as
large. The latter inequality is obtained from the original
$p^2$-torsion, even though the quotient comparison itself has only
level~$p$.
\end{remark}

\subsection{The supersingular range and the full-class upper bound}
We first make explicit the pointwise estimate used in the proof. Proposition~\ref{fc-thm:valuation} gives, for an unequal-character comparison at level $p^n$,
\begin{equation}\label{cr-eq:graph}
 v_L(T(P))\ge\frac{p-1}{p^n-1}
 \quad\text{for every nonzero $P\in G[p^n]$, and likewise for $H$.}
\end{equation}
Indeed its proof obtains $x^{p^n},y^{p^n}\in\pi^{p-1}I$ in the graph algebra. Evaluating at $(P,f(P))$ and taking the smaller coordinate value $r$ gives $p^nr\ge p-1+r$, proving \eqref{cr-eq:graph}.

For a two-slope model, the minimum nonzero $p^n$-torsion value is $z/p^{2n-1}$. At level one it is $B=z/p$. The inequality $v_L([p]x)>v_L(x)$ for a nonzero formal point with nonzero $p$-multiple shows that every nonzero torsion target has valuation at most $A<E$, so the precise range of Lemma~\ref{rc-lem:newton} applies throughout the inverse chain. At every further level the minimum targets have all $p^2$ divisions on the low branch, which divides that value by $p^2$. A target of larger value below $p^2B$ has divisions of strictly larger value; if its value is at least $p^2B$, all its divisions have value at least $B$. Division points of zero are first-level points and cannot change the minimum. Induction proves the assertion. It follows that
\begin{equation}\label{cr-eq:zbound}
 z\ge\frac{(p-1)p^{2n-1}}{p^n-1}>(p-1)p^{n-1}.
\end{equation}

\begin{proof}[Proof of Theorem~\ref{thm:ss-prime-range}]
Suppose the two special traces differ. Lemma~\ref{lem:one-slope-recovery}, which has no upper bound on $e$, excludes every case involving a one-slope model. Thus both models have two slopes. If both normalized Hasse values are prime to $p$, Theorem~\ref{rc-thm:primehasse} gives equal traces already at level $p$. We may therefore choose one model $G$ whose normalized Hasse value $z$ is divisible by $p$.

If $z>e$, then $G$ is central and $A=(p+1)e-z$ is prime to $p$. Apply Theorem~\ref{thm:central-p2-rigidity} to the restricted $p^2$-isomorphism, a contradiction. Equality $z=e$ is impossible because $p\nmid e$.

Suppose $z<e$. Quotient $G$ by its canonical line $C$, and quotient $H$ by the transported $K$-stable line $f(C)$. These paired quotients have isomorphic $p^{n-1}$-torsion by Lemma~\ref{lem:actual-quotients}. The first quotient has normalized Hasse value $z_1=pz$. Equations~\eqref{eq:ss-prime-range} and \eqref{cr-eq:zbound} give
\[
 e<\frac{(p-1)p^{2n}}{p^n-1}\le pz=z_1<pe.
\]
It is therefore central, $p\mid z_1$, and $p\nmid(p+1)e-z_1$. Since $n-1\ge2$, Theorem~\ref{thm:central-p2-rigidity} applies to the paired quotients and gives equal traces. Their isogenies are defined over $K$, so the original traces agree as well, the final contradiction.
\end{proof}

\begin{proof}[A second proof of Theorem~\ref{thm:prime-index-upper}]
Suppose $n\ge3$, $p\nmid e$, and $e<p^n$. This numerical range is contained in \eqref{eq:ss-prime-range}, so Theorem~\ref{thm:ss-prime-range} treats the comparison of two whole supersingular groups. For the full height-two filtration class, pass to the completed maximal unramified base. A determinant has the form $\chi_p^d\epsilon$, with $d\in\{0,1,2\}$ and finite $\epsilon\in\mu_{p-1}$. Since $p\nmid e$, the wild cyclotomic order modulo $p^n$ is $p^{n-1}$. Equality of the determinant truncations therefore forces equal $d$ and equal finite factors.

For dimension zero the special character is the Brauer lift of the semisimplified residual character, and Cartier duality gives dimension two. In dimension one, the connected--\'etale filtration of a potentially ordinary piece descends by uniqueness, so the possibilities are a whole supersingular group or two height-one graded pieces. For two filtered groups, $\chi_p\equiv1\pmod p$ makes the special character recoverable from the residual semisimplification. These are the structural reductions detailed in the proof of Theorem~\ref{thm:second-wild-exact}. A mixed filtered/supersingular comparison would require $e>p^{2n-3}$ by Proposition~\ref{prop:mixed-obstruction}; this contradicts $e<p^n\le p^{2n-3}$. The supersingular theorem handles the remaining case, proving $\nuMin(p,e,2)\le n-1$.

For $p<e<p^3$, Theorem~\ref{thm:general-arc} gives unequal special traces with isomorphic $p^2$-torsion, while the upper bound applies to $p^3$-torsion. For $p^3<e<p^4$, it gives failure at level $p^3$, while the upper bound applies at level $p^4$. These are exactly the two asserted index values.
\end{proof}

\begin{proof}[Proof of Corollary~\ref{cor:prime-recovery-intro}]
Write $S_n=(p-1)p^{2n}/(p^n-1)$. For $n\ge4$,
\[
 S_n<p^{n+1}\le p^{2n-3}.
\]
The first inequality is equivalent to $p^n>p$, and the second to $n\ge4$. Under $e<S_n$, Theorem~\ref{thm:ss-prime-range} excludes an unequal-character supersingular comparison. Proposition~\ref{prop:mixed-obstruction} excludes a mixed comparison because it would require $e>p^{2n-3}$. The determinant, dimension and filtered reductions in the preceding proof impose no further upper bound on $e$. Thus the same argument proves recovery throughout the full filtration class. For $n=3$ the preceding theorem gives recovery under $e<p^3$.

Since $e$ is a multiple of $p-1$, the identity
\[
 \frac{S_n}{p-1}=p^n+1+\frac{1}{p^n-1}
\]
shows that $e<S_n$ is equivalent to $e\le(p-1)(p^n+1)$. The hypothesis $p\nmid e$ remains in force. The index-two interval follows from the preceding proof. The exact mixed construction of Corollary~\ref{cor:sharp-mixed-prime}, at level $p^3$, gives unequal traces whenever $e>p^3$; the new full-class bound at level $p^4$ supplies the matching upper bound up to $(p-1)(p^4+1)$. At level $p^4$, the mixed construction gives unequal traces whenever $e>p^5$; recovery at level $p^5$ supplies the matching upper bound up to $(p-1)(p^5+1)$. These are the asserted index-three and index-four intervals. The lower constructions allow the field of index $e$ to vary, whereas both upper bounds hold over every field in the specified base class.
\end{proof}

For example, at $p=5$ the index is three for $125<e\le2504$ and four for $3125<e\le12504$, subject in both cases to $4\mid e$ and $5\nmid e$. At level $p^3$, the mixed obstruction occurs inside the larger supersingular range, so the distinction between the two classes is necessary there.

\section{Canonical quotients and the complete prime-to-$p$ index}\label{sec:canonical-iteration}
This section proves the complete index table by iterating central rigidity while retaining the actual torsion comparison at every quotient.

Throughout, $p\ge5$, $K$ is complete discretely valued of characteristic zero with perfect residue field, $\mu_p\subset K$, and $p\nmid e=e_K$. For potentially supersingular groups use the common tame good field $L/K$ of degree $p^2-1$ and the normalized Hasse parameter of Section~\ref{sec:central-ramification}. The inputs are the pointwise graph bound, the one-slope and prime-Hasse recovery statements, the actual quotient formulas, and central $p^2$ rigidity, each proved in the preceding sections.

\begin{lemma}[Iteration with transported kernels]\label{ci-lem:iteration}
Let $f:G[p^n]\simeq H[p^n]$ be an isomorphism of potentially supersingular height-two groups. Suppose $G$ has a two-slope parameter $z$ divisible by $p$. If an integer $c$ satisfies
\[
 0\le c\le n-2,\qquad e<p^c z<pe,
 \qquad p^jz<e\quad(0\le j<c),
\]
then $\theta_G=\theta_H$.
\end{lemma}

\begin{proof}
Put $G_0=G$, $H_0=H$ and $f_0=f$. For $j<c$, let $C_j$ be the canonical $p$-line of $G_j$. It is $K$-stable: the semilinear good-model action preserves formal valuations, and the high line consists of zero and all points of the larger first-level valuation. Define
\[
 D_j=f_j(C_j),\qquad
 G_{j+1}=G_j/C_j,\qquad H_{j+1}=H_j/D_j.
\]
The subgroup $D_j$ is $K$-stable because $f_j$ is Galois-equivariant. Lemma~\ref{lem:actual-quotients} gives the actual quotients over $K$ and preserves their special characters. No condition that $D_j$ be canonical in $H_j$ is required.

To check the retained truncation precisely, put $q=n-j$. As finite Galois modules,
\begin{equation}\label{ci-eq:transport}
 (G_j/C_j)[p^{q-1}]
 =\{x\in G_j[p^q]:p^{q-1}x\in C_j\}/C_j.
\end{equation}
Indeed a quotient point killed by $p^{q-1}$ has a lift $x$ with $p^{q-1}x\in C_j$, hence $p^qx=0$ since $pC_j=0$; the converse follows by multiplication. The map $f_j$ carries this exact subquotient to its partner for $H_j,D_j$. It therefore defines an isomorphism $f_{j+1}$ at level $p^{n-j-1}$. Each step loses exactly one level and requires no integral extension of the original torsion isomorphism.

For every step taken, $p^jz<e$. The canonical quotient formula of Proposition~\ref{rc-prop:quotients} gives the parameter $z_j=p^jz$ of $G_j$ by induction. The partner remains potentially supersingular because it is $K$-isogenous to $H$. At the final step, $n-c\ge2$ and
\[
 e<z_c<pe,\qquad p\mid z_c,
 \qquad p\nmid (p+1)e-z_c.
\]
The last assertion uses $p\nmid e$. Restrict $f_c$ to $p^2$ and apply Theorem~\ref{thm:central-p2-rigidity}. The terminal traces agree, and all intervening $K$-isogenies preserve those traces. Hence the original traces agree.
\end{proof}

\begin{proof}[Proof of Theorem~\ref{thm:ss-iterated-range}]
Suppose the special characters differ. Since $n\ge3$, Lemma~\ref{lem:one-slope-recovery} excludes every case involving a one-slope model. Both models are therefore two-slope. If both normalized Hasse values were prime to $p$, Theorem~\ref{rc-thm:primehasse} would give equal traces already from $p$-torsion. Choose the model $G$ so that its parameter $z$ is divisible by $p$.

The smallest nonzero $p^n$-torsion value of $G$ is $z/p^{2n-1}$. To recall the range of the division calculation, its first-level minimum is $B=z/p$. A nonzero torsion target of value $q\le A=(p+1)e-z$ has all $p$-division values $q/p^2$ when $q\le p^2B$; otherwise the division values are $(q-p(p-1)B)/p$ and $B$. Multiplication by $p$ strictly increases positive nonzero valuations, so every target stays within $q\le A$. The smallest targets always divide by $p^2$; larger targets below $p^2B$ yield larger divisions, and targets above that threshold give divisions at least $B$. Division points of zero are first-level points and do not lower the minimum. Induction proves the asserted minimum, as also established before~\eqref{cr-eq:zbound}.

The pointwise graph bound, Proposition~\ref{fc-thm:valuation}, now gives
\[
 z\ge Z_n:=\frac{(p-1)p^{2n-1}}{p^n-1}.
\]
The assumed bound on $e$ is $e<p^{n-2}Z_n$, so $p^{n-2}z>e$. Let $c$ be the least nonnegative integer with $p^cz>e$. It exists and $c\le n-2$. Equality $p^jz=e$ is impossible because $p\mid z$ and $p\nmid e$. If $c=0$, the two-slope range gives $z<pe$. If $c>0$, minimality gives $p^{c-1}z<e$, hence $p^cz<pe$. Lemma~\ref{ci-lem:iteration} applies and contradicts the unequal characters.

For the integral form of the bound, use
\[
 \frac{p^{3n-3}}{p^n-1}
 =p^{2n-3}+p^{n-3}+\frac{p^{n-3}}{p^n-1}.
\]
The final term is strictly between zero and one, and $e/(p-1)$ is an integer.
\end{proof}

\begin{proof}[Proof of Corollary~\ref{cor:complete-prime-recovery}]
The supersingular bound is larger than the asserted full-class threshold, since
\[
 \frac{(p-1)p^{3n-3}}{(p^n-1)p^{2n-3}}
 =\frac{(p-1)p^n}{p^n-1}>1.
\]
We give the structural reduction to make its lack of an additional upper bound on $e$ explicit. Determinants have the form $\chi_p^d\epsilon$, with $d\in\{0,1,2\}$ and $\epsilon$ of order dividing $p-1$. Because $p\nmid e$, the wild cyclotomic order modulo $p^n$ is $p^{n-1}$. Equality of determinant truncations forces $d=d'$, since $|d-d'|\le2<p$, and then equality of the finite factors.

In dimension zero, Brauer lifting of the semisimplified residual character recovers the special character; Cartier duality gives dimension two. In dimension one, the connected--\'etale filtration of a potentially ordinary piece descends by uniqueness. The possibilities are therefore a whole supersingular group or two height-one graded pieces. For two filtered groups, $\chi_p\equiv1\pmod p$ identifies their special traces with the Brauer lifts of their residual semisimplifications. A mixed comparison would require $e>p^{2n-3}$ by Proposition~\ref{prop:mixed-obstruction}, contrary to the hypothesis. The remaining supersingular comparison is covered by Theorem~\ref{thm:ss-iterated-range}. This proves the uniform full-class bound.
\end{proof}

\begin{proof}[Proof of Theorem~\ref{thm:complete-prime-index}]
At $e=p-1$, Theorem~\ref{thm:minimal-index} gives the exact index one and an actual pair witnessing failure at the first level. Every other positive multiple of $p-1$ is at least $2(p-1)>p$. Since $p\nmid e$, there is a unique $r\ge2$ such that
\[
 p^{2r-3}<e<p^{2r-1}.
\]
Corollary~\ref{cor:complete-prime-recovery} with $n=r+1$ gives $\nuMin(p,e,2)\le r$. Corollary~\ref{cor:sharp-mixed-prime} at level $p^r$ gives an actual unequal-character comparison over some field of index $e$, because $e>p^{2r-3}$. Thus $\nuMin(p,e,2)\ge r$. Both inequalities are attained at every permitted index, proving the exact table. Solving its strict inequalities yields~\eqref{eq:complete-prime-index}; no equality endpoint occurs because $p\nmid e$.
\end{proof}

For example, at $p=5$ the index is three for every allowed $125<e<3125$, and four for every allowed $3125<e<78125$. Here ``allowed'' means $4\mid e$, $5\nmid e$, with $\mu_5$ in the base. The lower constructions vary the field and are actual Barsotti--Tate comparisons; the upper assertions are uniform over every field in the specified class.

\section{Primitive functions and the excluded two-slope configurations}\label{sec:primitive-grading}
The primitive-function method gives a separate numerical reduction of a hypothetical unequal-character comparison. It describes the Frobenius weights and identifies canonical residual lines without assuming such an identification. The central criterion excludes each of the three configurations at the end of this section. The shorter proof in Section~\ref{sec:central-ramification} does not require this reduction.

\begin{theorem}\label{thm:three-family-reduction}
Let $p\ge5$, $n\ge3$, $p\nmid e$, $(p-1)\mid e$, and $e<p^n$. Let two potentially supersingular height-two groups over the same complete field $K$ have isomorphic $p^n$-torsion and unequal special inertia traces. After completed unramified base change, let $L/K$ be their common tame good field of degree $p^2-1$. Suppose both good models have two Newton slopes. Write their Hasse valuations as $h_i=(p-1)z_i$, normalized by $v_L(L^\times)=\Z$, and let $\alpha_i$ be their tangent descent scalars for a fixed generator of $\operatorname{Gal}(L/K)$. After interchanging the models if necessary, one of the following holds:
\begin{align}
 z_1&=(p-1)p^n+b,&z_2&=p^{n-1}b,
 &p+1\le b\le p^2-p-1,\quad p\nmid b;
 \label{eq:pg-family-one}\\
 z_1&=b+(p-1)p^{n-1}\ell,&z_2&=p^nb,
 &1\le b\le p-1,\quad b+1\le\ell\le p+1.
 \label{eq:pg-family-two}
\end{align}
The third possibility is
\begin{equation}\label{eq:pg-family-sign}
 z_1=z_2=\frac{p-1}{2}p^n,\qquad \alpha_2=-\alpha_1,
\end{equation}
with the further necessary congruence
\begin{equation}\label{eq:pg-family-congruence}
 (p-1)p^{n-2}\mid(p+1)e-\frac{p-1}{2}(p^n+1).
\end{equation}
In this third case the residual representations split into distinct characters, and the given isomorphism interchanges the canonical residual lines with the opposite character lines. Every family is still subject to the strict two-slope inequalities $0<z_i<pe$.
\end{theorem}

The first step is the classical description of a finite additive group by primitive functions \citep[\S\S2.1--2.2]{Koch2015}. Rescaling the torsion graph gives exactly $2n$ primitive dimensions. Counting two Frobenius chains gives a congruence, and a one-dimensional zero-weight space identifies the canonical lines. The resulting families are necessary intermediate configurations and are proved impossible in Corollary~\ref{pg-cor:exclusion}.

\subsection{The rescaled graph and its primitive functions}
We first identify the finite-dimensional object on which the two Frobenius chains can be compared.

Assume $p\ge5$, $n\ge3$, $p\nmid e$, $(p-1)\mid e$ and $e<p^n$. Suppose two potentially supersingular height-two groups have isomorphic $p^n$-torsion and unequal special inertia traces. We consider the two-slope case; Lemma~\ref{lem:one-slope-recovery} treats the other polygons. After completed unramified base change, take the common tame good field $L/K$ of degree $m=p^2-1$, with ring $R$, algebraically closed residue field $k$, and $E=v_L(p)=me$. Choose $\tau,\pi,\zeta$ as in \eqref{rc-eq:hasseweight}. Write
\[
 h_i=(p-1)z_i,\qquad s_i=v_p(z_i),\qquad 0<z_i<pe,
 \qquad i=1,2.
\]
The integrality of $z_i$ follows from the deformation weight. Let $\alpha_i$ be the tangent scalar of $\tau$. The two Frobenius orbits of these scalars are disjoint, and determinant comparison gives
\begin{equation}\label{pg-eq:determinants}
 \alpha_1^{p+1}=\alpha_2^{p+1}.
\end{equation}
We may use exact tame eigenparameters: applying $m^{-1}\sum_{j=0}^{m-1}\alpha_i^{-j}\tau^j$ to a parameter gives one with a unit linear coefficient and exact eigenvalue $\alpha_i$. Such changes preserve all positive point valuations.

\begin{lemma}[Primitive presentation]\label{pg-lem:primitives}
Let $B$ be a finite-dimensional commutative Hopf algebra over $k$, generated by primitive elements $X,Y$. The space
\[
 V=\operatorname{span}_k\{X,X^p,X^{p^2},\ldots,Y,Y^p,Y^{p^2},\ldots\}
\]
is its full primitive space, and $\dim_k B=p^{\dim_k V}$.
\end{lemma}
\begin{proof}
This is the field case of the classical additive-group/primitive-module correspondence; see \citet[\S\S2.1--2.2, pp.~3--4]{Koch2015}. We give the algebra proof needed here. Choose a basis $v_1,\ldots,v_d$ of $V$, write $v_i^p=\sum_j a_{ij}v_j$, and form
\[
 C=k[t_1,\ldots,t_d]/(t_i^p-\textstyle\sum_j a_{ij}t_j).
\]
The monomials with each exponent less than $p$ form a basis. Indeed reduction by the displayed relations decreases total degree; the overlap between the $i$th and $j$th relations reduces to the difference of the two products of their linear right sides, which is zero. Thus reductions have unique normal forms and $\dim C=p^d$.

Each $t_i$ is primitive, and the primitive relations give an additive Hopf structure on $C$. Sending $t_i$ to $v_i$ gives a surjective Hopf map $C\to B$. If its kernel were nonzero, choose a nonzero kernel element $P$ of least normal-form degree $q$. Augmentation and linear independence exclude $q=0,1$. Every term of
\[
 \Delta(P)-P\otimes1-1\otimes P
\]
has degree less than $q$ in each tensor factor. No reduction is needed when expanding a normal monomial, since its individual exponents are less than $p$. The degree-less-than-$q$ subspace injects into $B$ by minimality, so its tensor square also injects. The displayed expression must therefore vanish already in $C\otimes C$.

A primitive normal form with individual exponents less than $p$ is linear. For any monomial of degree at least two, split off one full positive variable exponent if two variables occur; if only one occurs, split off exponent one. The resulting coproduct coefficient is nonzero modulo $p$. Different original monomials have different sums of exponent vectors and cannot cancel that term. Thus $P$ would be linear, a contradiction. The map $C\to B$ is an isomorphism, and the same argument identifies its primitive elements with the span of the $t_i$.
\end{proof}

\begin{lemma}[Minimum torsion values and leading characters]\label{pg-lem:scaling}
Put $D=p^{2n-1}$ and choose $r$ with $rD\equiv1\pmod m$. For model $i$, the minimum nonzero $p^n$-torsion coordinate value is $z_i/D$, attained by $(p-1)D$ points. The rescaled graph has an additive special fibre whose primitive space has dimension $2n$. It carries a $\Z/D\Z$ grading in which $X^{p^j},Y^{p^j}$ have respective weights $-z_1p^j,-z_2p^j$. The tame scalars
\begin{equation}\label{pg-eq:lambda}
 \lambda_i=\alpha_i\zeta^{-rz_i}
\end{equation}
belong to $\F_p^\times$, and no Frobenius iterate of either coordinate is zero. Moreover,
\begin{equation}\label{pg-eq:bounds}
 (p-1)p^{n-1}<z_i<pe<p^{n+1},\qquad s_i\le n.
\end{equation}
\end{lemma}
\begin{proof}
At level one the minimum is $b=z_i/p$, attained by $p^2-p$ noncanonical points. Suppose the level-$j$ minimum is $b/p^{2j-2}$. Each minimum target has all $p^2$ divisions on the low branch, of value $b/p^{2j}$. A larger target value $q<p^2b$ has division value $q/p^2$, strictly larger than that minimum. If $q\ge p^2b$, every division value is at least $b$, also larger. Division points of zero are first-level points and satisfy the same strict inequality. These statements follow from Lemma~\ref{rc-lem:newton}; its range is valid because every nonzero torsion value is at most the maximum first-level value. Induction proves the minimum, its multiplicity $(p^2-p)p^{2n-2}=(p-1)D$, and that the remaining $D$ points, including zero, have larger value.

Let $A$ be the finite free graph algebra. Make the Eisenstein extension $R'=R[\varpi]/(\varpi^D-\pi)$ and, in its generic algebra, put
\[
 X=x/\varpi^{z_1},\qquad Y=y/\varpi^{z_2},\qquad B'=R'[X,Y].
\]
All geometric values of $X,Y$ are integral. Their characteristic polynomials therefore have coefficients in the integrally closed DVR $R'$, so they are integral over $R'$. Consequently $B'$ is finite free, with the same generic algebra as $A$ and rank $p^{2n}$. Rescaling the formal group laws makes every nonlinear term in the first coordinate divisible by $\varpi^{z_1(d-1)}$, where $d\ge2$ is its total degree, and similarly for the second coordinate. The resulting series converge in the finite complete algebras and give integral Hopf operations. Thus $B_0=B'/\varpi B'$ is generated by primitive $X,Y$. Lemma~\ref{pg-lem:primitives} gives primitive dimension $2n$.

The diagonalizable group scheme $\mu_D$ acts by $\varpi\mapsto t\varpi$, fixing $x,y$. It preserves $B'$ and its special fibre and gives the stated grading. This uses the character grading of $\mu_D$, which is valid even though $\mu_D$ is nonreduced in characteristic $p$. Extend the tame generator by $\tau(\varpi)=\zeta^r\varpi$. Its scalars on the normalized coordinates are \eqref{pg-eq:lambda}.

To determine these scalars, perform the same scaling on one entire $p^n$-torsion algebra. Its normalized coordinate has $p^{2n}$ distinct generic values and satisfies a monic degree-$p^{2n}$ polynomial with integral coefficients. The resulting monogenic algebra has that rank and is the whole scaled algebra. Its special polynomial $P_i(U)$ is additive: $P_i(U+V)-P_i(U)-P_i(V)$ vanishes in the tensor square of the quotient, and has degree less than $p^{2n}$ in each variable, so linear independence of normal monomials makes it identically zero. An additive polynomial in characteristic $p$ has only $p$-power exponents, by the binomial-coefficient criterion. The minimum-value count makes its zero root have multiplicity exactly $D$. Hence
\[
 P_i(U)=U^{pD}+c_iU^D,\qquad c_i\ne0.
\]
Its distinct roots form an $\F_p$-line. Tame action preserves that line and scales it by $\lambda_i$, proving $\lambda_i\in\F_p^\times$. A minimum-value point gives a nonzero specialization of each coordinate in the graph, so none of its Frobenius powers vanishes.

The proof of Proposition~\ref{fc-thm:valuation} bounds every nonzero torsion coordinate below by $(p-1)/(p^n-1)$. Applying this to the minimum gives
\[
 z_i\ge\frac{(p-1)p^{2n-1}}{p^n-1}>(p-1)p^{n-1}.
\]
Together with $z_i<pe<p^{n+1}$ this proves \eqref{pg-eq:bounds}.
\end{proof}

\subsection{Counting the two Frobenius chains}
We now obtain divisibility restrictions from the dimension of the primitive space.

\begin{proposition}\label{pg-prop:chains}
If $\lambda_1\ne\lambda_2$, then $s_1=s_2=n$. In every case $\max(s_1,s_2)\ge n-1$. If $\lambda_1=\lambda_2$, order $s=s_1\le s_2=t$ and write $z_1=p^s a$, $z_2=p^t b$ with $p\nmid ab$. Then $a\ne b$ and
\begin{equation}\label{pg-eq:chain-congruence}
 (p-1)p^{\max(0,2n-1-s-t)}\mid a-b.
\end{equation}
\end{proposition}
\begin{proof}
For coordinate $i$, the weights of its successive Frobenius powers are nonzero and distinct until exponent $d_i=2n-1-s_i$, when the weight first becomes zero. These nonzero vectors represent $d_i+1=2n-s_i$ different weights. All powers have the same tame character $\lambda_i$, since it belongs to $\F_p^\times$. If the two characters differ, their chains are in different tame eigenspaces, so $2n\ge(2n-s_1)+(2n-s_2)$. The bounds $s_i\le n$ force both to equal $n$.

Suppose the characters agree and put $d=t-s$. A nonzero weight in the second chain at exponent $j$ can equal a first-chain weight only at exponent $j+d$, by comparing their $p$-adic valuations modulo $D$. Let $j_0$ be the first equality
\[
 z_1p^{j_0+d}\equiv z_2p^{j_0}\pmod D.
\]
It occurs at latest when both weights are zero; all later weights then coincide. The union has $(2n-s)+j_0$ elements, so $j_0\le s$. Thus $p^{2n-1-s}$ divides $z_1p^{t-s}-z_2$, giving the $p$-power part of \eqref{pg-eq:chain-congruence}. Equality of \eqref{pg-eq:lambda}, together with \eqref{pg-eq:determinants}, gives $p-1\mid z_1-z_2$, and hence $p-1\mid a-b$.

If $a=b$, then $z_2=p^d z_1$ and $\alpha_2=\lambda_1\zeta^{rp^dz_1}=\alpha_1^{p^d}$, contradicting disjoint special character orbits. If $t\le n-2$, the nonzero difference in \eqref{pg-eq:chain-congruence} has absolute value at least $(p-1)p^{2n-1-s-t}\ge(p-1)p^{n+1-s}$. But both $a,b$ are positive and strictly less than $p^{n+1-s}$ by \eqref{pg-eq:bounds}. This contradiction proves $t\ge n-1$.
\end{proof}

\begin{proposition}[One prime-to-$p$ Hasse value]\label{pg-prop:prime-families}
If one normalized Hasse value $z_0$ is prime to $p$, the other $z_*$ belongs with it to one of the two families \eqref{eq:pg-family-one} and \eqref{eq:pg-family-two}.
\end{proposition}
\begin{proof}
Proposition~\ref{pg-prop:chains} gives equal leading characters and $v_p(z_*)\in\{n-1,n\}$. Corollary~\ref{rc-cor:strictordering} gives $z_*<z_0$.

If $z_*=p^{n-1}b$, congruence~\eqref{pg-eq:chain-congruence} says that $z_0-b$ is a nonzero multiple of $(p-1)p^n$. Here $b<p^2$, whereas $z_0>(p-1)p^{n-1}>p^2$, so the multiple is positive. Since $z_0<p^{n+1}$, it is exactly the first positive multiple. The lower bound for $z_*$ gives $b>p-1$, hence $b\ge p+1$. Also
\[
 b<\frac{(p-1)p^n}{p^{n-1}-1}
 =p(p-1)+\frac{p(p-1)}{p^{n-1}-1}<p(p-1)+1.
\]
Thus $b\le p^2-p-1$, since the endpoint $p(p-1)$ is divisible by $p$. This proves \eqref{eq:pg-family-one}.

If $z_*=p^n b$, then $1\le b\le p-1$, and the congruence gives $z_0=b+(p-1)p^{n-1}\ell$. The bounds on $z_0$ give $1\le\ell\le p+1$. The strict ordering implies $(p-1)\ell>pb-b/p^{n-1}$, hence $(p-1)\ell\ge pb$ by integrality. For the stated range of $b$, this is equivalent to $\ell\ge b+1$, proving \eqref{eq:pg-family-two}.
\end{proof}

\subsection{Zero weight, canonical lines and lower levels}
We identify when the primitive grading forces the canonical directions to match, and record the quotient congruence needed when they do not.

\begin{lemma}\label{pg-lem:zero-weight}
If the weight-zero part of the primitive space in Lemma~\ref{pg-lem:scaling} has dimension one, the given torsion isomorphism identifies the two canonical $p$-lines.
\end{lemma}
\begin{proof}
Frobenius eventually takes every weight to zero. Its stable image therefore has dimension at most one, and has dimension at least one because neither coordinate is Frobenius nilpotent. For a sufficiently large $N$, the semilinear Fitting decomposition is
\[
 V=\ker(F^N)\oplus\operatorname{im}(F^N).
\]
It follows by stabilization of kernels and images and rank-nullity over the perfect field $k$. The primitive presentation in Lemma~\ref{pg-lem:primitives} splits accordingly as a tensor product. The nilpotent factor has only its zero $k$-point. The one-dimensional bijective factor has equation $u^p=cu$ with $c\ne0$ and exactly $p$ roots. Thus the support of the graph is an additive group of order $p$ in $\mathbf G_a^2$, an $\F_p$-line. Each coordinate is nonzero somewhere on that support, so both coordinate projections are injective.

Consequently the two rescaled coordinates vanish on specialization at exactly the same generic graph points. By the minimum-value induction, a coordinate vanishes precisely when the $p^{n-1}$ multiple of that point is on the canonical $p$-line: the minimum points are exactly the successive low divisions of the noncanonical first-level points. The torsion isomorphism commutes with multiplication by $p^{n-1}$, which is surjective onto $p$-torsion. Equality of these vanishing kernels identifies the canonical lines.

In particular, the criterion applies when the union of the two chains already represents $2n$ weights. Every represented weight is nonzero as a vector; the primitive dimension is $2n$, so weight zero then has dimension one.
\end{proof}

\begin{lemma}[A lower-level congruence]\label{pg-lem:lower-level}
For an unequal-type two-slope comparison at level $p^N$, $N\ge1$, suppose the normalized Hasse values are $z$ and $p^t b$, with $p\nmid zb$. No bound on $e$ is needed. Then
\begin{equation}\label{pg-eq:lower-level}
 (p-1)p^{\max(0,2N-1-t)}\mid z-b.
\end{equation}
\end{lemma}
\begin{proof}
Rescale with $D_N=p^{2N-1}$. The prime-to-$p$ coordinate has $2N$ distinct weights, including zero, so its chain spans the whole primitive space. The second coordinate has the same leading tame character, and its weight must appear in that chain. If $t<2N-1$, its only possible nonzero matching weight is at exponent $t$, giving $zp^t\equiv p^tb\pmod{D_N}$. If $t\ge2N-1$, no $p$-power divisibility is required. Equality of the leading and determinant characters always gives $p-1\mid z-p^tb$, hence $p-1\mid z-b$. The argument is valid also at $N=1$: it uses only the primitive dimension and grading, not the bound on $e$ in \eqref{pg-eq:bounds}.
\end{proof}

Let $C_i$ denote the canonical residual line, with character $\kappa_i$, and let $\eta_i$ be the quotient character. The high-root Hensel argument of Lemma~\ref{fc-lem:rational} and determinant comparison make both characters trivial over $G_L$. On the tame generator,
\begin{equation}\label{pg-eq:residual-characters}
 \kappa_i(\tau)=\zeta^{(p+1)e-z_i}\alpha_i^{-1},
 \qquad \eta_i(\tau)=\lambda_i^{-1}=\zeta^{pz_i}\alpha_i^{-1}.
\end{equation}
The first formula is \eqref{rc-eq:kappa}. For the second, a minimum-value $p^n$-point has leading field multiplier $\zeta^{rz_i}$; descended point action divides this by $\alpha_i$. Its leading quotient is the noncanonical residual quotient by the minimum-point description. Since $D\equiv p\pmod m$ and $p^2\equiv1\pmod m$, we have $r\equiv p\pmod m$. This proves the formula, independently of the lift of $\tau$, since the quotient character is trivial over $G_L$. It gives
\begin{equation}\label{pg-eq:scalar-condition}
 \kappa_i=\eta_i\quad\Longleftrightarrow\quad z_i\equiv e\pmod{p-1}.
\end{equation}
If the $\lambda_i$ agree, the quotient characters agree, and determinant equality makes the canonical characters agree. Failure of canonical-line matching then forces these two common characters to coincide: otherwise the canonical-character eigenspace is unique. In the coincident case, the transported noncanonical stable line together with the canonical line makes each residual representation scalar.

We will also use the following consequence of the quotient formulas and markings of Proposition~\ref{rc-prop:quotients} and Lemma~\ref{fc-lem:rational}. In normalized values, the canonical quotient changes $z$ to $pz$ for $z<e$, to $(p+1)e-z$ for $z>e$, and to one slope for $z=e$. Its dual kernel is noncanonical in the first case and canonical in the second. A noncanonical stable quotient changes $z$ to $z/p$ and has canonical dual kernel. A paired quotient of a level-$p^n$ comparison retains level $p^{n-1}$ and the dual marking.

Call $z>e$ the central case and $z<e$ the outer case. If $p\mid z$, equality $z=e$ is impossible. A central canonical quotient has prime-to-$p$ Hasse value and nontrivial residual wild action with a unique canonical fixed line. An outer canonical quotient has a stable noncanonical dual line, so its residual action over $L$ is trivial. Hence matched canonical quotients cannot pair a central and an outer source when the central source has $p$-divisible Hasse value. Two central sources with $p$-divisible values have prime-to-$p$ canonical quotient values, so Theorem~\ref{rc-thm:primehasse} gives equal quotient types. The isogenies are defined over $K$ and preserve rational special inertia traces. We have therefore excluded a matched-canonical comparison with both Hasse values divisible by $p$ whenever at least one source is central.

\subsection{Eliminating positive Hasse valuations}
We now prove that the only configuration with two positive $p$-adic Hasse valuations is the quadratic-sign configuration in the main reduction.

\begin{proposition}\label{pg-prop:positive}
If $s_1,s_2>0$, then \eqref{eq:pg-family-sign} and \eqref{eq:pg-family-congruence} hold. The residual representations are split with distinct characters, and the torsion isomorphism interchanges their canonical and opposite residual lines.
\end{proposition}
\begin{proof}
Order $s=s_1\le s_2=t$. Proposition~\ref{pg-prop:chains} leaves $t=n-1$ or $n$.

First suppose $t=n-1$ and $s<n-1$. The leading characters agree. Write $z_1=p^sa$, $z_2=p^{n-1}b$. Congruence~\eqref{pg-eq:chain-congruence} gives
\[
 z_1=p^sb+(p-1)p^n\ell.
\]
Here $\ell\ne0$. A negative value would force $b-a\ge(p-1)p^{n-s}$, although $b<p^2$ and $n-s\ge2$. Thus $\ell>0$, and $z_1<p^{n+1}$ makes $\ell=1$. The lower bound on $z_2$ gives $b\ge p+1$. Both $z_i>p^n>e$ are central. The first chain coincidence is at $j_0=s$, because $v_p(a-b)=n-s$; thus the weight union has $2n$ elements. Lemma~\ref{pg-lem:zero-weight} matches the canonical lines, and the central-quotient exclusion above is a contradiction.

If $s=t=n-1$, the nonzero difference $a-b$ is divisible by $(p-1)p$. But the bounds and prime-to-$p$ condition give $p+1\le a,b\le p^2-1$, so $|a-b|<(p-1)p$. This is impossible. All cases with $t=n-1$ are excluded.

Now suppose $t=n$ and $1\le s<n$. Write $z_2=p^nb$, where $1\le b\le p-1$. The leading characters again agree, and
\begin{equation}\label{pg-eq:mixed-positive}
 z_1=p^sb+(p-1)p^{n-1}\ell,\qquad1\le\ell\le p+1.
\end{equation}
The lower bound on $z_1$ gives positivity, and its upper bound gives the displayed range. If $p\nmid\ell$, then $j_0=s$, so the canonical lines match. Model two is central; the matched-canonical exclusion applies whether model one is central or outer. Thus $\ell=p$, the only multiple of $p$ in its range. Both sources are now central. Their canonical lines cannot match, so equality of the leading characters forces scalar residual action by \eqref{pg-eq:scalar-condition}. Since $(p-1)\mid e$ and $z_i\equiv b\pmod{p-1}$, we obtain $b=p-1$. Hence
\[
 z_1=(p-1)(p^n+p^s),\qquad z_2=(p-1)p^n.
\]
Pair the canonical line of model two with its noncanonical inverse image in model one. The quotient values, at remaining level $p^{n-1}$, are
\[
 w_1=(p-1)(p^{n-1}+p^{s-1}),\qquad
 w_2=(p+1)e-(p-1)p^n.
\]
The second is prime to $p$. If $s=1$, both are, contradicting Theorem~\ref{rc-thm:primehasse}. If $s\ge2$, write $e=(p-1)u$ and apply Lemma~\ref{pg-lem:lower-level} at $N=n-1$; it gives
\begin{equation}\label{pg-eq:positive-arithmetic}
 p^{2n-s-2}\mid A_0:=(p+1)u-p^n-1-p^{n-s}.
\end{equation}
The two-slope and base bounds give
\[
 p^{n-1}+p^{s-1}<u<\frac{p^n}{p-1},\qquad
 0<A_0<\frac{2p^n}{p-1}<3p^{n-1}.
\]
For $s\le n-2$, the divisor in \eqref{pg-eq:positive-arithmetic} is at least $p^n$, whereas $A_0<2p^n/(p-1)<p^n$. For $s=n-1$, integrality gives $u\ge(p+1)p^{n-2}+1$, and therefore $A_0\ge(2p+1)p^{n-2}>2p^{n-1}$. There is no multiple of $p^{n-1}$ strictly between $2p^{n-1}$ and $3p^{n-1}$. Both cases contradict \eqref{pg-eq:positive-arithmetic}.

It remains that $s=t=n$. Write $z_i=b_ip^n$, $1\le b_i\le p-1$. Both sources are central, so their canonical lines cannot match. Pairing each canonical line with its noncanonical image in the other model and applying Lemma~\ref{pg-lem:lower-level} at level $n-1$ gives
\[
 \begin{split}
 (p-1)p^{n-2}&\mid(p+1)e-b_1p^n-b_2,\\
 (p-1)p^{n-2}&\mid(p+1)e-b_2p^n-b_1.
 \end{split}
\]
Subtracting gives $p^{n-2}\mid b_1-b_2$. Since $n\ge3$ and $|b_1-b_2|<p$, we have $b_1=b_2=b$. Reduction modulo $p-1$ gives $(p-1)\mid2b$, so $b=(p-1)/2$ or $p-1$. The latter would make both residual representations scalar. Their scalar characters, and hence their $\lambda_i$, would agree; equal $z_i$ in \eqref{pg-eq:lambda} would then imply $\alpha_1=\alpha_2$, contrary to unequal type.

Thus $b=(p-1)/2$. Write $\alpha_i=\zeta^{a_i}$. The deformation weight gives $z_i\equiv-a_i\pmod{p+1}$. Equal $z_i$ makes $\alpha_2/\alpha_1\in\mu_{p-1}$; determinant equality also puts it in $\mu_{p+1}$. Its order divides two. Equality is excluded by unequal type, so $\alpha_2=-\alpha_1$. The first quotient congruence is exactly \eqref{eq:pg-family-congruence}. Equation~\eqref{pg-eq:scalar-condition} now gives distinct residual characters. A transported noncanonical stable line splits the representation, and the failure of canonical matching forces interchange with the opposite character line. This proves every assertion.
\end{proof}

\begin{proof}[Proof of Theorem~\ref{thm:three-family-reduction}]
If a normalized Hasse value is prime to $p$, Proposition~\ref{pg-prop:prime-families} gives the first two families. If both are divisible by $p$, Proposition~\ref{pg-prop:positive} gives the third. The graph grading and quotient argument provide necessary conditions only. They neither construct a comparison in any family nor identify groups with opposite tangent characters as quadratic twists.
\end{proof}

\begin{corollary}\label{pg-cor:exclusion}
None of the three configurations in Theorem~\ref{thm:three-family-reduction} can occur for an unequal-character comparison under its hypotheses.
\end{corollary}
\begin{proof}
In the first family, $z_2=p^{n-1}b\ge(p+1)p^{n-1}>p^n>e$. In the second, $z_2=p^nb\ge p^n>e$. In the third, $z_1=z_2=((p-1)/2)p^n\ge2p^n>e$. Thus every family has a central model with $p$-divisible normalized Hasse value $z$. Since $p\nmid e$, the integer $(p+1)e-z$ is prime to $p$. The strict two-slope bound $z<pe$ is part of the hypotheses, and the available torsion level is at least $p^2$. Theorem~\ref{thm:central-p2-rigidity} therefore forces equal special traces, a contradiction in every case.
\end{proof}

For instance, the arithmetic quadratic-sign candidates at $p=5,n=3$ have $z_1=z_2=250$ and $e=52,72,92,112$. Each is central and has $5\nmid6e-250$, so the same criterion excludes all four. This exclusion uses actual point fields; it makes no identification of arbitrary opposite-sign groups as quadratic twists.

\section{Terminal images and exact point-field refinements}\label{sec:point-refinements}
This section records the point-field information behind the first-wild comparison. The terminal-image estimate works for every prime at least five. At five we determine the two point-field filtrations exactly and construct a base with the necessary first cyclotomic break. The latter construction concerns the base field alone; it does not construct two groups with unequal characters.

Throughout the first two subsections, the residue field is algebraically closed, $\mu_p\subset K$, $p\ge5$, $e_K=p(p-1)$, and $L/K$ is the tame extension of degree $N=p^2-1$. Put $s=pN$. A two-slope model has integral parameter $0<B<pe_K$ and first-level values
\[
                     \alpha=s-B,\qquad \beta=B/p.
\]
The high line $H$ is rational over $L$: scaling its degree-$(p-1)$ root factor leaves simple nonzero roots over the algebraically closed residue field, which lift by Hensel's lemma. Its quotient character is trivial, since the determinant is cyclotomic and $\mu_p\subset L$. Thus the residual image over $L$ is trivial or has order $p$, with unique fixed line $H$ in the latter case. Twice dividing a low point by $p$ gives third-level value $B/p^5$. These assertions also follow from the Newton polygons in Sections~\ref{sec:quotients} and~\ref{sec:restricted}.

\subsection{A terminal plane and a wild image bound}
We recover the high line and its parameter from a small full division field.
\begin{lemma}\label{lem:terminal-plane-refinement}
Suppose $B=pb$ with $p\nmid b$ and the full $p^3$-division field $F/L$ has degree $p^4$. In a basis with $H=\F_pe_1$, the terminal space
\[
 V_2=\{C\in M_2(\F_p):1+p^2C\in\operatorname{Gal}(F/L)\}
\]
has dimension two. Every $C\in V_2$ has $C_{21}=0$, and the kernel of $C\mapsto C_{22}$ is exactly $\F_pE_{12}$.
\end{lemma}
\begin{proof}
Normalize $w=p^4v_L$. Every primitive third-level point $P$ with $v=p^2P\notin H$ has $w(T(P))=b$. Its valuation over $L$ has denominator $p^4$, so $L(P)=F$. If a wild automorphism $g$ has lower break $t(g)$, expansion in a uniformizer with fixed Teichm\"uller coefficients gives
\[
 w(gx-x)\ge w(x)+t(g),
\]
with equality when $p\nmid w(x)$, and strict inequality when $p\mid w(x)$. Indeed the first monomial $c\pi^a$ changes to first order by $ac\pi^a(g\pi/\pi-1)$; later exponents have strictly larger valuation. Formal addition gives $w(T(P+R)-T(P))=w(T(R))$, because its difference is $T(R)$ times a unit. Consequently
\begin{equation}\label{eq:terminal-point-break}
                      t(g)=w(T(gP-P))-b.
\end{equation}
A nonzero terminal matrix cannot kill any $v\notin H$, since $P$ generates $F$. For $v=(a,1)$, a nonzero $C_{21}$ would make the second coordinate of $Cv$ vanish for exactly one $a$. At that value $Cv$ is nonzero and high, while at the other values it is low. Equation~\eqref{eq:terminal-point-break} would give two different breaks for the same automorphism. Thus $C_{21}=0$.

If also $C_{22}=0$, nonvanishing of $C_{11}a+C_{12}$ for every $a$ forces $C_{11}=0$. The kernel of $C\mapsto C_{22}$ is therefore contained in $\F_pE_{12}$, and $\dim V_2\le2$. The first congruence space has dimension at most $\dim V_2$, by the $p$th-power map. The residual image has order at most $p$. If $\dim V_2\le1$, the total order would be at most $p^3$, contrary to $[F:L]=p^4$. Hence its dimension is two and the asserted kernel is the full line.
\end{proof}

\begin{proposition}\label{prop:terminal-parameter-recovery}
Suppose two two-slope models have isomorphic $p^3$-torsion and common full division field of degree $p^4$ over $L$. If $B_1=pb_1$ with $p\nmid b_1$, then $H_1=H_2$, $B_1=B_2$, and their special characters agree.
\end{proposition}
\begin{proof}
The denominator of $B_2/p^5$ forces $B_2=pb_2$, with $1\le b_i\le e_K-1$. In the first model, the terminal matrices of the preceding lemma have exact breaks
\[
 t_{\rm low}=(p^4-1)b_1\quad(C_{22}\ne0),\qquad
 t_{\rm high}=p^5N-(p^5+1)b_1\quad(0\ne C\in\F_pE_{12}).
\]
If the high lines differed, choose a residual vector outside both. Its translate by $E_{12}$ is high in the first model and low in the second. The general ramification inequality in the proof of the lemma gives
\[
 t_{\rm high}\le(p^4-1)b_2\le(p^4-1)(e_K-1).
\]
But $p^5N=(p^5+p^4)e_K$, so
\[
 t_{\rm high}=(p^4-1)e_K+(p^5+1)(e_K-b_1)>(p^4-1)e_K.
\]
Thus the lines agree. A matrix with $C_{22}\ne0$ now gives $b_1\le b_2$ from the low displacement inequality; $E_{12}$ gives $b_1\ge b_2$ from the high one. Equality follows. The common high-line character $\zeta^{c_i-B_i}$ then gives $c_1=c_2\pmod N$, hence equality of special characters.
\end{proof}

\begin{corollary}\label{cor:wild-image-at-least-five}
For two potentially supersingular height-two groups with isomorphic $p^3$-torsion, unequal special characters and nontrivial residual wild inertia, the common full division field over $L$ has order between $p^5$ and $p^9$.
\end{corollary}
\begin{proof}
Nontrivial residual wild inertia excludes the one-slope case, whose first-level points are rational over $L$. The common residual transvection identifies both high lines. Write the full degree as $p^r$. If $p\nmid B_i$, the value $B_i/p^5$ forces $r\ge5$. If $v_p(B_i)=1$, it forces $r\ge4$; equality is excluded by Proposition~\ref{prop:terminal-parameter-recovery}. If $r\le4$ remains possible, both parameters must therefore be $B_i=p^2k_i$ with $1\le k_i\le p-2$. Equal high-line characters give $c_2-c_1\equiv B_2-B_1\pmod N$. Equal determinants give $c_2-c_1\equiv0\pmod{p-1}$. Hence $k_2-k_1$ is divisible by $p-1$; its absolute value is smaller than $p-1$, so it is zero. The characters again agree, a contradiction. Finally the residual group has order $p$ and the two higher matrix layers have dimension at most four each, giving the upper bound.
\end{proof}

\subsection{Exact filtrations for the exchanged parameters at five}
We determine the filtrations directly from point coordinates; their exact values contain more information than the conductor bounds used earlier.
\begin{proposition}\label{prop:exact-exchanged-point-profiles}
Consider two potentially supersingular height-two models with tame reducible residual inertia and unequal special characters, at $p=5$, $e_K=20$, with parameters $B_1=25$, $B_2=75$, exchanged high lines $\F_5e_1,\F_5e_2$, and an isomorphism of their $125$-torsion. Assume the residual tame weights are normalized as in the proof of Theorem~\ref{thm:tame-five-recovery}. The full image over $L$ is $1+5M_2(\Z/125)$. For $5P_i=e_i$, the fields $E_i=L(P_i)$ are Galois of degree $25$ and have lower breaks
\[
                  E_1/L:(6,336),\qquad E_2/L:(24,234).
\]
The full second division field has the following breaks:
\[
\begin{array}{c|rrrr}
\text{upper}&6&24&66&72\\
\text{lower}&6&96&1146&1896.
\end{array}
\] Its cyclic degree-five subfields have breaks $6,24,66,72$ with respective multiplicities $1,5,25,125$. The cyclotomic quotient $L(\mu_{25})/L$ has break $72$, and $K(\mu_{25})/K$ has break $3$.
\end{proposition}
\begin{proof}
The first congruence space $W_1$ contains both off-diagonal axes and has both diagonal projections nonzero, by Proposition~\ref{prop:opposite-quotient-image} and its proof. Thus each column evaluation is onto $\F_5^2$. The fields $E_i$ are the corresponding quotient fields of the elementary abelian second-level extension, so they are Galois of degree $25$. The choice of lift $P_i$ does not change its stabilizer.

The first-level high and low values are $(95,5)$ and $(45,15)$. Their second-level values are $(19/5,3/5)$ for $P_1$ and $(1/5,9/5)$ for $P_2$. Multiplication by $[E_i:L]=25$ and formal addition therefore give
\[
\begin{array}{c|c|r|r|r|c}
\text{field}&\text{model}&A&D_0&D_1&H\\\hline
E_1&1&95&125&2375&\text{diagonal}\\
E_1&2&15&375&1125&\text{off-diagonal}\\
E_2&1&5&125&2375&\text{off-diagonal}\\
E_2&2&45&375&1125&\text{diagonal}
\end{array}
\]
Here $A$ is the coordinate valuation, $D_1$ the displacement on the indicated line and $D_0$ the displacement off it. Every displacement is below the absolute index $12000$.

The tame residual eigencharacters are $1,\zeta^6$, up to a common quadratic twist. The diagonal, $21$ and $12$ quotient characters consequently have breaks congruent to $0,6,18$ modulo $24$. Indeed a tame lift on a cyclic degree-five quotient has leading multiplier $\rho$ with $\rho^5=\zeta$, and conjugation on a break-$b$ coefficient is $\rho^b$; multiplication by five fixes these three residues. The two distinct tame weights in each point field force two distinct breaks, because a ramification graded piece has a single tame scalar weight. The last line is the direction with larger cyclic quotient break. These statements can also be read from the quotient index formula of \citet[Proposition~3.7, pp.~14--15, and Proposition~3.13, p.~16]{ColmezLocal}.

Write $d_i,o_i$ for the diagonal and off-diagonal cyclic quotient breaks. In $E_1$, an off-diagonal last line would, by Lemma~\ref{lem:opposite-line-conductor} and the first coordinate row, force $o_1=(125-95)/5=6$. The smaller positive diagonal break would be divisible by $24$, an impossibility. The last line is diagonal, and the second row gives
\[
                         d_1=(375-15)/5=72,\qquad o_1<72.
\]
For $E_2$, a diagonal last line gives $d_2=(125-5)/5=24$. An off-diagonal last line gives $o_2=(375-45)/5=66$ and $d_2<66$. Thus $d_2$ cannot equal $72$.

If $\dim W_1=3$, its diagonal subspace is one dimensional with both projections nonzero. The two diagonal functionals then have the same kernel and define the same field, contradicting $d_1\ne d_2$. Hence $W_1=M_2(\F_5)$. The fifth-power inclusion gives $W_2=M_2(\F_5)$, and counting both layers gives the full congruence image. Its nonzero trace modulo $25$ also shows $\mu_{25}\not\subset L$.

We refine the two filtrations using the fixed-coefficient expansion. If $v_5(A)=1$ and $r$ is the first nonzero prime-to-five exponent of a coordinate, then below the absolute index its only possible leading changes at break $b$ are
\begin{equation}\label{eq:first-wild-leading-changes}
                              A+5b,\qquad r+b.
\end{equation}
To prove this, expand each monomial in $\pi(1+u)$, $v(u)=b$. An exponent divisible by five first contributes modulo five in degree at least five; the initial exponent $A$ has order exactly one and contributes in degree five. Every later exponent has greater value. The earliest exponent prime to five is $r$ and contributes in degree one. Binomial coefficients divisible by five produce terms above the absolute index. Unequal leading values cannot cancel. In particular a displacement greater than $A+5b$, but below the absolute index, forces $r=A+4b$; for any smaller break $t<b$ the displacement is then exactly $A+5t$.

For $E_1$, let $t=o_1$ be the smaller lower break. The first coordinate row and the strict inequality for a coordinate divisible by five give $125>95+t$, so $t<30$. Since $t\equiv6\pmod{24}$, it is $6$. If $u$ is the last lower break, $72=(u+4t)/5$, hence $u=336$.

For $E_2$, a diagonal last line would give larger quotient break $24$, smaller off-diagonal break $18$, and lower breaks $(18,48)$. In its second coordinate row, $1125>45+5\cdot48$ forces $r=45+4\cdot48$ in~\eqref{eq:first-wild-leading-changes}. The displacement at break $18$ must then be $45+5\cdot18=135$, contradicting $375$. Thus its last line is off-diagonal. The larger quotient break is $66$ and the smaller diagonal break $t$ is $24$ or $48$, giving last lower break $u=330-4t$, namely $234$ or $138$. In the first coordinate row, $2375>5+5u$ forces the lower displacement to be $5+5t=125$. Therefore $t=24$ and $u=234$.

The coordinate characters on $M_2(\F_5)$ now have upper breaks
\[
                    A_{21}:6,\quad A_{22}:24,\quad
                    A_{12}:66,\quad A_{11}:72.
\]
A nonzero linear combination has the largest break among its active coordinates, because the other terms vanish on the last nonzero ramification group. Herbrand's inverse function gives successively
\[
 6,\quad6+5(24-6)=96,\quad96+25(66-24)=1146,
 \quad1146+125(72-66)=1896.
\]
Normalizing the last nonzero coefficient of a projective character to one gives $1,5,25,125$ cyclic quotients at these four breaks. The determinant modulo $25$ is $A_{11}+A_{22}$, hence has break $72$. It is the cyclotomic character over $L$. A degree-$24$ tame base change multiplies a cyclic wild break by $24$, as in the uniformizer calculation in the preceding section; the break over $K$ is therefore $3$.
\end{proof}

The proposition is a conditional calculation of the point fields and does not assert that the unequal-character pair exists. Its proof determines the ramification filtration from coordinates and quotient fields, without deriving a table from the impossibility theorem.

\subsection{The exact third-point towers and terminal filtration}
We continue the direct coordinate calculation of Proposition~\ref{prop:exact-exchanged-point-profiles}. This supplies an alternative obstruction in the actual determinant quotient, without using the preceding character-recovery theorem.

Keep the mixed-characteristic base with algebraically closed residue field, $p=5$, $e_K=20$, $\mu_5\subset K$, and tame good extension $[L:K]=24$. Assume the entire comparison hypotheses of Proposition~\ref{prop:exact-exchanged-point-profiles}: two whole potentially supersingular height-two groups have unequal finite special characters, tame reducible residual inertia, isomorphic actual $125$-torsion, parameters $B_1=25,B_2=75$, and exchanged high axes $\F_5e_1,\F_5e_2$, with the residual tame characters normalized there. Put
\[
 F=L(M[125]),\qquad F_2=L(M[25]),\qquad
 W=1+5M_2(\Z/125),\qquad U=1+25M_2(\F_5).
\]
Proposition~\ref{prop:exact-exchanged-point-profiles} gives these full actual images. All lower breaks below use the integer valuation on the top field of the named extension.

\begin{proposition}\label{prop:exact-third-point-profiles}
One can choose a common basis $Q_1,Q_2$ of the actual $125$-torsion with $25Q_i=e_i$ and point fields stable under a tame lift. Set $P_i=5Q_i$, $E_i=L(P_i)$ and $E_i^+=L(Q_i)$. Then $E_i^+/E_i$ is Galois of degree $25$, and
\[
\begin{array}{c|c|r|r}
 &\text{lower breaks}&\text{diagonal quotient}&\text{off-diagonal quotient}\\\hline
 E_1^+/E_1&(606,6936)&1872&606\\
 E_2^+/E_2&(624,6834)&624&1866.
\end{array}
\]
The directions are the actual first-level translations on the two residual axes.
\end{proposition}
\begin{proof}
Choose a lift $\tau$ of the tame generator. Its two residual eigenvalues are distinct. Hensel lifting of the simple roots of its characteristic polynomial and the polynomial idempotents
$(\tau-\lambda_j)/(\lambda_i-\lambda_j)$ split the free $\Z/125$-module into two free rank-one eigensummands. Choose $Q_i$ on them. Then $\tau Q_i$ is a unit multiple of $Q_i$, so the point fields are $\tau$-stable. A finite-order choice of $\tau$ is unnecessary.

Let $S_i$ and $T_i$ be the stabilizers of $Q_i$ and $P_i$ in $W$. The full third-level column orbit has $625$ elements, and the second-level orbit has $25$ elements. For $g\in T_i$, write $gQ_i-Q_i=25b_i$ in the common torsion coordinates. The map $g\mapsto b_i\in\F_5^2$ is a homomorphism, because $g\equiv1\pmod5$. The terminal matrices already make it surjective, and its kernel is $S_i$. Hence $S_i$ is normal in $T_i$, proving the Galois assertion and the precise translation marking. We do not assert that $E_i^+/L$ is Galois.

The second-level point values are $19/5,3/5,1/5,9/5$ in the four field/model rows. All are below their corresponding first-level low slopes. Dividing once more therefore divides each valuation by $25$, while the field degree grows from $25$ to $625$. The normalized coordinate values remain $95,15,5,45$, and formal addition gives
\[
\begin{array}{c|c|r|r|r|c}
\text{field}&\text{model}&A&D_0&D_1&H\\\hline
E_1^+&1&95&3125&59375&\text{diagonal}\\
E_1^+&2&15&9375&28125&\text{off-diagonal}\\
E_2^+&1&5&3125&59375&\text{off-diagonal}\\
E_2^+&2&45&9375&28125&\text{diagonal}.
\end{array}
\]
Every $A$ has five-adic order one, and every displacement is below $v_{E_i^+}(5)=300000$. The tame weights are again zero and six for the first tower, and eighteen and zero for the second. Indeed conjugation on a translation divides by the eigenvalue of the source $Q_i$. A cyclic degree-five quotient over $E_i$ has total ramification degree $125$ over $L$, so its leading tame multiplier satisfies $\rho^{125}=\zeta$. Thus the diagonal quotient break is a positive multiple of $24$, and the off-diagonal quotient break is six or eighteen modulo $24$. There are two distinct lower breaks $t<u$, with an eigenaxis as last line.

For a row with coordinate valuation $A$, let $r$ be the first nonzero exponent prime to five in its fixed-coefficient expansion. The first possible displacement at break $b$, below the absolute index, is $\min(A+5b,r+b)$. In particular a displacement larger than $A+5b$ forces $r=A+4b$; at every smaller break $t$ the displacement is then exactly $A+5t$.

In the first tower, an off-diagonal last line would force the larger cyclic quotient break $(3125-95)/5=606$ by Lemma~\ref{lem:opposite-line-conductor}. Thus $t<606$ and $u=3030-4t$. The second model has $28125>15+5u$, so its lower displacement would be $15+5t=9375$, forcing $t=1872$, a contradiction. The last line is diagonal. The second row then gives larger quotient break $1872$, so $u=9360-4t$. The first row has $59375>95+5u$ and forces $3125=95+5t$. Hence $t=606$ and $u=6936$.

For the second tower, a diagonal last line would give larger quotient break $(3125-5)/5=624$, so $t<624$ and $u=3120-4t$. The fourth row has $28125>45+5u$, forcing $9375=45+5t$ and $t=1866$, a contradiction. Its last line is off-diagonal. The fourth row gives larger quotient break $1866$, hence $u=9330-4t$. The third row has $59375>5+5u$ and forces $3125=5+5t$. Thus $t=624$ and $u=6834$. The quoted cyclic breaks follow from $(u+4t)/5$.
\end{proof}

\begin{proposition}\label{prop:exact-terminal-filtration}
The terminal extension $F/F_2$ would have lower breaks
\[
                  7896,\quad11646,\quad769146,\quad937896,
\]
with coordinate directions removed in the order $B_{22},B_{21},B_{12},B_{11}$ for $1+25B\in U$. Outside $U$, the lower breaks in $F/L$ equal the first-layer breaks $6,96,1146,1896$, in the order $A_{21},A_{22},A_{12},A_{11}$.
\end{proposition}
\begin{proof}
Put $K_i=F_2E_i^+$. The terminal group maps onto $T_i/S_i$, so $T_i=S_iU$. Galois correspondence gives $E_i^+\cap F_2=E_i$ and
\[
 \operatorname{Gal}(K_i/E_i)
 =\operatorname{Gal}(E_i^+/E_i)\times\operatorname{Gal}(F_2/E_i).
\]
Both factors are actual Galois extensions of degree $25$. The group $\operatorname{Gal}(K_i/F_2)$ is the $i$th terminal column $U/(U\cap S_i)$. The two column kernels have trivial intersection, so $K_1K_2=F$.

Lower numbering is inherited by subgroups. In the first case the subgroup for $F_2/E_1$ has coordinates $A_{22},A_{12}$, and in the second it has $A_{21},A_{11}$. Their lower and upper break lists are
\[
\begin{array}{c|c|c}
 &\text{lower}&\text{upper}\\\hline
 F_2/E_1&(96,1146)&(96,306)\\
 F_2/E_2&(6,1896)&(6,384).
\end{array}
\]
The new tower upper breaks are $(606,1872)$ and $(624,1866)$. All four breaks in each elementary abelian compositum are distinct. The unequal-break sum rule and upper-quotient compatibility therefore give their union, and inverse Herbrand integration gives
\[
\begin{array}{c|c|c}
 &\text{upper for }K_i/E_i&\text{lower for }K_i/E_i\\\hline
 i=1&(96,306,606,1872)&(96,1146,8646,166896)\\
 i=2&(6,384,624,1866)&(6,1896,7896,163146).
\end{array}
\]
Restricting to the terminal column subgroups retains the last two lower values. Converting each column back to upper numbering gives
\[
 K_1/F_2:\ (8646,40296)\quad(B_{21},B_{11}),\qquad
 K_2/F_2:\ (7896,38946)\quad(B_{22},B_{12}).
\]
These markings are restrictions of the actual maps $g\mapsto gQ_i-Q_i$, so they cannot be changed by an abstract identification of Galois groups.

The two columns generate $F/F_2$. Their coordinate upper breaks are
\[
               7896,\quad8646,\quad38946,\quad40296.
\]
Integrating the indices $1,5,25,125$ yields
\[
\begin{aligned}
 b_1&=7896,& b_2&=7896+5\cdot750=11646,\\
 b_3&=11646+25\cdot30300=769146,&
 b_4&=769146+125\cdot1350=937896.
\end{aligned}
\]

Every nonidentity element of $U$ has lower break at least $7896$. In any nonzero coset of $U$, the quotient average shows that some lift has smaller break, because the quotient $F_2/L$ has last lower break $1896$. Multiplying that lift by an element of $U$ preserves the smaller break, so every lift has the same break. The quotient average then identifies it with the first-layer break. This proves the final assertion. The only ramification inputs are subgroup restriction and the norm-polynomial/Herbrand quotient formulas of \citet[Proposition~3.7, Theorem~3.10 and Proposition~3.13(ii), pp.~14--16]{ColmezLocal}.
\end{proof}

\begin{corollary}\label{cor:determinant-fibre-obstruction}
The preceding two exact filtrations cannot arise from the stated actual torsion comparison. They give a second proof of tame reducible character recovery at $p=5$, $e_K=20$.
\end{corollary}
\begin{proof}
The actual determinant on $W$ has image $1+5\Z/125$, of order $25$, and is cyclotomic over $L$. Fix a nonzero $s\in\F_5$ and its order-five element $1+25s$. Write $g=1+5A+25B$ with entries of $A,B$ chosen in $\{0,1,2,3,4\}$. The determinant condition requires $\operatorname{tr}A=0\pmod5$ and
\[
 \operatorname{tr}B=s-\det A-\operatorname{tr}A/5\pmod5.
\]
Thus every trace-zero $A$ has $125$ lifts in the fibre, with the integer trace carry included.

Among the $124$ nonzero trace-zero first-layer matrices, $100$ have $A_{21}\ne0$, $20$ have $A_{21}=0$ with nonzero diagonal, and four have only $A_{12}\ne0$. Their breaks are $6,96,1146$. For $A=0$, the condition is $\operatorname{tr}B=s$; its terminal coordinate counts are $100,20,4,1$ in the removal order. The complete fibre distribution is
\[
\begin{array}{c|rrrrrrr}
 b(g)&6&96&1146&7896&11646&769146&937896\\\hline
 \#g&12500&2500&500&100&20&4&1.
\end{array}
\]
It has $15625$ elements and weighted break sum $5925000$. The exact quotient-index formula of \citet[Proposition~3.7, pp.~14--15]{ColmezLocal}, after subtracting one, would give the nonidentity determinant quotient element lower break
\[
                  \frac{5925000}{15625}=\frac{1896}{5}.
\]
This is impossible for a normalized discrete valuation: a lower break is an integer. This contradiction uses the exact actual-field tables and determinant fibre, without appealing to a later recovery theorem or to Hasse--Arf.
\end{proof}


\subsection{A base with cyclotomic break three}
We show that the first cyclotomic condition in the preceding proposition can occur for a complete base field.
\begin{proposition}\label{prop:break-three-base}
There is a complete discretely valued field $K$ with algebraically closed residue field of characteristic five, absolute index $20$ and $\mu_5\subset K$, such that $K(\mu_{25})/K$ is cyclic of degree five with lower break three, and $K(\mu_{125})/K$ is cyclic of degree $25$. After a tame extension $L/K$ of degree $24$, the extension $L(\mu_{25})/L$ has lower break $72$.
\end{proposition}
\begin{proof}
Let $k$ be algebraically closed of characteristic five, put $K_0=W(k)[1/5]$, $K_1=K_0(\zeta_5)$ and $\pi_1=\zeta_5-1$. The shifted cyclotomic polynomial is Eisenstein, so $e_{K_1}=4$ and $\pi_1$ is a uniformizer. Set
\[
                  f(T)=(1+T)^5-\zeta_5+\pi_1^4T^2.
\]
Its constant term is $-\pi_1$, all its other nonleading coefficients have positive valuation, and it is monic. Thus it is Eisenstein. For a root $t$, the field $K=K_1(t)$ has index $20$ and uniformizer $t$.

Put $z=\zeta_{25}-(1+t)$. The equation $f(t)=0$ gives
\[
 z^5+5(1+t)z^4+10(1+t)^2z^3+10(1+t)^3z^2
           +5(1+t)^4z-\pi_1^4t^2=0.
\]
The constant coefficient has $K$-valuation $22$, while all four intermediate coefficients have valuation $20$. They lie strictly above the Newton segment from $(0,22)$ to $(5,0)$. Hence every root has valuation $22/5$, and $[K(\zeta_{25}):K]=5$. This extension is cyclic because $\mu_5\subset K$.

Normalize its valuation by $w=5v_K$. For a nonidentity automorphism,
\[
 w(z)=22,\qquad
 \sigma(z)-z=(\zeta_5^a-1)\zeta_{25},\qquad
 w(\sigma(z)-z)=25.
\]
Since $22$ is prime to five, the uniformizer-expansion equality gives lower break $25-22=3$.

The extension $K_1(\mu_{125})/K_1$ is cyclic of degree $25$ and has unique degree-five subfield $K_1(\mu_{25})$. A nontrivial intersection with the degree-five field $K$ would force equality with that subfield, contrary to $[K(\mu_{25}):K]=5$. The intersection is therefore $K_1$, proving the degree and cyclicity assertion. The same tame uniformizer calculation as above multiplies the degree-five break by $24$, giving $72$ over $L$. No formal-group existence is inferred from this field construction.
\end{proof}

\section{Remaining questions}\label{sec:open}
This section identifies the wild-index and rational-torsion questions left open by the proved results. Theorem~\ref{thm:complete-prime-index} determines every full-class height-two index with $p\nmid e$ and $\mu_p\subset K$. The optimal supersingular-only range beyond Theorem~\ref{thm:ss-iterated-range} is a separate question, since the sharp full-class lower examples can be mixed comparisons.

At indices divisible by $p$, Theorem~\ref{thm:first-wild-general} gives $\nuMin(p,p(p-1),2)=2$ for every $p\ge5$, and Theorem~\ref{thm:wild-bands-exact} determines every coefficient $2\le k\le(p^2-1)/2$ in $e=kp^r(p-1)$: the index is $r+2$ when $p\nmid k$ or $k=p$, and $r+3$ for the other multiples of $p$. Theorem~\ref{thm:all-index-upper} gives the uniform sufficient range $e\le p^{n-1}$ for every index, and Theorem~\ref{thm:mixed-all-indices} determines the exact mixed-comparison threshold at every index. The remaining full-class problem concerns comparisons between two whole supersingular groups outside these proved ranges. More precisely, Corollary~\ref{cor:all-index-two-slope-reduction} reduces recovery under
\[
 e\le p^{2n-3-v_p(e)}
\]
to two models with two first-level Newton slopes. Theorem~\ref{thm:all-index-central-reduction} further produces a central comparison at a remaining level $p^r$, where $r\ge v_p(e)+2$ and the high numerator has exact $p$-adic valuation $v_p(e)$. The partner polygon is unrestricted, and the original numerical bound is not automatically valid at level $r$. For $v_p(e)>0$, exclusion of these terminal configurations remains open. Theorem~\ref{thm:full-supersingular-maximum} expresses the full index exactly as the maximum of the supersingular-only index and the mixed lower threshold. Excluding the remaining terminal comparisons would give
\[
 \nuMin(p,e,2)=\left\lceil\frac{1+\log_p e+v_p(e)}2\right\rceil
\]
for every positive multiple $e$ of $p-1$. This is a proposed extension, not a proved formula in the remaining wild cases. For $r\ge1$, the next coefficient $k=(p^2+1)/2$ is prime to $p$ and remains outside the proved exact band. Theorem~\ref{thm:first-boundary-normal-form} forces any unequal comparison there to have $z_1=z_2=((p-1)/2)p^{r+3}$, opposite tangent scalars, and split residual representations with exchanged canonical lines. Corollary~\ref{cor:boundary-point-field-degrees} further forces canonical point-field degrees at least $p^{r+3}$ at level $p^{n-1}$ and $p^{r+5}$ at level $p^n$. The translation spaces persist along the actual point tower, so the last three relative steps all have degree $p^2$; when $r=1$, the successive total degrees are exactly $p^2,p^4,p^6$. Corollary~\ref{cor:boundary-nonsplit-cartan-image} further excludes containment of either actual finite image over the common tame good field in a principal nonsplit Cartan. Corollary~\ref{cor:boundary-nonsplit-self-twist} excludes the corresponding tame quadratic self-twist with irreducible intertwiner. Whether the remaining non-Cartan form occurs for actual full $p^{r+3}$-torsion remains open. The detailed reducible point-field profiles have their own explicitly stated local hypotheses and proofs; the full-class first-wild exclusion does not replace those local arguments. In particular the exchanged-high-line and large-wild-image restrictions remain useful as local comparison statements, while the graph valuation bound excludes the entire unequal-character comparison at $e=p(p-1)$.

For the universal level deformation at weights $(p^2+p+1,p^2)$, Proposition~\ref{prop:divisorial-full-cartan} gives the exact full principal Cartan invariants at every level: $e=f=p^{n-1}$ and residue field $k(t^{p^{n-1}})$. The separate scalar calculation at level four gives relative invariants $e=p^2,f=p$ between the scalar and full Cartan fixed fields. At $p=5,n=4$, the additional tame quotient of ramification index $12$ with unchanged residue field has absolute index exactly $6500$. Proposition~\ref{prop:cartan-torus-regularity} additionally proves that the actual quotient is regular, with regular Cartier reduced vertical divisor, at every nonzero point of the exceptional torus. Corollary~\ref{cor:cartan-torus-specialization} gives $78000\mid e_L$ for every local specialization of those level-four charts; if $e(L/K)\le12$, it follows that $e_K\ge6500$. These chart identities also obstruct an actual boundary comparison when its finite image is contained in a principal nonsplit Cartan: Corollary~\ref{cor:boundary-nonsplit-cartan-image} constructs an invariant-ring map and contradicts the required divisibility, including over arbitrary perfect residue fields. A general comparison need not have that image, and the divisorial value alone is not a bound for arbitrary normal models. At a normal surface point with a single vertical multiplicity $6500$, Lemma~\ref{lem:cartier-index-contact} shows that an index-$260$ specialization requires Cartier index divisible by $25$. The known cyclic quotient in Example~\ref{ex:noncartier-quotient-specialization} exhibits the corresponding denominator phenomenon at vertical multiplicity $1300$ and Cartier index $5$.

Under the actual-comparison and unit-preserving chart hypotheses of Proposition~\ref{prop:boundary-specialization-contact}, an upstairs specialization has contact divisible by $5$. A downstairs transverse curve may have precisely such contact. What remains is to analyze the non-Cartan images and formal models compatible with the actual finite-level isomorphism and unequal special characters. Proposition~\ref{prop:cartan-orbit-data} proves that the point-orbit sizes alone do not force Cartan containment; the full congruence image can be strictly larger and noncommutative. The exchanged residual lines do not themselves identify the two deformations as quadratic twists. The remaining degree-$p^6$ boundary configuration is neither constructed nor excluded.

For the first numerical case, $p=5$ and $e=260$, the index convention gives
\[
                         3\le\nuMin(5,260,2)\le4.
\]
Indeed, Theorem~\ref{thm:mixed-all-indices} supplies an unequal-character
$p^3$ comparison since $260>5^{6-3-1}=25$, whereas
Theorem~\ref{thm:all-index-upper} gives recovery from $p^5$-torsion
since $260\le5^4$. Recovery from actual $p^4$-torsion would give
index three; the existence of an unequal-character $p^4$ comparison
over an allowed field of index $260$ would give index four.
The nonsplit-Cartan exclusion restricts that undecided comparison
without deciding between these two values.

The tame all-prime and wild characteristic-three existence criteria allow the base field to vary. Corollary~\ref{cor:cyclotomic-full-p} settles the particular fields $\Q_p(\zeta_p)$, but these criteria do not classify which other fixed fields of a given index carry each example, nor do they determine all higher-height analogues. Proposition~\ref{prop:height-one-index} settles height one for every prime; the remaining higher-height questions have a different scope. Likewise, the character index concerns a group with potentially good graded pieces, whereas full rational torsion is a property of one actual representation. Keeping these quantifiers distinct is necessary when applying the theorems to a particular elliptic curve or Barsotti--Tate group.


The three necessary configurations produced by the primitive-grading reduction are all excluded by Corollary~\ref{pg-cor:exclusion}. The central criterion therefore closes that comparison argument as well as the shorter graph-valuation proof. The remaining index questions in this section concern absolute ramification divisible by $p$; the full-class prime-to-$p$ table has no unresolved interval.

\section*{Method and records}
The mathematics and text were produced and checked by GPT-6 Astra (ChatGPT Pro mode). No human has checked this paper. Checks include reconstruction from statements, complete proof comparison and direct verification of hypotheses. All checks used the same model; reasoning-effort metadata is not exposed, and no different-model independence is claimed. The workbench records the precise boundaries of each check.

Separate reconstructions cover central rigidity, tame elliptic criteria, ordinary reduction and the first-level elliptic models. The complete prime-to-$p$ index, all-index mixed criterion, exact wild bands, arbitrary-index primitive comparisons, the structural maximum identity, the first-boundary normal form, proper canonical point-field bounds, persistence of point-field translations, the full principal Cartan at every level, its regular exceptional-torus quotient, the scalar subgroup calculation and the height-one assertion have statement-led reconstructions using the preceding theorems as disclosed inputs. The elementary orbit and quadratic self-twist diagnostics also have a separate reconstruction. The actual nonsplit-Cartan exclusion over arbitrary perfect residue fields has a separate statement-led reconstruction from the preceding normal-form, universal-frame and invariant-ring theorems. The Hesse defect has two complete field constructions. Actual point fields, quotient transport, integral descent and extension lifting are proved in the text. Independently written arithmetic programs and unchanged original-program reruns supplement the arguments; finite computations do not prove the field or geometric assertions. Remaining questions are stated in Section~\ref{sec:open}.

The text is dedicated to the public domain under CC0 1.0. Editable sources and exact-arithmetic programs are prepared at \url{\workbenchurl}; full mathematical and checking records are available at \url{\recordurl}.

\begin{thebibliography}{99}
\bibitem[Berger(2004)]{Berger2004}
Berger, L. (2004). An introduction to the theory of $p$-adic representations. In \emph{Geometric aspects of Dwork theory} (Vol.~I, pp.~255--292). Walter de Gruyter. Exact example and Sen criterion: Section~II.1.2, p.~9 of the 2002 author version, \url{https://arxiv.org/abs/math/0210184}. Author publication list: \url{https://perso.ens-lyon.fr/laurent.berger/articles.php}.

\bibitem[Edixhoven, van der Geer, and Moonen(n.d.)]{EdixhovenGeerMoonen}
Edixhoven, B., van der Geer, G., \& Moonen, B. (n.d.). \emph{Abelian varieties (preliminary version of the first chapters)}. Author manuscript. Frobenius, Verschiebung and geometric torsion: Propositions~5.20 and~5.22, pp.~80--81. Author-hosted Chapter~5: \url{https://www.math.ru.nl/~bmoonen/BookAV/Isogs.pdf}.

\bibitem[Milne(2008)]{Milne2008AV}
Milne, J. S. (2008). \emph{Abelian varieties} (Version~2.00). Author course notes. Weil pairing: Section~I.13, p.~57; N\'eron mapping property: Theorem~17.1 and Remark~17.2, p.~70. \url{https://www.jmilne.org/math/CourseNotes/AV.pdf}.

\bibitem[Anema, Top, and Tuijp(2018)]{AnemaTopTuijp2018}
Anema, A. S. I., Top, J., \& Tuijp, A. (2018). Hesse pencils and 3-torsion structures. \emph{SIGMA, 14}, 102. \href{https://doi.org/10.3842/SIGMA.2018.102}{doi:10.3842/SIGMA.2018.102}.
\bibitem[Beaudry et al.(2019)]{BDMMRR2019}
Beaudry, A., Downey, N., McCranie, C., Meszar, L., Riddle, A., \& Rock, P.
(2019).
\newblock Computations of orbits for the Lubin--Tate ring.
\newblock \emph{Journal of Homotopy and Related Structures, 14}(3), 691--718.
\newblock \url{https://doi.org/10.1007/s40062-018-00228-7}.
\newblock Author manuscript: \url{https://arxiv.org/abs/1801.07559}.

\bibitem[Bondarko(2006)]{Bondarko2006}
Bondarko, M. V. (2006). The generic fibre of finite group schemes; a ``finite wild'' criterion for good reduction of Abelian varieties. \emph{Izvestiya: Mathematics, 70}(4), 661--691. \href{https://doi.org/10.1070/IM2006v070n04ABEH002323}{doi:10.1070/IM2006v070n04ABEH002323}.
\bibitem[Carter and Strauch(2021)]{CarterStrauch2021}
Carter, A., \& Strauch, M. (2021). Lubin--Tate deformation spaces and fields of norms. \emph{Journal de Th\'eorie des Nombres de Bordeaux, 33}(2), 433--458. \href{https://doi.org/10.5802/jtnb.1166}{doi:10.5802/jtnb.1166}.
\bibitem[Demb\'el\'e, Dieulefait, Freitas, and Pacetti(2026)]{DembeleEtAl2026}
Demb\'el\'e, L., Dieulefait, L. V., Freitas, N., \& Pacetti, A. (2026). \emph{Serre modularity over real quadratic fields} (PDF dated 16 November 2026, version~10). Hexagon 2611.00071v3. \url{https://hexagonmath.org/pdf/2611.00071v3}.
\bibitem[Drinfeld(1974)]{Drinfeld1974}
Drinfeld, V. G. (1974). Elliptic modules. \emph{Mathematics of the USSR-Sbornik, 23}(4), 561--592. \href{https://doi.org/10.1070/SM1974v023n04ABEH001731}{doi:10.1070/SM1974v023n04ABEH001731}. Original deformation and full-level theorems: Propositions~4.2--4.3, pp.~570--573.
\bibitem[Grothendieck(ca.~1969)]{Grothendieck1969}
Grothendieck, A. (ca.~1969). \emph{Groupe de Barsotti--Tate et cristaux (jusqu'en \'et\'e 1969)} [Unpublished notes, \emph{Cristaux}, cote~9, sheets~10--13]. Universit\'e de Montpellier, Institut Montpelli\'erain Alexander Grothendieck. \url{https://grothendieck.umontpellier.fr/9.pdf}.
\bibitem[Hattori(2008)]{Hattori2008}
Hattori, S. (2008). Tame characters and ramification of finite flat group schemes. \emph{Journal of Number Theory, 128}(5), 1091--1108. \href{https://doi.org/10.1016/j.jnt.2007.05.006}{doi:10.1016/j.jnt.2007.05.006}.
\bibitem[Hopkins and Gross(1994)]{HopkinsGross1994}
Hopkins, M. J., \& Gross, B. H. (1994). Equivariant vector bundles on Lubin--Tate moduli space. In \emph{Topology and representation theory} (Contemporary Mathematics, Vol.~158, pp.~23--88). American Mathematical Society. \href{https://doi.org/10.1090/conm/158/01453}{doi:10.1090/conm/158/01453}.
\bibitem[Katz(1981)]{Katz1981}
Katz, N. M. (1981). Serre--Tate local moduli. In \emph{Surfaces alg\'ebriques} (Lecture Notes in Mathematics, Vol.~868, pp.~138--202). Springer. \href{https://doi.org/10.1007/BFb0090648}{doi:10.1007/BFb0090648}.
\bibitem[Kisin(2006)]{Kisin2006}
Kisin, M. (2006). Crystalline representations and $F$-crystals. In \emph{Algebraic geometry and number theory} (Progress in Mathematics, Vol.~253, pp.~459--496). Birkh\"auser. \href{https://doi.org/10.1007/978-0-8176-4532-8_7}{doi:10.1007/978-0-8176-4532-8\_7}.
\bibitem[Koch(2015)]{Koch2015}
Koch, A. (2015). \emph{Primitively generated Hopf orders in characteristic $p$}. arXiv:1509.07393. \url{https://arxiv.org/abs/1509.07393}.

\bibitem[Liu(2007)]{Liu2007}
Liu, T. (2007).
\newblock Potentially good reduction of Barsotti--Tate groups.
\newblock \emph{Journal of Number Theory, 126}(2), 155--184.
\newblock \url{https://doi.org/10.1016/j.jnt.2006.11.008}.

\bibitem[Liu(2013)]{Liu2013}
Liu, T. (2013).
\newblock The correspondence between Barsotti--Tate groups and Kisin modules
when $p=2$.
\newblock \emph{Journal de Th\'eorie des Nombres de Bordeaux, 25}(3), 661--676.
\newblock \url{https://www.numdam.org/article/JTNB_2013__25_3_661_0.pdf}.

\bibitem[Lubin and Tate(1965)]{LT1965}
Lubin, J., \& Tate, J. (1965).
\newblock Formal complex multiplication in local fields.
\newblock \emph{Annals of Mathematics, 81}(2), 380--387.
\newblock \url{https://doi.org/10.2307/1970622}.
\newblock Primary scan:
\url{https://www.sas.rochester.edu/mth/sites/doug-ravenel/otherpapers/lubin-tate.pdf}.

\bibitem[Lubin and Tate(1966)]{LT1966}
Lubin, J., \& Tate, J. (1966).
\newblock Formal moduli for one-parameter formal Lie groups.
\newblock \emph{Bulletin de la Soci\'et\'e Math\'ematique de France, 94}, 49--59.
\newblock \url{https://doi.org/10.24033/bsmf.1633}.


\bibitem[Lubin(1967)]{Lubin1967}
Lubin, J. (1967).
\newblock Finite subgroups and isogenies of one-parameter formal Lie groups.
\newblock \emph{Annals of Mathematics, 85}(2), 296--302.

\bibitem[Melistas(2022)]{Melistas2022}
Melistas, M. (2022). Purely additive reduction of abelian varieties with torsion. \emph{Journal of Number Theory, 239}, 21--39. \href{https://doi.org/10.1016/j.jnt.2021.10.015}{doi:10.1016/j.jnt.2021.10.015}.
\bibitem[OpenAI(2026)]{OpenAI2026}
OpenAI. (2026, October 6). \emph{Fontaine--Mazur modularity at the prime 2} [Preprint, commit \texttt{4eaee8754825}]. \url{https://github.com/openai/math/tree/4eaee87548250b2682affda7d05d947c05554bae/preprints/Fontaine-Mazur-modularity-at-the-prime-2-October-6-2026}.
\bibitem[Ozeki and Yoshida(2026)]{OzekiYoshida2026}
Ozeki, Y., \& Yoshida, M. (2026). \emph{Fields of definition of $p$-torsion points of elliptic curves and their ramification}. arXiv:2609.06554. \url{https://arxiv.org/abs/2609.06554}.
\bibitem[Serre and Tate(1968)]{SerreTate1968}
Serre, J.-P., \& Tate, J. (1968). Good reduction of abelian varieties. \emph{Annals of Mathematics, 88}(3), 492--517. \href{https://doi.org/10.2307/1970722}{doi:10.2307/1970722}.
\bibitem[Serre(1979)]{Serre1979}
Serre, J.-P. (1979). \emph{Local fields} (M. J. Greenberg, Trans.; Graduate Texts in Mathematics, Vol.~67). Springer. \href{https://doi.org/10.1007/978-1-4757-5673-9}{doi:10.1007/978-1-4757-5673-9}.
\bibitem[Smith(2023)]{Smith2023}
Smith, H. (2023). Ramification in division fields and sporadic points on modular curves. \emph{Research in Number Theory, 9}, Article~17. \href{https://doi.org/10.1007/s40993-023-00424-2}{doi:10.1007/s40993-023-00424-2}. Author version: \url{https://arxiv.org/abs/1810.04809}.
\bibitem[Strauch(2006)]{Strauch2006}
Strauch, M. (2006). \emph{Deformation spaces of one-dimensional formal modules and their cohomology} [Author preprint, arXiv:math/0611109v1]. \url{https://arxiv.org/abs/math/0611109v1}. Regular parameters and compatible level presentation: Theorem~2.1.2 and proof, pp.~9--11; level action: \S2.2, pp.~12--14.
\bibitem[Tate(1967)]{Tate1967}
Tate, J. T. (1967). $p$-divisible groups. In T. A. Springer (Ed.), \emph{Proceedings of a conference on local fields} (pp.~158--183). Springer. \href{https://doi.org/10.1007/978-3-642-87942-5_12}{doi:10.1007/978-3-642-87942-5\_12}.




\bibitem[Colmez(n.d.)]{ColmezLocal}
Colmez, P. (n.d.). \emph{Corps locaux, notes du cours de M2}. Author course notes. \url{https://webusers.imj-prg.fr/~pierre.colmez/CL.pdf}. Quotient ramification formula: Proposition~3.7, pp.~14--15; Herbrand theorem and upper quotient compatibility: Theorem~3.10 and Proposition~3.13, pp.~15--16.
\bibitem[The Stacks Project Authors(2026)]{StacksCohen}
The Stacks Project Authors. (2026). \emph{The Cohen structure theorem}, Theorem~10.160.8 (Tag~032A). \url{https://stacks.math.columbia.edu/tag/032A}.

\bibitem[The Stacks Project Authors(2026)]{StacksDVR}
The Stacks Project Authors. (2026).
\emph{Extensions of discrete valuation rings}, Section~15.113,
especially Lemma~15.113.3 and Remark~15.113.6 (Tag~0EXQ).
\url{https://stacks.math.columbia.edu/tag/0EXQ}.

\bibitem[Bisatt et~al.(2026)]{BisattEtAl2026}
Bisatt, M., Furio, L., \& Lombardo, D. (2026). \emph{Explicit $p$-adic Hodge theory for elliptic curves and non-split Cartan images} (arXiv:2603.04021v2, revised June~8). \url{https://arxiv.org/abs/2603.04021v2}.

\bibitem[The Stacks Project Authors(2026)]{StacksRegularNormal}
The Stacks Project Authors. (2026).
\emph{A regular ring is normal}, Lemma 10.157.5 (Tag 0567).
\url{https://stacks.math.columbia.edu/tag/0567}.

\bibitem[The Stacks Project Authors(2026)]{StacksHeightOneIntersection}
The Stacks Project Authors. (2026).
\emph{Height-one intersection for a Noetherian normal domain},
Lemma 10.157.6 (Tag 031T).
\url{https://stacks.math.columbia.edu/tag/031T}.

\bibitem[Lorenzini(2013)]{Lorenzini2013WildQuotients}
Lorenzini, D. (2013). Wild quotient singularities of surfaces.
\emph{Mathematische Zeitschrift, 275}, 211--232.
\url{https://doi.org/10.1007/s00209-012-1132-7}.

\bibitem[Lorenzini and Schr\"oer(2023)]{LorenziniSchroer2023WildQuotients}
Lorenzini, D., \& Schr\"oer, S. (2023).
Discriminant groups of wild cyclic quotient singularities.
\emph{Algebra \& Number Theory, 17}(5), 1017--1068.
\url{https://doi.org/10.2140/ant.2023.17.1017}.

\bibitem[Liedtke and Matsumoto(2018)]{LiedtkeMatsumoto2018}
Liedtke, C., \& Matsumoto, Y. (2018). Good reduction of K3 surfaces.
\emph{Compositio Mathematica, 154}(1), 1--35.
\url{https://doi.org/10.1112/S0010437X17007400}.

\bibitem[Kir\'aly and L\"utkebohmert(2013)]{KiralyLutkebohmert2013}
Kir\'aly, F., \& L\"utkebohmert, W. (2013).
Group actions of prime order on local normal rings.
\emph{Algebra \& Number Theory, 7}(1), 63--74.
\url{https://doi.org/10.2140/ant.2013.7.63}.

\end{thebibliography}

\end{document}
